% =====================================================================
%  Data-Center GPU Platforms Compared: Lecture Notes
%  Companion notes to the 33-slide deck "Data-Center GPU Platforms Compared".
%  One subsection per slide: what the slide shows, the details behind
%  it, definitions of every term, and the intuition.
%
%  Compile:  pdflatex, makeindex, pdflatex, pdflatex
%  Needs only standard packages (mathptmx, tcolorbox, imakeidx, ...).
% =====================================================================
\documentclass[11pt,letterpaper]{article}
\usepackage[T1]{fontenc}
\usepackage{mathptmx}
\usepackage[scaled=0.9]{helvet}
\usepackage{courier}
\usepackage[margin=1in,headheight=14pt]{geometry}
\usepackage{amsmath,amssymb,bm}
\usepackage{booktabs,array}
\usepackage{xcolor}
\usepackage{microtype}
\usepackage{enumitem}
\usepackage{titlesec}
\usepackage[most]{tcolorbox}
\usepackage{imakeidx}
\usepackage{fancyhdr}
\PassOptionsToPackage{hyphens}{url}
\usepackage[hidelinks,bookmarksnumbered]{hyperref}

% ---------- palette (same as the deck) --------------------------------
\definecolor{ink}{HTML}{1F2937}
\definecolor{muted}{HTML}{6B7280}
\definecolor{tok}{HTML}{5F6F85}
\definecolor{emb}{HTML}{0CA678}
\definecolor{posc}{HTML}{7048E8}
\definecolor{attn}{HTML}{E8590C}
\definecolor{mlp}{HTML}{1C7ED6}
\definecolor{outc}{HTML}{D6336C}
\definecolor{trn}{HTML}{9C6644}
\definecolor{vc}{HTML}{1098AD}
\colorlet{stagec}{ink}
\newcommand{\setstage}[1]{\colorlet{stagec}{#1}}

% ---------- headings ---------------------------------------------------
\titleformat{\section}{\sffamily\bfseries\Large\color{ink}}{\thesection}{0.8em}{}[{\color{stagec}\titlerule[1.2pt]}]
\titleformat{\subsection}{\sffamily\bfseries\large\color{ink}}{\thesubsection}{0.7em}{}
\titleformat{\subsubsection}{\sffamily\bfseries\color{ink}}{\thesubsubsection}{0.6em}{}
\titlespacing*{\section}{0pt}{20pt}{10pt}
\titlespacing*{\subsection}{0pt}{16pt}{6pt}

\pagestyle{fancy}
\fancyhf{}
\fancyhead[L]{\sffamily\small\color{muted}\nouppercase{\leftmark}}
\fancyhead[R]{\sffamily\small\color{muted}\thepage}
\renewcommand{\headrulewidth}{0pt}

\makeindex[title=Index of terms,intoc,columns=2]

% ---------- slide-note machinery -------------------------------------
% \slide{number}{title}{subtitle}
\newcommand{\slide}[3]{%
  \subsection[Slide #1: #2]{Slide #1: #2}%
  {\color{muted}\itshape #3\par}\medskip}
\newcommand{\para}[1]{\medskip\noindent\textbf{#1}\quad}

\tcbset{
  termsbox/.style={breakable,enhanced,colback=stagec!4,colframe=stagec!60,
    boxrule=0.5pt,arc=2pt,left=6pt,right=6pt,top=4pt,bottom=4pt,
    title=Terms,fonttitle=\sffamily\bfseries\small,coltitle=white,colbacktitle=stagec!80,
    attach boxed title to top left={yshift=-2mm,xshift=4mm},boxed title style={arc=1.5pt,boxrule=0pt}},
  intuitionbox/.style={breakable,enhanced,colback=white,colframe=white,
    borderline west={3pt}{0pt}{stagec},left=8pt,right=2pt,top=2pt,bottom=2pt,arc=0pt,boxrule=0pt,
    title=Intuition,fonttitle=\sffamily\bfseries\small,coltitle=stagec,colbacktitle=white,
    attach boxed title to top left={yshift=-1mm,xshift=0mm},boxed title style={boxrule=0pt,left=8pt}},
  formulabox/.style={enhanced,colback=ink!4,colframe=ink!30,boxrule=0.4pt,arc=2pt,
    left=6pt,right=6pt,top=4pt,bottom=4pt}
}
\newenvironment{terms}{\begin{tcolorbox}[termsbox]\begin{description}[leftmargin=1.1em,itemsep=3pt,topsep=2pt,font=\normalfont]}{\end{description}\end{tcolorbox}}
\newcommand{\term}[2]{\item[\textbf{#1}]\index{#1} #2}
\newenvironment{intuition}{\begin{tcolorbox}[intuitionbox]}{\end{tcolorbox}}
\newenvironment{formulas}{\begin{tcolorbox}[formulabox]}{\end{tcolorbox}}

\setlist[itemize]{leftmargin=1.4em,itemsep=2pt}
\setlength{\emergencystretch}{1.5em}
\setlist[enumerate]{leftmargin=1.6em,itemsep=2pt}
\newcommand{\paper}[1]{\textit{#1}}
\newcommand{\passk}{\mathrm{pass}@k}
\newcommand{\E}{\mathbb{E}}
\DeclareMathOperator{\clip}{clip}






\begin{document}

\begin{titlepage}
\thispagestyle{empty}
\vspace*{2.5cm}
{\color{attn}\rule{3cm}{3pt}}\par\vspace{0.6cm}
{\sffamily\bfseries\color{muted}\large LECTURE NOTES \ $\cdot$ \ HARDWARE, 2026\par}
\vspace{0.4cm}
{\sffamily\bfseries\fontsize{30}{34}\selectfont\color{ink} Data-Center GPU\\Platforms Compared\par}
\vspace{0.3cm}
{\sffamily\fontsize{18}{22}\selectfont\color{ink} Companion notes to the slides\par}
\vspace{1.2cm}
{\large\color{muted}\parbox{0.82\textwidth}{NVIDIA from the HGX A100 to the Vera Rubin NVL72, AMD Instinct from the MI100 to the MI455X and the Helios rack, and Intel from the Gen9 integrated graphics through Xe, Ponte Vecchio and Aurora to Gaudi and the GPU D-series: the compute units (SM, CU, Xe-core), the packages (dies, chiplets, tiles), memory per GPU, FP64 against FP4, the scale-up fabrics (NVLink, Infinity Fabric, Xe Link, UALink), the racks, the software stacks (CUDA, ROCm, oneAPI) and what the TOP500 and MLPerf measure. One section per slide, every term defined, and a reading list for each part.}\par}
\vspace{1.5cm}
{\color{muted}For advanced undergraduates and beginning graduate students\par}
\vspace{0.3cm}
{\color{muted}November 2026\par}
\vfill
{\small\color{muted}\parbox{0.9\textwidth}{These notes follow the deck \textit{Data-Center GPU Platforms Compared} (33 slides) slide by slide. Numbers quoted from an architecture whitepaper or a paper were checked against that document; numbers from product datasheets, product pages, press releases and press coverage are marked $\dagger$; numbers marked ``derived'' were computed from quoted ones and the formulas are given. Tensor peaks are dense unless marked sparse. Where a vendor has not published a figure the notes say so (``n.p.'') rather than estimate it. The account runs to November 2026: Rubin and the MI455X were shipping, Helios racks had begun to ship, and Intel's GPU D-series was sampling.}\par}
\end{titlepage}

\tableofcontents
\clearpage

\section*{How to use these notes}
\addcontentsline{toc}{section}{How to use these notes}
Each section of these notes corresponds to one section of the deck, and each subsection to one slide. A subsection has four parts:
\begin{itemize}
  \item \textbf{What the slide shows}: a description of the figure or table, so the notes can be read without the slides.
  \item \textbf{Details}: the material the slide compresses, with the arithmetic written out and the numbers given their sources.
  \item A \textbf{Terms} box defining every technical term the slide uses, the first time it appears. All terms are collected in the index at the end.
  \item An \textbf{Intuition} box: the picture to keep in mind, and the question to ask when the picture seems to fail.
\end{itemize}
Conventions. A \emph{GPU} here is one data-center accelerator package as sold: an NVIDIA SXM module (one or two dies), an AMD Instinct OAM module (two to twelve dies), an Intel Max 1550 (47 active tiles) or a Gaudi 3. A \emph{platform} is the package on its board, in its rack, with its fabric and software. The compute units are the NVIDIA \emph{SM}, the AMD \emph{CU} (compute unit; the MI455X counts \emph{WGPs}, work-group processors) and the Intel \emph{Xe-core} (\emph{EU} and \emph{subslice} in the Gen families); matrix units are \emph{tensor cores}, \emph{matrix cores} and \emph{XMX} engines respectively. Peak FLOP/s counts a fused multiply-add as two operations; tensor peaks are \emph{dense} unless marked sparse (NVIDIA headlines 2:4-sparse figures, twice the dense ones; AMD and Intel usually quote dense). FP64 is given as vector / matrix where the two differ. Bandwidths use decimal prefixes (1 TB/s $=10^{12}$ B/s). A dagger ($\dagger$) marks any number whose source is a datasheet, product page, press release, benchmark list or press coverage rather than an architecture whitepaper or a paper; ``n.p.'' means not published by the vendor as of November 2026.

A reading plan for a single lecture: Sections 1--2 (the overview, the NVIDIA line) take about a fifth of the session; Section 3 (AMD) a fifth; Sections 4--5 (Intel's Gen and Xe families, Ponte Vecchio and Aurora) a quarter; Sections 6--7 (fabrics, racks, standards, software) a fifth; Sections 8--10 (systems, benchmarks, wrap-up) the rest. The exercises in Appendix~C are arithmetic on the deck's own numbers and can be set as homework. Appendix~D is the reading list, section by section, and matches the paper kit that accompanies these notes.
% =====================================================================
\setstage{ink}
\section{Overview: three vendors, six levels (slides 1--3)}

\slide{1}{Data-Center GPU Platforms Compared}{NVIDIA, AMD and Intel, 2020--2026}

\para{What the slide shows} The cover carries the title, a one-paragraph description of the lecture, and a drawing of three accelerator packages seen from above, one per vendor, each with its memory stacks along the sides. At the top, NVIDIA's Blackwell: two large gray rectangles, the two compute dies, with a thin dark bar between them for the die-to-die link and four HBM3e stacks on each side. In the middle, AMD's MI300X: four pale green base rectangles, the I/O dies, with eight narrower bright green strips across them, the compute dies, and four HBM3 stacks on each side. At the bottom, Intel's Ponte Vecchio: two orange base tiles, each carrying eight compute tiles in two rows with a strip of cache tiles between the rows, and four HBM2e stacks on each side. The caption calls the three drawings ``three answers to the reticle limit'' and dates them: Blackwell 2024, MI300X 2023, Ponte Vecchio 2023. The description paragraph lists what the lecture covers: NVIDIA from the A100 to Vera Rubin, AMD Instinct from the MI100 to the MI455X and the Helios rack, Intel from the Gen graphics families through Xe and Ponte Vecchio to what followed; the compute units, the packages, the memory, the fabrics, the racks and the software.

\para{Details} This lecture is the third in a sequence. The first, \emph{Introduction to GPU Architectures}, explained how a graphics chip became a throughput processor and ended with NVIDIA's Volta (2017). The second, \emph{NVIDIA GPU Architectures after Volta}, followed one vendor's compute dies from Turing (2018) to Rubin (2026). This one widens the view in two directions at once. It adds the two other companies that build data-center GPUs, AMD and Intel, and it moves the unit of comparison up from the die to the \emph{platform}: the package as sold, the board it sits on, the rack it fills, the fabric that joins the GPUs, and the software that programs them. The thesis, stated on slide 32, is that the dies converged while the platforms did not; the purpose of the lecture is to show where each level is decided and by what.

The three packages on the cover are the lecture's first comparison, and they illustrate a problem all three vendors faced at the same moment. A lithography stepper exposes one rectangular field at a time, and the largest field, the reticle limit, is about 830 mm$^2$ (858 mm$^2$ on the most common definition, less once the scribe lines are removed). NVIDIA's A100 (826 mm$^2$) and H100 (814 mm$^2$) were monolithic dies at that limit; so were the largest dies of every earlier generation back to Pascal. A product that wanted more transistors than fit in one field had to be built from several pieces of silicon, and the three vendors chose three different ways of doing it, in the reverse order of the drawing.

Intel went first and furthest. Ponte Vecchio, announced in August 2021 and launched as the Data Center GPU Max series on January 10, 2023, stacks 47 active tiles from five process nodes and two foundries: 16 compute tiles on TSMC N5, 8 ``RAMBO'' cache tiles on Intel 7, 2 base tiles on Intel 7 of 646 mm$^2$ each, 2 Xe Link tiles on TSMC N7, 8 HBM2e stacks, joined by 11 embedded silicon bridges (EMIB); 16 thermal tiles on top make 63 pieces in all$^\dagger$. The compute and cache tiles sit face-down on the base tiles (Intel's Foveros 3D stacking); the two halves and the memory are bridged by EMIB (2.5D). The whole carries over 100 billion transistors and 2{,}330 mm$^2$ of active silicon on a 4{,}844 mm$^2$ package$^\dagger$. Slide 17 takes the package apart.

AMD's MI300X, launched December 6, 2023, is the second design. Four I/O dies (IODs) on TSMC N6 form the base; each carries 64 MB of a 256 MB Infinity Cache, two HBM3 controllers and the Infinity Fabric; eight accelerator complex dies (XCDs) on TSMC N5, each with 38 active compute units of 40, are hybrid-bonded on top, two per IOD; eight HBM3 stacks of 24 GB ring the package for 192 GB at 5.3 TB/s$^\dagger$. The transistor count is 153 billion across twelve dies. The same base with six XCDs and three Zen 4 CPU dies is the MI300A, the processor of El Capitan (slide 9).

NVIDIA's Blackwell, announced March 18, 2024, is the simplest of the three in construction and the largest in transistor count: two reticle-sized dies on TSMC 4NP, 208 billion transistors together, joined edge to edge by the NV-HBI interface at 10 TB/s so that the pair presents itself to software as one GPU with one L2; eight stacks of HBM3e give 180 to 288 GB depending on the product (slide 4)$^\dagger$. There are no separate I/O or cache dies: everything but memory is on the two compute dies. Rubin (2026) keeps this construction with two 3 nm dies and 336 billion transistors; AMD's MI455X (2026) goes the other way to twelve dies (8 compute on TSMC N2, 2 cache, 2 I/O on N3P); Intel built no successor to Ponte Vecchio's package.

Why the three differ is a question of what each company already knew how to do. Intel owns Foveros and EMIB and had reasons to show them; AMD had shipped chiplet CPUs since 2017 and brought the same partition of compute dies on an I/O die to the GPU; NVIDIA, with the highest-volume product and the most to lose from yield surprises, chose the fewest pieces. The cost of each choice appears later in the lecture: Ponte Vecchio's 47 tiles came a year late and ran at 600 W (slide 17); the MI300X's I/O dies made its 192 GB possible but its all-to-all GPU links share one fabric (slide 21); Blackwell's two dies shipped nine months after announcement, after a mask change, and then in volume (slide 4).

The reading list for each section is in Appendix~D and matches the paper kit that accompanies these notes. The architecture whitepapers (AMD's CDNA 1--4, Intel's Gen9, Gen11 and Xe documents, NVIDIA's through Hopper) and the conference papers (Ponte Vecchio at ISSCC 2022, the MI300X at ISSCC and Hot Chips 2024, Vera Rubin at Hot Chips 2026) are the primary sources; product pages, datasheets and press releases are marked with the dagger throughout.

\begin{terms}
\term{platform}{The accelerator package together with its board, rack, scale-up fabric and software stack: the unit a data center actually buys. The lecture compares platforms, not only dies.}
\term{package}{The assembly of one or more silicon dies and memory stacks on a substrate that is sold as one GPU: an SXM or OAM module.}
\term{reticle limit}{The largest area a lithography stepper exposes in one shot, about 830 square millimeters; a die cannot be larger, so a product that needs more silicon must use several dies.}
\term{die}{One piece of silicon cut from a wafer; a chiplet or tile is a die designed to be combined with others in one package.}
\term{chiplet}{AMD's term for a small die that is combined with others in one package; on the MI300X the compute chiplets (XCDs) sit on I/O dies (IODs).}
\term{tile}{Intel's term for a die in a multi-die package; Ponte Vecchio has compute, cache (RAMBO), base and Xe Link tiles.}
\term{XCD}{Accelerator complex die: an AMD compute chiplet of 40 CDNA 3 compute units (38 enabled on the MI300X), eight per package.}
\term{IOD}{I/O die: the AMD base die carrying the Infinity Cache, the memory controllers and the Infinity Fabric; four per MI300X.}
\term{base tile}{The large Intel 7 die at the bottom of each Ponte Vecchio stack, carrying the L2 cache, memory controllers and the interconnect that the compute tiles sit on.}
\term{HBM}{High-bandwidth memory: DRAM dies stacked vertically and placed beside the GPU on the package, giving terabytes per second over a wide, short bus. Generations HBM2, HBM2e, HBM3, HBM3e, HBM4 appear in this lecture.}
\term{NV-HBI}{NVIDIA High-Bandwidth Interface: the 10 TB/s connection between Blackwell's two dies that lets them act as one GPU.}
\term{Foveros}{Intel's 3D stacking technology, in which active dies are bonded face to face on a base die; used in Ponte Vecchio.}
\term{EMIB}{Embedded multi-die interconnect bridge: a small piece of silicon embedded in the package substrate that joins the edges of two dies; Ponte Vecchio uses eleven.}
\term{hybrid bonding}{Joining two dies with direct copper-to-copper connections at very fine pitch rather than solder bumps; AMD's XCD-on-IOD stacking and its MI455X use it.}
\end{terms}

\begin{intuition}
The cover is the lecture in one picture: three companies hit the same wall, the size of a single die, within two years of each other, and each climbed it with the packaging it already owned. Everything that follows, from memory per GPU to the fabric to the software, is a consequence of choices made at this level and at the five above it.
\end{intuition}

% ---------------------------------------------------------------------
\slide{2}{Three vendors, one decade: the shape of the race}{Per-GPU tensor peaks from the three vendors since 2020; the lines converge on paper, the market does not.}

\para{What the slide shows} The chart on the left plots, on a logarithmic vertical axis from 100 TFLOP/s to 100 PFLOP/s, the fastest dense tensor peak of each vendor's newest data-center GPU against its year of first shipment. NVIDIA's line (gray circles) runs from the A100 at 312 TFLOP/s in 2020 through the H100 at 1{,}979 in 2022 to the GB200 at 10{,}000 in 2024, the GB300 at 15{,}000 in 2025 and Rubin at 35{,}000 in 2026. AMD's line (green squares) starts lower, at the MI100's 185 in late 2020 and the MI250X's 383 in 2021, reaches 2{,}615 with the MI300X in late 2023, meets NVIDIA at 10{,}000 with the MI355X in 2025, and ends above it at 40{,}260 with the MI455X in 2026. Intel appears as two orange triangles, Ponte Vecchio at 839 in 2023 and Gaudi 3 at 1{,}835 in 2024, and a hollow triangle at 1{,}300 in 2027 for the announced GPU D-series. The labels name the precision at which each peak is reached: FP16 for the first points, FP8 from 2022--23, FP4 from 2024--25. Two stat tiles on the right give the number of vendors (three, one of them at about 80\% of the market) and the number of exascale systems on the June 2026 TOP500 (five). The panel ``What the curves say'' states the three facts the chart and the tiles put together: the peaks converge on paper, the market has not followed, and HPC is the exception.

\para{Details} The chart compresses a decade of three product lines into one measure, so the first thing to understand is what the measure is and is not. Each point is the highest throughput the vendor states for its matrix units at the narrowest precision the part supports, counted without structured sparsity. For the A100 that is FP16 or BF16 with FP32 accumulation, 312 TFLOP/s; for the H100 it is FP8, 1{,}979; for Blackwell and Rubin it is 4-bit floating point (NVFP4), 10, 15 and 35 PFLOP/s; for AMD, FP16 on the MI100 and MI250X, FP8 on the MI300X, MXFP4 on the MI355X and MI455X. A line that rises by 112$\times$ from 2020 to 2026 (35{,}000 divided by 312; derived) is therefore not measuring one quantity: about 4$\times$ of it is the move from 16-bit to 4-bit operands, which quadruples the number of multiplies a fixed array of transistors does per clock, and the rest is more transistors, more matrix units per compute unit and higher clocks. The FP64 rates, which cannot be narrowed, rose far less over the same years and in NVIDIA's case fell (slide 12); the FP16 rates, which can be compared across the whole period, rose from 312 to 4{,}000 TFLOP/s, 13$\times$ (derived).

The vendor figures also need two conventions kept straight, and the lecture keeps them on every slide. First, dense against sparse. NVIDIA's headline numbers assume 2:4 structured sparsity, two zeros in every four weights, which doubles the throughput of a matrix multiply whose weights have been pruned to that pattern; the dense figure is half. The H100's often-quoted 3{,}958 TFLOP/s of FP8 is sparse; the dense figure plotted here is 1{,}979. AMD and Intel usually quote dense figures and state the sparse ones separately; the MI355X's 10 PFLOP/s of FP4 is dense (20 with sparsity), as is the MI455X's 40.26. Second, announced against shipped. Every point is placed at the year the part first shipped to customers, not the year it was announced, because the gap has grown: the MI455X was unveiled in July 2026 and shipped in volume in the third quarter; Rubin was declared in ``full production'' in January 2026 and the first racks shipped in August; Intel's GPU D-series was announced in October 2025, sampled from October 2026 and is placed at 2027 as a hollow marker because it has not shipped$^\dagger$.

With those conventions the chart's first message is that the two lines converge. In 2020 the A100 had 1.7$\times$ the MI100's dense FP16 peak; in 2022--23 the H100's FP8 and the MI300X's FP8 were within a third of each other (1{,}979 against 2{,}615, in AMD's favor, a year later); in 2025 the MI355X and the GB200 both stated 10 PFLOP/s of dense FP4; in 2026 the MI455X's 40.26 PFLOP/s exceeded Rubin's 35. Slide 6 lays the numbers out in a table and slide 32 draws the conclusion: on peak compute the two vendors' dies are, as of 2026, the same kind of object built from the same TSMC processes with comparable transistor counts (320 billion against 336 billion), and the difference between a 35 and a 40 is smaller than the difference between a peak and what a library achieves (slide 26).

Intel's two solid points are the exception that shows the rule. Ponte Vecchio's 839 TFLOP/s of FP16 (a vendor figure consistent with its confirmed 419 TFLOP/s of TF32) sat between the A100 and the H100 when it shipped in January 2023, a year later than planned and three months after the H100; it supported no FP8. Gaudi 3, a different architecture bought with Habana Labs in 2019, reached 1{,}835 TFLOP/s of BF16 or FP8 in 2024, close to the H100's dense FP8, but with no FP4 and no successor. The hollow 2027 marker, 1.3 PFLOP/s of FP8 on the GPU D-series, is an inference part on commodity memory, and the lecture treats it as what it is: a restart, not a continuation (slide 20).

The second message, carried by the first stat tile and the panel, is that the convergence of peaks has not been matched by any convergence of sales. The market-share figures are estimates by analysts, not vendor disclosures, and the deck marks them as secondary: NVIDIA held about 87\% of data-center AI accelerator revenue in 2024 and about 80\% in 2025, AMD between 6 and 8\%, Intel under 1\%$^\dagger$. The vendors' own numbers are consistent with that order of magnitude: AMD's data-center segment, which includes its EPYC CPUs, reported \$16.6 billion for 2025 and \$8.5 billion for the third quarter of 2026 alone, up 97\% in a year, while NVIDIA guided to about \$108 billion of revenue for a single quarter ending in October 2026$^\dagger$. Why a die that matches on paper sells a tenth as much is the question slides 21 to 26 answer: the fabric that joins GPUs into a training machine (slide 21), the rack it comes in (slide 22), the library that reaches the peak (slide 26) and the economics of who buys and when (slide 25).

The third message, the second stat tile, is that one market behaves differently. Of the five systems on the June 2026 TOP500 that exceed an exaflop on the HPL benchmark, three use AMD or Intel GPUs: El Capitan (AMD MI300A), Frontier (AMD MI250X) and Aurora (Intel Max 1550). One, JUPITER, uses NVIDIA GH200. The new number one, LineShine, uses Chinese CPUs and no accelerator at all$^\dagger$. The slide's panel counts the exascale systems by vendor; the exact split, three AMD-or-Intel, one NVIDIA and one without GPUs, is given on slide 28. In high-performance computing, which buys FP64 and bandwidth per dollar and per watt under multi-year contracts and ports its codes once, the vendor with the best CUDA libraries has not been the vendor of choice; slide 25 explains why the two markets differ.

\begin{terms}
\term{tensor peak}{The highest matrix-multiply throughput a vendor states for a GPU, in FLOP/s at a named precision; the matrix units (tensor cores, matrix cores, XMX engines) reach it, the vector lanes do not.}
\term{dense peak}{A matrix peak counted with every multiply-add performed; the figure plotted throughout this lecture.}
\term{structured sparsity (2:4)}{A pruning pattern in which two of every four weights are zero; hardware that skips the zeros doubles throughput, and NVIDIA's headline peaks assume it.}
\term{FP16 and BF16}{Sixteen-bit floating-point formats; FP16 has a 10-bit mantissa and a small range, BF16 an 8-bit exponent (the range of FP32) and a 7-bit mantissa. Both run at the same matrix rate on every GPU here.}
\term{FP8}{Eight-bit floating point (E4M3 or E5M2), first in hardware on the H100 (2022) and the MI300X (2023); twice the FP16 rate.}
\term{FP4}{Four-bit floating point (E2M1) with a shared scale per small block of values; NVIDIA's NVFP4 and the OCP MXFP4 format. Twice the FP8 rate; first shipped in 2024--25 on Blackwell and CDNA 4.}
\term{MXFP4 and MXFP8}{The Open Compute Project's microscaling formats: blocks of 32 values sharing one 8-bit power-of-two scale; AMD states its MI355X and MI455X peaks in them.}
\term{NVFP4}{NVIDIA's 4-bit format: blocks of 16 values with an FP8 scale and a second FP32 scale per tensor; the format of Blackwell's and Rubin's highest peaks.}
\term{year of first shipment}{The year a GPU first reached customers in volume, as distinct from the year it was announced; the lecture plots and tabulates by shipment.}
\term{TOP500}{The twice-yearly list of the 500 fastest computers ranked by the HPL benchmark (solving a dense linear system in FP64); exascale means more than $10^{18}$ FLOP/s on it.}
\term{HPL}{High-Performance Linpack, the TOP500 benchmark: FP64 dense LU factorization, rewarding matrix throughput and memory bandwidth per watt.}
\term{market share}{The fraction of data-center accelerator revenue or units a vendor sells; no vendor publishes it, so every figure in the lecture is an analyst's estimate marked as secondary.}
\end{terms}

\begin{intuition}
Read the chart as two statements that are both true: the engineering gap closed, the commercial gap did not. The rest of the lecture is a search for where the second gap lives, and the answer is never the die.
\end{intuition}

% ---------------------------------------------------------------------
\slide{3}{How to compare a platform, and the decade's timeline}{A GPU is sold as a package on a board in a rack with a fabric and a software stack; each level is a separate contest.}

\para{What the slide shows} A table on the upper left lists six levels at which the three vendors' platforms can be compared and the name each vendor gives its own: the compute unit (NVIDIA's SM with a warp of 32, AMD's CU with a wavefront of 64, Intel's Xe-core with a subgroup of 8 to 32); the package (GB200's two dies, the MI300X's eight XCDs on four IODs, Ponte Vecchio's 47 tiles); the board (SXM and the 8-GPU HGX, OAM and the 8-GPU UBB, OAM boards of four to eight GPUs); the scale-up fabric (NVLink with NVSwitch, Infinity Fabric, Xe Link); the rack (NVL72, Helios with 72 GPUs, none shipped by Intel); and the software (CUDA; ROCm and HIP; oneAPI and SYCL). To its right a timeline from 2015 to 2028 marks seven moments: Intel's Gen9 (2015) and AMD's MI25 (2017); the A100, MI100 and Xe-LP with oneAPI 1.0 (2020); the MI250X in Frontier, the H100 and Arc Alchemist (2021--22); Ponte Vecchio in Aurora, the MI300X and MI300A, and the cancellation of Rialto Bridge (2023); Blackwell, El Capitan at number one, Gaudi 3 and the forming of UALink (2024); the cancellation of Falcon Shores, the MI355X, the GB300 and the OpenAI--AMD 6 GW deal (2025); and, in color, Rubin, the MI455X and Helios shipping in 2026 with the MI430X and Jaguar Shores beyond. Two panels below give the lecture's method (``Why compare at every level'') and its summary of each vendor's decade in one sentence.

\para{Details} The table is the lecture's organizing device, and each of its rows is a section. The first row is the subject of the two earlier lectures for NVIDIA and of slides 7, 13 and 14 here for AMD and Intel. All three vendors build a GPU from a few hundred copies of one unit that holds vector lanes, a matrix unit, registers, a scratchpad and schedulers, and all three run groups of threads in lockstep on it; the names and the widths differ. NVIDIA's streaming multiprocessor (SM) has four partitions of 32 lanes and schedules warps of 32 threads; AMD's compute unit (CU) has four SIMD units of 16 lanes, each of which runs a wavefront of 64 work-items over four clocks; Intel's Xe-core has 8 or 16 vector engines, and the width of the group it schedules, the subgroup, is chosen by the compiler from 8, 16 or 32. Slide 14 puts the three side by side. The point for comparison is that a peak FLOP/s figure is the product of the number of these units, the lanes in each, two operations per fused multiply-add and the clock, so that two GPUs with the same peak may reach it with very different counts of units and clocks, and software tuned for one width will not be tuned for another.

The second row, the package, is the cover's subject and slide 17's for Ponte Vecchio, slide 9's for the MI300X, and slide 11's for the twelve-die MI455X. The third row, the board, is where a vendor's choice becomes a customer's: NVIDIA's SXM module and its 8-GPU HGX baseboard are NVIDIA designs that only NVIDIA GPUs fit, while AMD's and Intel's data-center GPUs ship on the Open Compute Project's OAM module and UBB baseboard, so that a server built for one fits the other (slide 23). The fourth row, the scale-up fabric, is the subject of slide 21: how many GPUs can read each other's memory at full rate, and through what. The fifth, the rack, is slide 22: NVIDIA has sold 72-GPU liquid-cooled racks since late 2024, AMD's first (Helios) shipped in the third quarter of 2026, and Intel has shipped none. The sixth, software, is slides 24 to 26.

The panel ``Why compare at every level'' gives the reason for the method in two examples. A faster die behind a slower fabric loses the training run, because a large model is trained on hundreds or thousands of GPUs exchanging gradients and activations every step, and the time spent waiting on the slowest link is paid by every GPU; slide 21 gives the measured case, the MI300X's 64 GB/s between a pair of GPUs against 450 GB/s on NVLink in 2024, and slide 26 the result, multi-node training 10 to 25\% slower on a GPU with a higher peak$^\dagger$. A better rack with a thinner library loses the deployment, because a cloud customer buys throughput on a working stack by the quarter and cannot wait for the port. The contests are decided separately, and the panel names the winners the lecture will argue for: Intel in packaging (the most tiles ever shipped), AMD in memory per GPU and FP64, NVIDIA in the fabric and the software. Slide 31 collects the scorecard in one table.

The timeline supplies the dates the rest of the lecture assumes. Its first marks are the starting points: Intel's Gen9 graphics in Skylake (2015), the integrated GPU whose execution unit and slice are the ancestors of Xe (slide 13), and AMD's Radeon Instinct MI25 (2017), the first product of a data-center line that until 2020 used the same Vega architecture as AMD's gaming cards. In 2020 three things happened within six months: NVIDIA's A100 (May), AMD's MI100 (November), the first GPU of the CDNA architecture split off from the graphics line, and Intel's Xe-LP graphics in Tiger Lake (September) with the oneAPI specification 1.0 (September 28)$^\dagger$. In 2021--22 AMD's MI250X built Frontier, the first exascale system (June 2022); NVIDIA announced and shipped the H100; and Intel shipped its first discrete gaming GPU in 24 years, Arc Alchemist (October 2022). In 2023 Ponte Vecchio finally launched (January) and Aurora was installed (June); AMD launched the MI300X and MI300A (December); and Intel cancelled Rialto Bridge, Ponte Vecchio's planned successor (March). In 2024 NVIDIA announced Blackwell (March) and shipped the first racks (November); El Capitan took the number one position (November); Intel launched Gaudi 3 (April); and eight companies including AMD and Intel formed the UALink consortium (October) to standardize a scale-up fabric. In 2025 Intel cancelled Falcon Shores as a product (January 30); AMD announced the MI355X (June) and shipped it in the third quarter; NVIDIA shipped the GB300 (July); and OpenAI committed to 6 GW of AMD Instinct GPUs (October 6). The colored final mark is the year of this lecture: Rubin racks from August, the MI455X and Helios from the third quarter, and beyond them AMD's MI430X for HPC (first half of 2027) and Intel's Jaguar Shores, which has no date$^\dagger$.

The second panel's three sentences are the lecture's thesis by vendor, and each section will test one of them. NVIDIA turned an architecture lead into a platform: a fabric (NVLink and NVSwitch), a CPU (Grace, then Vera), a rack (NVL72) and libraries, so that its advantage no longer rests on the die. AMD built the GPUs of two of the first three exascale machines (Frontier in 2022 and El Capitan in 2024; Intel's Aurora was the second, in 2024) and, from 2023, a credible AI line that reached rack scale in 2026. Intel shipped one exascale GPU, cancelled its two successors, and is restarting from inference parts. The sentences are deliberately unqualified; the qualifications are the next thirty slides.

\begin{terms}
\term{compute unit (CU)}{AMD's name for the unit of a GPU that holds four 16-lane SIMD units, matrix cores, a scalar unit, registers and a local data share; the counterpart of NVIDIA's SM and Intel's Xe-core.}
\term{streaming multiprocessor (SM)}{NVIDIA's compute unit: four partitions of 32 lanes with a tensor core each, 256 KB of registers and a shared L1 and shared memory.}
\term{Xe-core}{Intel's compute unit since Xe (2020): 8 (HPC) or 16 (HPG) vector engines, as many XMX matrix engines, and a shared L1 and local memory; the renamed subslice of the Gen families.}
\term{warp, wavefront, subgroup}{The group of threads scheduled together in lockstep: 32 on NVIDIA (warp), 64 on AMD (wavefront), 8 to 32 on Intel (subgroup, chosen by the compiler).}
\term{SXM}{NVIDIA's mezzanine module for data-center GPUs, with NVLink pins; four or eight SXM modules make an HGX baseboard.}
\term{HGX}{NVIDIA's baseboard of 4 or 8 SXM GPUs with NVSwitch chips, sold to server vendors; the unit of an H100 or B200 server.}
\term{OAM}{Open Accelerator Module, the Open Compute Project's mezzanine form factor (102 by 165 mm) used by AMD Instinct and Intel Max GPUs.}
\term{UBB}{Universal Baseboard, the OCP baseboard that carries eight OAM modules; one chassis design fits any OAM vendor.}
\term{scale-up fabric}{The interconnect over which GPUs in one system read and write each other's memory directly: NVLink, Infinity Fabric (xGMI), Xe Link, UALink. Distinct from the scale-out network (Ethernet, InfiniBand) between servers.}
\term{NVSwitch}{NVIDIA's switch chip for NVLink, which lets every GPU in a board or rack reach every other at full rate rather than over a fixed set of point-to-point links.}
\term{Infinity Fabric}{AMD's interconnect family; between GPUs it is also called xGMI, and it joins the eight MI300X GPUs of a UBB directly to one another.}
\term{Xe Link}{Intel's GPU-to-GPU link on Ponte Vecchio: 16 links per package, joining up to eight GPUs without a switch.}
\term{NVL72}{NVIDIA's rack of 72 GPUs and 36 CPUs on one NVLink fabric, shipped in GB200, GB300 and Vera Rubin versions.}
\term{Helios}{AMD's 72-GPU rack of MI455X GPUs, EPYC Venice CPUs and UALink over Ethernet, shipping from the third quarter of 2026.}
\term{CUDA, ROCm, oneAPI}{The three vendors' software platforms: NVIDIA's CUDA (2007), AMD's ROCm with the HIP language (2016), Intel's oneAPI with SYCL (2020).}
\term{CDNA}{AMD's data-center GPU architecture line (Compute DNA), split from the RDNA graphics line in 2020: CDNA 1 (MI100), 2 (MI250X), 3 (MI300), 4 (MI355X), 5 (MI455X).}
\term{UALink}{Ultra Accelerator Link, the open scale-up fabric standard formed in October 2024 by AMD, Intel, AWS, Google, Meta, Microsoft and others; specification 1.0 in April 2025.}
\end{terms}

\begin{intuition}
Six rows, three columns, eighteen cells: fill the table for any accelerator and its place in the market follows, because a product that wins four rows and loses the fabric and the software has, so far, lost. Keep the timeline beside it; most of what looks like a design decision was a date.
\end{intuition}
% =====================================================================
\setstage{tok}
\section{The NVIDIA line as platforms (slides 4--6)}

\slide{4}{The NVIDIA line as platforms: HGX A100 to Vera Rubin NVL72}{Six platforms in six years: the board became a rack, memory grew 7$\times$, the scale-up link 5$\times$, FP64 went up and down.}

\para{What the slide shows} A table with six columns, one per NVIDIA data-center platform from 2020 to 2026, fills the left two thirds of the slide: the HGX A100 and HGX H100 and H200 eight-GPU boards, then the GB200, GB300 and Vera Rubin NVL72 racks. Its fourteen rows give the year of first shipment, the GPU, HBM capacity and bandwidth per GPU, FP64 as vector / tensor rate, the dense FP16, FP8 and FP4 tensor peaks, the scale-up bandwidth per GPU, the number of GPUs in one scale-up domain, the CPU, the power per GPU and the power of the board or rack. A caption states the conventions: datasheet values, dense tensor peaks (NVIDIA's headlines are the 2:4-sparse figures, twice these), power for the DGX systems and the racks, Rubin from the specification page revised in September 2026, and ``n.p.'' for figures not published. The panel ``Reading the line'' draws four conclusions: the unit grew from a board to a rack, memory grew 7$\times$ and bandwidth 12$\times$, precision went down while peaks went up 110$\times$, and FP64 rose and fell; it ends by naming the comparisons the next sections make.

\para{Details} The previous lecture told this story die by die; this slide retells it as the six products a customer could buy, because that is what the AMD and Intel sections will be compared with. Reading down the columns, the six are as follows.

The HGX A100 (May 2020) is a baseboard of eight A100 SXM4 modules and six NVSwitch 2 chips, sold to server vendors and, as the DGX A100, by NVIDIA itself with two 64-core AMD EPYC 7742 host processors$^\dagger$. Each A100 has 40 GB of HBM2 in five stacks at 1{,}555 GB/s (the 80 GB version of November 2020 uses HBM2e at 2{,}039 GB/s), 12 NVLink 3 links at 50 GB/s for 600 GB/s per GPU, and a 400 W rating; the DGX A100 draws 6.5 kW at maximum$^\dagger$. The 8-GPU NVLink domain is the unit of scale-up: beyond it, servers talk over 200 Gb/s InfiniBand. Its peaks are the baseline for everything that follows: 9.7 TFLOP/s of FP64 on the vector lanes, 19.5 on the tensor cores, 312 TFLOP/s of dense FP16 or BF16.

The HGX H100 (shipping from the fourth quarter of 2022) keeps the shape and changes every number: eight H100 SXM5 modules, four NVSwitch 3 chips, 18 NVLink 4 links at 50 GB/s for 900 GB/s per GPU, 80 GB of HBM3 in five stacks at 3.35 TB/s, 700 W per GPU and 10.2 kW for the DGX H100 with two Intel Xeon host processors$^\dagger$. FP64 rose to 34 TFLOP/s on the vector lanes and 67 on the tensor cores; dense FP16 to 989 TFLOP/s; and FP8, new with Hopper's Transformer Engine, reached 1{,}979 TFLOP/s dense (3{,}958 with sparsity, the figure usually quoted)$^\dagger$. The HGX H200 (announced November 13, 2023; shipped in the second quarter of 2024) is the same GH100 die with the memory replaced: six stacks of HBM3e for 141 GB at 4.8 TB/s, the same peaks, links and power$^\dagger$. It exists because AMD's MI300X, launched a month after the H200 was announced, carried 192 GB (slide 5).

The GB200 NVL72 (first racks November 2024, volume in 2025) is where the unit changes. The product is a rack, not a board: 18 compute trays, each holding two GB200 superchips, where a superchip is one Grace CPU (72 Arm Neoverse V2 cores, up to 480 GB of LPDDR5X at 512 GB/s) joined by the 900 GB/s NVLink-C2C to two Blackwell GPUs; and 9 switch trays, each with two NVLink Switch chips of 72 ports and 7.2 TB/s. The 18 switch chips give every one of the 72 GPUs 18 NVLink 5 links at 100 GB/s, 1.8 TB/s per GPU and 130 TB/s for the rack, and the rack is one scale-up domain$^\dagger$. Per GPU: 186 GB of HBM3e at 8 TB/s (NVIDIA states 372 GB and 16 TB/s per superchip), FP64 40 TFLOP/s (2{,}880 per rack, with no separate tensor-core rate), dense FP16 2.5 PFLOP/s, dense FP8 5 PFLOP/s and dense FP4 10 PFLOP/s (NVIDIA's rack figures are 1{,}440 PFLOP/s of FP4 and 720 of FP8 with sparsity), 1{,}200 W; the rack draws about 120 kW, a figure from Dell and the press rather than NVIDIA's pages$^\dagger$. The x86 host is gone: the CPU in the table is NVIDIA's own.

The GB300 NVL72 (first racks at CoreWeave, July 3, 2025) is more than the H200 was to the H100: the H200 was a memory refresh of an unchanged die, while Blackwell Ultra rebalances the same silicon, as the FP4 and FP64 rows record. Blackwell Ultra (B300) carries eight 12-high stacks of HBM3e for 288 GB at the same 8 TB/s, raises the dense FP4 rate to 15 PFLOP/s (the sparse figure stays at 20; the increase comes from a redesigned tensor-core datapath, not from more SMs), doubles the special-function rate for attention softmax, moves to PCIe 6, and is rated at 1{,}400 W$^\dagger$. In exchange it removes nearly all of the FP64 datapath: 1.4 TFLOP/s per GPU, 100 TFLOP/s per rack of 72, against 2{,}880 on the GB200. The rack power is about 130 kW by press estimates (NVIDIA has published none).

The Vera Rubin NVL72 (first revenue racks August 2026) is the first new silicon in two years and the first whose specification NVIDIA revised after announcement. Rubin is two 3 nm dies with 336 billion transistors; Vera is an 88-core Arm CPU of NVIDIA's own Olympus design; the rack has 72 Rubin and 36 Vera, NVLink 6 at 3.0 TB/s per GPU and 216 TB/s per rack, 288 GB of HBM4 per GPU at 19.2 TB/s, FP64 33 TFLOP/s, dense FP16 4 PFLOP/s, FP8 17.5 PFLOP/s and dense NVFP4 35 PFLOP/s (50 PFLOP/s for inference, the sparse figure)$^\dagger$. The HBM bandwidth was 22 TB/s and the NVLink 3.6 TB/s and 260 TB/s per rack when announced at CES in January 2026; the product page was edited in early September 2026 and NVIDIA said on October 29 that the earlier numbers were ``preliminary architectural targets'' and the new ones ``what the system sustains at qualified, volume production''$^\dagger$. The GPU's power is not published; supply-chain reports put it between 1.8 and 2.3 kW$^\dagger$.

The panel's arithmetic is in the table. Memory per GPU went from 40 GB to 288 GB, $7.2\times$; bandwidth from 1.555 to 19.2 TB/s, $12.3\times$; the scale-up link from 600 to 3{,}000 GB/s, $5\times$; the number of GPUs sharing it from 8 to 72, $9\times$; the fastest dense tensor peak from 312 TFLOP/s (FP16) to 35 PFLOP/s (FP4), $112\times$, of which $4\times$ is the narrowing of the operand from 16 bits to 4 and the rest is silicon (derived). FP64 went 9.7, 34, 34, 40, 1.4, 33 on the vector lanes and 19.5, 67, 67, 40, 1.4, 33 on the tensor cores: up by $3.4\times$ from A100 to H100, flat to the GB200, cut by a factor of nearly 30 on the GB300, and back to half the H100's tensor-core rate on Rubin, as a deliberate decision about who the product is for (slide 12). The power per GPU went 400, 700, 700, 1{,}200, 1{,}400 W and beyond, and the unit of power from a 6.5 kW box to a rack of about 120 kW, both of which the rack comparison on slide 22 takes up.

The last point of the panel sets up the sections that follow. AMD matched this line product for product from 2020: the MI100 against the A100, the MI250X against the A100 80 GB, the MI300X against the H100 and H200, the MI325X against the H200, the MI355X against the GB200 and GB300, the MI455X against Rubin (slide 7). Intel entered once, with Ponte Vecchio against the H100 in 2023 (slide 17), and its announced return is an inference card (slide 20).

\begin{terms}
\term{HGX A100 / HGX H100}{NVIDIA's eight-GPU baseboards of 2020 and 2022, with NVSwitch chips joining the GPUs; sold to server vendors and as the DGX A100 and DGX H100.}
\term{DGX}{NVIDIA's own server built around an HGX board: DGX A100 (6.5 kW), DGX H100 (10.2 kW).}
\term{H200}{The H100's GH100 die with six stacks of HBM3e (141 GB at 4.8 TB/s) in place of five of HBM3; announced November 2023, shipped 2024.}
\term{GB200 superchip}{One Grace CPU and two Blackwell GPUs on one board, joined by NVLink-C2C at 900 GB/s; 36 of them make an NVL72.}
\term{Grace}{NVIDIA's first data-center CPU (2023): 72 Arm Neoverse V2 cores with up to 480 GB of LPDDR5X at 512 GB/s.}
\term{Vera}{NVIDIA's second CPU (2026): 88 Olympus cores of NVIDIA's own design, 176 threads, 1.8 TB/s of NVLink-C2C.}
\term{NVLink-C2C}{The chip-to-chip version of NVLink that joins an NVIDIA CPU to its GPUs with cache coherence; 900 GB/s on Grace, 1.8 TB/s on Vera.}
\term{Blackwell Ultra (B300)}{The 2025 refresh of Blackwell: 288 GB of HBM3e, 15 PFLOP/s of dense FP4, FP64 cut to 1.4 TFLOP/s, 1,400 W.}
\term{Rubin}{NVIDIA's 2026 GPU: two 3 nm dies, 336 billion transistors, 288 GB of HBM4 at 19.2 TB/s, NVLink 6 at 3 TB/s, 35 PFLOP/s of dense NVFP4.}
\term{specification revision}{A change to a vendor's published figures between announcement and shipment; Rubin's HBM and NVLink bandwidths were lowered in September 2026, and Blackwell's memory and FP64 figures had been in 2024.}
\term{n.p.}{Not published: the vendor has stated no figure as of November 2026. Used rather than an estimate throughout.}
\term{LPDDR5X}{Low-power DDR5 memory, used in large capacity on NVIDIA's Grace and Vera CPUs (and on Intel's GPU D-series) in place of DDR5 or HBM.}
\end{terms}

\begin{intuition}
Read the table left to right as one object changing shape: a board with eight GPUs and someone else's CPU became a rack with seventy-two GPUs and NVIDIA's own, and every number in between grew with it except the one NVIDIA decided its customers no longer needed. The AMD and Intel tables that follow are built to the same rows so that the shapes can be laid on top of each other.
\end{intuition}

% ---------------------------------------------------------------------
\slide{5}{Memory per GPU: the contest AMD kept winning on paper}{HBM per GPU, 2020--2026: AMD led from the MI250X on and tied in 2025; Ponte Vecchio matched it in 2023.}

\para{What the slide shows} Two small charts share the left half. The first plots HBM capacity per GPU in gigabytes, on a logarithmic axis from 32 to 512 GB, against year of shipment: NVIDIA's gray line steps from 40 GB (A100, 2020) to 80 (A100 80 GB and H100), 141 (H200, 2024), 186 (GB200), 288 (GB300, 2025) and 288 again (Rubin, 2026); AMD's green line runs from 32 GB (MI100, 2020) to 128 (MI250X, 2021), 192 (MI300X, 2023), 256 (MI325X, 2024), 288 (MI355X, 2025) and 432 (MI455X, 2026); two orange triangles mark Intel's Ponte Vecchio (128 GB, 2023) and Gaudi 3 (128 GB, 2024). The second chart plots bandwidth per GPU in TB/s on a logarithmic axis from 1 to 32: NVIDIA 1.6, 2.0, 3.35, 4.8, 8, 8, 19.2; AMD 1.23, 3.2, 5.3, 6.0, 8.0, 23.3; Intel 3.28 and 3.67. The panel on the right walks through the generations, 2020--22 to 2026 (the MI100's 32 GB trailing the A100, then the MI250X taking the lead, Ponte Vecchio matching the MI250X in 2023, a tie at 288 GB in 2025, the MI455X ahead in 2026), and ends with why the lead did not decide the market.

\para{Details} High-bandwidth memory is the one component the three vendors buy from the same three suppliers (SK hynix, Samsung and Micron), and the capacity and bandwidth of a GPU are set by how many stacks the package has room for, how many DRAM dies each stack carries, and which generation of the HBM interface the GPU's memory controllers speak. The arithmetic is worth doing once. A stack's capacity is its die count times the die capacity (8 GB dies through HBM3, then 16 and 24 Gb dies allowing 24 and 36 GB in 8- and 12-high stacks); its bandwidth is its 1{,}024-bit interface (2{,}048 bits on HBM4) times the per-pin data rate. So the A100's 40 GB is five 8 GB HBM2 stacks at 1{,}555 GB/s, 311 GB/s per stack; the H100's 80 GB is five 16 GB HBM3 stacks at 3.35 TB/s, 670 GB/s per stack; the H200's 141 GB is six 24 GB HBM3e stacks (144 GB nominal) at 4.8 TB/s, 800 GB/s per stack; Blackwell's 186 or 192 GB is eight 24 GB HBM3e stacks at 8 TB/s, 1 TB/s per stack; the B300's and the MI355X's 288 GB is eight 12-high 36 GB stacks at the same 8 TB/s; Rubin's 288 GB is eight 36 GB HBM4 stacks at 19.2 TB/s, 2.4 TB/s per stack, the doubled interface width of HBM4 showing; the MI455X's 432 GB is twelve 36 GB HBM4 stacks at 23.3 TB/s, 1.94 TB/s per stack (derived throughout)$^\dagger$. Two numbers set everything: how many stacks fit around the dies, and which HBM generation is in production.

The generations in the panel can now be read as packaging decisions. In 2020 the MI100's single die carried four stacks of HBM2 for 32 GB, less than the A100's 40 or 80, so AMD's lead begins only with the MI250X. In 2021--22 the MI250X carried 128 GB against the A100's 80 because it had two compute dies, each with four stacks of 16 GB HBM2e, eight stacks in all at 3.2 TB/s, where the A100's single die had room for six stack sites and used five$^\dagger$. In 2023--24 the MI300X's 192 GB at 5.3 TB/s came from eight 24 GB HBM3 stacks around its four I/O dies, against the H100's five stacks and 80 GB; NVIDIA's answer was the H200 with six stacks of HBM3e (141 GB, 4.8 TB/s), and AMD's answer to that was the MI325X (October 2024) with eight 32 GB HBM3e stacks for 256 GB at 6.0 TB/s on the same silicon$^\dagger$. Ponte Vecchio's 128 GB in eight stacks of HBM2e at 3.28 TB/s matched the MI250X's 128 GB and exceeded the H100's 80 from its January 2023 launch until the MI300X's December one, a parity the Intel part rarely gets credit for; Gaudi 3 (2024) carried the same 128 GB of HBM2e at 3.67 TB/s$^\dagger$. In 2025 the two lines met: the B300 and the MI355X both carry eight 12-high HBM3e stacks for 288 GB at 8 TB/s. In 2026 they diverged again on stack count: Rubin has eight HBM4 stacks (288 GB, 19.2 TB/s after the September revision), the MI455X twelve (432 GB, 23.3 TB/s; AMD had announced 19.6 TB/s in 2025 and raised it at the July 2026 launch)$^\dagger$. AMD's twelve stacks are possible because its package puts the memory controllers on separate I/O dies at the edge; Rubin's two large compute dies leave room for eight.

Why does memory per GPU matter enough to fight over? For training, capacity sets how many GPUs a model must be split across and bandwidth sets how fast each GPU can stream its share of the weights, activations and optimizer state; both are large but neither is usually the bound, because the matrix multiplies are. For inference, memory is the bound. A 70-billion-parameter model in 16-bit weights occupies 140 GB; it fits on one MI300X and needs two H100s. During decode, each generated token requires reading every weight once, so the tokens per second a single user sees is approximately the bandwidth divided by the bytes of weights, 5.3 TB/s over 140 GB gives about 38 tokens per second on an MI300X and 3.35 over 140 about 24 on one H100 before any batching (derived); and the key-value cache of every concurrent request has to live in the same memory. This is why the first MLPerf Inference results AMD submitted (slide 30) were on Llama 2 70B, and why NVIDIA shipped one memory refresh of an unchanged die (H200) and one memory-led rebalancing of the same silicon (B300).

The panel's last paragraph is the lecture's recurring argument. The per-GPU memory lead did not decide the market for three reasons. First, the unit of memory a customer sees is the scale-up domain, not the GPU: a GB200 NVL72 rack has 13.4 TB of HBM in one NVLink domain, against 8 times 192 GB, 1.5 TB, on an MI300X UBB, so a model that needs 2 TB of memory is served from one NVLink rack and split across two Infinity Fabric boards with a network between them (slide 21). Second, the achieved bandwidth depends on the kernels: the 2024 measurements on slide 26 found the MI300X's GEMM throughput at 47\% of its peak against the H100's 73\%, which erases a memory-bandwidth lead of 1.6$\times$. Third, HBM is the shared bottleneck of the whole industry: all three vendors buy from the same three makers, whose 12-high HBM3e and HBM4 output was sold out through 2026, so a design with twelve stacks per GPU ships as many GPUs as there are stacks for$^\dagger$.

\begin{formulas}
Per-GPU HBM capacity and bandwidth from the stack parameters, with $n$ stacks of $h$ dies of capacity $c$ each, interface width $w$ bits and per-pin data rate $r$:
\[
\text{capacity} = n\,h\,c,\qquad \text{bandwidth} = n\,\frac{w\,r}{8}.
\]
H100: $5\times 8\times 2\,\text{GB}=80$ GB; $5\times 1{,}024\times 5.2\,\text{Gb/s}/8 = 3.3$ TB/s. Rubin: $8\times 12\times 3\,\text{GB}=288$ GB; $8\times 2{,}048\times 9.4\,\text{Gb/s}/8\approx 19.2$ TB/s (derived; the per-pin rate inferred from the stated total). Decode tokens per second for one user $\approx$ bandwidth / weight bytes: $5.3\,\text{TB/s}/140\,\text{GB}\approx 38$ for a 70B model in FP16 on an MI300X.
\end{formulas}

\begin{terms}
\term{HBM2, HBM2e, HBM3, HBM3e, HBM4}{Successive generations of stacked DRAM: 8 GB per stack at about 0.3 TB/s (HBM2), 16 GB at 0.4--0.5 (HBM2e), 16--24 GB at 0.67 (HBM3), 24--36 GB at 0.8--1.0 (HBM3e), 36 GB at 2--2.4 TB/s on a doubled 2,048-bit interface (HBM4).}
\term{12-high stack}{An HBM stack of twelve DRAM dies rather than eight; 36 GB with 24 Gb dies. The B300, MI355X, Rubin and MI455X use them.}
\term{stack site}{A position on the package substrate for one HBM stack; the number of sites around the compute dies bounds capacity (5 used of 6 on A100 and H100, 8 on Blackwell and Rubin, 12 on the MI455X).}
\term{memory refresh}{A product that keeps the compute die and changes only the memory: H200 (2024), MI325X (2024); the B300 (2025) added 288 GB but also rebalanced the die (FP4 up, FP64 cut).}
\term{decode}{The phase of language-model inference that generates one token at a time; each token reads all the weights once, so its speed is bounded by memory bandwidth.}
\term{KV cache}{The key and value tensors of every previous token in a request, kept in GPU memory during decode; it grows with context length and concurrent requests and competes with the weights for capacity.}
\term{MI325X}{AMD's October 2024 memory refresh of the MI300X: 256 GB of HBM3e at 6.0 TB/s, 1,000 W, the same compute dies.}
\term{HBM supply}{The output of the three HBM makers, SK hynix, Samsung and Micron; the binding constraint on how many GPUs any vendor could ship in 2024--26.}
\end{terms}

\begin{intuition}
Memory is the contest decided by the package: whoever leaves room for more stacks wins it, and AMD's I/O-die design left room. It is also the contest that mattered least on its own, because customers buy memory by the rack and bandwidth by the kernel, not by the GPU.
\end{intuition}

% ---------------------------------------------------------------------
\slide{6}{Compute peaks per GPU: three vendors, FP64 to FP4}{Dense vendor peaks; AMD doubled FP64 on matrix cores until CDNA 4; Intel's PVC ran FP64 at full rate.}

\para{What the slide shows} A table of twelve GPUs in three blocks: NVIDIA (A100, H100 SXM, GB200, GB300, Rubin), AMD (MI100, MI250X, MI300X, MI355X, MI455X) and Intel (Ponte Vecchio Max 1550, Gaudi 3). For each it gives the year, the process node, FP64 as vector / matrix rate, the dense FP16, FP8 and FP4 matrix peaks in TFLOP/s, and the power in watts. A caption lists the exceptions: the MI100's 185 is FP16 (its BF16 rate is 92); the MI455X's figures are MXFP8 and MXFP4; Ponte Vecchio's compute tiles are TSMC N5, its base tiles Intel 7 and its link tiles TSMC N7; its FP16 is the vendor figure, consistent with its confirmed 419 TFLOP/s of TF32; Gaudi 3 states one figure for BF16 and FP8; Rubin is from the September 2026 specification page. The panel ``How to read it'' makes four points: low precision converged; FP64 diverged; Intel bet on full-rate FP64; and a peak is not what is achieved.

\para{Details} Every number in the table is the product of four factors: the number of compute units, the operations each unit performs per clock at that precision, the clock, and two for the multiply-add. Writing the products out shows where each vendor's peaks come from and why they differ. The clocks and per-unit rates are from the architecture whitepapers; the peaks are the vendors' stated figures, which the products reproduce to within rounding$^\dagger$.

NVIDIA's A100 has 108 SMs at 1.41 GHz. Each SM does 32 FP64 fused multiply-adds per clock on its vector lanes, so $108\times32\times2\times1.41=9.7$ TFLOP/s, and twice that, 19.5, on the tensor cores' double-precision matrix instruction; each SM does 1{,}024 FP16 multiply-adds per clock on its tensor cores, $108\times1{,}024\times2\times1.41=312$ TFLOP/s. The H100 SXM has 132 SMs with 64 FP64 lanes each at up to 1.98 GHz: $132\times64\times2\times1.98=33.4$, the datasheet's 34; the tensor cores give 67 for FP64, 989 for dense FP16 (at a slightly lower clock, as the previous lecture discussed) and 1{,}979 for FP8, the first 8-bit format in NVIDIA hardware. Blackwell's GB200 states 40 TFLOP/s of FP64 with no faster tensor-core path, 2.5 PFLOP/s of dense FP16, 5 of FP8 and 10 of FP4; the GB300 keeps the three matrix rates except FP4, which rises to 15 PFLOP/s, and cuts FP64 to 1.4; Rubin gives 33, 4{,}000, 17{,}500 and 35{,}000$^\dagger$. The FP16 rate per SM per clock doubled from Ampere to Hopper and again to Blackwell; the FP8 and FP4 rates double again with each halving of the operand.

AMD's CDNA parts follow the same construction with the compute unit (CU) in place of the SM, and the table shows the one place AMD chose differently. The MI100 (120 CUs, 1.5 GHz) does 64 FP64 FLOP per CU per clock, half its FP32 rate: $120\times64\times1.5=11.5$ TFLOP/s; its matrix cores do 1{,}024 FP16 FLOP per CU per clock, $120\times1{,}024\times1.5=185$ TFLOP/s, but only 512 for BF16, hence the caption. CDNA 2 raised FP64 to the full FP32 rate, 128 FLOP per CU per clock on the vector units and 256 on the matrix cores: the MI250X's 220 CUs at 1.7 GHz give $220\times128\times1.7=47.9$ TFLOP/s of vector FP64 and 95.7 of matrix FP64, and 383 of FP16. CDNA 3 kept both ratios and added FP8: the MI300X's 304 CUs at 2.1 GHz give $304\times128\times2.1=81.7$ TFLOP/s vector and $304\times256\times2.1=163.4$ matrix FP64, $304\times2{,}048\times2.1=1{,}307$ FP16 and $304\times4{,}096\times2.1=2{,}615$ FP8 (derived; the whitepaper's figures). CDNA 4 then changed the balance: the MI355X's 256 CUs at 2.4 GHz double the FP16 rate per CU to 4{,}096 FLOP per clock ($256\times4{,}096\times2.4=2{,}516$, stated 2.5 PFLOP/s), add FP6 and FP4 at 16{,}384 FLOP per CU per clock (10 PFLOP/s dense), and halve the matrix FP64 rate per CU from 256 to 128 FLOP per clock, so that matrix and vector FP64 are both $256\times128\times2.4=78.6$ TFLOP/s$^\dagger$. CDNA 5's MI455X states 5.03 PFLOP/s of FP16, 20.13 of MXFP8 and 40.26 of MXFP4 (and 315 TFLOP/s of FP32) for its 256 work-group processors, and no FP64 figure at all; AMD has said the FP64 part of the generation is the MI430X with 288 TFLOP/s, for 2027$^\dagger$.

Intel's Ponte Vecchio is the third column of the panel's argument. Its 128 Xe-cores hold 8 vector engines each of 16 FP32 lanes, 16{,}384 lanes, and each 512-bit engine performs 16 FP64 multiply-adds per clock, the same count as its FP32 ones (Intel's per-core figure of 256 FP32 and 256 FP64 operations per clock), so FP64 and FP32 share one peak: $16{,}384\times2\times1.6\,\text{GHz}=52.4$ TFLOP/s. Its XMX matrix engines, 1{,}024 of them, do 128 TF32 multiply-adds per clock each, $1{,}024\times128\times2\times1.6=419$ TFLOP/s, and double that for FP16 and again for INT8 (839 TFLOP/s and 1{,}678 TOPS, Intel's figures, which no independent source has confirmed but which follow from the confirmed TF32 rate); there is no FP8 and no FP4$^\dagger$. In January 2023 that made it the fastest FP64 GPU shipping (52.4 against the H100's 34 vector / 67 tensor and the MI250X's 47.9 / 95.7) and a generation behind in matrix work (839 against 989 at FP16, nothing against 1{,}979 at 8 bits). Gaudi 3, Intel's other accelerator, states 1{,}835 TFLOP/s for both BF16 and FP8 from its eight matrix engines, close to the H100's dense FP8, and has no FP64 at all; slide 16 covers it.

The panel's first conclusion follows from the table: on the low-precision rates the two main vendors converged. The MI355X and GB200 both state 10 PFLOP/s of dense FP4, 5 of FP8 and 2.5 of FP16; the MI455X's 40.26 and Rubin's 35 are within 15\%. The reason is in the process column: from 2023 both bought the same TSMC nodes (N5 and 4N, then N3P and 3 nm, then N2 against a 3 nm pair), and a matrix unit's throughput per watt and per square millimeter at a given node and precision is set by physics more than by design. The second conclusion is the FP64 divergence, the subject of slide 12: AMD held matrix FP64 at twice the vector rate from 2021 to 2024 and reached 163 TFLOP/s, then halved the matrix rate on CDNA 4 for the AI market and moved full FP64 to a separate HPC part; NVIDIA peaked at 67 on the H100, dropped the tensor-core advantage on Blackwell, cut FP64 almost entirely on the B300, and restored 33 on Rubin with emulation on the integer tensor cores offered for the rest. The third is Intel's bet: full-rate FP64 at the cost of a matrix engine a generation behind, which won Aurora and nothing else. The fourth, which slide 26 quantifies, is that the table's numbers are ceilings: in December 2024 independent measurements put large matrix multiplies on the MI300X at 47\% of its stated peak and on the H100 at 73\%, so the right comparison is measured throughput at the same date, and a 25\% peak lead can be a 20\% deficit in practice$^\dagger$.

\begin{formulas}
Peak at precision $p$:
\[
P_p = N_{\text{units}}\times \mathrm{ops}_p/\text{clk}\times f .
\]
MI300X: $304\times 256\times 2.1\,\text{GHz} = 163.4$ TFLOP/s matrix FP64; $304\times 4{,}096\times 2.1 = 2{,}615$ TFLOP/s FP8. MI355X: $256\times 128\times 2.4 = 78.6$ (FP64, vector and matrix); $256\times 16{,}384\times 2.4 = 10{,}066$ TFLOP/s FP4. Ponte Vecchio: $16{,}384\times 2\times 1.6 = 52.4$ TFLOP/s FP64 $=$ FP32; $1{,}024\times 128\times 2\times 1.6 = 419$ TFLOP/s TF32. Halving the operand width doubles $\mathrm{ops}_p/\text{clk}$ on every vendor's matrix unit from FP16 down to FP4.
\end{formulas}

\begin{terms}
\term{vector FP64 and matrix FP64}{Double-precision throughput on the ordinary lanes and on the matrix units; AMD's matrix cores and NVIDIA's Ampere and Hopper tensor cores doubled the vector rate, Blackwell's and Rubin's do not, and Intel's XMX has no FP64 mode.}
\term{matrix core}{AMD's name for the matrix multiply-add unit in a CDNA compute unit, four per CU; the counterpart of a tensor core.}
\term{XMX}{Xe Matrix Extensions: Intel's systolic matrix engine, one per vector engine in an Xe-core; 4,096-bit (128 TF32 multiply-adds per clock) on Xe-HPC.}
\term{full-rate FP64}{A design in which a lane performs FP64 at the same FLOP rate as FP32: CDNA 2 and 3 (vector), Ponte Vecchio. NVIDIA's compute dies run FP64 at half the FP32 rate.}
\term{FLOP per clock per unit}{The number of floating-point operations one compute unit completes per cycle at a given precision; the architectural constant that, times unit count and clock, gives the peak.}
\term{TF32}{A 19-bit tensor format (8-bit exponent, 10-bit mantissa) used for FP32 inputs on matrix units; Ponte Vecchio's confirmed 419 TFLOP/s and the H100's 494 are TF32 rates.}
\term{work-group processor (WGP)}{The compute unit of AMD's CDNA 5 (MI455X), a pair of CUs sharing front-end and local memory; AMD counts 256 of them and publishes no lane count.}
\term{MI430X}{AMD's HPC variant of the CDNA 5 generation, with 288 TFLOP/s of FP64 and 432 GB of HBM4, due in the first half of 2027 for Discovery, Alice Recoque and LUMI-AI.}
\term{Gaudi 3}{Intel's 2024 Habana-architecture accelerator: 1,835 TFLOP/s of BF16 or FP8 from eight matrix engines, 128 GB of HBM2e, Ethernet on the die, no FP64.}
\term{achieved fraction}{Measured throughput divided by the stated peak; 47\% (MI300X) and 73\% (H100) on large BF16 matrix multiplies in December 2024.}
\end{terms}

\begin{intuition}
Every peak is a count of units times a rate per clock times a clock: write the product and the vendors' choices appear as the middle factor. AMD spent it on FP64 until 2025, Intel spent it on FP64 once, and NVIDIA spent it on the narrowest format it could sell. By 2026 the low-precision products agree to within rounding, which is why the next twenty slides are about everything except the die.
\end{intuition}

% =====================================================================
\setstage{emb}
\section{AMD Instinct and CDNA (slides 7--12)}

\slide{7}{AMD Instinct and CDNA: the compute unit, five generations}{A GCN-derived CU with 64-wide wavefronts (32 on CDNA 5) and matrix cores, from one die to a dozen chiplets.}

\para{What the slide shows} The upper left is a block diagram of one CDNA 3 compute unit (CU), drawn from AMD's 2023 whitepaper$^\dagger$: a scheduler and scalar unit across the top left, the 512 KB vector register file across the top right, four SIMD units of 16 lanes each in the middle (each drawn as two rows of eight lane boxes), and along the bottom the four matrix cores (FP64 down to FP8 and INT8) beside the 64 KB local data share (LDS; 160 KB in CDNA 4) and the 32 KB L1 cache. The caption under it states the execution rule that distinguishes AMD's design from NVIDIA's: a wavefront of 64 work-items issues on a 16-lane SIMD over four cycles, so the four SIMDs give 64 FP32 lanes and 128 FLOP per clock (256 in CDNA 3, which doubled the FP32 rate per CU), with the matrix cores running at 2--16$\times$ that; the comparison is to the NVIDIA SM's four 32-lane partitions. Below the diagram a table runs the five CDNA generations across the page (MI100, MI250X, MI300X, MI355X, MI455X) with rows for shipping date, compute dies and process, base dies, transistors, compute units, HBM capacity and bandwidth, and power. The right panel, ``The lineage, in one paragraph each,'' summarizes what each generation changed and points to the slide that treats it.

\para{Details} AMD's compute GPUs descend from Graphics Core Next (GCN), the architecture AMD introduced in 2012 and used in every Radeon until 2019. In GCN a compute unit holds four SIMD units of 16 lanes each; a \emph{wavefront} is a group of 64 work-items (threads), and a wavefront's instruction is issued to one 16-lane SIMD and executes over four consecutive cycles, a quarter of the wavefront per cycle. Each SIMD therefore owns a set of wavefronts (up to ten in GCN and the MI100, eight from CDNA 2) and round-robins among them, hiding memory latency the same way the NVIDIA warp scheduler does, but with half the issue rate per lane: an NVIDIA SM partition issues one 32-thread warp instruction per cycle onto 32 lanes, a CDNA SIMD issues one 64-work-item instruction every four cycles onto 16 lanes. The arithmetic per CU per clock comes out the same, 64 FP32 FMAs or 128 FLOP, because four SIMDs of 16 lanes make 64 lanes; what differs is granularity. A branch that diverges within a wavefront idles lanes for all 64 work-items, and a kernel that cannot fill 64 work-items per wavefront wastes more of the machine than one that cannot fill 32.

In November 2020 AMD split the line: RDNA (Radeon DNA) for graphics, which moved to 32-wide wavefronts and a dual-CU ``work-group processor,'' and CDNA (Compute DNA) for data-center compute, which kept the GCN compute unit, the 64-wide wavefront and the four SIMD16 units, dropped the graphics fixed-function hardware (rasterizers, render back-ends, display) and added one \emph{matrix core} per CU$^\dagger$. The first CDNA part, the Instinct MI100, is a single 7 nm die of 120 CUs in eight shader engines of 15, 7{,}680 lanes, 32 GB of HBM2 in four stacks at 1.23 TB/s, three Infinity Fabric links of 92 GB/s for GPU-to-GPU traffic (276 GB/s aggregate), a PCIe 4.0 card at 300 W$^\dagger$. Its peaks at 1.5 GHz were 11.5 TFLOP/s FP64 and 23.1 TFLOP/s FP32 on the vector units, 46.1 TFLOP/s FP32 on the matrix cores and 184.6 TFLOP/s FP16 (92.3 BF16) on the matrix cores$^\dagger$. The arithmetic is $120 \times 1.502\,\text{GHz} \times 64 = 11.5$ TFLOP/s for FP64 (64 FLOP per CU per clock, half the FP32 vector rate) and $120 \times 1.502 \times 1{,}024 = 184.6$ for FP16 (derived); the per-CU-per-clock figures in this paragraph and the next come from AMD's ROCm microarchitecture pages$^\dagger$.

CDNA 2 (MI250X, November 2021, slide 8) made two changes to the CU and one to the package. The FP64 vector rate was raised to the FP32 rate (128 FLOP per CU per clock), so double precision ran at the single-precision rate, which no GPU had done before; FP32 gained a packed mode in which two FP32 values ride in one 64-bit register and execute together, doubling the FP32 rate to 256 FLOP per CU per clock when the compiler can pair operations; and the matrix cores grew to four per CU with an FP64 matrix mode at 256 FLOP per CU per clock. The package carried two graphics compute dies (GCDs) on TSMC N6, each a complete GPU with its own four HBM2e stacks and L2, joined by in-package Infinity Fabric; 220 of 224 CUs were enabled. CDNA 3 (MI300X and MI300A, December 2023, slide 9) kept the CU's structure but doubled the FP16 and added FP8 matrix rates (2{,}048 and 4{,}096 FLOP per CU per clock), doubled the vector FP32 rate to 256 FLOP per CU per clock (so that FP32 vector equals FP64 matrix on the MI300X: 163.4 TFLOP/s each), enlarged the L1 to 32 KB and the vector register file to 512 KB per CU, and moved to the 3D-stacked package of eight accelerator complex dies (XCDs, TSMC N5, 38 of 40 CUs each) on four I/O dies (IODs, N6) with a 256 MB Infinity Cache. CDNA 4 (MI355X, 2025, slide 10) moved the XCDs to N3P, enabled 32 of 36 CUs per XCD (256 in all), grew the LDS from 64 to 160 KB per CU, added the block-scaled MXFP6 and MXFP4 matrix formats, doubled FP16 and FP8 rates per CU again (4{,}096 and 8{,}192 FLOP per CU per clock), and halved the FP64 matrix rate per CU per clock from 256 to 128, equal to the vector rate. CDNA 5 (MI455X, 2026, slide 11) moved to TSMC N2 compute dies, replaced the CU with a work-group processor (WGP) as its unit of count, 256 of them, with 320 KB of LDS per WGP and, by the launch reporting, a 32-wide wavefront in place of GCN's 64$^\dagger$, dropped the Infinity Cache for a 192 MB shared L2 on two cache dies, went to HBM4, and replaced Infinity Fabric between GPUs with UALink carried over Ethernet.

The table on the slide gives the die census generation by generation$^\dagger$. The transistor counts rise from 25.6 billion (MI100; a third-party database figure, not in AMD's whitepaper$^\dagger$) through 58.2 billion for the two MI250X dies (a launch-slide figure$^\dagger$) to 153 billion for the twelve dies of the MI300X, 185 billion for the ten of the MI355X and 320 billion for the twelve of the MI455X; the HBM goes 32, 128, 192, 288, 432 GB at 1.2, 3.2, 5.3, 8.0 and 23.3 TB/s; the power goes 300 W (PCIe card), 560 W, 750 W, 1{,}400 W and, for the MI455X, not published. The MI455X's compute-die row reads ``8 XCD (N2)'' and its base-die row ``2 IOD, 2 FCD'': two I/O dies and two cache dies, all on N3P. Read against the NVIDIA table on slide 4 the pattern is that AMD moved to multiple compute dies three generations before NVIDIA's two-die Blackwell, and that AMD's HBM capacity led at every generation (slide 5).

Two cautions on reading the CU against the SM. First, the counts are not commensurable: a CDNA 3 CU has 64 FP32 lanes and four matrix cores, a Hopper SM has 128 FP32 lanes and four tensor cores, so ``304 CUs against 132 SMs'' overstates the MI300X's advantage by a factor of two at the lane level (19{,}456 lanes against 16{,}896, a ratio of 1.15, is the fair comparison; derived). Second, the wavefront width matters to programmers porting code: a CUDA kernel written for warps of 32 that uses warp-level primitives (shuffles, ballots) must be re-examined for wavefronts of 64, and HIP exposes the wavefront size as a compile-time constant for this reason (slide 24).

\begin{terms}
\term{CDNA}{Compute DNA, AMD's data-center GPU architecture line since 2020: CDNA 1 (MI100), CDNA 2 (MI200), CDNA 3 (MI300), CDNA 4 (MI350) and CDNA 5 (MI400). Graphics parts use the separate RDNA line.}
\term{GCN}{Graphics Core Next, AMD's GPU architecture of 2012--2019, from which CDNA inherits the compute unit, the 64-wide wavefront and the four-cycle SIMD issue.}
\term{compute unit (CU)}{AMD's counterpart of the NVIDIA SM: four 16-lane SIMD units, a scalar unit, a 512 KB vector register file (CDNA 3), four matrix cores, the local data share and an L1 cache.}
\term{wavefront}{AMD's group of 64 work-items that execute one instruction together, over four cycles on a 16-lane SIMD; the counterpart of NVIDIA's 32-thread warp.}
\term{work-item}{OpenCL's and HIP's word for one thread of a kernel; 64 of them make a wavefront on CDNA.}
\term{SIMD16}{A 16-lane single-instruction multiple-data unit; four per CDNA compute unit.}
\term{matrix core}{AMD's matrix-multiply unit, the counterpart of NVIDIA's tensor core; one per CU on CDNA 1, four per CU from CDNA 2, with FP64, FP32, FP16, BF16, FP8 and INT8 modes (FP6 and FP4 from CDNA 4).}
\term{local data share (LDS)}{AMD's software-managed on-chip memory per CU, the counterpart of CUDA shared memory: 64 KB per CU through CDNA 3, 160 KB on CDNA 4, 320 KB per WGP on CDNA 5.}
\term{XCD}{Accelerator complex die: a CDNA 3 or 4 compute die of 40 (CDNA 3) or 36 (CDNA 4) CUs, eight of which sit on the I/O dies of an MI300X or MI355X.}
\term{IOD}{I/O die: the base die of the MI300 and MI350 packages carrying the Infinity Cache, the HBM controllers and the Infinity Fabric; four per MI300X on N6, two per MI355X.}
\term{GCD}{Graphics compute die: one of the two complete GPU dies in an MI250X module; software sees each as a separate device.}
\term{work-group processor (WGP)}{The dual-CU unit that RDNA uses and that CDNA 5 adopts as its unit of count, with 32-wide waves; the MI455X has 256 of them.}
\term{Infinity Fabric}{AMD's coherent interconnect, used on-package between dies, between GPUs on a board (xGMI links) and between CPU and GPU in Frontier and El Capitan.}
\term{Instinct}{AMD's brand for data-center GPUs: MI100, MI200, MI300, MI350 and MI400 series.}
\end{terms}

\begin{intuition}
A CDNA compute unit and an NVIDIA SM are the same object drawn at different widths: both hold four schedulers, four vector units and four matrix units, but AMD's vector unit is 16 lanes fed by a 64-wide wavefront over four cycles and NVIDIA's is 32 lanes fed by a 32-wide warp each cycle. Hold that in mind when comparing counts: a CU is roughly half an SM in lanes, and a wavefront is twice a warp in threads.
\end{intuition}

% ---------------------------------------------------------------------
\slide{8}{MI250X and Frontier: two dies per module and the first exaflop}{Each MI250X holds two compute dies that software sees as two GPUs; 37,632 of them made Frontier (2022).}

\para{What the slide shows} Two diagrams side by side under the heading ``an MI250X module (left) and a Frontier node (right).'' The module is drawn as two tall green dies, GCD 0 and GCD 1 with 110 CUs each, joined by a double-headed arrow for the in-package Infinity Fabric, with four purple HBM2e stacks along each outer edge; the caption gives the in-package link (400 GB/s), the memory (four stacks of HBM2e per GCD, 64 GB each) and the eight external Infinity Fabric links. The node is drawn as one CPU box, ``EPYC 7A53 `Trento', 64 cores,'' above four MI250X boxes joined in a chain, with a caption for the coherent CPU--GPU links (36 $+$ 36 GB/s per link), the GPU--GPU links (50 $+$ 50 GB/s), the four Slingshot-11 NICs at 200 Gb/s and the 512 GB of DDR4. Below the diagrams, a panel ``The MI250X, by the numbers'' gives the module's peaks and transistor count; the right panel ``Frontier and its siblings'' gives the machine's size, its HPL history, the other systems built from the same blade, why AMD won the contracts, and what the two-GPU module cost its programmers.

\para{Details} The MI250X, launched November 8, 2021, is the first data-center GPU built from two compute dies in one package$^\dagger$. Each graphics compute die (GCD), on TSMC N6, carries 112 CUs of which 110 are enabled, 8 MB of L2, four HBM2e stacks (64 GB at 1.6 TB/s) and its own Infinity Fabric and PCIe 4.0 interfaces; the two GCDs sit side by side on an interposer with four in-package Infinity Fabric links between them, 400 GB/s bidirectional, and eight external Infinity Fabric links of 100 GB/s each leave the module (800 GB/s aggregate)$^\dagger$. The module is an OCP Accelerator Module (OAM, slide 23) rated 500 W, or 560 W with the higher power limit; the lower-clocked MI250 (208 CUs, 362 TFLOP/s FP16) and the single-GCD PCIe card MI210 (104 CUs, 64 GB, 300 W, March 2022) complete the series$^\dagger$. AMD's launch slide gave 58.2 billion transistors for the two dies; AMD's whitepaper does not repeat the figure, so the deck marks it as the launch-slide value$^\dagger$.

The peaks follow from the CDNA 2 compute unit (slide 7) at the 1.7 GHz peak clock: $220 \times 1.7\,\text{GHz} \times 128 = 47.9$ TFLOP/s FP64 vector, and the same 47.9 for FP32 vector because CDNA 2 runs FP64 at the FP32 rate; $220 \times 1.7 \times 256 = 95.7$ TFLOP/s for FP64 on the matrix cores and for packed FP32; $220 \times 1.7 \times 1{,}024 = 383$ TFLOP/s FP16 and BF16 on the matrix cores, and 383 TOPS INT8 (derived; all vendor peaks$^\dagger$). Against the contemporary A100 (slide 4), the MI250X had 2.5$\times$ the FP64 vector rate (47.9 against 19.5 on A100's tensor cores, 4.9$\times$ against its 9.7 vector rate), 1.6$\times$ the memory (128 GB against 80) and 1.6$\times$ the bandwidth (3.2 TB/s against 2.0), but only 1.2$\times$ the FP16 matrix rate (383 against 312) and no TF32, sparsity or INT4 modes (derived). It was a part designed for the FP64 contracts it had won, not for the deep-learning market, and its two-die construction had a consequence for software: ROCm presents each GCD as a separate device, so a node of four modules is eight GPUs, each with 64 GB, and an application that needs more than 64 GB per rank or wants to use the full module's bandwidth must itself distribute work across two devices joined by a 400 GB/s link.

Frontier, at Oak Ridge National Laboratory, is the machine the MI250X was built for$^\dagger$. It is an HPE Cray EX system of 9{,}408 nodes in 74 cabinets, each node one 64-core AMD EPYC 7A53 ``Trento'' processor (a third-generation EPYC with an Infinity Fabric interface to the GPUs) with 512 GB of DDR4, four MI250X modules (eight GCDs) with 512 GB of HBM2e, and four HPE Slingshot-11 network interfaces at 200 Gb/s, one attached to each GPU so that a GPU can send to the network without crossing the CPU; the nodes are joined in a dragonfly topology. The CPU and each GPU are connected by coherent Infinity Fabric at 36 GB/s in each direction, so the CPU and GPU share one address space with cache coherence, and the GPUs are connected to each other by Infinity Fabric links of 50 GB/s in each direction. The counts: 37{,}632 MI250X modules, 75{,}264 GCDs, 9{,}408 CPUs. On the June 2022 TOP500 list Frontier was the first system to exceed one exaflop on the High-Performance Linpack benchmark, 1.102 EFLOP/s; the run was tuned over the following years to 1.194 EFLOP/s (May 2023) and 1.353 EFLOP/s (November 2024), where it has stayed; in June 2026 it was third behind LineShine and El Capitan at 24.6 MW (slide 27)$^\dagger$. The HPL figure is 66\% of the list's 2.06 EFLOP/s peak for the machine (derived; slide 27).

The same blade design, the EX235a, built several other systems$^\dagger$: LUMI at CSC in Kajaani, Finland (2{,}978 GPU nodes, 11{,}912 MI250X, 380 PFLOP/s, eleventh in June 2026), Setonix at the Pawsey Centre in Australia, Adastra at CINES in France (1{,}424 MI250X plus later MI300A nodes) and Eni's HPC6 in Italy (478 PFLOP/s, eighth in June 2026). Together these made the MI250X the most widely deployed HPC GPU of 2022--23 outside NVIDIA's line.

Why AMD won these: the U.S. Department of Energy's CORAL-2 procurement, announced in 2018 and awarded in 2019, bought two exascale systems, Frontier (Oak Ridge, May 2019) and El Capitan (Livermore, August 2019; AMD's processors were confirmed in March 2020), from Cray with AMD CPUs and GPUs; the third CORAL-2 system, Aurora at Argonne, went to Intel (slide 19)$^\dagger$. The criteria were FP64 performance per dollar and per watt, a coherent CPU--GPU link from one vendor, and an open software stack that the laboratories could extend; AMD's bid offered full-rate FP64, Infinity Fabric between its own CPU and GPU, and ROCm, which is open source (slide 12 takes up the FP64 argument). The cost of the choice was software: in 2022 ROCm was less mature than CUDA, each application had to be ported to HIP or to a portability layer, and the two-GPUs-per-module structure complicated memory management. The Oak Ridge Center for Accelerated Application Readiness ported the first applications over 2019--22, and slide 25 reports what the portability studies found.

\begin{terms}
\term{MI250X}{AMD's CDNA 2 flagship (November 2021): two GCDs, 220 CUs, 128 GB HBM2e at 3.2 TB/s, 47.9 / 95.7 TFLOP/s FP64 vector / matrix, 500--560 W OAM; the Frontier GPU.}
\term{Frontier}{The HPE Cray EX system at Oak Ridge National Laboratory: 9,408 nodes of one EPYC 7A53 and four MI250X, 37,632 GPUs, the first exascale HPL result (1.102 EFLOP/s, June 2022), 1.353 EFLOP/s since November 2024.}
\term{HPL}{High-Performance Linpack, the dense LU factorization benchmark that ranks the TOP500; it measures FP64 matrix throughput and the network, and is what ``exaflop'' refers to in this deck.}
\term{TOP500}{The twice-yearly list (June and November) of the 500 fastest HPL results; slide 28 gives the June 2026 top ten.}
\term{CORAL-2}{The Department of Energy procurement that bought the three U.S. exascale systems: Frontier and El Capitan from Cray with AMD processors (2019), Aurora from Intel.}
\term{EPYC 7A53 Trento}{The 64-core AMD CPU of Frontier, a third-generation EPYC variant with coherent Infinity Fabric to the GPUs.}
\term{Slingshot}{HPE Cray's Ethernet-based high-performance network; Slingshot-11 provides 200 Gb/s per NIC in Frontier, Aurora and El Capitan.}
\term{dragonfly}{A network topology of fully connected groups joined by global links, used by Slingshot; it reaches tens of thousands of nodes in a few hops.}
\term{OAM}{OCP Accelerator Module, the Open Compute Project's mezzanine form factor for accelerators; MI250X, MI300X, MI355X, Ponte Vecchio and Gaudi 3 ship on it (slide 23).}
\term{packed FP32}{A CDNA 2 mode in which two FP32 values occupy one 64-bit register and one instruction operates on both, doubling the FP32 vector rate when the compiler can pair operations.}
\term{coherent CPU--GPU link}{A connection over which CPU and GPU share one address space with hardware cache coherence, so either can read the other's memory without explicit copies; Infinity Fabric in Frontier, NVLink-C2C in Grace Hopper.}
\end{terms}

\begin{intuition}
Frontier is an argument made in hardware: one vendor's CPU and GPU joined coherently, double precision at full rate, and everything bought on FP64 per watt. The price was two GPUs in every package and a software stack the laboratories had to help finish. Keep both halves in view, because the same trade reappears in every later AMD design win.
\end{intuition}

% ---------------------------------------------------------------------
\slide{9}{MI300X and MI300A: 3D chiplets and the exascale APU}{Eight compute dies stacked on four I/O dies with a 256 MB cache; with CPU cores it became El Capitan.}

\para{What the slide shows} A plan view of the MI300X package fills the upper left: four tall pale rectangles are the I/O dies, each labeled with its 64 MB share of the Infinity Cache, and on them stand eight narrow green rectangles, the XCDs of 38 CUs each, two per IOD; four purple HBM3 stacks stand at each side. The caption explains that the XCDs are hybrid-bonded onto the IODs, which hold the 256 MB Infinity Cache (17.2 TB/s), the memory controllers and the Infinity Fabric, and that the MI300A swaps two XCDs for three Zen 4 CPU dies (24 cores) sharing the same 128 GB. Two stat tiles give 153 billion transistors on 12 dies and 5.3 TB/s from 192 GB of HBM3 in 8 stacks. The right panel, ``MI300X, MI300A, and their machines,'' gives the MI300X's figures and its ratios to the H100, the MI300A's figures, El Capitan's size and HPL history, and the MI325X memory refresh of October 2024.

\para{Details} CDNA 3 is the generation in which AMD's chiplet experience from its CPUs (the Zen 2 EPYC of 2019 put eight compute dies around one I/O die) was applied to a GPU in three dimensions$^\dagger$. The MI300X package, announced in June 2023 and launched on December 6, 2023, holds twelve dies: four I/O dies (IODs) on TSMC N6, each about 370 mm$^2$, laid side by side on a silicon interposer and joined by Infinity Fabric, and eight accelerator complex dies (XCDs) on TSMC N5, each about 115 mm$^2$, stacked face-to-face on the IODs by hybrid bonding (copper-to-copper bonding without solder bumps, at a pitch fine enough for tens of thousands of connections per square millimeter), two XCDs per IOD; eight HBM3 stacks sit on the interposer at the edges, two per IOD$^\dagger$. Each XCD carries 40 CUs of which 38 are enabled, with 4 MB of L2; the eight XCDs give 304 CUs, 19{,}456 lanes and 1{,}216 matrix cores, and the package's 153 billion transistors (146 billion on the MI300A) made it the largest chip AMD had built$^\dagger$.

The IODs do the work that a monolithic die's uncore does. They hold the 256 MB Infinity Cache, a memory-side last-level cache of 64 MB per IOD that sits in front of the HBM controllers and serves every XCD with an aggregate bandwidth of 17.2 TB/s (the whitepaper figure; a Hot Chips 2024 slide gave 14.7 TB/s, and the deck follows the whitepaper)$^\dagger$; the eight HBM3 controllers for 192 GB at 5.3 TB/s; the Infinity Fabric network that connects the XCDs to each other and to the cache; and the external interfaces, seven Infinity Fabric links of 128 GB/s each (896 GB/s aggregate) that connect the GPU directly to each of the other seven on an eight-GPU Universal Baseboard (UBB 2.0), and one PCIe 5.0 x16 host link. Because every XCD can reach every HBM stack through the IODs and the cache, the MI300X presents itself to software as a single GPU with 192 GB, unlike the MI250X's two devices (it can also be partitioned into up to eight, one per XCD, for small workloads). The module is an OAM at 750 W$^\dagger$.

The peaks at the 2.1 GHz peak clock: $304 \times 2.1\,\text{GHz} \times 128 = 81.7$ TFLOP/s FP64 vector; $304 \times 2.1 \times 256 = 163.4$ TFLOP/s FP64 matrix, and the same 163.4 for FP32 vector and matrix (CDNA 3 doubled the vector FP32 rate per CU); $304 \times 2.1 \times 1{,}024 = 653.7$ TFLOP/s TF32; $304 \times 2.1 \times 2{,}048 = 1{,}307.4$ TFLOP/s FP16 and BF16; $304 \times 2.1 \times 4{,}096 = 2{,}614.9$ TFLOP/s FP8 and TOPS INT8, all dense; AMD's datasheet also lists sparse figures of twice each matrix rate (derived from the per-CU rates; peaks$^\dagger$). The slide's comparison with the H100 SXM (slide 4) is 192 GB against 80 GB, 2.4$\times$; 5.3 against 3.35 TB/s, 1.6$\times$; 2{,}615 against 1{,}979 TFLOP/s dense FP8, 1.3$\times$; 163.4 against 67 TFLOP/s FP64 matrix, 2.4$\times$ (derived). On paper the MI300X led the H100 in every row of the table; slide 26 reports what the achieved fraction looked like a year later, and slide 30 what MLPerf measured.

The MI300A is the same package with the two XCDs on one side replaced by three Zen 4 core complex dies (CCDs), the same CCDs used in EPYC ``Genoa'' processors, eight cores each, 24 cores in all; the six remaining XCDs give 228 CUs$^\dagger$. CPU and GPU share the four IODs, the Infinity Cache and the 128 GB of HBM3 (eight stacks of 16 GB) in one physical address space, so there is no separate CPU memory to copy to and from: an accelerated processing unit (APU) in AMD's term, the first data-center APU. Its peaks are 61.3 / 122.6 TFLOP/s FP64 vector / matrix (228 CUs at 2.1 GHz), 980.6 TFLOP/s FP16, 1{,}961 TFLOP/s FP8; its power is 550 W air-cooled or 760 W liquid-cooled; it ships in an SH5 socket rather than an OAM, and four of them connect through Infinity Fabric on an El Capitan node$^\dagger$.

El Capitan, at Lawrence Livermore National Laboratory, is the MI300A machine$^\dagger$: an HPE Cray EX system of about 11{,}100 compute nodes of four MI300A each (11{,}136 on the November 2024 list; LLNL's current page counts 11{,}424 batch nodes plus debug and login nodes, 46{,}080 APUs in all), with Slingshot-11 and no discrete CPUs. Its HPL run of 1.742 EFLOP/s put it first on the November 2024 list; it was dedicated on January 9, 2025; the run was tuned to 1.809 EFLOP/s for November 2025, where it was first again, and in June 2026 it was second behind the CPU-only LineShine at 29.7 MW (slide 27)$^\dagger$. The deck's ``about 44{,}500'' is the November 2024 configuration of $11{,}136 \times 4 = 44{,}544$ APUs (derived). El Capitan's HPL per megawatt, 61 PFLOP/s, is the highest of the DOE machines (slide 27), which is the APU argument made in measurement: no CPU memory, no PCIe hop, 128 GB of HBM shared.

The MI325X (October 10, 2024) is the MI300X silicon with the HBM3 replaced by eight stacks of HBM3e, 256 GB at 6.0 TB/s, at 1{,}000 W$^\dagger$; the compute peaks are unchanged. It answered NVIDIA's H200 (141 GB at 4.8 TB/s), which had done the same to the H100 (slide 4): both vendors learned in 2024 that a memory refresh of an existing die could be brought to market in a year and that inference customers would pay for capacity.

\begin{terms}
\term{MI300X}{AMD's CDNA 3 GPU (December 2023): 8 XCDs on 4 IODs, 304 CUs, 192 GB HBM3 at 5.3 TB/s, 256 MB Infinity Cache, 81.7 / 163.4 TFLOP/s FP64, 1,307 TFLOP/s dense FP16, 750 W OAM.}
\term{MI300A}{The APU variant of the MI300: 6 XCDs (228 CUs) and 3 Zen 4 CCDs (24 cores) on the same IODs, sharing 128 GB of HBM3; the El Capitan processor.}
\term{APU}{Accelerated processing unit, AMD's term for a package with CPU cores and a GPU sharing one memory; the MI300A is the first data-center APU.}
\term{hybrid bonding}{Die stacking by direct copper-to-copper bonding of the two dies' top metal layers, without solder bumps; it gives connection pitches of about 10 micrometers or less and joins the XCDs to the IODs in MI300.}
\term{Infinity Cache}{AMD's memory-side last-level cache on the I/O dies: 256 MB at 17.2 TB/s on MI300X and MI355X, serving all compute dies in front of the HBM; dropped on MI455X in favor of a 192 MB L2.}
\term{UBB 2.0}{The Universal Baseboard, the OCP eight-module board on which MI300X and MI355X OAMs are mounted and connected by Infinity Fabric; the AMD counterpart of NVIDIA's HGX board.}
\term{Zen 4 CCD}{AMD's eight-core CPU chiplet of 2022, used in EPYC Genoa; three of them give the MI300A its 24 cores.}
\term{El Capitan}{The HPE Cray EX system at Lawrence Livermore National Laboratory: about 11,100 nodes of four MI300A, 1.742 EFLOP/s on HPL in November 2024 (first), 1.809 EFLOP/s from November 2025, 29.7 MW.}
\term{MI325X}{The October 2024 refresh of the MI300X with 256 GB of HBM3e at 6.0 TB/s and 1,000 W; same compute.}
\term{HBM3e}{The extended HBM3 generation, with higher per-pin rate and 12-high stacks of up to 36 GB; used on MI325X, MI355X, H200 and Blackwell.}
\term{memory-side cache}{A cache placed in front of the memory controllers rather than next to the compute units, so that it serves all compute dies and caches by address, not by requester.}
\end{terms}

\begin{intuition}
The MI300 is AMD's CPU chiplet strategy turned vertical: put the uncore (cache, memory controllers, fabric) on large, cheap base dies and stack small compute dies on top, so that the compute can move to a new process node without redesigning the memory system. Swap two compute dies for CPU dies and the GPU becomes a processor with 128 GB of HBM as its only memory; that one substitution is El Capitan.
\end{intuition}

% ---------------------------------------------------------------------
\slide{10}{CDNA 4: MI355X, 288 GB and FP4, matrix FP64 halved}{The same recipe as Blackwell Ultra one quarter later: more HBM, 4-bit formats, less double precision, 1,400 W.}

\para{What the slide shows} A horizontal bar chart on the left, ``what CDNA 4 changed (MI300X $=$ 1),'' plots the MI355X relative to the MI300X, per GPU, on nine rows: dense FP4 (new, no MI300X figure, drawn at 3.4 to mark it), dense FP8 (2.6 $\to$ 5 PFLOP/s, 1.9$\times$), dense FP16 (1.3 $\to$ 2.5 PFLOP/s, 1.9$\times$), HBM capacity (192 $\to$ 288 GB, 1.5$\times$), HBM bandwidth (5.3 $\to$ 8 TB/s, 1.5$\times$), Infinity Fabric per GPU (896 $\to$ 1{,}008 GB/s, 1.1$\times$), FP64 vector (81.7 $\to$ 78.6 TFLOP/s, 0.96$\times$), FP64 matrix (163 $\to$ 78.6 TFLOP/s, 0.48$\times$) and power (750 $\to$ 1{,}400 W, 1.9$\times$); a dashed line marks 1. The right panel, ``CDNA 4 (announced June 12, 2025; shipping Q3 2025),'' gives the silicon, the formats and their rates, the first measured results and the first deployment.

\para{Details} CDNA 4 is AMD's generation of 2025, announced at the Advancing AI event on June 12, 2025 and shipping in the third quarter$^\dagger$. The package keeps the MI300's construction but changes both kinds of die: eight XCDs on TSMC N3P, each with 36 CUs of which 32 are enabled (256 CUs, 16{,}384 lanes, 1{,}024 matrix cores), stacked on two I/O dies on N6 (the MI300X's four IODs merged into two larger ones, which removes the die-to-die crossings on the base layer), with the 256 MB Infinity Cache, eight stacks of 12-high HBM3e for 288 GB at 8 TB/s, and 185 billion transistors in all$^\dagger$. Two products share the silicon: the MI355X, liquid-cooled at 1{,}400 W with a 2.4 GHz peak clock, and the MI350X, air-cooled at 1{,}000 W at 2.2 GHz; both are OAMs on UBB 2.0 boards of eight with seven Infinity Fabric links per GPU, now 144 GB/s each (1{,}008 GB/s aggregate)$^\dagger$.

The compute unit changed in three ways. The LDS grew from 64 KB to 160 KB per CU, the first increase since GCN, bringing it closer to the 228 KB of shared memory an H100 SM can configure (slide 14 compares the hierarchies). The matrix cores gained two block-scaled formats from the Open Compute Project's Microscaling (MX) specification, MXFP6 and MXFP4, in which a block of 32 values shares one 8-bit exponent scale, so that each value costs 6 or 4 bits plus a quarter of a bit of shared scale; FP8 was already supported in both the OCP FP8 form and the MXFP8 form. And the FP64 matrix rate per CU per clock was halved, from 256 to 128 FLOP, so that the matrix cores do FP64 no faster than the vector units$^\dagger$. The per-CU-per-clock matrix rates doubled for the 16-bit and 8-bit formats, from 2{,}048 to 4{,}096 FLOP for FP16 and BF16 and from 4{,}096 to 8{,}192 for FP8, with 16{,}384 for MXFP6 and MXFP4.

The peaks for the MI355X at 2.4 GHz: FP64, $256 \times 2.4\,\text{GHz} \times 128 = 78.6$ TFLOP/s, vector and matrix alike; FP32, $256 \times 2.4 \times 256 = 157.3$ TFLOP/s; FP16 and BF16, $256 \times 2.4 \times 4{,}096 = 2{,}516$ TFLOP/s, quoted as 2.5 PFLOP/s; FP8, 5.0 PFLOP/s; MXFP6 and MXFP4, 10.1 PFLOP/s, quoted as 10; INT8 5.0 POPS (derived from the per-CU rates; peaks$^\dagger$). The MI350X at 2.2 GHz scales each by $2.2/2.4$: 72.1 TFLOP/s FP64, 2.3 PFLOP/s FP16, 9.2 PFLOP/s FP4 (derived). AMD's whitepaper lists the FP16 and FP8 figures with structured-sparsity doubles (5.0 and 10 PFLOP/s) but gives no sparse figure for FP6 and FP4 (the datasheet lists N/A, though AMD's MLPerf Inference 5.1 blog quotes ``up to 20 petaflops of FP4'' for the MI355X), and the deck says so; the 10 PFLOP/s of FP4 is a dense number, the same dense FP4 rate as NVIDIA's GB200 (slide 4), which is the convergence slide 6 describes.

The FP64 decision deserves its own paragraph because it reverses two generations of AMD's design. From the MI250X to the MI300X the matrix cores did FP64 at twice the vector rate, and 163.4 TFLOP/s of FP64 matrix on the MI300X was the figure that made it the fastest HPC GPU of 2023 (slide 12). On CDNA 4 the matrix FP64 is 78.6 TFLOP/s, half the MI300X's, and only 4\% below the MI300X's vector rate (81.7) because the clock went up while the CU count went down ($256 \times 2.4 = 614$ against $304 \times 2.1 = 638$ CU-GHz; derived). The transistors saved went to the low-precision matrix paths. The slide puts this beside NVIDIA's parallel decision: Blackwell Ultra (the B300, announced March 2025, shipping from July 2025) kept 288 GB of HBM3e at 8 TB/s, raised dense FP4 to 15 PFLOP/s, and cut FP64 to about 1.4 TFLOP/s per GPU (slide 12 and the previous deck). The MI355X's 78.6 TFLOP/s is $78.6/1.4 = 56\times$ the B300's (derived), so the two vendors moved in the same direction by very different amounts: AMD halved FP64 on its AI part and kept a separate HPC part (the MI430X, slide 11), NVIDIA removed it from the AI part and restored a middle figure on Rubin (33 TFLOP/s).

What was measured. The MI355X's first public results were in MLPerf Inference 5.1 (September 2025), which also contained the first MXFP4 submissions by any vendor; on Llama 2 70B the MI355X delivered 2.7$\times$ the MI325X$^\dagger$. In MLPerf Inference 6.0 (April 2026) AMD's own comparison put the MI355X at 100\% of the B200 per GPU on the Llama 2 70B offline scenario, 97\% on the server scenario and 119\% on the interactive scenario (slide 30), and a 87-GPU multi-node run passed one million tokens per second$^\dagger$. The first deployment named on the slide is Oak Ridge's Lux, an AI cluster of MI355X nodes built by HPE and operated with Oracle, announced in October 2025 for early 2026$^\dagger$; cloud instances were announced by Oracle Cloud Infrastructure among others$^\dagger$.

Read the bar chart as a statement about where the transistors went. Memory and bandwidth rose 1.5$\times$ with the move to 12-high HBM3e; the low-precision matrix rates rose about 1.9$\times$ per GPU on 16\% fewer CUs at a 14\% higher clock, which means the per-CU rate doubled; FP64 stood still or fell; and power rose 1.9$\times$, so that performance per watt at FP16 and FP8 barely changed from the MI300X (2.5 PFLOP/s at 1{,}400 W against 1.3 at 750 W: 1.8 against 1.7 TFLOP/s per watt; derived) and the gain is in FP4, where there was no MI300X figure to compare.

\begin{terms}
\term{CDNA 4}{AMD's 2025 compute architecture: N3P compute dies, 160 KB LDS per CU, MXFP6 and MXFP4 matrix formats, 288 GB HBM3e, FP64 matrix rate halved per CU.}
\term{MI355X}{The CDNA 4 flagship: 256 CUs at 2.4 GHz, 288 GB HBM3e at 8 TB/s, 78.6 TFLOP/s FP64, 2.5 / 5 / 10 PFLOP/s dense FP16 / FP8 / FP4, 1,400 W liquid-cooled OAM.}
\term{MI350X}{The air-cooled 1,000 W sibling of the MI355X at 2.2 GHz, with peaks about 8\% lower.}
\term{N3P}{TSMC's performance-tuned 3 nm process, used for the MI355X's compute dies and for the MI455X's cache and I/O dies.}
\term{MXFP4}{The Open Compute Project Microscaling 4-bit floating-point format: E2M1 values in blocks of 32 sharing one 8-bit power-of-two scale; supported by CDNA 4, CDNA 5 and Blackwell.}
\term{MXFP6}{The 6-bit member of the same family (E2M3 or E3M2 values, blocks of 32, shared 8-bit scale); CDNA 4 runs it at the FP4 rate.}
\term{block scaling}{Sharing one scale factor among a small block of values (32 in MX formats, 16 in NVIDIA's NVFP4) so that a narrow format can cover a wide dynamic range tensor; the successor to per-tensor scaling.}
\term{12-high HBM3e}{HBM3e stacks of twelve DRAM dies, 36 GB per stack, giving 288 GB from eight stacks on the MI355X, B300 and Rubin.}
\term{FP64 matrix rate}{The double-precision throughput of the matrix units, which on CDNA 2 and 3 was twice the vector rate and on CDNA 4 equals it.}
\term{Blackwell Ultra}{NVIDIA's 2025 refresh of Blackwell (B300, GB300): 288 GB HBM3e, 15 PFLOP/s dense FP4, FP64 cut to about 1.4 TFLOP/s, 1,400 W; the NVIDIA counterpart of CDNA 4.}
\term{Lux}{Oak Ridge National Laboratory's AI cluster of MI355X nodes (HPE, with Oracle), announced October 2025 for early 2026; the first named MI355X deployment on the slide.}
\end{terms}

\begin{intuition}
CDNA 4 is CDNA 3 with the dial turned from HPC toward inference: more memory, 4-bit formats, double the 16- and 8-bit matrix rates per compute unit, half the FP64 matrix rate, and nearly twice the power. That AMD and NVIDIA made the same move in the same quarter tells you the market, not the architecture, chose the direction.
\end{intuition}

% ---------------------------------------------------------------------
\slide{11}{MI400, MI455X and Helios: AMD's rack (2026)}{CDNA 5 on TSMC N2, 432 GB of HBM4, UALink over Ethernet, 72 GPUs per rack; shipping from Q3 2026.}

\para{What the slide shows} A front elevation of a Helios rack on the left, under the heading ``Helios: 72 MI455X, 18 EPYC `Venice' CPUs, 18 DPUs, 800G NICs, in one OCP Open Rack Wide'': nine green compute trays at the top and nine at the bottom, each drawn with four GPU bars, and six purple switch trays in the middle; labels to the right of the rack give the tray contents (18 compute trays of 4 MI455X, 1 Venice and 1 Salina DPU each), the rack totals (72 GPUs, 31 TB of HBM4 at 1.7 PB/s, 260 TB/s of UALink-over-Ethernet scale-up, 3.6 TB/s per GPU over 72 lanes), the switch trays (UALink over Ethernet), the scale-out (Vulcano 800G NICs, 43 TB/s), the power (225--245 kW), the cooling (liquid) and the form factor (double-wide for serviceability), and the rack peaks (2.9 EFLOP/s FP4, 1.4 EFLOP/s FP8, dense). The right panel, ``MI455X and the deals,'' gives the GPU's silicon and peaks, its shipping status as of November 2026, the customer commitments announced in 2025--26, and the next products.

\para{Details} The MI400 series was previewed on June 12, 2025, with the MI355X: a 2026 GPU with 432 GB of HBM4 at 19.6 TB/s, 40 PFLOP/s of FP4 and 20 PFLOP/s of FP8, to ship in a 72-GPU rack called Helios$^\dagger$. The agreements of late 2025 (below) name the series ``MI450,'' and the part itself, the MI455X, was unveiled at Advancing AI on July 23, 2026, with the HBM4 bandwidth raised to 23.3 TB/s$^\dagger$. CDNA 5's package is the third form of AMD's stacking: eight accelerator complex dies on TSMC N2, two cache dies (FCDs) and two I/O dies on N3P, assembled with CoWoS-L packaging and hybrid bonding, 320 billion transistors, twelve HBM4 stacks of 36 GB$^\dagger$. The compute is counted in 256 work-group processors (WGPs), the RDNA-style dual-CU unit, each with 320 KB of LDS; the lane count has not been published. The Infinity Cache of the two previous generations is gone, replaced by a 192 MB shared L2 split across the two cache dies (96 MB each), a change that makes the memory hierarchy look more like NVIDIA's (a large L2, no memory-side cache) and less like the MI300's$^\dagger$.

The peaks, all dense and vendor-stated$^\dagger$: 40.26 PFLOP/s MXFP4, 20.13 PFLOP/s MXFP6 and MXFP8, 5.03 PFLOP/s FP16 and BF16, 315 TFLOP/s FP32; the clock and the per-WGP rates were not given. Two ratios are new: MXFP8 is four times FP16, not two as on CDNA 4 (so the 8-bit matrix path was widened relative to the 16-bit one), and MXFP4 is twice MXFP8 as before (derived). FP64 has no published figure; AMD's answer for FP64 is the separate MI430X (below), and the slide marks the MI455X's FP64 and power as not published. Against Rubin (slide 4, as revised in September 2026): 432 against 288 GB (1.5$\times$), 23.3 against 19.2 TB/s (1.2$\times$), 40.3 against 35 PFLOP/s dense FP4 (1.15$\times$), 20.1 against 17.5 PFLOP/s FP8 (1.15$\times$), 5.0 against 4.0 PFLOP/s FP16 (derived). The scale-up link is the largest change: each MI455X has 72 lanes of UALink carried over Ethernet (UALoE), 3.6 TB/s bidirectional, 3.6$\times$ the MI355X's 1{,}008 GB/s of Infinity Fabric and 1.2$\times$ Rubin's 3.0 TB/s of NVLink 6 (derived; slide 21), plus a 256 GB/s host link and a Pensando Vulcano network interface for scale-out at 800 Gb/s per port$^\dagger$. The module is a new form factor AMD calls the Enhanced Accelerator Module, and its power rating has not been published; the rack figure of 225--245 kW divided by 72 GPUs puts the GPU share at roughly 2 kW with everything else included, but that is a reader's bound, not a specification.

Helios is the rack$^\dagger$. Eighteen compute trays each hold four MI455X, one EPYC ``Venice'' (sixth-generation EPYC, Zen 6) and one Pensando Salina data processing unit: 72 GPUs, 18 CPUs, 18 DPUs, a ratio of one CPU per four GPUs against NVIDIA's one per two (slide 22). The switch trays carry UALink over Ethernet switches, so that every GPU reaches every other at 3.6 TB/s, 260 TB/s across the rack, the figure AMD compares with the NVL72's 216 TB/s after NVIDIA's September 2026 revision (slide 22). The totals are the per-GPU figures times 72: $72 \times 432\,\text{GB} = 31.1$ TB of HBM4; $72 \times 23.3 = 1{,}678$ TB/s, quoted as 1.7 PB/s (the Supermicro datasheet, written before the bandwidth increase, still says 1.4 PB/s); $72 \times 40.26 = 2{,}899$ PFLOP/s, 2.9 EFLOP/s of FP4; $72 \times 20.13 = 1{,}449$, 1.4 EFLOP/s of FP8 (derived). Scale-out is 43 TB/s aggregate through the Vulcano 800G NICs$^\dagger$. The rack is liquid-cooled, draws 225--245 kW by press accounts, and is built to the Open Compute Project's Open Rack Wide specification, a double-width rack AMD chose for serviceability (slide 23); it was shown at the OCP Global Summit in October 2025$^\dagger$.

Status as of November 2026$^\dagger$. AMD's third-quarter 8-K of November 3, 2026 says the MI450 series ``launched in the quarter'' with Helios racks running workloads at Anthropic, Meta, Microsoft, OpenAI and Oracle Cloud (the deck's table and panel quote the filing); on the call AMD spoke of shipping ``in volume'' and of first Helios rack revenue, a ``big step up'' expected in the fourth quarter and a ramp through 2027; the data-center segment's revenue was \$8.5 billion, up 97\% from a year earlier (slide 29). The commitments that precede the ramp: OpenAI agreed on October 6, 2025 to deploy 6 gigawatts of Instinct GPUs beginning with 1 GW of MI450 in the second half of 2026, with a warrant for up to 160 million AMD shares that vests as the deployments are met; Oracle announced about 50{,}000 MI450 for its cloud from the third quarter of 2026; Meta agreed on February 24, 2026 to deploy a custom MI450-based part, up to 6 GW, with an identical warrant for up to 160 million shares; Anthropic committed on July 22, 2026 to up to 2 GW of MI450-series GPUs from the first half of 2027, with AMD investing up to \$5 billion in Anthropic; and Microsoft Azure offers Helios racks. Slide 29 puts these beside NVIDIA's and notes that a gigawatt is a unit of power, not a count of GPUs.

Next. The MI430X, announced in November 2025, is the HPC member of the MI400 series: the same 432 GB of HBM4 with ``true hardware FP64'' at 288 TFLOP/s (stated in July 2026), shipping in the first half of 2027, and already the GPU of three announced systems: Discovery at Oak Ridge (HPE Cray GX5000, with Venice CPUs, delivery 2028), Alice Recoque at CEA/GENCI in France (Eviden, 94 racks, over one exaflop on HPL) and LUMI-AI (EuroHPC, second half of 2027) (slide 29)$^\dagger$. Beyond it AMD's roadmap names the MI500 for 2027, with a successor Helios rack and EPYC ``Verano'' CPUs, and the MI600 for 2028$^\dagger$.

\begin{terms}
\term{CDNA 5}{AMD's 2026 compute architecture (MI400 series): N2 compute dies, work-group processors, 192 MB L2 on cache dies, HBM4, UALink over Ethernet.}
\term{MI455X}{The CDNA 5 AI GPU unveiled July 23, 2026: 256 WGPs, 320 billion transistors, 432 GB HBM4 at 23.3 TB/s, 40.3 / 20.1 / 5.0 PFLOP/s dense MXFP4 / MXFP8 / FP16, 3.6 TB/s UALoE; FP64 and power not published.}
\term{MI450}{The series name AMD and its customers used for the MI455X generation in the 2025 agreements (OpenAI, Oracle, Meta).}
\term{MI430X}{The HPC member of the MI400 series: 432 GB HBM4, 288 TFLOP/s FP64, first half of 2027; the GPU of Discovery, Alice Recoque and LUMI-AI.}
\term{Helios}{AMD's 72-GPU rack: 18 trays of 4 MI455X, 1 Venice and 1 Salina DPU, UALink-over-Ethernet switch trays, 31 TB HBM4, 2.9 EFLOP/s dense FP4, 260 TB/s scale-up, 225--245 kW, OCP Open Rack Wide; first revenue in Q3 2026.}
\term{FCD}{The cache die of the MI455X package: two on N3P, each carrying 96 MB of the 192 MB shared L2.}
\term{HBM4}{The fourth HBM generation, with a 2,048-bit interface per stack (twice HBM3's); twelve stacks give the MI455X 432 GB at 23.3 TB/s, eight give Rubin 288 GB at 19.2 TB/s.}
\term{UALink over Ethernet (UALoE)}{Carrying the UALink memory-semantic protocol on Ethernet physical links and switches; Helios's scale-up fabric, 72 lanes and 3.6 TB/s per GPU.}
\term{EPYC Venice}{AMD's sixth-generation EPYC server CPU (Zen 6), one per compute tray in Helios and the CPU of Discovery and Alice Recoque.}
\term{Pensando Salina and Vulcano}{AMD's data processing unit (Salina, 400G) and network interface (Vulcano, 800G), one DPU per tray and one NIC per GPU in Helios.}
\term{CoWoS-L}{TSMC's chip-on-wafer-on-substrate packaging with local silicon interconnect bridges in an organic interposer, used for very large multi-die packages (MI455X, Blackwell).}
\term{Open Rack Wide}{The Open Compute Project's double-width rack specification, to which Helios is built.}
\term{Enhanced Accelerator Module}{AMD's name for the MI455X's module form factor, the successor of the OAM.}
\end{terms}

\begin{intuition}
With Helios AMD stopped selling a GPU and started selling what NVIDIA sells: a rack with its own CPU, DPU, fabric and switch trays, delivered as a unit. The silicon numbers favor AMD by 15 to 50 percent on paper; the open questions, which slide 32 lists, are whether the Ethernet-based fabric and ROCm hold up at that scale, and those are measured, not announced.
\end{intuition}

% ---------------------------------------------------------------------
\slide{12}{FP64 across the three vendors: who still builds for HPC}{AMD kept doubling it, NVIDIA cut it, Intel ran it at full rate once; three of five exascale machines are AMD or Intel.}

\para{What the slide shows} A bar chart, ``peak FP64 per GPU, matrix rate where it differs,'' on a logarithmic axis from 1 to several hundred TFLOP/s, with ten bars in three colors: NVIDIA's A100 (19.5), H100 (67), GB200 (40), GB300 (1.4) and Rubin (33) in gray; AMD's MI250X (95.7), MI300X (163.4), MI355X (78.6) and MI430X (288) in green; Intel's Ponte Vecchio (52.4) in orange; each bar carries its value. A caption notes that the MI430X is a 2027 part with a vendor-stated figure, that NVIDIA's Blackwell and Rubin have no separate matrix FP64 rate, and that Ponte Vecchio's FP64 equals its FP32 rate. The right panel, ``Why the lines diverged,'' gives four paragraphs: AMD designed to contracts, NVIDIA designed to inference, Intel's one entry, and the scoreboard of exascale systems.

\para{Details} Double precision is the arithmetic of simulation: a climate model, a reactor code or a molecular dynamics run accumulates over millions of steps and needs the 53-bit mantissa of FP64 to keep its error below the physics it computes, and the High-Performance Linpack benchmark that ranks the TOP500 is an FP64 LU factorization. It is also arithmetic that deep learning does not use at all. Over 2020--2026 the three vendors answered that split in three ways, and the chart is the record.

AMD's line rises and then bends. The MI250X's 95.7 TFLOP/s of FP64 on the matrix cores (47.9 vector) in 2021 was already more than any NVIDIA part has offered since; the MI300X's 163.4 (81.7 vector) in 2023 is the highest FP64 rate of any shipped GPU on the chart; the MI355X's 78.6 (vector and matrix, slide 10) halved the matrix rate in 2025; and the MI430X's 288 TFLOP/s, a vendor statement of July 2026 for a part shipping in 2027, would be the highest again$^\dagger$. The reason is the one slide 8 gave: AMD's compute GPUs were specified by the contracts they were built to win. The CORAL-2 awards of 2019 for Frontier and El Capitan set FP64 targets, and so did Europe's EuroHPC procurements (LUMI, Adastra), and the MI250X and MI300A were those machines' processors; the MI300X carried the MI300A's FP64 into the AI market because it shared the dies; CDNA 4 could halve matrix FP64 because the FP64 customers were being moved to a separate SKU, and the MI430X is that SKU, sold to Discovery at Oak Ridge (2028), Alice Recoque in France and LUMI-AI with ``true hardware FP64'' as its selling point against emulation$^\dagger$.

NVIDIA's line rises to 67 and falls to 1.4. The A100's 19.5 TFLOP/s (its tensor cores doubled the 9.7 vector rate) and the H100's 67 (34 vector) were the HPC GPUs of 2020--22, and the H100, in its GH200 superchip form, is the GPU of JUPITER, Europe's exascale machine (slide 27). Blackwell's GB200 gives 40 TFLOP/s with no separate matrix rate (the tensor cores no longer accelerate FP64), Blackwell Ultra's B300 about 1.4 TFLOP/s, a figure read from the GB300 NVL72's 100 TFLOP/s per rack of 72 GPUs (derived; the Hot Chips 2025 table gives 1.39), and Rubin's spec page gives 33 TFLOP/s, half the H100's matrix rate$^\dagger$. The B300 removed most of the FP64 units to make room for the FP4 path and the larger special-function units that inference needs; what NVIDIA offers HPC instead is emulation, in which an FP64 matrix multiply is decomposed into INT8 tensor-core products (the Ozaki scheme) with enough slices to reproduce 53-bit accuracy, shipped in cuBLAS from CUDA 13.0 Update 2 (October 2025) and reported by NVIDIA at up to 2.3$\times$ the native FP64 rate on a GB200 NVL72$^\dagger$. The scheme depends on INT8 tensor throughput, which Blackwell has (4.5 POPS dense on the B200) and Blackwell Ultra largely gave up along with FP64 (the GB300 NVL72 page lists 24 POPS of sparse INT8 per rack, about 0.17 POPS dense per GPU; derived), so for a 2025 NVIDIA AI GPU neither path to FP64 is fast; Rubin's 33 TFLOP/s of native FP64 is NVIDIA's correction. The practical consequence for a buyer is on the slide's last line: ask whether a quoted FP64 figure is the vector or the matrix rate, native or emulated, and what fraction a real code achieves (slide 26).

Intel's single bar, Ponte Vecchio's 52.4 TFLOP/s, is a different design choice: the Xe-HPC vector engine does eight FP64 FMAs per clock, half its FP32 lane count, so FP64 runs at the full FP32 rate with no separate matrix path (slide 18). In 2023 that was more than the H100's vector rate and less than its matrix rate, and it is what gave Aurora its exaflop (slide 19). Intel's successors, Rialto Bridge and Falcon Shores, were cancelled (slide 20), so the bar has no successor on the chart.

The scoreboard the panel gives is the June 2026 TOP500's five exascale systems (slide 28): El Capitan (MI300A, 1.809 EFLOP/s), Frontier (MI250X, 1.353), Aurora (Max 1550, 1.012) and JUPITER (GH200, 1.000), and LineShine (2.198), which has no accelerator at all and reaches its result on Chinese CPUs$^\dagger$. Counted by GPU vendor, that is two AMD machines, one Intel, one NVIDIA and one with none; counted by architecture, four of the five are not NVIDIA, which is the sense in which the AMD and Intel designs ``won'' the exascale decade. The two columns of the comparison, the HPC list and the AI clusters of slide 29, are bought on different arithmetic, and the FP64 chart is the clearest single picture of why the same vendors rank differently on the two.

\begin{formulas}
FP64 peak from the compute unit: $P_{64} = N_{\text{CU}} \times f \times R_{64}$, with $R_{64}$ the FP64 FLOP per CU per clock (vector: 64 on CDNA 1, 128 on CDNA 2--4; matrix: 256 on CDNA 2--3, 128 on CDNA 4). MI300X: $304 \times 2.1\,\text{GHz} \times 256 = 163.4$ TFLOP/s. MI355X: $256 \times 2.4 \times 128 = 78.6$. Per-GPU figure from a rack: GB300 NVL72, $100\,\text{TFLOP/s} / 72 = 1.39$ TFLOP/s. Ratio of the two 2025 AI parts: $78.6 / 1.4 = 56$.
\end{formulas}

\begin{terms}
\term{FP64}{IEEE double precision: 1 sign, 11 exponent and 52 stored mantissa bits (53 with the hidden bit); the arithmetic of most simulation codes and of HPL.}
\term{vector FP64 rate}{Double-precision throughput of the SIMD or CUDA-core lanes; the rate a non-matrix code sees.}
\term{matrix FP64 rate}{Double-precision throughput of the tensor or matrix cores for matrix multiply; twice the vector rate on A100, H100 and CDNA 2--3, equal to it on Blackwell, Rubin and CDNA 4.}
\term{FP64 emulation}{Computing a double-precision matrix multiply from several lower-precision tensor-core products (the Ozaki scheme on INT8) with error compensation; offered in cuBLAS from CUDA 13.0 Update 2 (October 2025).}
\term{Ozaki scheme}{An error-free transformation that splits FP64 operands into slices representable in a narrow integer format, multiplies the slices on integer matrix units and recombines them exactly; the basis of NVIDIA's FP64 emulation.}
\term{MI430X}{AMD's 2027 HPC GPU: 432 GB HBM4 and a stated 288 TFLOP/s of hardware FP64; the processor of Discovery, Alice Recoque and LUMI-AI.}
\term{exascale system}{A machine with an HPL result of at least one exaflop (10 to the 18 FP64 FLOP/s); five existed on the June 2026 TOP500.}
\term{LineShine}{The Chinese system that led the June 2026 TOP500 at 2.198 EFLOP/s on CPUs alone, the first Chinese number one since 2017.}
\term{JUPITER}{Europe's first exascale system, at Forschungszentrum J\"ulich: about 24,000 Grace Hopper superchips, 1.000 EFLOP/s on HPL, fifth in June 2026.}
\term{Aurora}{Argonne National Laboratory's exascale system of 63,744 Ponte Vecchio GPUs, 1.012 EFLOP/s on HPL, fourth in June 2026 (slide 19).}
\end{terms}

\begin{intuition}
Follow the FP64 bars and you follow the customers: AMD's rise is Frontier and El Capitan, NVIDIA's fall is the inference market, Intel's single bar is Aurora. When a vendor cuts double precision it has decided who it is selling to; when it promises 288 TFLOP/s for 2027 it has decided to keep selling to the laboratories. Ask which rate, native or emulated, and measure your own code.
\end{intuition}

% =====================================================================
\setstage{attn}
\section{Intel: the Gen families, Xe and Gaudi (slides 13--16)}

\slide{13}{Intel's integrated GPUs: Gen9 to Gen12, the EU and the slice}{For a decade Intel's GPU shipped inside every laptop CPU; its EU and slice are what Xe grew from.}

\para{What the slide shows} The left half draws one Gen9 slice as Intel's 2015 whitepaper describes it: three subslices side by side, each a block of eight small orange squares (the execution units, EUs) above a bar labeled ``sampler; 64 KB SLM'', and under all three a long teal bar for the 768 KB L3 cache and the fixed-function geometry and raster units that the slice shares. A caption below the diagram gives the EU's internals: seven hardware threads, each with 128 registers of 32 bytes, two SIMD-4 floating-point units issuing every clock for 16 FP32 FLOP per clock per EU, and the product tiers built from slices (GT2 with one slice of 24 EUs and 384 FLOP per clock, GT4 with three slices and 72 EUs). A small table then follows the integrated line through four generations: Gen9 (Skylake, 2015, 14 nm, 24 EUs, about 0.4 TFLOP/s), Gen11 (Ice Lake, 2019, 10 nm, 64 EUs, about 1 TFLOP/s), Xe-LP in Tiger Lake (2020, 10 nm SuperFin, 96 EUs, 1{,}536 FLOP per clock, about 2.1 TFLOP/s) and the Xe-LP DG1 discrete card (Iris Xe MAX, 2020, 96 EUs, about 2.5 TFLOP/s). The right panel, ``What the integrated years built'', explains the EU, the hierarchy, the scale and the software that Intel carried into its data-center attempt.

\para{Details} From the i740 of 1998, Intel's only earlier discrete graphics card, until the Arc cards of 2022, Intel's GPUs lived inside its CPUs. That history matters for this lecture because the data-center GPU of 2023, Ponte Vecchio (slide 17), was not designed from a blank sheet: its compute unit, its memory hierarchy and its programming model descend from the integrated graphics in every Core laptop processor of the 2010s. Intel documented that lineage in a series of architecture whitepapers, of which Stephen Junkins's ``The Compute Architecture of Intel Processor Graphics Gen9'' (2015) is the clearest, and the vocabulary on this slide is his.

The unit of execution is the \emph{execution unit} (EU). An EU is a small in-order processor that holds seven hardware threads. Each thread owns a general register file of 128 registers of 32 bytes, so 4 KB per thread and 28 KB per EU, and a thread's registers hold its vectors: a 32-byte register is eight FP32 values, so a SIMD-8 operation reads and writes whole registers and a SIMD-16 operation uses pairs of them. The EU has two floating-point units, each four lanes wide, each able to complete a fused multiply-add on all four lanes every clock; together they deliver 16 FP32 FLOP per clock, and one of the two also executes the transcendental (``extended math'') instructions at a lower rate while integer operations run on both. The EU's thread arbiter can issue up to four instructions per clock from different threads to different units, which is how it keeps both FPUs busy from code that stalls on memory.

Here is the first difference from NVIDIA and AMD. On an SM a warp is 32 threads and on a CU a wavefront is 64 work-items, fixed by the hardware. On an EU a single hardware thread executes a SIMD instruction whose width the \emph{compiler} chooses: SIMD-8, SIMD-16 or SIMD-32, meaning that one EU thread carries 8, 16 or 32 work-items of the kernel. A SIMD-16 FP32 instruction on a four-wide FPU takes four clocks to retire, so a wider SIMD choice trades latency for fewer instructions and fewer threads. The compiler picks a width from register pressure: a kernel with many live values compiles to SIMD-8 so that 128 registers suffice for eight work-items, a lean kernel compiles to SIMD-32. In OpenCL and later SYCL this choice surfaces as the \emph{subgroup size}, which is why slide 14 lists Intel's thread group as ``8--32'' where NVIDIA's is 32 and AMD's 64.

Above the EU is the \emph{subslice}: eight EUs sharing an instruction cache, a thread dispatcher, a texture sampler and a data port, the load/store path to the memory hierarchy. The 64 KB of \emph{shared local memory} (SLM) per subslice is the counterpart of CUDA's shared memory and OpenCL's local memory: a software-managed scratchpad that a work-group's items use to exchange data, allocated in Gen9 from a region of the L3. Above the subslice is the \emph{slice}: three subslices (24 EUs) plus the fixed-function hardware of a graphics pipeline, the geometry and rasterization units, and 768 KB of L3 cache shared by the slice. A GPU is built from one, two or three slices, and Intel's product tiers name the count: GT2 is one slice (24 EUs, the HD Graphics 530 of a desktop Skylake), GT3 two (48 EUs) and GT4 three (72 EUs, the Iris Pro 580), the larger tiers adding 64 or 128 MB of eDRAM as a last-level cache on the package$^\dagger$.

The arithmetic of a Gen9 GT2 is $24 \text{ EUs} \times 16 \text{ FLOP/clk} = 384$ FLOP per clock; at the 1.15 GHz the part typically ran this is $384 \times 1.15\,\text{GHz} = 442$ GFLOP/s, the ``about 0.4 TFLOP/s'' in the table (derived). For comparison, the NVIDIA GTX 980 of the same year delivered about 4.6 TFLOP/s at 165 W, so the integrated GPU was an order of magnitude smaller, as its 15 W share of a laptop's power budget required.

Gen11, in the 10 nm Ice Lake processors of 2019, grew the slice to eight subslices and 64 EUs, with 3 MB of L3, and kept the EU: $64 \times 16 = 1{,}024$ FLOP per clock, about 1 TFLOP/s at the roughly 1.1 GHz the product ran, a figure Intel advertised as the first integrated GPU to pass a teraflop$^\dagger$. The whitepaper for Gen11 also describes tile-based rendering and coarse pixel shading, graphics features outside this lecture. Gen12, which Intel renamed \emph{Xe-LP} (low power) when it introduced the Xe brand in 2020, is the Tiger Lake processor's GPU and the first of the Xe family: 96 EUs in six \emph{dual subslices} of 16, on the 10 nm SuperFin process. Its EU was rebuilt, with a single eight-wide FP32/INT ALU and a two-wide extended-math unit in place of the two SIMD-4 FPUs, so an EU is now eight lanes and the GPU $96 \times 8 = 768$ lanes, 1{,}536 FLOP per clock, and about 2.07 TFLOP/s at 1.35 GHz (derived). Xe-LP left out FP64 hardware altogether, a decision Xe-HPC reversed (slide 18). The same Xe-LP design, at 1.65 GHz with 4 GB of LPDDR4X on a card, was sold in late 2020 as the DG1, the Iris Xe MAX, Intel's first discrete GPU in 22 years, at about 2.5 TFLOP/s (derived); it mattered as a proof that Intel could build and ship a discrete product, not as a competitor to anything.

Three things from these years carried into Xe and Ponte Vecchio. The hierarchy survived with new names (slide 14). The compile-time SIMD width survived as the subgroup size. And the software did: Intel had shipped OpenCL drivers for its integrated GPUs since 2011 and built a mature compiler for them; the Level Zero low-level driver interface and the oneAPI programming model (announced December 2018, specification 1.0 in September 2020) were designed to run on the integrated GPUs first$^\dagger$. The business argument Intel made for a data-center GPU in 2019 rested on exactly this: hundreds of millions of integrated GPUs already ran the architecture, so a scaled-up version would have drivers, compilers and an installed base from day one. The rest of this section is about why that argument did not carry the market, and what it did achieve.

\begin{terms}
\term{execution unit (EU)}{Intel's smallest processor in the Gen architectures: seven hardware threads, a 28 KB register file, two SIMD-4 FPUs (Gen9) delivering 16 FP32 FLOP per clock; renamed Xe Vector Engine in Xe.}
\term{hardware thread (Intel)}{One of the seven contexts an EU holds; each executes a SIMD instruction that carries 8, 16 or 32 work-items of the kernel.}
\term{general register file (GRF)}{The EU's registers: 128 registers of 32 bytes per thread on Gen9, 4 KB per thread.}
\term{SIMD width (compile-time)}{On Intel GPUs the number of work-items one instruction carries (8, 16 or 32) is chosen by the compiler per kernel, not fixed by the hardware as NVIDIA's warp of 32 and AMD's wavefront of 64 are.}
\term{subgroup}{The OpenCL and SYCL name for the set of work-items executing one SIMD instruction together; on Intel GPUs its size is the compile-time SIMD width.}
\term{subslice}{Eight EUs with a shared instruction cache, thread dispatcher, texture sampler and data port; the Gen counterpart of an SM or CU, renamed Xe-core in Xe.}
\term{shared local memory (SLM)}{Intel's name for the software-managed scratchpad a work-group shares, 64 KB per subslice on Gen9; the counterpart of CUDA shared memory and AMD's LDS.}
\term{slice}{Three subslices (24 EUs on Gen9) with the fixed-function geometry and raster units and 768 KB of L3; GPUs are built from one to three slices.}
\term{GT2, GT3, GT4}{Intel's integrated-graphics tiers by slice count: one, two or three slices (24, 48, 72 EUs on Gen9), the larger tiers with eDRAM on the package.}
\term{eDRAM}{Embedded DRAM, a 64 or 128 MB last-level cache on the package of Intel's higher integrated-graphics tiers (2013--2019).}
\term{Gen9, Gen11, Gen12}{Intel's integrated GPU generations in Skylake (2015), Ice Lake (2019) and Tiger Lake (2020); Gen12 was renamed Xe-LP.}
\term{Xe-LP}{The low-power member of the Xe family: 96 eight-wide EUs, 768 lanes, about 2 TFLOP/s, no FP64 hardware; the GPU of Tiger Lake and the DG1 card.}
\term{dual subslice}{Xe-LP's grouping of 16 EUs sharing thread control and a 128 KB L1/SLM, the step between the Gen subslice and the Xe-core.}
\term{DG1 (Iris Xe MAX)}{Intel's first discrete GPU since the i740 of 1998: a 96-EU Xe-LP on a card with 4 GB of LPDDR4X (2020).}
\term{Level Zero}{Intel's low-level driver interface for its accelerators (1.0 in 2020), the layer under oneAPI that corresponds to the CUDA driver API.}
\term{oneAPI}{Intel's programming model and library collection for CPUs, GPUs and other accelerators, announced in 2018 with SYCL-based Data Parallel C++ as its language; specification 1.0 in September 2020.}
\end{terms}

\begin{intuition}
Picture an Intel GPU of the 2010s as a stack of three boxes: an EU is a tiny seven-thread processor with two four-wide FPUs; a subslice is eight of them with a sampler and 64 KB of scratchpad; a slice is three subslices with an L3. Everything Intel built afterward, up to the 47-tile Ponte Vecchio, is this stack with wider units and new names. The one idea with no counterpart at NVIDIA or AMD is that the compiler, not the hardware, decides how many work-items an instruction carries.
\end{intuition}

% ---------------------------------------------------------------------
\slide{14}{From EU to Xe-core: the Xe hierarchy against SM and CU}{Xe (2020) renamed the units, doubled the vector width for HPC and added the XMX matrix engine.}

\para{What the slide shows} Two small block diagrams sit side by side at the top left. The left one is an Xe-HPC core from Ponte Vecchio: a row of eight purple blocks labeled ``8 XVE, 512 b: 16 FP32 and 16 FP64 FMA/clk each'', a row of eight orange blocks labeled ``8 XMX, 4096-bit systolic arrays'', and a bar for the ``512 KB L1 / shared local memory''. The right one is an Xe-HPG core from the Arc graphics cards: sixteen narrower orange blocks (``16 XVE, 256-bit: 8 FP32 lanes each''), sixteen XMX blocks (``16 XMX, 1024-bit; 1 ray-tracing unit'') and a ``192 KB L1 / SLM'' bar. Below them a table lines up the three vendors' hierarchies level by level: thread group (warp of 32; wavefront of 64; subgroup of 8--32), vector unit (an SM partition of 32 lanes on Hopper; a 16-lane SIMD that takes four cycles per wavefront; a 16-lane XVE), core (SM with 4 partitions and 4 tensor cores; CU with 4 SIMDs and 4 matrix cores; Xe-core with 8 XVE and 8 XMX), cluster (GPC of 18 SMs; XCD of 38--40 CUs; Xe-slice of 16 cores) and die (132 SMs on the H100, 148 across the B200's two dies; 304 CUs on 8 XCDs; a stack of 64 cores). The right panel explains the renaming, the two Xe variants, the XMX engine and what the table means across vendors.

\para{Details} When Intel introduced the Xe brand in 2020 it renamed every level of the Gen hierarchy and announced that one architecture would span four markets: Xe-LP for integrated and entry graphics, Xe-HPG (high-performance graphics) for discrete gaming cards, Xe-HP for data-center media and visualization (a product line that was cancelled before shipping, in 2021), and Xe-HPC for supercomputers. The names map as follows. The EU became the \emph{Xe Vector Engine} (XVE). The subslice became the \emph{Xe-core}, Intel's counterpart of the SM and the CU. The slice became a \emph{render slice} in Xe-HPG (four Xe-cores with four ray-tracing units and the graphics fixed-function hardware) or an \emph{Xe-slice} in Xe-HPC (sixteen Xe-cores, no graphics hardware). Xe-HPC adds a level above the slice, the \emph{stack}: a die of four Xe-slices (64 Xe-cores) with four HBM2e memory controllers, eight Xe Link ports and a large L2 cache; Ponte Vecchio is two stacks in one package (slide 17).

The two variants share the hierarchy but not the widths, and the slide's two diagrams show the difference. An Xe-HPG core has sixteen XVEs, each 256 bits wide, so each XVE executes eight FP32 lanes per clock and a core has $16 \times 8 = 128$ FP32 lanes. An Xe-HPC core has eight XVEs, each 512 bits wide (sixteen FP32 lanes), so also 128 FP32 lanes per core, from half as many, twice as wide engines. Intel's architecture material gives the Xe-HPC core's vector throughput as 256 FP32 operations per clock (128 fused multiply-adds) and, unusually, the same 256 operations per clock for FP64: Xe-HPC runs double precision at the full single-precision rate, where the H100 runs it at half (slide 18 does the arithmetic). FP16 runs at 512 operations per clock on the vector engines. Per XVE this is 16 FMAs per clock at FP32 and 16 at FP64; the 512-bit figure describes the operand width of one SIMD instruction, and the full-rate FP64 figure means the FP64 pipeline completes a SIMD-16 FP64 instruction at the same rate as a SIMD-16 FP32 one.

The second difference is the matrix engine. Intel's \emph{XMX} (Xe Matrix eXtensions) is a systolic array that performs a block of multiply-adds on a small matrix tile per clock, the counterpart of NVIDIA's tensor core and AMD's matrix core; it first appeared in Xe-HPG and Xe-HPC in 2022 and is not present in Xe-LP. An Xe-HPG core has sixteen XMX engines of 1{,}024 bits each; an Xe-HPC core has eight engines of 4{,}096 bits each, four times the width, so that a Ponte Vecchio core does 4{,}096 FP16 or BF16 operations per clock on its XMX units against 2{,}048 for an Arc core (derived from the vendor peaks: 839 TFLOP/s over 128 cores at 1.6 GHz; 138 TFLOP/s over 32 cores at 2.1 GHz). The XMX engines handle INT8 (at twice the FP16 rate), FP16, BF16 and TF32 (at half); they have no FP64 mode, which is why Ponte Vecchio's FP64 is a vector-unit figure only (slide 12), and no FP8, which arrived in the competing H100 and MI300X and is the main reason PVC's matrix peaks look a generation old (slide 18).

The third difference is memory per core. The Xe-HPC core has a 512 KB L1 data cache and shared local memory, configurable between the two uses; the Xe-HPG core has 192 KB. Against the H100 SM's 256 KB of combined L1 and shared memory and the MI300X CU's 64 KB LDS plus 32 KB L1, Ponte Vecchio's core has the largest local store of the three, which slide 18 counts as one of its strengths and the Aurora porting experience (slide 25) counts as one of its tuning problems: a kernel tiled for 48 or 96 KB of shared memory leaves most of 512 KB unused.

The comparison table is the heart of the slide, and it is worth reading row by row because it is the Rosetta stone for the rest of the lecture. \emph{Thread group}: a warp of 32 threads, a wavefront of 64 work-items, an Intel subgroup of 8, 16 or 32 work-items chosen by the compiler (slide 13). \emph{Vector unit}: on Hopper one SM partition has 32 FP32 lanes and executes a warp instruction in one clock; on CDNA a SIMD has 16 lanes and executes a 64-wide wavefront instruction over four clocks; on Xe-HPC an XVE has 16 lanes and executes a SIMD-16 instruction in one clock (or a SIMD-32 one in two). \emph{Core}: the SM has four partitions and four tensor cores; the CU four SIMDs and four matrix cores; the Xe-core eight XVEs and eight XMX engines. All three cores have 128 FP32 lanes on their 2022--23 parts, which is not a coincidence but a convergence: 128 lanes with a 256 KB-class register file is what fits in a core that can still be fed by one scheduler. \emph{Cluster}: the GPC of 18 SMs on Hopper, the XCD die of 38 or 40 CUs on MI300, the Xe-slice of 16 cores. \emph{Die}: 132 enabled SMs on H100 and 148 on a B200 die pair; 304 CUs across eight XCDs on MI300X; 64 cores per Ponte Vecchio stack, 128 in the package.

Two lessons follow. First, an Xe-core, an SM and a CU are the same kind of object, a group of about 128 lanes with matrix units and a local store, and a program written as work-groups over local memory maps to all three; the portability studies on slide 25 depend on this. Second, the places where the vendors differ, lane counts per instruction, thread-group width, matrix-unit width and local-store size, are exactly the parameters a kernel is tuned to, so the same code reaches peak on one vendor and not another until it is retuned. Intel's compile-time SIMD width adds a parameter the other two do not have, and Intel's compilers choose it per kernel.

\begin{terms}
\term{Xe}{Intel's GPU architecture brand from 2020, spanning Xe-LP (integrated), Xe-HPG (gaming), Xe-HP (cancelled) and Xe-HPC (supercomputing); Xe2 and Xe3 are its later generations.}
\term{Xe Vector Engine (XVE)}{The renamed EU: a 256-bit (Xe-HPG, 8 FP32 lanes) or 512-bit (Xe-HPC, 16 lanes) SIMD engine.}
\term{Xe-core}{The renamed subslice and Intel's counterpart of the SM and CU: 16 XVE and 16 XMX (Xe-HPG) or 8 XVE and 8 XMX (Xe-HPC), 128 FP32 lanes either way, with a 192 or 512 KB L1/SLM.}
\term{XMX (Xe Matrix eXtensions)}{Intel's systolic matrix engine for INT8, FP16, BF16 and TF32, the counterpart of tensor and matrix cores; 1,024 bits wide in Xe-HPG, 4,096 in Xe-HPC; no FP64 or FP8 mode.}
\term{render slice}{Xe-HPG's grouping of four Xe-cores with four ray-tracing units and the graphics fixed-function hardware.}
\term{Xe-slice}{Xe-HPC's grouping of sixteen Xe-cores.}
\term{stack (Xe-HPC)}{One Ponte Vecchio compute die: four Xe-slices (64 Xe-cores), four HBM2e controllers, eight Xe Links and a large L2; a Max 1550 has two.}
\term{Xe-HPG}{The gaming and workstation variant of Xe (Arc Alchemist, 2022): 256-bit XVEs, ray-tracing units, 1,024-bit XMX.}
\term{Xe-HPC}{The supercomputing variant of Xe (Ponte Vecchio): 512-bit XVEs, full-rate FP64, 4,096-bit XMX, 512 KB L1/SLM per core, no graphics hardware.}
\term{Xe-HP}{A planned data-center media and visualization Xe variant, cancelled in 2021 before shipping.}
\term{systolic array}{A grid of multiply-add units through which operands flow in lockstep so that a matrix tile is multiplied in a fixed number of clocks; the structure of XMX, of Gaudi's MME and, in a different arrangement, of tensor cores.}
\term{full-rate FP64}{Double-precision throughput equal to single-precision throughput; Xe-HPC and AMD's CDNA 2 offer it on their vector units, NVIDIA runs FP64 at half the FP32 rate.}
\end{terms}

\begin{intuition}
The three vendors arrived at the same core, about 128 lanes, four to eight matrix engines and a few hundred kilobytes of local store, and then named every level differently. Read the table as a dictionary: warp, wavefront, subgroup; SM, CU, Xe-core; GPC, XCD, Xe-slice. When a paper reports results on one vendor's hardware, translate the units before judging whether the result should carry over.
\end{intuition}

% ---------------------------------------------------------------------
\slide{15}{Xe-HPG, Xe2, Xe3: Arc Alchemist to Panther Lake}{Three discrete generations in four years; Intel's data-center restart is built from this line.}

\para{What the slide shows} A table at the left compares four products from Intel's graphics line: the Arc A770 (Xe-HPG, October 2022: ACM-G10 die on TSMC N6, 21.7 billion transistors, 406 mm$^2$, 32 Xe-cores, 4{,}096 FP32 lanes, 512 XMX engines, 17.2 TFLOP/s FP32, 16 GB of GDDR6 at 560 GB/s, 225 W), the Arc B580 (Xe2, December 2024: BMG-G21 on TSMC N5, 19.6 billion transistors, 272 mm$^2$, 20 Xe-cores, 2{,}560 lanes, 160 XMX, 14.6 TFLOP/s, 233 INT8 TOPS, 12 GB at 456 GB/s, 190 W), the Arc Pro B60 (Xe2, 2025: the same die with 24 GB at 456 GB/s, 197 TOPS, 120--200 W) and the Panther Lake integrated GPU (Xe3, January 2026, on Intel 18A: up to 12 Xe-cores, 1{,}536 lanes, 96 XMX, about 120 TOPS, LPDDR5X, 25 W). A caption adds the A770's 138 TFLOP/s FP16 XMX peak, the Xe2 change to SIMD16 engines with half as many XMX per core, Panther Lake's clocks and NPU, and Lunar Lake as the first Xe2 part. The right panel, ``Three generations, and the restart'', tells the story of Alchemist, Battlemage and Celestial and why a lecture on data-center GPUs includes them.

\para{Details} Intel's discrete graphics cards are in this lecture for one reason, stated in the subtitle: after the cancellation of Falcon Shores in January 2025 (slide 20), Intel's remaining data-center GPU plan is built from the graphics line, as NVIDIA's L40S is built from the Ada graphics die. To judge the GPU D-series of 2026 one has to know what Xe2 and Xe3 are.

\emph{Alchemist} (Xe-HPG), launched on October 12, 2022, was the first high-end discrete Intel GPU since 1998 and the first product with XMX matrix engines and ray-tracing units$^\dagger$. The large die, ACM-G10, is built on TSMC's N6 process, not Intel's own, with 21.7 billion transistors on 406 mm$^2$: eight render slices of four Xe-cores, 32 Xe-cores in all, 512 XVEs (16 per core) and 512 XMX engines, 32 ray-tracing units and a 16 MB L2$^\dagger$. The Arc A770 runs it at 2.1 GHz with 16 GB of GDDR6 on a 256-bit bus at 560 GB/s for 225 W and launched at \$329. Its FP32 peak is $32 \times 128 \times 2 \times 2.1\,\text{GHz} = 17.2$ TFLOP/s (derived, matching Intel's figure), and its XMX FP16 peak of 138 TFLOP/s implies 2{,}048 FP16 operations per clock per core (derived). As a graphics card it landed between NVIDIA's RTX 3060 and 3070 at a lower price; as a product it was known for drivers that took most of a year to reach competitive performance in older games, which is the recurring Intel story of this section in miniature: capable silicon, immature software.

\emph{Battlemage} (Xe2), launched in December 2024 with the Arc B580 at \$249, moved to TSMC N5 and changed the core$^\dagger$. The BMG-G21 die has 19.6 billion transistors on 272 mm$^2$ and 20 Xe2-cores. Each Xe2-core has eight XVEs of SIMD16 width instead of sixteen of SIMD8, so the lane count per core stays at 128 ($20 \times 128 = 2{,}560$ lanes) while the number of XMX engines per core halves to eight (160 in all), each engine doing more work per clock; the L2 grew to 18 MB. Intel quotes 14.6 TFLOP/s FP32 and 233 TOPS of INT8 on the XMX engines, with 12 GB of GDDR6 on a 192-bit bus at 456 GB/s and 190 W$^\dagger$. The move to SIMD16 brings Xe2 closer to the Xe-HPC core of Ponte Vecchio (eight 512-bit engines per core) and away from the Xe-HPG core, a sign that the graphics and compute lines were converging. The first Xe2 part was in fact the integrated GPU of Lunar Lake (September 2024)$^\dagger$.

The workstation cards of 2025 put the same die in front of a different buyer. The Arc Pro B50 (16 Xe2-cores, 128 XMX, 16 GB at 224 GB/s, 70 W, 170 TOPS, \$299) and Arc Pro B60 (20 Xe2-cores, 160 XMX, 24 GB at 456 GB/s, 120--200 W, 197 TOPS) were announced at Computex on May 19, 2025, and Intel's ``Project Battlematrix'' paired a Xeon W workstation with up to eight B60 cards, 192 GB of GDDR6 in all, as a local inference box for models that fit in that memory$^\dagger$. Memory capacity at a low price, rather than peak throughput, was the pitch, and it is the pitch the GPU D-series repeats at larger scale (slide 20).

\emph{Celestial} (Xe3) arrived first as the integrated GPU of Panther Lake, the Core Ultra Series 3 processors shown at CES on January 5, 2026, and the first high-volume product on Intel's 18A process$^\dagger$. The top configuration has 12 Xe3-cores at 2.5 GHz ($12 \times 128 = 1{,}536$ lanes, 96 XMX engines), which Intel rates at about 120 TOPS of INT8 on the GPU alongside about 50 TOPS from the processor's NPU, with LPDDR5X-9600 memory shared with the CPU$^\dagger$. The discrete derivative, Xe3P, is the architecture of the GPU D-series that Intel announced in October 2025 under the code name Crescent Island and showed as silicon in 2026; slide 20 gives its specification.

Why this line and not Xe-HPC? Three reasons appear in Intel's statements and in the parts themselves. Volume: a graphics architecture is built in tens of millions of units and its drivers and compilers are exercised by a large market, which is what Intel lacked for Ponte Vecchio. Cost: the Xe-HPC package of 47 tiles (slide 17) was expensive to build and to cool, while a graphics die on a standard process and a PCIe card is cheap. And the market: by 2025 the growth in data-center GPUs was inference, which needs INT8 and FP8 matrix throughput and memory capacity rather than FP64, and a graphics core with XMX engines has the first two. NVIDIA had made the same move in the other direction years earlier: the L40S (Ada) is a data-center card built from the RTX 4090's die, and the GPU D-series is Intel's L40S.

\begin{terms}
\term{Arc}{Intel's brand for discrete graphics cards from 2022: Alchemist (Xe-HPG), Battlemage (Xe2) and, in integrated form, Celestial (Xe3).}
\term{Alchemist}{The first Arc generation (October 2022): the ACM-G10 die on TSMC N6 with 32 Xe-cores, XMX engines and ray-tracing units; the Arc A770 is its largest card.}
\term{ACM-G10}{Alchemist's large die: 21.7 billion transistors, 406 square millimeters, TSMC N6, 32 Xe-cores.}
\term{Battlemage}{The second Arc generation (December 2024), the Xe2 architecture: the BMG-G21 die on TSMC N5 with 20 Xe2-cores of eight SIMD16 engines each.}
\term{Xe2}{Intel's second Xe generation: eight SIMD16 XVEs and eight XMX engines per core (128 lanes, as before), first in Lunar Lake (2024) and the Arc B580.}
\term{Arc Pro B50 and B60}{Workstation cards on the Battlemage die (2025) with 16 or 24 GB of GDDR6, sold on memory capacity for local inference.}
\term{Project Battlematrix}{Intel's 2025 workstation configuration of a Xeon W with up to eight Arc Pro B60 cards, 192 GB of GDDR6, for running models locally.}
\term{Celestial (Xe3)}{Intel's third Xe generation, first as the integrated GPU of Panther Lake (2026) on Intel 18A; Xe3P is its discrete derivative.}
\term{Panther Lake}{Intel's Core Ultra Series 3 processors (CES 2026), the first high-volume product on Intel 18A, with up to 12 Xe3-cores.}
\term{Intel 18A}{Intel's 2 nm-class process node with gate-all-around transistors and backside power delivery; Panther Lake is its first volume product.}
\term{NPU}{Neural processing unit: a low-power matrix engine for inference inside a client processor, separate from the GPU; Panther Lake's is rated at about 50 TOPS.}
\term{TOPS}{Tera-operations per second, the unit vendors use for INT8 matrix throughput; counts a multiply-add as two operations.}
\term{L40S}{NVIDIA's data-center card built from the Ada graphics die (2023); the pattern Intel's GPU D-series follows with Xe3P.}
\end{terms}

\begin{intuition}
Alchemist, Battlemage and Celestial look like a story about gaming cards, but they are the only Intel GPU architecture still shipping in volume, and volume is what pays for compilers and drivers. When Intel's data-center GPU plan collapsed in 2025, the graphics core was what remained to build on, exactly as NVIDIA's data-center inference cards are built from its graphics dies.
\end{intuition}

% ---------------------------------------------------------------------
\slide{16}{Gaudi 2 and Gaudi 3: Intel's other accelerator}{The Habana architecture: matrix engines, tensor cores and Ethernet on the die; strong on paper, weak in the market.}

\para{What the slide shows} A block diagram of the Gaudi 3 package fills the top left: two compute dies side by side, each carrying four pink ``MME'' blocks (matrix multiplication engines), a field of 32 small orange squares labeled ``32 TPCs, 48 MB SRAM'' (tensor processor cores) and a teal bar for ``12 $\times$ 200 GbE RoCE'' ports, with four HBM2e stacks drawn along each outer edge. The caption defines the MME as a large systolic array and the TPC as a VLIW SIMD engine for everything else, and notes that the chip's scale-up and scale-out links are both standard Ethernet from the die, with no proprietary fabric. Two stat tiles read 1{,}835 TFLOP/s (BF16 and FP8, dense) and 128 GB (HBM2e, 3.67 TB/s). The right panel, ``The line and its outcome'', gives the Habana acquisition, Gaudi 2, Gaudi 3 and the market result.

\para{Details} Intel's second data-center accelerator line did not come from the Gen or Xe families at all. Habana Labs, an Israeli startup founded in 2016, had shipped the Goya inference chip in 2018 and announced the Gaudi training chip in 2019; Intel bought the company in December 2019 for about \$2 billion$^\dagger$. For the next six years Intel sold two AI accelerator lines with different instruction sets, different software stacks and different arguments: Ponte Vecchio and its planned successors for HPC and AI under oneAPI, and Gaudi for deep-learning training and inference under Habana's own SynapseAI stack. The slide's last line, that a second ISA with its own software divided Intel's effort, is the summary; the details show what the Gaudi design was and what it achieved.

The architecture has three kinds of unit. The \emph{matrix multiplication engine} (MME) is a large systolic array that performs the matrix multiplies of a neural network, the convolutions and the fully connected and attention layers, in BF16 or FP8 with FP32 accumulation; Gaudi 2 has two and Gaudi 3 eight. The \emph{tensor processor core} (TPC) is a VLIW (very long instruction word) processor with wide SIMD vector units and its own local memory, programmable in C, that executes everything the MMEs do not: activations, normalizations, element-wise operations, reductions and data movement. Gaudi 2 has 24 TPCs and Gaudi 3 64. The third kind of unit is the network: RDMA-over-Converged-Ethernet (RoCE v2) engines integrated on the die, 24 ports of 100 Gb/s on Gaudi 2 and 24 of 200 Gb/s on Gaudi 3. This is the distinctive choice. NVIDIA scales up over NVLink and NVSwitch, AMD over Infinity Fabric and later UALink (slide 21); Gaudi uses the same standard Ethernet for both the links between the eight accelerators on a baseboard and the links out of the server, so that a Gaudi cluster is built from merchant Ethernet switches. On the eight-accelerator baseboard, 21 of each chip's 24 ports connect it to the other seven (three links per pair, an all-to-all) and 3 ports leave the chassis for scale-out$^\dagger$.

\emph{Gaudi 2} (May 10, 2022) is built on TSMC's 7 nm process with 24 TPCs, two MMEs, 48 MB of on-die SRAM, 96 GB of HBM2e at 2.45 TB/s, 24 ports of 100 Gb Ethernet and a 600 W OAM module$^\dagger$. In 2023 it was the only merchant accelerator from a vendor other than NVIDIA with competitive MLPerf Training results (Google's TPUs and Graphcore's IPUs had submitted earlier, but the TPU is not sold as a chip), and it is the chip Intel's early Gaudi customers deployed. \emph{Gaudi 3}, announced in April 2024 and presented by Roman Kaplan at Hot Chips in August 2024, moves to a 5 nm process and a two-die package: 64 TPCs and 8 MMEs (32 and 4 per die), 96 MB of SRAM (48 per die), 128 GB of HBM2e in eight stacks at 3.67 TB/s, 24 ports of 200 Gb Ethernet, and a rated power of 900 W on an air-cooled OAM module or 1{,}200 W liquid-cooled$^\dagger$. Intel's quoted matrix peak is 1{,}835 TFLOP/s, one figure for both BF16 and FP8, dense: the MME runs FP8 at the same rate as BF16 rather than at twice, which is the opposite of the NVIDIA and AMD choice and reflects a design in which the systolic array is sized for BF16 training$^\dagger$. If each MME is a 256 $\times$ 256 array of multiply-adds, as the Hot Chips presentation describes, then $8 \times 65{,}536 \times 2\,\text{FLOP} = 1{,}048{,}576$ FLOP per clock and the quoted peak corresponds to a clock of 1.75 GHz (derived). There is no FP64 at all and no FP4.

How does Gaudi 3 compare with its contemporaries? Against the H100 (slide 18) its FP8 peak is 1{,}835 against 1{,}979 TFLOP/s dense, its memory 128 GB against 80 at 3.67 against 3.35 TB/s, and its scale-up links 24 $\times$ 200 Gb/s $=$ 600 GB/s of aggregate Ethernet per chip against 900 GB/s of NVLink (the figure plotted on slide 21). On paper it was an H100-class training part with more memory, and Intel priced it below the H100. Against the MI300X (slide 9) it had less memory (128 against 192 GB) and less FP8 (1{,}835 against 2{,}615), and against the Blackwell parts shipping from late 2024 it had no FP4 and a third or less of the matrix throughput.

The market result is the subtitle's second half. Intel set a target of \$500 million of Gaudi revenue for 2024, modest against NVIDIA's data-center revenue of tens of billions per quarter, and on October 31, 2024, Pat Gelsinger said on the third-quarter call that Intel would not reach it, citing software readiness and the transition from Gaudi 2 to Gaudi 3$^\dagger$. Analysts put Gaudi's share of the accelerator market below one percent; Intel's roadmap statements of late 2024 and 2025 folded Gaudi's successor into the Falcon Shores and then Jaguar Shores plans (slide 20); and Intel's third-quarter 2026 results mention only the GPU D-series among accelerators$^\dagger$. MLPerf tells the same story: Gaudi 2 and 3 appear in the 2023--24 inference and training rounds and in none since (slide 30).

Why did a competent design fail commercially? The slide's own explanation is software. Gaudi ran SynapseAI, a graph compiler and runtime of Habana's own, with PyTorch support through a plug-in; CUDA code did not port to it, oneAPI did not target it, and the kernels that defined performance in 2023--24 (FlashAttention, the vLLM serving engine, quantization libraries) arrived on Gaudi late or not at all. A buyer of a second-source accelerator in those years could choose AMD's MI300X, which ran ported CUDA code through HIP with ROCm's growing library set (slide 24), or Gaudi, which ran a different stack; most who left NVIDIA chose AMD. Two further reasons sit in the hardware: the Ethernet scale-up fabric, though standard and cheap, carried higher latency than NVLink or Infinity Fabric for the collective operations of training, and the absence of FP8 at double rate and of any FP4 left Gaudi 3 behind on the inference workloads that grew fastest in 2025. The lesson the slide draws is organizational as much as technical: two accelerator lines with two software stacks meant that neither received the whole of Intel's effort.

\begin{terms}
\term{Habana Labs}{The Israeli accelerator startup (founded 2016) that designed Goya and Gaudi; bought by Intel in December 2019 for about two billion dollars.}
\term{Gaudi}{Habana's training-and-inference accelerator family: Gaudi (2019), Gaudi 2 (2022) and Gaudi 3 (2024), built around matrix engines, tensor processor cores and on-die Ethernet.}
\term{matrix multiplication engine (MME)}{Gaudi's large systolic array for matrix multiplies in BF16 or FP8 with FP32 accumulation; two on Gaudi 2, eight on Gaudi 3.}
\term{tensor processor core (TPC)}{Gaudi's programmable VLIW SIMD processor with local memory that executes the non-matrix operations of a network; 24 on Gaudi 2, 64 on Gaudi 3.}
\term{VLIW}{Very long instruction word: an instruction format that packs several operations for different units into one word, scheduled by the compiler rather than by hardware.}
\term{RoCE}{RDMA over Converged Ethernet: a protocol for direct memory-to-memory transfers over Ethernet, used by Gaudi for both scale-up and scale-out and by Ultra Ethernet's predecessors.}
\term{SynapseAI}{Habana's software stack for Gaudi, a graph compiler and runtime with PyTorch support; separate from oneAPI.}
\term{Gaudi 3}{Intel's 2024 accelerator: two 5 nm dies, 64 TPCs, 8 MMEs, 96 MB SRAM, 128 GB HBM2e at 3.67 TB/s, 1,835 TFLOP/s BF16 or FP8, 24 ports of 200 Gb Ethernet, 900 to 1,200 W.}
\term{OAM}{Open Compute Project Accelerator Module: the open form factor for accelerator modules used by AMD, Intel and others, the counterpart of NVIDIA's SXM (slide 23).}
\term{second-source accelerator}{A buyer's alternative to the dominant vendor's part; in 2023--25 AMD's Instinct line took that role and Gaudi did not.}
\end{terms}

\begin{intuition}
Gaudi 3 was an H100-class training chip with more memory, standard Ethernet instead of a proprietary fabric, and a lower price, and it failed to sell. The chip was not the problem; the stack was. A buyer leaving CUDA could choose a GPU that ran ported CUDA code (AMD) or an accelerator with its own software (Gaudi), and nearly all chose the former. The deeper point for the whole lecture: the second place in this market goes to whoever makes leaving CUDA cheapest.
\end{intuition}

% =====================================================================
\setstage{mlp}
\section{Xe-HPC: Ponte Vecchio, Aurora and after (slides 17--20)}

\slide{17}{Ponte Vecchio: 47 active tiles from five process nodes}{Intel's exascale GPU (announced 2021, launched 2023) is the most complex package any vendor has shipped.}

\para{What the slide shows} A plan view of the Data Center GPU Max 1550 package fills the top left. Two large pale rectangles are the base tiles (``base tile 0, Intel 7, 646 mm$^2$'' and ``base tile 1''), each carrying eight purple compute tiles in two rows, four orange RAMBO cache tiles below them, and a teal ``Xe Link tile (TSMC N7): 8 links'' bar at the bottom; a thin orange strip between the two base tiles is labeled EMIB, and four HBM2e stacks stand along each outer edge. The caption counts the tiles: per stack, 8 compute tiles on TSMC N5, 4 RAMBO cache tiles on Intel 7, 1 Xe Link tile and 4 HBM2e stacks; 11 EMIB bridges; 47 active tiles in all, over 100 billion transistors and 2{,}330 mm$^2$ of active silicon. Two stat tiles read 128 (Xe-cores; 1{,}024 XVE, 1{,}024 XMX) and 3.28 TB/s (128 GB HBM2e, 8 stacks). The right panel, ``What the package is'', explains Foveros and EMIB, the cache, the dates and the SKUs.

\para{Details} Ponte Vecchio, sold as the Intel Data Center GPU Max series, is the GPU of the Aurora supercomputer (slide 19) and the only Xe-HPC product Intel shipped. It is in this lecture for two reasons. It is the one Intel entry in the exascale-GPU contest of 2022--23 against the H100 and the MI250X (slide 18), and it is the most aggressive piece of packaging any GPU vendor has sold: where NVIDIA's Blackwell of 2024 joins two dies and AMD's MI300X of 2023 stacks eight compute dies on four I/O dies (slide 9), Ponte Vecchio in 2023 assembled 47 active tiles from five process technologies into one 4{,}844 mm$^2$ package$^\dagger$.

Count the tiles with the diagram. The package has two halves, each a \emph{stack} in Intel's vocabulary (slide 14). A stack is built on one \emph{base tile}, a 646 mm$^2$ die on Intel's own Intel 7 process (the renamed 10 nm Enhanced SuperFin) that carries the memory controllers, the PCIe and inter-stack links, and a large part of the L2 cache$^\dagger$. On top of the base tile, bonded face to face by Intel's \emph{Foveros} 3D stacking, sit eight \emph{compute tiles}, each of eight Xe-cores, made on TSMC's N5 process, and four \emph{RAMBO} tiles (``random access memory, bandwidth optimized'') of SRAM cache on Intel 7. Beside the base tile sits one \emph{Xe Link} tile on TSMC's N7 process with eight serial link ports, and four HBM2e stacks. That is $8 + 4 + 1 + 1 + 4 = 18$ tiles per stack, 36 for two, plus 11 \emph{EMIB} bridges (embedded multi-die interconnect bridges, small silicon dies set into the package substrate that carry dense wiring between the two base tiles and from each base tile to its HBM and Xe Link neighbors): 47 active tiles$^\dagger$. Sixteen passive thermal tiles on top of the stack, there to give the lid a flat surface to cool, bring the count that Intel sometimes quoted to 63$^\dagger$. The five process technologies are TSMC N5 (compute), TSMC N7 (Xe Link), Intel 7 (base and RAMBO), the DRAM process of the HBM2e and the process of the EMIB bridges. The transistor count exceeds 100 billion; the active silicon area is 2{,}330 mm$^2$, nearly three reticle-limit dies' worth, on a package of 4{,}844 mm$^2$$^\dagger$.

Two packaging technologies make this possible and are worth separating. \emph{Foveros} is 3D: a die is thinned and bonded face down onto a base die with a fine pitch of micro-bumps, so that every compute tile talks to the base tile through thousands of short vertical connections, as AMD's XCDs do to their IODs on MI300X. \emph{EMIB} is 2.5D: two dies sit side by side on the substrate and a small bridge die beneath their adjoining edges carries the wires between them, giving interposer-class wiring density without a full silicon interposer under the whole package (NVIDIA's CoWoS puts the entire GPU and its HBM on one interposer). Ponte Vecchio uses Foveros within a stack and EMIB between the stacks and to the memory; the industry's current parts use one or the other (MI300X: hybrid bonding on an interposer; Blackwell: two dies on CoWoS-L with local bridges), and Intel's claim to have led in packaging rests on having shipped both at once, two years before Blackwell's two dies and in the same year as MI300X.

The compute is the sum of the tiles: 16 compute tiles of 8 Xe-cores make 128 Xe-cores, each with 8 XVEs and 8 XMX engines, so 1{,}024 of each, plus 128 ray-tracing units that the data-center product keeps but does not advertise$^\dagger$. The memory is 128 GB of HBM2e in eight stacks at 3{,}276.8 GB/s, the slide's 3.28 TB/s$^\dagger$. The cache is the package's most-quoted and least-consistent figure: Intel's specification gives 408 MB of L2, spread across the base tiles and the RAMBO tiles, against 50 MB on the H100; the split between the two kinds of tile is reported differently in the ISSCC paper, the 2021 Architecture Day material and the ALCF documentation, so the slide gives only the total$^\dagger$. In addition there are 64 MB of L1 (512 KB per Xe-core, slide 14) and 64 MB of register file across the package.

The dates tell the project's history. Intel described the architecture at its Architecture Day on August 19, 2021; Wilfred Gomes and colleagues presented ``Ponte Vecchio: a multi-tile 3D stacked processor for exascale computing'' at ISSCC in February 2022; Hong Jiang presented it at Hot Chips 34 in August 2022; Intel announced the Data Center GPU Max series on November 9, 2022, and launched it on January 10, 2023, alongside the Sapphire Rapids Xeons$^\dagger$. Against the plan announced in 2019, which had Aurora delivered in 2021, this was about a year late, and the slip was mostly the 2020 delay of Intel's 7 nm process, which pushed the base tiles to Intel 7 and the compute tiles to TSMC (slide 19).

Three products were announced. The Max 1550 is the full part: 128 Xe-cores, 128 GB, 600 W, in the OAM form factor for liquid-cooled servers, the Aurora GPU. The Max 1100 is a PCIe card with 56 Xe-cores, 48 GB of HBM2e and 300 W, the part most Max-series customers outside Aurora bought and the one Reguly's portability study used (slide 25)$^\dagger$. A Max 1350 (112 cores, 96 GB, 450 W) was announced in November 2022 and dropped in 2023 in favor of a planned Max 1450 that was reported to be an export-compliant variant; neither reached volume$^\dagger$.

\begin{terms}
\term{Ponte Vecchio}{Intel's Xe-HPC GPU (2021--23), sold as the Data Center GPU Max series: 47 active tiles, 128 Xe-cores, 128 GB HBM2e, 600 W; the GPU of Aurora.}
\term{Data Center GPU Max 1550}{The full Ponte Vecchio product: 128 Xe-cores, 128 GB of HBM2e at 3.28 TB/s, 408 MB L2, 600 W OAM module.}
\term{Data Center GPU Max 1100}{The PCIe Ponte Vecchio: 56 Xe-cores, 48 GB HBM2e, 300 W.}
\term{tile}{Intel's word for a die in a multi-die package; Ponte Vecchio has compute, RAMBO, base, Xe Link and HBM tiles.}
\term{base tile}{Ponte Vecchio's 646 square millimeter Intel 7 die carrying memory controllers, links and L2, on which compute and cache tiles are stacked.}
\term{compute tile}{A Ponte Vecchio die of eight Xe-cores on TSMC N5; eight per stack, sixteen per package.}
\term{RAMBO tile}{``Random access memory, bandwidth optimized'': a Ponte Vecchio SRAM cache tile on Intel 7, four per stack, part of the 408 MB L2.}
\term{Xe Link}{Intel's GPU-to-GPU serial link, eight ports per Xe Link tile, sixteen per Ponte Vecchio, 53 Gb/s per lane; joins the six GPUs of an Aurora node.}
\term{Foveros}{Intel's 3D die-stacking technology: a thinned die bonded face down onto a base die through fine-pitch micro-bumps.}
\term{EMIB}{Embedded multi-die interconnect bridge: a small silicon die set into the package substrate carrying dense wiring between the edges of two adjacent dies (2.5D).}
\term{CoWoS}{TSMC's chip-on-wafer-on-substrate packaging, in which dies and HBM sit on a silicon interposer; used by NVIDIA and AMD.}
\term{Intel 7}{Intel's renamed 10 nm Enhanced SuperFin process, used for Ponte Vecchio's base and RAMBO tiles and for Sapphire Rapids.}
\term{thermal tile}{A passive die placed on top of a stack to give the lid a flat surface; sixteen on Ponte Vecchio, bringing the quoted total to 63 tiles.}
\term{Architecture Day}{Intel's periodic technical disclosure event; the August 2021 edition described Ponte Vecchio and Xe-HPC.}
\end{terms}

\begin{intuition}
Ponte Vecchio is three reticles' worth of silicon from two foundries, stacked and bridged into one 600 W package in 2023, a year before Blackwell's two dies. It proved that a GPU could be assembled from many small dies on mixed processes; it also showed the cost, in yield, power and schedule, of being first. When the deck says Intel led in packaging, this is the exhibit.
\end{intuition}

% ---------------------------------------------------------------------
\slide{18}{The Xe-HPC core; the Max 1550 against its contemporaries}{Full-rate FP64, a 512 KB local store per core, a matrix engine a generation behind.}

\para{What the slide shows} A table down the left compares the three compute GPUs of 2022--23 row by row: Max 1550 (shipped January 2023; N5 $+$ Intel 7 $+$ N7; over 100 billion transistors; 128 Xe-cores; 16{,}384 FP32 lanes; 1.6 GHz; FP64 vector 52.4 TFLOP/s and no separate matrix FP64; FP32 52.4; TF32 matrix 419; FP16/BF16 matrix 839; INT8 1{,}678 TOPS; 408 MB L2; 128 GB HBM2e at 3.28 TB/s; 16 Xe Links; 600 W), H100 SXM (Q4 2022; TSMC 4N; 80 billion; 132 SMs; 16{,}896 lanes; 1.98 GHz boost; FP64 34 vector / 67 matrix; FP32 67; TF32 494; FP16 989; INT8 1{,}979; 50 MB L2; 80 GB HBM3 at 3.35 TB/s; NVLink 4 at 900 GB/s; 700 W) and MI250X (November 2021; TSMC N6; 58.2 billion; 220 CUs; 14{,}080 lanes; 1.7 GHz; FP64 47.9 / 95.7; FP32 47.9; FP16 383; INT8 383; $2 \times 8$ MB L2; 128 GB HBM2e at 3.2 TB/s; 8 Infinity Fabric links at 800 GB/s; 560 W). A caption notes that the figures are dense, that the Max 1550's FP16 and INT8 are Intel's figures and follow from its per-core table (4{,}096 FP16 and 8{,}192 INT8 operations per clock per Xe-core), that the H100's tensor peaks imply a clock of about 1.83 GHz against the 1.98 GHz its FP32 figure implies, that Xe Link runs 53 Gb/s per lane with the lanes per link unpublished, and that the MI250X's two dies appear as two GPUs. The right panel, ``The core, and the arithmetic'', derives the lane, FP64 and XMX peaks and gives the verdict of 2023.

\para{Details} Start with the lanes, because every peak on the slide follows from them. A Max 1550 has 128 Xe-cores, each with 8 XVEs of 16 FP32 lanes (slide 14): $128 \times 8 \times 16 = 16{,}384$ lanes. At the product's maximum clock of 1.6 GHz and two FLOP per fused multiply-add, the FP32 peak is $16{,}384 \times 2 \times 1.6\,\text{GHz} = 52.4$ TFLOP/s, Intel's figure$^\dagger$. The H100 SXM has $132 \times 128 = 16{,}896$ lanes, three percent more, and reaches 67 TFLOP/s because its clock is higher, about 1.98 GHz for the FP32 figure; the MI250X has $220 \times 64 = 14{,}080$ lanes at 1.7 GHz for 47.9 TFLOP/s. On the vector units, then, the three parts are within a factor of 1.4 of one another, and the differences are clock and lane count, not architecture.

Double precision is where Ponte Vecchio stands out. Intel's Xe-HPC core executes 256 FP64 operations per clock, the same as its FP32 rate, so the FP64 peak is also 52.4 TFLOP/s: every XVE completes 16 double-precision fused multiply-adds per clock, as many as single-precision ones (derived from Intel's per-core figures). The H100's SM has 64 FP64 lanes against 128 FP32 lanes, so its vector FP64 is half its FP32, 34 TFLOP/s; its tensor cores double that to 67 for matrix multiplication. The MI250X runs FP64 at the full FP32 rate on its vector units (47.9) and doubles it on the matrix cores (95.7). So in January 2023 a Max 1550 had the highest vector FP64 of any shipping GPU, 52.4 against 47.9 and 34, and the lowest matrix FP64, because the XMX engines have no FP64 mode and the vector figure is all there is. For a code whose double-precision work is dense matrix multiplication, which is what HPL measures, the MI250X's 95.7 and the H100's 67 beat Ponte Vecchio's 52.4; for a code whose FP64 is vector arithmetic in stencils, sparse solvers or particle pushes, Ponte Vecchio was the fastest GPU of its year. This is the sense in which Aurora's design, with its emphasis on FP64 vector throughput and bandwidth, was a bet on classical simulation at the moment the market turned to low-precision AI.

The matrix engines tell the other half. Intel's confirmed XMX figure is 419 TFLOP/s of TF32: $1{,}024$ engines $\times 128$ FMA per clock $\times 2 \times 1.6\,\text{GHz} = 419$ TFLOP/s (derived), equivalently 2{,}048 TF32 operations per clock per Xe-core. BF16 and FP16 run at twice that, 839 TFLOP/s (4{,}096 operations per clock per core), and INT8 at twice again, 1{,}678 TOPS; these two figures are Intel's and follow directly from the per-Xe-core throughput table in Intel's architecture material (4{,}096 FP16 and 8{,}192 INT8 operations per clock), which is why the deck quotes them as vendor figures rather than as independently measured ones$^\dagger$. Against the H100's 494 TF32, 989 FP16 and 1{,}979 INT8 (dense), Ponte Vecchio's matrix peaks are 85\% of the H100's at each of these precisions, a respectable showing for a chip designed in 2019. The gap is FP8. Hopper's tensor cores run FP8 at twice the FP16 rate, 1{,}979 TFLOP/s; XMX has no FP8 mode, so an FP8 workload has to run on Ponte Vecchio in FP16 at 839 TFLOP/s, 42\% of the H100's FP8 rate (at INT8 the two are close, 1{,}678 against 1{,}979 TOPS), and it cannot run the FP8 training and inference recipes that became standard in 2023--24 at all. That is what the subtitle means by ``a matrix engine a generation behind'': XMX matches the H100's older precisions and lacks its new one.

Memory is a strength. The 512 KB of L1 and shared local memory per Xe-core is twice the H100 SM's 256 KB and eight times the MI250X CU's 64 KB LDS; the 408 MB of L2 is eight times the H100's 50 MB and 25 times the MI250X's two 8 MB caches; the 128 GB of HBM2e matches the MI250X's and is 1.6 times the H100's 80 GB; and the 3.28 TB/s of bandwidth is within 2\% of the H100's 3.35 TB/s and above the MI250X's 3.2$^\dagger$. A ridge-point calculation (peak over bandwidth, slide 15 of the previous lecture) gives $52.4 / 3.28 = 16$ FLOP per byte for FP64 on Ponte Vecchio against $34 / 3.35 = 10$ on the H100, so the Intel part can sustain more arithmetic per byte fetched before it becomes compute-bound (derived); in practice the large caches matter more, since a working set that fits in 408 MB never touches HBM.

The links are a weakness. Xe Link provides sixteen ports per GPU (eight per Xe Link tile) with serializers at 53 Gb/s per lane, and Intel did not publish the lanes per port or the aggregate bandwidth, which is why the slide lists ``16 Xe Links'' where it gives 900 GB/s for NVLink 4 and 800 GB/s for the MI250X's eight Infinity Fabric links$^\dagger$. In Aurora the links join the six GPUs of a node (slide 19); there is no Xe Link switch, so the domain never grew past a node (slide 21). Power: 600 W on the liquid-cooled OAM module, against 700 W for the H100 SXM and 560 W for the MI250X$^\dagger$.

A caveat on the MI250X column: the module holds two graphics compute dies that the software sees as two separate GPUs with their own 64 GB each (slide 8), so the ``per GPU'' figures in the table are per module, and a program must use both dies to reach them; the H100 and Max 1550, despite the latter's two stacks, present themselves as one device by default (the Max 1550 can also be exposed as two, one per stack, a mode Aurora's users exploit, slide 25).

The verdict of 2023, as the panel states it: competitive FP64 and bandwidth, a generation behind on low-precision matrix work, 600 W, a year late. In the market that bought GPUs in 2023, the one that bought H100s for large language models, the last three items decided; in the market that bought Aurora, the first two did. The next slide is about that machine.

\begin{terms}
\term{vector FP64 and matrix FP64}{Double-precision throughput on the SIMD lanes (vector) and on the matrix units (matrix); Ponte Vecchio has only the former, the H100 and MI250X double the former on the latter.}
\term{TF32}{NVIDIA's 19-bit tensor format (8-bit exponent, 10-bit mantissa) used for FP32-range matrix multiplication at reduced precision; XMX supports it at half the FP16 rate.}
\term{ridge point}{Peak FLOP/s divided by memory bandwidth, the arithmetic intensity (FLOP per byte) above which a kernel is compute-bound rather than memory-bound.}
\term{Infinity Fabric link}{AMD's GPU-to-GPU and CPU-to-GPU interconnect (xGMI); the MI250X has eight links of 100 GB/s.}
\term{graphics compute die (GCD)}{One of the MI250X's two compute dies, each presented to software as a separate GPU with its own 64 GB of HBM2e.}
\term{OAM module}{An accelerator in the Open Compute Accelerator Module form factor, as the Max 1550, MI250X and MI300X are; liquid or air cooled on a universal baseboard.}
\term{two-stack mode}{Exposing each of a Max 1550's two stacks as a separate device, so that a six-GPU Aurora node appears as twelve; used to improve locality of memory access.}
\end{terms}

\begin{intuition}
Line up the three 2022--23 parts and the Max 1550 wins the rows HPC cares about (vector FP64, caches, memory capacity) and loses the rows AI cares about (FP8, matrix FP64, links). The chip was designed for a 2019 view of what a supercomputer would run. Whether that was a mistake depends on which buyer you ask, and the two buyers' answers are slides 19 and 20.
\end{intuition}

% ---------------------------------------------------------------------
\slide{19}{Aurora: 63,744 GPUs, 1.012 exaflops, 38.7 MW}{Planned in 2015 for 2018; delivered in 2023 with Ponte Vecchio; exascale HPL on Intel GPUs in 2024.}

\para{What the slide shows} A diagram of one Aurora node occupies the top left: two gray boxes for the Xeon Max 9470C processors (52 cores, 64 GB of HBM2e each) joined by a line, each connected downward to three of six purple ``Max 1550'' boxes that are themselves chained by Xe Link; a caption states that Xe Link joins the six GPUs all-to-all, that PCIe 5 connects them to the CPUs with 512 GB of DDR5 per CPU, that each node has eight Slingshot-11 NICs at 200 Gb/s, and that 10{,}624 such nodes in 166 HPE Cray EX racks hold 63{,}744 GPUs. Below it a table follows Aurora's TOP500 entries: November 2023 (5{,}439 nodes in the run, 0.585 EFLOP/s, rank 2, 24.7 MW), June 2024 (9{,}234 nodes, 1.012 EFLOP/s, rank 2, HPL-MxP 10.6 EFLOP/s, 38.7 MW), November 2024 through 2025 (the same run, rank 3, HPL-MxP 11.6) and June 2026 (rank 4, HPL-MxP third). The right panel, ``The machine, its history, and its science'', gives the peak and efficiency, the eight-year history, the 2025 Gordon Bell finalists and the other Ponte Vecchio systems.

\para{Details} Aurora is the Argonne Leadership Computing Facility's exascale system and the only one of the five exascale machines on the June 2026 TOP500 built on Intel GPUs. Its node is the slide's diagram. Two Xeon CPU Max 9470C processors (Sapphire Rapids with HBM: 52 cores each and 64 GB of HBM2e on the CPU package, plus 512 GB of DDR5 per socket) connect over PCIe 5 to six Data Center GPU Max 1550s, three per CPU; the six GPUs are joined to one another by Xe Link in an all-to-all; and eight HPE Slingshot-11 network interfaces at 200 Gb/s each give the node 1.6 Tb/s into the system's dragonfly network$^\dagger$. Across the machine there are 10{,}624 nodes in 166 HPE Cray EX cabinets, 21{,}248 CPUs and 63{,}744 GPUs, with 20.4 petabytes of memory counting the GPUs' HBM, the CPUs' HBM and the DDR5$^\dagger$. The peak is 1.980 EFLOP/s of FP64, essentially $63{,}744 \times 52.4\,\text{TFLOP/s} = 3.34$ EFLOP/s of GPU vector FP64 on paper reduced by the clock the machine actually sustains, plus the CPUs; the TOP500 Rpeak figure of 1.980 EFLOP/s is what the HPL run was measured against$^\dagger$.

The design choices differ from Frontier's and El Capitan's in instructive ways (slide 27). Aurora has two CPUs per node to Frontier's one and six GPUs to Frontier's four modules; it puts HBM on the CPU, which no other exascale machine does; its CPU-to-GPU link is PCIe 5, not a coherent fabric like Frontier's Infinity Fabric or El Capitan's shared APU memory; and its GPU-to-GPU link is Xe Link within the node only. Eight NICs per node, against four on Frontier, give it the highest injection bandwidth per node of the three, a choice that suits the data-intensive and AI workloads Argonne planned for it as much as the simulation codes.

The history is the subtitle's first clause. In April 2015 the Department of Energy awarded Intel and Cray a contract for Aurora as a 180 PFLOP/s system (180--200 in contemporary accounts) built on Intel's Knights Hill many-core processors, to be delivered in 2018$^\dagger$. Knights Hill was cancelled in 2017 after Intel ended the Xeon Phi line, and in September 2017 the contract was re-planned as an exascale system for 2021. In March 2019 Intel and Argonne announced that the exascale Aurora would use a new Intel GPU, Ponte Vecchio, with the Sapphire Rapids Xeon, for 2021$^\dagger$. In July 2020 Intel disclosed that its 7 nm process was about a year behind, which moved Ponte Vecchio's base tiles to Intel 7 and its compute tiles to TSMC and the machine to 2022, then 2023 (slide 17). Installation of all 10{,}624 nodes was completed in June 2023; a partial HPL run on 5{,}439 nodes gave 585.34 PFLOP/s and second place on the November 2023 list, at 24.69 MW; the run on 9{,}234 nodes in May 2024 gave 1.012 EFLOP/s and second place in June 2024, at 38.70 MW, the figure that has stood since$^\dagger$. The eight years from contract to exaflop, against three for Frontier, are the cost of two changes of processor and one change of process node; HPCwire's headline in 2023 was ``Aurora the survivor.''

The HPL figures deserve a careful reading. The 1.012 EFLOP/s was achieved on 9{,}234 of the 10{,}624 nodes, 87\% of the machine, and Argonne and Intel have not submitted a full-system run; the ratio of HPL to Rpeak is $1.012 / 1.980 = 51\%$, the lowest of the exascale systems (Frontier reaches $1.353 / 2.06 = 66\%$ and El Capitan $1.809 / 2.82 = 64\%$; derived from the June 2026 list). At 38.7 MW the machine delivers 26 PFLOP/s per megawatt against 55 for Frontier and 61 for El Capitan (derived), which is why the slide and slide 27 call Aurora the least efficient of the exascale group; the June 2026 list places it fourth behind LineShine, El Capitan and Frontier, having been second in 2023--24 and third in 2024--25$^\dagger$. The mixed-precision benchmark tells a kinder story: Aurora's HPL-MxP, which runs the same linear solve in low precision with iterative refinement to FP64 accuracy and so exercises the XMX engines, reached 10.6 EFLOP/s in June 2024 and 11.6 EFLOP/s (on 9{,}500 nodes) in November 2024, and stood third on the June 2026 HPL-MxP list behind El Capitan and Frontier$^\dagger$. Goto and colleagues' 2026 paper on the HPL and HPL-MxP campaigns on Aurora describes the tuning that these runs required, from the two-stack exposure of each GPU to the handling of 63{,}744 devices' failure modes$^\dagger$.

The science is the reason the machine exists, and 2025 was its year. Argonne teams were on four of the Gordon Bell Prize finalists at SC25; two of the four ran on Aurora itself, one (CRK-HACC) on Frontier, and one (BerkeleyGW) on Frontier, Aurora and Perlmutter$^\dagger$. AERIS, a weather and climate foundation model trained at scale, sustained 10.21 EFLOP/s in mixed precision on 10{,}080 nodes (11.21 EFLOP/s peak), the largest such run reported on any system to that date; a light-matter dynamics code from the University of Southern California reached 1.87 EFLOP/s with more than a trillion atoms; and a BerkeleyGW many-body calculation ran at 0.7 EFLOP/s$^\dagger$. These results use the parts of Ponte Vecchio the datasheet comparison favored, FP64 vector throughput, large caches and node bandwidth, together with the XMX engines at BF16 for the AI model, and they are the strongest evidence that the machine does what it was bought for.

Two other systems used Ponte Vecchio: Dawn at the University of Cambridge (256 Dell nodes, 1{,}024 Max 1550s, 19.46 PFLOP/s, 41st in November 2023) and SuperMUC-NG Phase 2 at the Leibniz Supercomputing Centre (240 Lenovo nodes, 960 Max GPUs, 17.19 PFLOP/s, 52nd)$^\dagger$. No system with Intel GPUs has entered the list since 2023 (slide 28). Aurora will run to the end of the decade; its successor at Argonne is not an Intel machine (slide 29).

\begin{terms}
\term{Aurora}{Argonne's exascale system: 10,624 HPE Cray EX nodes of two Xeon Max CPUs and six Max 1550 GPUs, 63,744 GPUs, 1.012 EFLOP/s HPL at 38.7 MW; fourth on the June 2026 TOP500.}
\term{Xeon CPU Max 9470C}{Aurora's processor: a 52-core Sapphire Rapids Xeon with 64 GB of HBM2e on the package, two per node.}
\term{Slingshot-11}{HPE's 200 Gb/s Ethernet-derived HPC network; Aurora has eight NICs per node, Frontier four.}
\term{dragonfly network}{A network topology of fully connected groups joined by global links, giving low diameter at exascale; used by Slingshot on Aurora and Frontier.}
\term{HPE Cray EX}{HPE's liquid-cooled supercomputer cabinet line, the platform of Aurora, Frontier and El Capitan.}
\term{Knights Hill}{The cancelled Xeon Phi processor on which the 2015 Aurora contract was based.}
\term{HPL}{The High Performance Linpack benchmark, a dense FP64 linear solve; its result Rmax ranks the TOP500.}
\term{Rpeak}{A system's theoretical FP64 peak as listed on the TOP500; Aurora's is 1.980 EFLOP/s.}
\term{HPL efficiency}{Rmax divided by Rpeak; 51 percent for Aurora, 64 to 66 percent for El Capitan and Frontier.}
\term{HPL-MxP}{The mixed-precision Linpack: the solve runs in low precision on matrix units and is refined to FP64 accuracy; Aurora reached 11.6 EFLOP/s.}
\term{Gordon Bell Prize}{The ACM award for the largest-scale scientific computing achievement of the year; Argonne teams were on four 2025 finalists, two of which ran on Aurora.}
\term{AERIS}{A weather and climate foundation model trained on Aurora at 10.21 EFLOP/s sustained mixed precision on 10,080 nodes, a 2025 Gordon Bell finalist.}
\term{Dawn and SuperMUC-NG Phase 2}{The two other Ponte Vecchio systems on the TOP500 (Cambridge, 1,024 GPUs; LRZ, 960 GPUs), both from 2023.}
\end{terms}

\begin{intuition}
Aurora is the machine that proves and indicts Ponte Vecchio at once: it reached an exaflop and ran the largest mixed-precision science to date, and it did so at the lowest efficiency and after the longest delay of any exascale system. Judge the GPU by what the machine computes and the program by how long it took; both judgments are fair, and they point in opposite directions.
\end{intuition}

% ---------------------------------------------------------------------
\slide{20}{After Ponte Vecchio: cancellations and a restart from Xe3}{Rialto Bridge and Falcon Shores were cancelled; Jaguar Shores has no date; the GPU D-series is sampling.}

\para{What the slide shows} A four-column table lists Intel's data-center accelerator roadmap after Ponte Vecchio: product, when announced, what it was to be, and what happened. Rialto Bridge (2022) was to be Ponte Vecchio's successor with more Xe-cores in an 800 W OAM module, planned for 2023; it was cancelled in March 2023 in what Intel called a ``streamlined and simplified'' roadmap. Falcon Shores (2022) was first an XPU of CPU and GPU tiles, then a GPU-only part for 2025 with Gaudi's networking; it was cancelled as a product on January 30, 2025, and became an ``internal test chip.'' Gaudi 3 (April 2024), Habana's accelerator, shipped, fell short of its \$500 million 2024 target and is to be folded into Jaguar Shores. Jaguar Shores (2025), a ``rack-scale'' AI system to succeed both, has no date and is absent from Intel's third-quarter 2026 results. The GPU D-series, code-named Crescent Island (October 2025), is an inference GPU on Xe3P with 32 cores, 256 XMX engines, 1.3 PFLOP/s of FP8, 160 GB of LPDDR5X (480 GB on partner cards) at 1.5 TB/s and 350 W on a PCIe card; silicon was shown at Hot Chips in August 2026, the specification given at OCP in October 2026, and it is sampling with no stated general availability. Two panels below explain why the data-center GPU faltered and what the restart and the NVIDIA deal are.

\para{Details} The slide is a chronology of a retreat and the start of a return, and the dates matter because each entry changed the plan that preceded it.

\emph{Rialto Bridge} was announced at ISC in May 2022 as the direct successor to Ponte Vecchio: the same Xe-HPC architecture with more Xe-cores (up to 160 against 128), in the same OAM form factor at up to 800 W, with sampling planned for 2023, so that Aurora-style systems could be refreshed without a new node design$^\dagger$. On March 3, 2023, less than two months after Ponte Vecchio launched, Intel's Jeff McVeigh announced that Rialto Bridge was cancelled and that the data-center GPU line would move to a two-year cadence, skipping to Falcon Shores in 2025, in a ``streamlined and simplified'' roadmap that would do ``fewer things better''$^\dagger$. The practical meaning was that Ponte Vecchio would have no successor for at least two years, which ended the prospect of further Xe-HPC systems after Aurora, Dawn and SuperMUC-NG Phase 2.

\emph{Falcon Shores} was announced in February 2022 as an \emph{XPU}: a package combining x86 CPU tiles and Xe GPU tiles with shared memory, Intel's answer to the Grace Hopper superchip and to AMD's MI300A, for 2024$^\dagger$. In March 2023, with the Rialto Bridge cancellation, Intel said the XPU design was premature for a market whose workloads were changing too fast and redefined Falcon Shores as a GPU-only part for 2025; in late 2023 it said Falcon Shores would also incorporate Gaudi's Ethernet networking and matrix IP, merging the two accelerator lines into one product$^\dagger$. On January 30, 2025, on Intel's fourth-quarter call, interim co-chief executive Michelle Johnston Holthaus said that Falcon Shores would not be sold: ``we plan to leverage Falcon Shores as an internal test chip only'' and concentrate on a rack-scale successor, Jaguar Shores$^\dagger$. Three years after announcement, the part that was to replace both Ponte Vecchio and Gaudi was gone, and with it Intel's only planned data-center GPU with HBM.

\emph{Gaudi 3} is the one product on the table that shipped (slide 16): the 2024 accelerator sold below its \$500 million target, has not appeared in MLPerf since 2024, and was to be absorbed into Jaguar Shores$^\dagger$. \emph{Jaguar Shores}, named in late 2024 and described in 2025 only as a ``rack-scale'' AI solution, has no announced architecture, date or specification; Intel's third-quarter 2026 results, released October 29, 2026, mention the GPU D-series and do not mention Jaguar Shores or Gaudi$^\dagger$. Whether it still exists as a program is not public.

\emph{The GPU D-series} is the restart. Intel announced it on October 14, 2025, at the OCP Global Summit under the code name Crescent Island, as an inference GPU for data centers built on the Xe3P architecture, the discrete derivative of the Panther Lake graphics core (slide 15), with 160 GB of LPDDR5X memory, air-cooled, for sampling in the second half of 2026$^\dagger$. Intel showed the silicon at Hot Chips in August 2026 and gave the specification and the product name at the OCP Summit in October 2026: 32 Xe3P cores with 256 XMX engines, 1.3 PFLOP/s of dense FP8, 160 GB of LPDDR5X on Intel's own card and up to 480 GB on partner cards, 1.5 TB/s of memory bandwidth, 350 W on a PCIe card, sampling, with general availability not stated (a render label reads ``GPU D270'')$^\dagger$. Three things about this part are new for Intel and one is old. It is built from the graphics line, not Xe-HPC. It has no HBM: LPDDR5X, the memory of laptops and phones, gives large capacity at a fraction of HBM's cost and about a fifth of its bandwidth, which suits the decode phase of language-model inference, where a model's weights must fit in memory but each token touches them once. It is a PCIe card for air-cooled servers rather than an OAM module for liquid-cooled racks. And, as with every entry on this table, it is a year or more from announcement to volume. At 1.3 PFLOP/s of FP8 the D-series sits between an H100 (1.98) and a Gaudi 3 (1.84) on matrix throughput and well below an MI355X (5.0) or a B300 (5.0); its argument is memory per dollar and per watt, which is the Arc Pro argument of 2025 at data-center scale.

The two panels give the diagnosis and the prognosis. Why the data-center GPU faltered: the 7 nm delay of 2020 pushed Ponte Vecchio a year into Hopper's market; a 47-tile package cost yield and power; oneAPI was young against a CUDA ecosystem fifteen years old; and the market moved from FP64 simulation to low-precision inference, the precise point where Ponte Vecchio was weakest (slide 18). Each of the four is a reason a well-designed chip can fail commercially, and together they explain why Intel, having shipped one exascale GPU, cancelled its next two. The restart: an inference card on commodity memory for the workload that is growing, built from an architecture that ships in volume, sold through the standard PCIe server channel. And the \emph{NVIDIA deal}: on September 18, 2025, NVIDIA announced a \$5 billion investment in Intel at \$23.28 per share, with Intel to build custom x86 CPUs that integrate NVLink for NVIDIA's AI platforms and x86 processors with NVIDIA RTX graphics chiplets for PCs$^\dagger$. For this lecture the deal's meaning is that Intel's most concrete data-center role in NVIDIA's racks is as a CPU supplier joined to NVIDIA's fabric (slide 23), not as a GPU competitor. Intel's data-center and AI revenue grew 78\% year over year in the third quarter of 2026, to \$7.3 billion, on Xeon demand for AI hosts$^\dagger$.

Where does that leave the comparison this lecture makes? Through November 2026 Intel has one data-center GPU family in the field (Max, in Aurora and two European systems), one accelerator that shipped and faded (Gaudi), one inference GPU sampling (D-series), and no announced part with HBM or a rack-scale fabric. Against NVIDIA's Rubin and AMD's MI455X, both shipping with HBM4 and 72-GPU racks (slides 4 and 11), Intel is not in the 2026 contest for training systems. Its remaining assets are the ones the first slides of this section described: a GPU architecture in hundreds of millions of client processors, a compiler and driver stack that runs on all of them, oneAPI and SYCL as an open programming model (slide 24), and packaging technology that the rest of the industry now licenses or imitates. Whether those are enough to re-enter is the open question slide 32 leaves.

\begin{terms}
\term{Rialto Bridge}{Ponte Vecchio's planned successor (announced 2022, up to 160 Xe-cores, 800 W, for 2024), cancelled in March 2023.}
\term{Falcon Shores}{Intel's planned 2024--25 data-center part, first an XPU of CPU and GPU tiles, then a GPU with Gaudi networking; cancelled as a product in January 2025 and kept as an internal test chip.}
\term{XPU}{Intel's term for a package combining CPU and accelerator tiles with shared memory, the counterpart of AMD's APU and NVIDIA's superchip; Falcon Shores was to be one.}
\term{Jaguar Shores}{The rack-scale AI system Intel named in 2024--25 as the successor to Falcon Shores and Gaudi; no architecture, date or specification has been published.}
\term{GPU D-series (Crescent Island)}{Intel's Xe3P inference GPU announced in October 2025: 32 cores, 256 XMX, 1.3 PFLOP/s FP8, 160 to 480 GB of LPDDR5X at 1.5 TB/s, 350 W PCIe; sampling in late 2026.}
\term{Xe3P}{The discrete derivative of the Xe3 graphics architecture, the core of the GPU D-series.}
\term{LPDDR5X}{Low-power DDR5 DRAM, the memory of phones and laptops; gives an accelerator large capacity at low cost and about a fifth of HBM's bandwidth.}
\term{decode-bound inference}{The token-generation phase of language-model serving, in which every output token reads the whole set of weights once, so memory capacity and bandwidth, not arithmetic, set throughput.}
\term{OCP Global Summit}{The Open Compute Project's annual conference, where Intel announced Crescent Island (2025) and specified the GPU D-series (2026).}
\term{Intel--NVIDIA agreement}{The September 2025 deal in which NVIDIA invested five billion dollars in Intel and Intel agreed to build x86 CPUs with NVLink for NVIDIA's platforms and x86 processors with RTX chiplets for PCs.}
\term{two-year cadence}{The product rhythm Intel announced in March 2023 for data-center GPUs, against the yearly cadence NVIDIA and AMD adopted in 2024.}
\end{terms}

\begin{intuition}
Read the table as a sequence of plans each replaced before it shipped: Rialto Bridge by Falcon Shores, Falcon Shores by Jaguar Shores, Jaguar Shores by silence, and in their place a graphics-derived inference card with laptop memory. The restart is modest on purpose: a part Intel can build from what it ships in volume, for the workload that is growing, sold where it already sells servers. Whether a company can return to a market from that position is the question the next slides on fabrics and software help answer.
\end{intuition}

% =====================================================================
\setstage{posc}
\section{Fabrics, racks and standards (slides 21--23)}

\slide{21}{Scale-up fabrics: NVLink, Infinity Fabric, Xe Link, UALink}{The link between GPUs is NVIDIA's widest lead; AMD answered with a standard, Intel with Xe Link.}

\para{What the slide shows} A log-scale chart on the left plots the vendor-stated scale-up bandwidth per GPU, in GB/s, against the year of first shipment, from 2020 to 2027. The NVIDIA line (circles) runs 600 (A100, NVLink 3, 2020), 900 (H100, NVLink 4, 2022), 1{,}800 (B200, NVLink 5, late 2024; the B300 of 2025 keeps it) and 3{,}000 (Rubin, NVLink 6, 2026); the labels ``NVLink 5'' and ``NVLink 6'' sit by the last two points. The AMD line (squares) runs 276 (MI100, 2020), 800 (MI250X, 2021), 896 (MI300X, 2023; labeled ``MI300X: 7 $\times$ 128''), 1{,}008 (MI355X, 2025) and 3{,}600 (MI455X with UALink over Ethernet, 2026; labeled ``MI455X UALoE''). A single orange triangle, labeled ``Intel Gaudi 3 (24 $\times$ 200 GbE)'', marks Intel's Gaudi 3 at 1{,}200 GB/s (2024), its twenty-four 200 Gb/s ports counted bidirectionally like the NVLink and Infinity Fabric figures, the only Intel part with a published per-accelerator scale-up figure; Ponte Vecchio's Xe Link has no comparable published number and is not plotted. The two lines are flat for four years at a factor of two or three apart and then both jump by a factor of three to four in 2026. The panel ``Four fabrics'' summarizes NVLink (switched since 2018, in-switch reductions since Hopper, a 72-GPU domain since 2024), Infinity Fabric (direct links on an eight-GPU board; SemiAnalysis's measured 64 GB/s per pair against 450 for NVLink), Xe Link (sixteen links joining six GPUs in an Aurora node, never scaled to a switch) and UALink (the open standard, 200 Gb/s per lane, up to 1{,}024 accelerators; Helios runs it over Ethernet at 3.6 TB/s per GPU).

\para{Details} A scale-up fabric is the network that joins the GPUs of one node or one rack into a single memory-coherent or load--store domain, as distinct from the scale-out network (InfiniBand or Ethernet) that joins nodes. The distinction matters because the two carry different traffic. Tensor-parallel training and multi-GPU inference split one layer's matrices across GPUs and exchange partial results after every layer, tens of times per token for a large model; these transfers are latency-sensitive and need bandwidth comparable to the local memory system. Data-parallel training exchanges gradients once per step, a larger but far less frequent transfer. The scale-up domain is sized to the first kind of traffic, which is why every vendor's figure is given per GPU and in hundreds of gigabytes per second, and why the size of the domain (8 GPUs or 72) decides how large a model can be served from one set of coherent memories.

NVIDIA's numbers are the datasheet figures for the SXM modules, each the sum of both directions over all links$^\dagger$. NVLink 3 on the A100 has 12 links at 50 GB/s for 600 GB/s; NVLink 4 on the H100 has 18 links at 50 GB/s for 900 GB/s; NVLink 5 on Blackwell has 18 links at 100 GB/s for 1.8 TB/s; NVLink 6 on Rubin was announced at CES in January 2026 at 3.6 TB/s and appears on the September 2026 spec revision at 3.0 TB/s per GPU (slide 4)$^\dagger$. A GB200 or GB300 NVL72 rack therefore has $72 \times 1.8 = 129.6$ TB/s of NVLink bandwidth, which NVIDIA rounds to 130; the Vera Rubin NVL72 has $72 \times 3.0 = 216$ TB/s (derived; the CES figure was 260). Since the DGX-2 of 2018 these links have terminated in switch chips rather than in other GPUs: NVSwitch 3 on the Hopper HGX board has 64 ports; the Blackwell-generation NVLink Switch chip has 72 ports and 7.2 TB/s of bandwidth, and eighteen of them in nine switch trays join the 72 GPUs of an NVL72$^\dagger$. The consequence of switching is that any pair of GPUs can use the whole per-GPU bandwidth at once, and that collectives can be executed partly in the switch: NVSwitch 3 introduced SHARP in-network reductions and the NVLink 6 switch carries 130 TFLOP/s of in-network arithmetic, a figure NVIDIA gave at Hot Chips 2026$^\dagger$. The 72-GPU domain has existed since the first GB200 NVL72 racks shipped in November 2024 (slide 22).

AMD's Infinity Fabric figures are also datasheet sums over all links$^\dagger$. The MI100 has three links at 92 GB/s for 276 GB/s; the MI250X eight external links at 100 GB/s for 800 GB/s aggregate (the two dies in a module are joined in-package at 400 GB/s bidirectional); the MI300X seven links at 128 GB/s for 896 GB/s; the MI355X seven links at 144 GB/s for 1{,}008 GB/s. The topology is the thing to notice. Seven links on an eight-GPU UBB board means that every GPU has exactly one direct link to each of the other seven; there is no switch, and the traffic between a pair is confined to the one 128 GB/s link that joins them. On an NVSwitch board the same pair could use the whole 900 GB/s (450 GB/s per direction) of the GPU's NVLink ports, routed over as many switch paths as needed. That difference is what SemiAnalysis measured in December 2024: about 64 GB/s per GPU pair on the MI300X against about 450 GB/s on the H100, for a one-direction point-to-point test$^\dagger$. For an all-to-all collective across eight GPUs the two topologies are closer, because the seven MI300X links are all busy at once; for a pair or for a tensor-parallel group of two or four they are not, and the deck's panel treats this as the structural weakness of a point-to-point fabric rather than a defect of the MI300X.

AMD's 2026 answer is to switch the fabric and to make the protocol a standard. The MI455X exposes 72 lanes of UALink carried over Ethernet physical layers and switches (``UALoE''), 3.6 TB/s bidirectional per GPU, and the Helios rack joins its 72 GPUs through Ethernet switch trays for 260 TB/s of scale-up bandwidth in total ($72\times 3.6 = 259$, derived)$^\dagger$. The per-GPU number is above Rubin's revised 3.0 TB/s; the rack number is above the Vera Rubin NVL72's 216 TB/s. Neither says anything yet about achieved collective bandwidth or latency, which depend on the switch silicon and the software, and for Helios no independent measurement existed in November 2026 (slide 22).

UALink itself is the open standard the deck plots as the AMD line's last point. The Ultra Accelerator Link consortium was incorporated in October 2024 by AMD, Intel, AWS, Google, Meta, Microsoft, Cisco, HPE, Apple, Alibaba, Astera Labs and Synopsys; its 200G 1.0 specification of April 8, 2025 defines a memory-semantic (load, store, atomic) link at 200 Gb/s per lane, in one-, two- or four-lane stations, 800 Gb/s per direction for a four-lane station, with a domain of up to 1{,}024 accelerators, a round-trip latency target under one microsecond and copper reach under four meters$^\dagger$. The protocol is derived from Infinity Fabric's memory semantics; the physical layer borrows from Ethernet so that merchant switch vendors can build the switches, which is what ``UALink over Ethernet'' in Helios means. Slide 23 places it among the other standards.

Intel's fabrics are the chart's two gaps. Ponte Vecchio's Xe Link tile provides sixteen links per GPU at 53 Gb/s per lane; Intel never published a lanes-per-link figure or a per-GPU aggregate, so the deck does not plot it$^\dagger$. In an Aurora node the six Max 1550s are joined all-to-all over Xe Link, and that is as far as it went: there was no Xe Link switch and no eight-GPU or rack-scale Xe Link domain (slide 19). Gaudi 3 took the opposite approach and put twenty-four 200 Gb/s Ethernet ports on the package, 4.8 Tb/s or 600 GB/s per direction (1{,}200 GB/s bidirectional, the convention the chart uses for all three vendors), used for both scale-up (typically twenty-one ports in an eight-card system) and scale-out (the remaining three), with standard Ethernet switches$^\dagger$; it is the one Intel point on the chart. Intel's later plan is to adopt UALink and Ultra Ethernet, of which it is a founding member, for Jaguar Shores, a product without a date (slide 20).

Two derived observations close the slide. First, the ratio of the two lines: NVIDIA's per-GPU scale-up bandwidth exceeded AMD's by $600/276 = 2.2\times$ in 2020, $900/896 \approx 1.0\times$ in 2023 (but switched against point-to-point), $1{,}800/1{,}008 = 1.8\times$ in 2025, and in 2026 AMD's figure is higher, $3{,}600/3{,}000 = 1.2\times$ (derived). Second, the ratio of fabric bandwidth to memory bandwidth, which says how expensive remote data is relative to local: for the H100, $900/3{,}350 = 0.27$; for the MI300X, $896/5{,}300 = 0.17$; for Rubin, $3{,}000/19{,}200 = 0.16$; for the MI455X, $3{,}600/23{,}300 = 0.15$ (derived). Both vendors' fabrics grew about as fast as HBM did, so the cost of reaching a neighbor's memory, in units of local memory bandwidth, stayed at roughly six to seven times, and the software's job of keeping communication off the critical path did not get easier.

\begin{terms}
\term{scale-up fabric}{The network joining the GPUs of one node or rack into one load--store or coherent domain; measured in bandwidth per GPU and in the number of GPUs it joins.}
\term{scale-out network}{The InfiniBand or Ethernet network joining nodes or racks; carries data-parallel and pipeline traffic rather than tensor-parallel traffic.}
\term{NVLink}{NVIDIA's GPU-to-GPU link: generations 3 (600 GB/s, A100), 4 (900, H100), 5 (1,800, Blackwell) and 6 (3,000 on the 2026 spec sheet, Rubin), each figure the sum of both directions over all links.}
\term{NVSwitch and NVLink Switch}{NVIDIA's NVLink crossbar chips: NVSwitch 3 (64 ports, Hopper), the Blackwell NVLink Switch (72 ports, 7.2 TB/s; eighteen per NVL72) and the NVLink 6 switch with 130 TFLOP/s of in-network arithmetic.}
\term{in-network reduction}{Computing the sums of an all-reduce inside the switch as data passes through (NVIDIA's SHARP), removing a round of communication from the collective.}
\term{Infinity Fabric (xGMI)}{AMD's coherent interconnect; between Instinct GPUs it is a set of direct point-to-point links, seven of 128 GB/s on the MI300X, seven of 144 on the MI355X.}
\term{point-to-point versus switched}{In a point-to-point fabric each pair of GPUs shares one link; in a switched fabric any pair can use the GPU's whole port bandwidth. The eight-GPU UBB board is the first, the HGX board the second.}
\term{UALink}{Ultra Accelerator Link, the open memory-semantic scale-up standard (consortium October 2024; 1.0 specification April 2025): 200 Gb/s per lane, up to 1,024 accelerators.}
\term{UALoE}{UALink over Ethernet: AMD's implementation in the MI455X and Helios, with UALink's protocol carried on Ethernet physical layers and switches; 3.6 TB/s per GPU.}
\term{Xe Link}{Intel's GPU-to-GPU link on Ponte Vecchio: sixteen links per GPU on a dedicated tile, 53 Gb/s per lane, used to join six GPUs in an Aurora node.}
\term{bidirectional bandwidth}{The convention, used by NVIDIA and AMD, of quoting a link as the sum of both directions; a 900 GB/s NVLink moves 450 GB/s each way.}
\end{terms}

\begin{intuition}
Think of the fabric as the second memory system: a GPU reads its own HBM at 20 TB/s and its neighbor's at 3, so the program must arrange that most bytes are local and that the remote ones are in flight while local work proceeds. NVIDIA's lead was never in the raw link speed, where AMD kept pace, but in switching it (any pair at full rate), scaling it (72 GPUs) and computing in it; UALink and Helios are AMD adopting all three.
\end{intuition}

% ---------------------------------------------------------------------
\slide{22}{Rack-scale: the NVL72 family against Helios}{Both sell 72-GPU liquid-cooled racks with their own CPUs; AMD's is wider, holds more HBM, shipped later.}

\para{What the slide shows} A table compares four racks column by column: GB200 NVL72, GB300 NVL72, Vera Rubin NVL72 and AMD Helios. The rows are first shipments (November 2024, July 2025, August 2026, Q3 2026), GPUs (72 B200, 72 B300, 72 Rubin, 72 MI455X), CPUs (36 Grace, 36 Grace, 36 Vera, 18 Venice), total HBM (13.4, 20, 20.7 and 31 TB), HBM bandwidth (0.58, 0.58, 1.4 and 1.7 PB/s), dense FP4 (0.72, 1.08, 2.52 and 2.9 EFLOP/s), dense FP8 (0.36, 0.36, 1.26 and 1.4 EFLOP/s), FP64 (2.9 PFLOP/s, 0.1, 2.4, not published), scale-up fabric (NVLink 5, NVLink 5, NVLink 6, UALink over Ethernet), total scale-up bandwidth (130, 130, 216 and 260 TB/s), scale-out per GPU (400, 800, 1{,}600 and 2{,}400 Gb/s), power (about 120 kW, about 130 kW, not published, 225--245 kW) and form factor (one standard rack three times; a double-wide rack for Helios). A footnote gives the sources: NVIDIA product pages with the Vera Rubin figures as revised in September 2026; AMD's July 2026 launch and the Supermicro Helios datasheet, which still lists the earlier 1.4 PB/s; the GB200 dense FP8 derived from the sparse 720 PFLOP/s; the Helios power from press coverage. The right panel reads the two designs against each other: the same shape (72 GPUs, CPUs in the tray, 216--260 TB/s of scale-up in the 2026 racks, liquid cooling), a different fabric (NVIDIA's own switch silicon with in-network compute against UALink on merchant Ethernet switches), different ratios (one Grace or Vera per two GPUs against one Venice per four; one ConnectX-9 per GPU against three Vulcano 800G NICs; one rack against a double-wide one), and the caution that delivered racks and measured throughput were not yet comparable in November 2026.

\para{Details} The rack is the unit in which both vendors now sell their largest accelerators, and the table is built so that every row is a quantity a buyer would compare. It is worth going through how each figure arises, because several are products of the per-GPU numbers of earlier slides and one or two are not.

The NVIDIA racks are the GB200 NVL72, announced at GTC in March 2024, with the first complete racks delivered by Dell to CoreWeave on November 18, 2024 (Microsoft Azure had run GB200 servers in October) and shipments in volume from December; the GB300 NVL72, announced at GTC in March 2025 with first racks at CoreWeave on July 3, 2025; and the Vera Rubin NVL72, announced at CES on January 5, 2026, declared in full production in May and shipped as first revenue racks in August 2026$^\dagger$. Each holds 18 compute trays of two superchips, each superchip one Grace or Vera CPU with two GPUs, so 36 CPUs and 72 GPUs; 9 switch trays of two NVLink Switch chips; and liquid cooling throughout. The memory rows are $72\times$ the per-GPU HBM: $72 \times 186$ GB $= 13.4$ TB for the GB200 (NVIDIA's page gives 13.4 TB for the rack and 372 GB per superchip, so 186 GB per GPU), $72 \times 288 = 20.7$ TB for both the GB300 and Rubin (NVIDIA rounds the GB300 page to 20 TB and the Rubin page to 20.7 TB); the bandwidth rows are $72 \times 8 = 576$ TB/s, rounded to 0.58 PB/s, for the two Blackwell racks and $72 \times 19.2 = 1{,}382$ TB/s, which NVIDIA states as 1{,}400 TB/s, for Rubin (derived)$^\dagger$. The compute rows are $72\times$ the dense per-GPU peaks of slide 6: FP4 $72 \times 10 = 720$ PFLOP/s, $72 \times 15 = 1{,}080$ and $72 \times 35 = 2{,}520$; FP8 $72 \times 5 = 360$ twice and $72 \times 17.5 = 1{,}260$; FP64 $72 \times 40 = 2{,}880$ TFLOP/s (2.9 PFLOP/s), $72 \times 1.4 \approx 100$ TFLOP/s and $72 \times 33 = 2{,}376$, stated as 2{,}400 (derived). NVIDIA's rack pages headline the sparse figures (1{,}440 PFLOP/s FP4 for both Blackwell racks, 3{,}600 for Rubin); the GB200 page gives FP8 only as 720 PFLOP/s sparse, so the table's 0.36 EFLOP/s dense is derived by halving. The scale-up row is $72\times$ the NVLink figure: $72 \times 1.8 = 130$ TB/s and $72 \times 3.0 = 216$. The scale-out row is per GPU: one ConnectX-7 at 400 Gb/s in the GB200 rack, one ConnectX-8 at 800 in the GB300 rack, one ConnectX-9 at 1.6 Tb/s in the Rubin rack, and three Vulcano 800G NICs (2.4 Tb/s) per MI455X in Helios$^\dagger$. The power row is not an NVIDIA figure: the ``about 120 kW'' for the GB200 NVL72 comes from Dell and press coverage of the first installations, the GB300's ``about 130'' from the same sources, and Rubin's is not published; the panel notes that these are the first racks for which a single rack's power was itself a datacenter design problem$^\dagger$.

Helios is AMD's first rack-scale product. It was previewed at Advancing AI in June 2025, shown as a mechanical design at the OCP Global Summit in October 2025, and launched with the MI455X at Advancing AI on July 23, 2026; AMD's third-quarter 2026 report of November 3 said the first Helios revenue shipments had occurred in the quarter and that racks were running at Anthropic, Meta, Microsoft, OpenAI and Oracle$^\dagger$. The rack holds 18 compute trays, each with four MI455X, one EPYC ``Venice'' CPU (the sixth-generation EPYC on Zen 6) and one Pensando Salina DPU, so 72 GPUs, 18 CPUs and 18 DPUs; scale-up switch trays running UALink over Ethernet; and Pensando Vulcano 800G NICs for scale-out$^\dagger$. The memory row is $72 \times 432 = 31.1$ TB; the bandwidth row $72 \times 23.3 = 1{,}678$ TB/s, stated as 1.7 PB/s (the Supermicro datasheet, written against the earlier 19.6 TB/s per GPU, still says 1.4 PB/s); the compute rows $72 \times 40.26 = 2{,}899$ PFLOP/s MXFP4 and $72 \times 20.13 = 1{,}449$ MXFP8, stated as 2.9 and 1.4 EFLOP/s (derived)$^\dagger$. FP64 is not published because the MI455X's FP64 rate is not published (slide 11). Scale-up is 260 TB/s ($72 \times 3.6 = 259$); scale-out is stated by AMD as 43 TB/s for the rack and 2.4 Tb/s per GPU: twelve Vulcano 800G NICs per four-GPU tray, 216 per rack, three per GPU, which is the table's 2{,}400 Gb/s ($72 \times 2.4\,\text{Tb/s} \approx 21.6$ TB/s per direction, 43 bidirectional; derived)$^\dagger$. The power range of 225--245 kW is from press coverage of the launch, not from an AMD datasheet$^\dagger$. The rack is built to the OCP Open Rack Wide specification, a double-wide 21-inch rack AMD chose, it said, for serviceability: the trays slide out from the front, the switch trays are separate, and the width allows cable-free blind-mate liquid and power connections at the back$^\dagger$.

The panel's comparisons follow from the rows. On memory Helios holds $31/20.7 = 1.5\times$ the HBM of a Vera Rubin NVL72 and $1.7/1.4 = 1.2\times$ the bandwidth; on dense FP4 it is $2.9/2.52 = 1.15\times$; on dense FP8 $1.4/1.26 = 1.1\times$; on scale-up bandwidth $260/216 = 1.2\times$ (derived). These are all the per-GPU ratios of slides 5, 6 and 21 multiplied by the same 72. The two rows where the designs differ in kind are the fabric and the CPU ratio. NVIDIA's NVLink 6 switch is its own silicon, with the in-network arithmetic of slide 21 and a protocol tuned over six generations; Helios's switches are Ethernet switches from merchant vendors running UALink's memory semantics over them, which is open and cheaper but inherits Ethernet's framing and latency. NVIDIA pairs one Grace or Vera with two GPUs, so that each GPU has 240 GB of LPDDR5X (half of Grace's 480 GB) or up to 750 GB (half of Vera's 1.5 TB) of CPU memory reachable at NVLink-C2C speed (derived), and the CPU can hold the key-value cache or the optimizer state that does not fit in HBM; AMD pairs one Venice with four MI455X over PCIe-class host links of 256 GB/s per GPU, and relies on the larger HBM (432 GB) to need the CPU less$^\dagger$. Which is right depends on the workload: for decode-bound inference with long contexts the CPU memory matters; for training at rack scale it matters less.

The panel's last point is the one to carry away. In November 2026 the Vera Rubin NVL72 had shipped as revenue racks for three months and the GB200 and GB300 racks had two years of deployment behind them, with MLPerf results at 8{,}192 GPUs (slide 30); Helios had shipped its first racks in the quarter and had no published benchmark. A table of datasheet products says what the two companies designed. What they delivered, at what fraction of these peaks, would be known only from measurements that did not yet exist.

\begin{terms}
\term{NVL72}{NVIDIA's rack of 72 GPUs and 36 CPUs in 18 compute trays, joined by 18 NVLink Switch chips in 9 switch trays into one NVLink domain; GB200 (2024), GB300 (2025) and Vera Rubin (2026) versions.}
\term{Helios}{AMD's 72-GPU rack (2026): 18 trays of four MI455X, one EPYC Venice and one Salina DPU; UALink over Ethernet for scale-up; Vulcano 800G NICs; OCP Open Rack Wide.}
\term{superchip}{NVIDIA's board of one Grace or Vera CPU and two Blackwell or Rubin GPUs joined by NVLink-C2C; two per compute tray.}
\term{Grace and Vera}{NVIDIA's Arm server CPUs: Grace (72 Neoverse V2 cores, up to 480 GB LPDDR5X) in the Blackwell racks; Vera (88 Olympus cores, up to 1.5 TB) in the Rubin rack.}
\term{EPYC Venice}{AMD's sixth-generation EPYC server CPU on Zen 6 cores; one per four GPUs in Helios.}
\term{DPU}{Data processing unit, a NIC with its own cores for networking, storage and security offload; AMD's Pensando Salina in Helios, NVIDIA's BlueField in its racks.}
\term{ConnectX and Vulcano}{The scale-out NICs: NVIDIA's ConnectX-7 (400 Gb/s), -8 (800) and -9 (1.6 Tb/s), one per GPU; AMD's Pensando Vulcano at 800 Gb/s, three per GPU in Helios.}
\term{OCP Open Rack Wide}{The Open Compute Project's double-wide (21-inch) rack specification, adopted by AMD for Helios.}
\term{rack power}{The electrical load of one rack, 120 kW to 245 kW for these designs against 10--20 kW for a conventional server rack; the reason all of them are liquid-cooled and sited in purpose-built halls.}
\term{PB/s}{Petabytes per second, $10^{15}$ B/s; the rack-level unit for HBM bandwidth (0.58 to 1.7 PB/s here).}
\end{terms}

\begin{intuition}
A rack is 72 of the per-GPU numbers multiplied out, plus three things that do not multiply: the fabric that joins them, the CPU that feeds them and the power that limits them. Helios and the NVL72 have converged on the same shape because the physics of a 100--250 kW liquid-cooled cabinet allows little else; the fabric and the software are where they still differ, and those are measured, not read from a datasheet.
\end{intuition}

% ---------------------------------------------------------------------
\slide{23}{Standards: OAM, UBB, UALink, Ultra Ethernet, NVLink Fusion}{Everyone but NVIDIA builds on open form factors and fabrics; NVIDIA opened its fabric to partners instead.}

\para{What the slide shows} A horizontal timeline across the top carries eight dated events, alternating above and below the line: 2019, OCP OAM 1.0, the accelerator module specification contributed by Facebook, Microsoft and Baidu; 2022--23, UBB 2.0 baseboards, with the note that the MI300X, Gaudi 3 and Ponte Vecchio all ship on OAM; July 2023, the Ultra Ethernet Consortium formed by AMD, Intel, Meta, Microsoft, Cisco, HPE, Broadcom and others; October 2024, the UALink consortium incorporated by AMD, Intel, AWS, Google, Meta, Microsoft, Cisco and others; April 2025, the UALink 200G 1.0 specification (200 Gb/s per lane, 1{,}024 accelerators); May 2025, NVLink Fusion, NVIDIA's offer of NVLink IP for partners' CPUs and accelerators; June 2025, the Ultra Ethernet 1.0 specification; and October 2025, highlighted, Helios shown on OCP Open Rack Wide and the OpenAI--AMD 6 GW agreement. Two panels below read the stack layer by layer: the module and board (OAM at 102 $\times$ 165 mm, 700 W and later more, and the eight-module UBB, which let a server vendor build one chassis for AMD, Intel and others, against NVIDIA's own SXM and HGX), the scale-up layer (UALink, derived from Infinity Fabric) and the scale-out layer (Ultra Ethernet, RDMA over Ethernet against InfiniBand); and what the standards have and have not achieved: they made the second source possible (an MI355X or Gaudi 3 drops into a UBB chassis; Helios's fabric is bought from Ethernet switch vendors) but had not, by November 2026, produced a shipped multi-vendor scale-up domain, UALink silicon having been announced for 2026 and Helios being AMD-only; NVIDIA's answer, NVLink Fusion, admits a partner's chip to an NVLink domain on NVIDIA's terms.

\para{Details} Standards in this industry exist because the buyers want a second source and the second-place vendors want to be it. The hyperscalers, who buy most of the accelerators, have no wish to design a new chassis, a new fabric and a new network for every vendor; the vendors other than NVIDIA cannot each afford to build the whole platform NVIDIA built. The result since 2019 has been a stack of open specifications, one per level of slide 3, each contributed by a different coalition and each with a different degree of success.

The module and the board came first. The OCP Accelerator Module (OAM) specification, version 1.0 in 2019, defines a mezzanine module of 102 by 165 millimeters with a standardized connector carrying power, PCIe to the host and up to seven high-speed interconnect ports to other modules; the initial power envelope was 700 W and later revisions allowed more$^\dagger$. It was contributed to the Open Compute Project by Facebook with Microsoft and Baidu, and its purpose was exactly the hyperscaler's: one socket into which any vendor's accelerator could be dropped. The Universal Baseboard (UBB) is the companion specification for the board that holds eight OAM modules, routes their interconnect ports among themselves (the seven links of slide 21 map onto the seven OAM ports) and presents the result to a host; UBB 2.0, from 2022--23 (the MI300X shipped on it in December 2023), raised the power and signaling rates for the MI300 generation$^\dagger$. AMD's MI250X, MI300X, MI325X and MI355X, Intel's Max 1550 and Gaudi 2 and 3, and several startups' parts are all OAM modules on UBB boards, which is why Supermicro, Dell, HPE and Lenovo could offer an ``eight-GPU AI server'' for all of them with one mechanical design. NVIDIA never joined: its SXM module and HGX baseboard are proprietary, with the NVSwitch chips on the board, and its market position has let it keep them so. The MI455X departs from OAM in the other direction, to an ``Enhanced Accelerator Module'' of AMD's own for the higher power and the UALink lanes of Helios, so that from 2026 the two leading vendors both ship rack products on modules of their own design, with the open module remaining the form for eight-GPU servers$^\dagger$.

The scale-up layer is UALink, described on slide 21. Its genealogy explains both its promise and its lateness. AMD had the only mature non-NVIDIA coherent GPU fabric, Infinity Fabric, and contributed its memory-semantic protocol; Broadcom, Cisco and the Ethernet switch vendors wanted the physical layer to be Ethernet so that they could sell the switches; the hyperscalers wanted a fabric they could put under their own accelerators (AWS Trainium, Google TPU, Microsoft Maia, Meta MTIA) as well as under AMD's and Intel's. The consortium was incorporated in October 2024, the 200G 1.0 specification followed in April 2025, and the first UALink switch silicon was announced by several vendors for 2026$^\dagger$. By November 2026 the one shipping UALink system was Helios, and it runs AMD GPUs through switches configured by AMD; a rack mixing two vendors' accelerators on one UALink domain, which is what the standard is for, had not been built.

The scale-out layer is Ultra Ethernet. The Ultra Ethernet Consortium was formed under the Linux Foundation in July 2023 by AMD, Arista, Broadcom, Cisco, Eviden, HPE, Intel, Meta and Microsoft, and its 1.0 specification was released on June 11, 2025$^\dagger$. The problem it addresses is that InfiniBand, which NVIDIA owns since its 2020 purchase of Mellanox, had become the default scale-out network for AI clusters, with RDMA, lossless flow control and in-network reductions, while standard Ethernet with RoCE (RDMA over Converged Ethernet) suffered from congestion collapse at the scale of thousands of nodes. Ultra Ethernet specifies a new transport layer with packet spraying over multiple paths, selective retransmission and congestion control designed for the all-to-all traffic of training, on standard Ethernet physical layers and switches. AMD's Pensando Pollara 400 (2024) was announced as the first UEC-ready NIC and the Vulcano in Helios follows the 1.0 specification; Broadcom's Thor Ultra (2025) is another; NVIDIA's Spectrum-X Ethernet is its proprietary counterpart, and NVIDIA sells both InfiniBand and Spectrum-X with its racks$^\dagger$.

NVLink Fusion is NVIDIA's response to the whole stack, announced at Computex on May 18, 2025$^\dagger$. It licenses the NVLink chip-to-chip interface and the NVLink switch connectivity to partners so that a custom CPU (Fujitsu's and Qualcomm's were named) or a custom accelerator (built with MediaTek, Marvell, Alchip, Astera Labs, Synopsys or Cadence as design partners) can sit in an NVLink domain alongside NVIDIA GPUs or, in the CPU case, replace Grace or Vera as the host of NVIDIA GPUs; Intel's September 2025 agreement to build x86 CPUs with NVLink is the largest instance (slide 20)$^\dagger$. The difference from UALink is governance: a UALink member can build a whole system without NVIDIA; an NVLink Fusion partner builds a component that NVIDIA's switches and NVIDIA's software admit to NVIDIA's domain. The deck's phrase is that NVIDIA opened its fabric to partners rather than to a standard.

The second panel's verdict is deliberately measured, and students should be able to restate it. What the standards achieved is real: a second source at the server level, because an MI355X or a Gaudi 3 drops into the same UBB chassis as its predecessor and a server vendor's engineering is done once; a second source at the network level, because Helios buys its scale-up and scale-out switches from the Ethernet industry rather than building them; and a credible threat that keeps NVIDIA's prices and openness under pressure, of which NVLink Fusion is evidence. What they had not achieved by November 2026 is the thing a standard is finally for: interoperation. No shipped rack joined two vendors' accelerators on one scale-up fabric; no shipped UALink switch had been measured against an NVLink switch on collective bandwidth and latency; Ultra Ethernet's first deployments were young. The October 2025 pairing on the timeline, Helios on an open rack and the OpenAI--AMD agreement for 6 GW, is highlighted because it was the moment the open stack acquired a buyer large enough to make it matter; whether it delivers is the question of slide 32.

\begin{terms}
\term{OAM}{OCP Accelerator Module: the Open Compute Project's 102 by 165 mm mezzanine module specification (1.0 in 2019) with standardized power, host and interconnect connectors; the form factor of AMD Instinct, Intel Max and Gaudi parts.}
\term{UBB}{Universal Baseboard: the OCP board holding eight OAM modules and routing their interconnect ports among them; UBB 2.0 (2022--23) for the MI300 generation.}
\term{SXM and HGX}{NVIDIA's proprietary module and eight-GPU baseboard, with the NVSwitch chips on the board; the closed counterpart of OAM and UBB.}
\term{Open Compute Project}{The Facebook-founded foundation (2011) that publishes open hardware specifications for data centers, including OAM, UBB and the Open Rack family.}
\term{Ultra Ethernet}{The Ultra Ethernet Consortium's (July 2023) scale-out transport for AI and HPC on standard Ethernet, with multipath packet spraying and new congestion control; 1.0 specification June 11, 2025.}
\term{RoCE}{RDMA over Converged Ethernet: remote direct memory access carried on Ethernet, the pre-Ultra-Ethernet way to run AI traffic on Ethernet; prone to congestion at scale.}
\term{InfiniBand}{The lossless RDMA network owned by NVIDIA since its purchase of Mellanox (2020); the default scale-out network of NVIDIA clusters, with in-network reductions.}
\term{NVLink Fusion}{NVIDIA's program (May 2025) licensing NVLink interfaces to partners so that custom CPUs and accelerators can join NVLink domains; Intel's x86-with-NVLink CPUs are its largest instance.}
\term{second source}{A buyer's ability to obtain an equivalent part from a second vendor without redesigning the system; the purpose of OAM, UBB and UALink.}
\term{interoperation}{Two vendors' parts working together on one fabric or board; the test a standard must eventually pass, not yet passed by UALink in November 2026.}
\end{terms}

\begin{intuition}
Each open standard sits at one level of the platform and was written by whoever needed a second source at that level: the hyperscalers for the module, AMD and the switch vendors for the scale-up fabric, everyone but NVIDIA for the network. Judge a standard by what has shipped across vendors on it, not by its membership list; by that test the module standards have succeeded, the fabric standards are one rack old, and NVIDIA's counter-offer is to let partners in without letting the standard in.
\end{intuition}

% =====================================================================
\setstage{trn}
\section{Software stacks and portability (slides 24--26)}

\slide{24}{CUDA, ROCm and HIP, oneAPI and SYCL: three stacks}{The same layers in each stack: driver, compiler, libraries, frameworks; what differs is age and breadth.}

\para{What the slide shows} Three columns of five stacked boxes, one column per vendor in its color, draw the software stacks layer by layer from the top down. Frameworks: PyTorch, JAX, TensorRT-LLM and vLLM for NVIDIA; PyTorch, JAX, vLLM and SGLang for AMD; PyTorch (the XPU back end) and OpenVINO for Intel. Communication: NCCL and NVSHMEM; RCCL; oneCCL. Math and deep-learning libraries: cuBLAS, cuDNN, CUTLASS and TensorRT; rocBLAS, hipBLASLt, MIOpen and Composable Kernel; oneMKL and oneDNN. Language and compiler: CUDA C++, \texttt{nvcc} and PTX, with Triton; HIP C++, \texttt{hipcc} and LLVM, with Triton; SYCL, DPC++ and OpenMP, with SYCLomatic. Runtime and driver: the CUDA runtime and driver; the HSA runtime, ROCr and the \texttt{amdgpu} kernel driver; Level Zero and Intel's compute runtime. A caption notes that the portability layers (Kokkos, RAJA, OpenMP offload, SYCL with Codeplay's plugins for NVIDIA and AMD, and Triton, whose compiler targets all three) sit above all three stacks. A small panel gives dates: CUDA 1.0 in June 2007; ROC(m) 1.0 in April 2016 (the April 19, 2016 release was branded ``ROC 1.0'', Radeon Open Compute; the ROCm name followed), 6.0 in December 2023, 7.0 in September 2025 and 10.0 in August 2026 (versions 8 and 9 were skipped); the oneAPI 1.0 specification in September 2020, SYCL 2020 in February 2021 and the UXL Foundation in September 2023. The right panel, ``What differs'', compares age and breadth, the porting paths (HIP's \texttt{hipify} and Intel's SYCLomatic, which Intel says converts 80--90\% of CUDA automatically), framework coverage (PyTorch on all three; vLLM and SGLang on NVIDIA and AMD; the performance-setting kernels arriving first for CUDA), Intel's split between oneAPI for Xe and SynapseAI for Gaudi, and governance (CUDA closed, ROCm open source, oneAPI's specification in the Linux Foundation's UXL).

\para{Details} The slide's claim is structural: the three stacks have the same layers because the problem is the same. A GPU program needs a driver that owns the device and its memory, a runtime that launches kernels and moves data, a compiler from some kernel language to the device's instruction set, libraries for the operations everyone needs (dense linear algebra, convolutions, attention, collectives), and frameworks in which the application is written and from which the libraries are called. Every vendor has had to build all five, and the layers map onto each other one for one. What differs is how long each layer has existed, how many people have optimized it, and how much of the world's code was written against it.

\emph{NVIDIA.} CUDA 1.0 was released in June 2007 (slide 29 of the previous deck traces its history). Nineteen years later the stack has a layer of libraries for every workload that has mattered: cuBLAS and cuBLASLt for matrix multiply, cuDNN for convolutions and attention, CUTLASS as the open-source template library from which the fastest kernels are built, TensorRT for inference graphs, NCCL for collectives across NVLink and InfiniBand, NVSHMEM for one-sided communication, and above them TensorRT-LLM, Dynamo and the vendor-maintained back ends of PyTorch and JAX. The compiler path is CUDA C++ through \texttt{nvcc} to PTX, a virtual instruction set, and then to the device's SASS; PTX is the stable target that lets a binary built for one generation run on the next, and it is also the target of Triton and of the other compilers that bypass CUDA C++. The runtime and driver are closed source, as is most of the library layer; CUTLASS and some tools are open. The breadth has a cost that slide 26 measures: even on NVIDIA's own hardware, a kernel reaches the peak only when it is one of the kernels NVIDIA's engineers tuned.

\emph{AMD.} ROCm (Radeon Open Compute) began in April 2016 as an open-source stack for GCN GPUs on Linux$^\dagger$. Its design decision was HIP: a C++ dialect and runtime API deliberately shaped like CUDA's, so that a CUDA program becomes a HIP program by a mechanical renaming (\texttt{cudaMalloc} to \texttt{hipMalloc}, \texttt{<<<grid, block>>>} launches kept as they are) that the \texttt{hipify} tools perform automatically, and so that one HIP source compiles for AMD through \texttt{hipcc} and LLVM and for NVIDIA by passing the calls through to CUDA. The library layer mirrors NVIDIA's name by name: rocBLAS and hipBLASLt for matrix multiply, MIOpen for deep-learning primitives, Composable Kernel as the template library, RCCL as the collective library compatible with NCCL's API, and the \texttt{amdgpu} kernel driver with the ROCr runtime implementing the HSA (Heterogeneous System Architecture) interfaces below. The versions on the slide mark the stack's phases: 5.0 in February 2022 with the MI200 generation; 6.0 on December 15, 2023, released with the MI300X; 7.0 on September 16, 2025, the release that rewrote the AI libraries after the 2024 criticism described on slide 26; and 10.0 on August 26, 2026, after a six-week release cadence through 7.1, 7.2 and up to 7.14 in 2026 and a renumbering that skipped 8 and 9$^\dagger$. The whole stack is open source on GitHub, which is both its advantage (a user can read and fix a kernel) and historically its weakness (AMD's own engineering was spread across more code than it had staff to tune).

\emph{Intel.} oneAPI was announced at Intel's Architecture Day in December 2018 and specified at version 1.0 on September 28, 2020, with the first production toolkits in December 2020$^\dagger$. Its language is SYCL, the Khronos Group's single-source C++ standard for heterogeneous computing, whose SYCL 2020 version (final specification February 9, 2021) added unified shared memory, a simpler kernel syntax and a formal notion of back ends; Intel's implementation is DPC++ (Data Parallel C++), built on LLVM and open source. Below it sits Level Zero, a low-level runtime API with the role of the CUDA driver API, and Intel's compute runtime; alongside it OpenMP offload for Fortran and C codes. The libraries are oneMKL (BLAS, LAPACK, FFT), oneDNN (deep-learning primitives, also used inside PyTorch on CPUs), oneCCL for collectives, and above them OpenVINO for inference and the Intel XPU back end of PyTorch. SYCLomatic (the open-source form of Intel's \texttt{dpct} tool) translates CUDA source to SYCL, and Intel's claim is that it migrates 80--90\% of a typical code automatically, leaving the rest to hand-editing$^\dagger$. In September 2023 the oneAPI specification was moved into the Unified Acceleration (UXL) Foundation under the Linux Foundation, with Arm, Fujitsu, Google, Imagination, Qualcomm and Samsung as founding members, so that its governance is formally not Intel's$^\dagger$. Intel's second accelerator line, Gaudi, has a separate stack, SynapseAI, with its own graph compiler and kernel library (slide 16); the slide's panel counts that split as a cost, because Intel's limited software effort was divided between two instruction sets that shared no code below the framework layer.

The caption's point about portability layers is the reason the three stacks are not three islands. Kokkos and RAJA, the C++ performance-portability libraries developed at Sandia and Livermore for the Department of Energy's machines, have back ends for CUDA, HIP and SYCL (and OpenMP for CPUs), so that one source compiles for all three; OpenMP offload does the same for directive-based codes; SYCL itself targets NVIDIA and AMD through the plugins Codeplay released on December 16, 2022 (Codeplay was acquired by Intel in 2022), which compile DPC++ kernels to PTX or to AMD's GCN/CDNA code; and Triton, the Python-embedded kernel language from OpenAI that PyTorch's compiler uses to generate fused kernels, has back ends for NVIDIA PTX, AMD (through LLVM's AMDGPU target) and Intel (through Intel's Triton fork)$^\dagger$. Slide 25 is about how well these layers perform; this slide only places them.

The panel's comparison of porting paths deserves a precise statement. HIP's path is syntactic: because the API is CUDA-shaped, most CUDA code compiles for AMD after \texttt{hipify}, and the work that remains is performance (warp size 64 against 32, different cache sizes, different library behavior) rather than correctness. SYCL's path is semantic: SYCL is a different language with a different memory model, and SYCLomatic produces a correct translation of the common patterns but leaves the uncommon ones and any inline PTX to a human. Both are far easier than the ports of a decade ago, and both still leave the performance tuning, which is where the time goes (slide 25).

The panel's framework row is the one that decides the AI market. PyTorch supports all three vendors, with the NVIDIA back end maintained jointly by NVIDIA and Meta, the ROCm back end by AMD, and the XPU back end by Intel; JAX supports NVIDIA and AMD. The serving systems vLLM and SGLang support NVIDIA and AMD. But the kernels that set the performance of a training run or a serving deployment (FlashAttention, the fused MoE kernels, the quantized GEMMs, speculative decoding) are written first for CUDA by their authors, who have NVIDIA hardware, and ported to ROCm later, by AMD or by the community; the deck's phrase is that ``the kernels that set performance arrive first for CUDA.'' AMD's response has been to fund the ports (AITER, its AI tensor library of 2025, is a collection of such kernels) and to make its hardware available to the authors, and the gap in 2026 is months rather than years. For Intel the frameworks are present but the serving systems and the kernel ecosystem largely are not, which is one reason the GPU D-series of slide 20 is positioned as an inference card with OpenVINO rather than as a training platform.

\begin{terms}
\term{software stack}{The layers between an application and a GPU: driver, runtime, compiler, libraries and frameworks; each vendor has one of each.}
\term{CUDA}{NVIDIA's platform (June 2007): the CUDA C++ language, the nvcc compiler, the PTX virtual instruction set, the runtime and driver, and the library layer above them.}
\term{PTX}{Parallel Thread Execution, NVIDIA's virtual instruction set; the stable compiler target that lets binaries carry forward across GPU generations and that Triton and other compilers emit.}
\term{ROCm}{Radeon Open Compute, AMD's open-source GPU software stack (ROC 1.0 April 2016; 6.0 December 2023; 7.0 September 2025; 10.0 August 2026).}
\term{HIP}{Heterogeneous-compute Interface for Portability: AMD's CUDA-shaped C++ dialect and runtime API, compiled by hipcc for AMD GPUs or passed through to CUDA for NVIDIA ones.}
\term{hipify}{AMD's source-translation tools (hipify-clang, hipify-perl) that rename CUDA API calls and types to their HIP equivalents.}
\term{HSA and ROCr}{The Heterogeneous System Architecture interfaces and AMD's runtime implementing them, the layer below HIP that talks to the amdgpu kernel driver.}
\term{oneAPI}{Intel's programming model and library specification (1.0 September 2020), built on SYCL with Level Zero beneath and oneMKL, oneDNN and oneCCL above; governed since 2023 by the UXL Foundation.}
\term{SYCL}{The Khronos Group's single-source C++ standard for heterogeneous computing; SYCL 2020 (February 2021) is the current version. Intel's implementation is DPC++.}
\term{Level Zero}{Intel's low-level device API, the counterpart of the CUDA driver API, beneath oneAPI's runtimes.}
\term{SYCLomatic}{The open-source CUDA-to-SYCL source translator (Intel's dpct); Intel claims 80--90\% automatic migration.}
\term{UXL Foundation}{The Unified Acceleration Foundation under the Linux Foundation (September 2023), which holds the oneAPI specification.}
\term{Kokkos and RAJA}{Department of Energy C++ performance-portability libraries (Sandia; Livermore) with back ends for CUDA, HIP, SYCL and OpenMP.}
\term{Triton}{OpenAI's Python-embedded language and compiler for GPU kernels, used by PyTorch's compiler; back ends for NVIDIA, AMD and Intel.}
\term{Codeplay plugins}{The December 2022 oneAPI plugins that let DPC++ compile SYCL for NVIDIA (through PTX) and AMD GPUs.}
\term{SynapseAI}{Intel's separate software stack for the Gaudi accelerators, with its own graph compiler and kernel library.}
\end{terms}

\begin{intuition}
Picture the three columns as the same building drawn three times, one nineteen years old, one ten, one six: the floors are the same, but the oldest has furniture in every room. Porting is no longer the hard part; the hard part is that the fastest kernels are written where the hardware is, and in 2026 that is still CUDA first.
\end{intuition}

% ---------------------------------------------------------------------
\slide{25}{Portability: what the SYCL and Kokkos studies measure}{HPC codes run on all three vendors through portability layers; the gaps are in maturity, not the model.}

\para{What the slide shows} A table of four studies occupies the upper left, with columns for the study, what was run and the finding. Reguly (SC Workshops 2023) ran SYCL against native code on the Intel Max 1100, the NVIDIA A100 and the AMD MI250X and found DPC++ within 1.2\% of CUDA on the A100, 15.9\% behind HIP on the MI250X, and weak on CPUs. Davis et al.\ (ICS 2025) ran seven programming models over five proxy applications on five GPUs and found that the portability layers reach most of native speed, with the spread in compilers and runtimes. ORNL/TM-2025/4398 ran Kokkos with SYCL back ends on AMD GPUs and found DPC++ less mature on AMD than on Intel, with HIP remaining the reference. The Aurora early-science projects (2023--25) ported DOE codes with SYCL, OpenMP and Kokkos and found that codes written for Summit and Frontier ran, with the tuning for the 512 KB shared local memory and the two-stack GPU taking the time. A panel below the table draws the lesson: the abstractions (a parallel-for over work-groups, local memory, atomics) map to all three vendors, and what differs is compiler maturity, library coverage and vendor tuning; one source runs everywhere and reaches peak where its author tuned. The right panel, ``How HPC and AI differ here'', explains why HPC codes are portable by construction (a decade of DOE funding for Kokkos, RAJA, OpenMP and SYCL ports because its machines use three vendors), why AI code is CUDA-first (written to CUDA kernels and ported later; a FlashAttention port to ROCm lags by months), the economics behind both (an HPC center buys for seven years and pays for a port once; an AI lab buys capacity by the quarter), and what might change it (Triton, PyTorch's compiler and vLLM making the vendor a back end; AMD's 2026 deals assume that is far enough along).

\para{Details} Performance portability is the property that one source code runs on several architectures at a useful fraction of each one's achievable performance, and the studies on the slide are attempts to measure it rather than assert it. The measurements matter for this lecture because they are the only evidence about whether the software gap of slide 24 is a property of the programming models or of their implementations, and the four studies agree that it is the latter.

\emph{Reguly (2023).} István Reguly's ``Evaluating the performance portability of SYCL across CPUs and GPUs'' (SC Workshops 2023, arXiv 2309.10075) took a set of benchmarks and mini-applications with existing CUDA, HIP and OpenMP versions, wrote SYCL versions, and ran them on an Intel Data Center GPU Max 1100 (the 56-core PCIe Ponte Vecchio of slide 17), an NVIDIA A100 and an AMD MI250X, as well as on Intel and AMD CPUs. The headline numbers the deck quotes are averages over the suite: the DPC++ SYCL code ran within 1.2\% of the native CUDA versions on the A100, through Codeplay's PTX plugin, and 15.9\% behind the native HIP versions on the MI250X, through the AMD plugin; on CPUs the SYCL versions were well behind OpenMP, because the SYCL compilers' CPU back ends vectorized poorly$^\dagger$. The reading is that SYCL on NVIDIA had, by 2023, a compiler path as good as CUDA's for these codes (both end in PTX and the same back-end compiler), while SYCL on AMD was still paying for a younger plugin and for the differences in how the two runtimes schedule work; and that portability to CPUs was the open problem, not portability between GPUs.

\emph{Davis et al.\ (2025).} ``Taking GPU programming models to task for performance portability'' by Joshua Davis and colleagues (ICS 2025, arXiv 2402.08950) is the broadest of the four: seven programming models, the vendor-native ones (CUDA, HIP) and the portability layers (Kokkos, RAJA, SYCL, OpenMP offload and OpenACC), applied to five proxy applications, on NVIDIA V100, A100 and H100 and AMD MI50 and MI250X GPUs$^\dagger$. The paper's method is the point: rather than report one portability number it reports the distribution of achieved performance for each model across the applications and machines, and asks where the losses come from. Its finding, as the deck summarizes it, is that the portability layers reach most of native performance in most cases, and that where they do not the cause is in the compiler (a missed unrolling, a register-spilling kernel) or the runtime (launch overhead, memory-allocation behavior) rather than in anything the abstraction prevents the programmer from expressing. The practical corollary is that a loss seen through a portability layer is usually recoverable by a compiler fix or a tuning flag, which is a different situation from a loss that would require a rewrite.

\emph{ORNL/TM-2025/4398.} The Oak Ridge technical report ``Experiences with SYCL on AMD GPUs with Kokkos,'' written with Sandia's Kokkos team, asks a narrower question that mattered to the DOE: if a code is written in Kokkos and Kokkos has a SYCL back end (built for Aurora), can that SYCL back end be used on Frontier's AMD GPUs instead of Kokkos's native HIP back end, so that one back end serves two machines? The answer was not yet. Both SYCL implementations tested, Intel's DPC++ through the Codeplay plugin and the independent AdaptiveCpp (formerly hipSYCL), ran the Kokkos tests on AMD GPUs, but DPC++ was less mature on AMD than on Intel, with gaps in features and performance, and the HIP back end remained the reference for production use on Frontier$^\dagger$. This is the maturity gap stated most concretely: the same compiler that matched CUDA on an A100 in Reguly's study was, two years later, still catching up on AMD.

\emph{Aurora early science.} The fourth row is not one paper but the experience of the Aurora Early Science Program and the Exascale Computing Project's application teams between 2023 and 2025, as reported in their papers and at ALCF workshops$^\dagger$. The codes in question had already been ported once, to Summit's V100s (CUDA or OpenACC) and then to Frontier's MI250Xs (HIP, or Kokkos and RAJA with HIP back ends), and the Aurora port was the third. The finding was that the codes written in the portability layers ran on Aurora with modest effort, because Kokkos, RAJA and OpenMP had SYCL and Intel back ends by then; the codes written directly in CUDA went through SYCLomatic and then through hand work. What took the time was performance: Ponte Vecchio's 512 KB of L1 and shared local memory per Xe-core, four times an MI250X compute unit's and twice an H100 SM's (slide 18), rewarded different blocking; its two stacks, exposed either as one device or as two, required a decision about how to decompose work; and the Xe-core's XVE needed SIMD-width choices (slide 14) that the other architectures fix in hardware. These are tuning problems, not portability problems, and the teams that had them also had working codes.

The lesson panel generalizes. Every one of these programming models offers a parallel loop over work-items grouped into work-groups (CUDA's threads and blocks, HIP's work-items and work-groups, SYCL's the same, Kokkos's parallel\_for over a team policy), a fast per-group memory (shared memory, LDS, shared local memory) and atomics, because every GPU of the three vendors has those things (slides 7, 13 and 14). The abstractions therefore map; what differs is how good the compiler is at turning them into each vendor's instruction set, how many of the libraries the code calls exist and are tuned on each vendor, and how much the vendor's own engineers have worked on the code in question. One source runs everywhere; it reaches peak where its author, or the vendor, tuned.

The right panel turns this into the lecture's central contrast. HPC codes are portable by construction because the DOE made them so: its leadership machines since 2017 have been NVIDIA (Summit, Perlmutter), AMD (Frontier, El Capitan) and Intel (Aurora), and the Exascale Computing Project spent seven years and about 1.8 billion dollars, a large part of it on applications and software, funding Kokkos, RAJA, OpenMP and SYCL ports precisely so that the laboratories could buy from whichever vendor offered the most FP64 per dollar and per watt (slide 12). The same applications run on all five machines. AI code is CUDA-first because it was written at companies and universities that bought NVIDIA GPUs; its kernels are written in CUDA or Triton against CUDA hardware, and the port to ROCm follows by months, as it did for FlashAttention and the fused mixture-of-experts kernels, and to oneAPI mostly not at all. The economics explain the difference without any judgment about the engineers: an HPC center procures a machine for seven years and pays for the port once, so a 15\% software gap at delivery is a nuisance that the center's own staff will close; an AI laboratory rents or buys capacity by the quarter and runs the same few models, so a 15\% gap today is 15\% of today's bill and the time to close it is time its competitors do not spend. That, the panel says, is why the TOP500 is AMD and Intel while the AI clusters are NVIDIA, and it is a more exact statement than any about the hardware.

What would change it is the last point. If the application is written in PyTorch and compiled through Triton, and served through vLLM, then the vendor is a back end chosen at install time, and the kernels the author wrote are generated rather than ported. PyTorch's compiler and Triton have made this true for a growing share of training and serving code, and AMD's 2026 commitments from OpenAI, Meta, Oracle and Anthropic (slide 29) are bets by buyers with the largest CUDA code bases that the share is large enough. Slide 26 is what was measured before those bets were made.

\begin{terms}
\term{performance portability}{The property that one source runs on several architectures at a useful fraction of each one's achievable performance; measured, not assumed.}
\term{proxy application}{A small program that reproduces the computational pattern of a large scientific code, used to compare machines and programming models.}
\term{DPC++}{Data Parallel C++, Intel's open-source LLVM-based SYCL compiler, the implementation used in the Reguly and ORNL studies.}
\term{AdaptiveCpp}{The independent SYCL implementation formerly called hipSYCL, with back ends for AMD, NVIDIA and Intel GPUs.}
\term{work-group and work-item}{SYCL's and HIP's names for a group of threads scheduled together on one compute unit and for one thread within it; CUDA's block and thread.}
\term{shared local memory (SLM)}{Intel's name for the per-work-group on-chip memory; 512 KB per Xe-core on Ponte Vecchio, the figure the Aurora tuning turned on.}
\term{Exascale Computing Project}{The DOE program (2016--2024) that funded the applications, libraries and programming models for Frontier, Aurora and El Capitan, including the ports that make HPC codes portable.}
\term{Aurora Early Science Program}{The ALCF program that gave application teams early access to Aurora to port and tune their codes (2023--25).}
\term{maturity gap}{A performance loss attributable to a young compiler, runtime or library rather than to the programming model; recoverable without rewriting the code.}
\term{CUDA-first}{The pattern in AI software of writing and tuning kernels for CUDA and porting to other stacks afterwards.}
\end{terms}

\begin{intuition}
The studies all say the same thing in different units: the abstractions are portable and the implementations are uneven, so a port costs tuning, not rewriting. Whether anyone pays for the tuning is an economic question, and the answer differs for a laboratory that keeps a machine seven years and a company that rents one by the quarter.
\end{intuition}

% ---------------------------------------------------------------------
\slide{26}{Achieved against peak: the software gap in numbers}{In 2024 independent tests put MI300X matrix multiply at half its peak, the H100 at three quarters.}

\para{What the slide shows} A grouped bar chart on the left compares, for the H100, the H200 and the MI300X, the datasheet dense BF16 peak (gray) with the BF16 matrix-multiply throughput SemiAnalysis measured in December 2024 (brown): 989 against about 720 TFLOP/s for the H100 and again for the H200, and 1{,}307 against about 620 for the MI300X, with the values printed above the bars and a legend inside the plot. The right panel gives the context: SemiAnalysis's December 2024 study found the H100 and H200 at about 73\% of peak on large BF16 matrix multiplies and the MI300X at about 47\%, multi-node H100 training 10--25\% faster, and GPU-to-GPU bandwidth of about 64 GB/s per pair on Infinity Fabric against 450 on NVLink, under the title ``The CUDA moat is still alive'': competitive hardware, immature software. AMD's response was ROCm 6.3 through 7.0, which rewrote the matrix-multiply and attention libraries (hipBLASLt, Composable Kernel, AITER), a six-week release cadence and the MLPerf submissions of slide 30. The rule the panel draws: a peak is a property of the die, the achieved fraction is a property of the libraries, and it rises over a product's life, so measured throughput should be compared at the same date.

\para{Details} The slide takes the most-cited independent measurement of the 2024 AMD--NVIDIA contest and uses it to make a general point about peaks. The source is SemiAnalysis's ``MI300X vs H100 vs H200 benchmark part 1: training,'' published on December 22, 2024 by Dylan Patel, Daniel Nishball and Reyk Knuhtsen after five months of testing on hardware provided by several cloud operators$^\dagger$. The study ran microbenchmarks and training workloads on eight-GPU systems of each type and on multi-node clusters, with the vendor-recommended public software and, for AMD, also with development builds the company supplied during the study.

The matrix-multiply result is the one the chart shows. On large BF16 general matrix multiplies (GEMMs) with FP32 accumulation, the shapes that dominate transformer training, the H100 and H200 achieved about 720 TFLOP/s against the datasheet's 989 dense, a fraction of $720/989 = 73\%$; the MI300X achieved about 620 TFLOP/s against its 1{,}307, a fraction of $620/1{,}307 = 47\%$ (derived)$^\dagger$. The MI300X's measured throughput was therefore below the H100's despite a peak 32\% higher, and the gap between the two fractions, 26 points, is the software gap in its purest form, because a GEMM is one kernel with no communication and no framework overhead; it measures how well the vendor's BLAS library maps the operation onto the matrix units. The H100 and H200 have the same die and the same peak and achieved the same throughput, as expected; the H200's larger memory does not change a compute-bound kernel.

The other two numbers in the panel widen the picture. Multi-node training of real models ran 10--25\% faster on the H100 clusters than on the MI300X clusters, a smaller gap than the GEMM figures alone would predict, because training also spends time in memory-bound operations where the MI300X's 5.3 TB/s helps and in communication$^\dagger$. The communication measurement, about 64 GB/s between a pair of MI300X GPUs over Infinity Fabric against about 450 GB/s between a pair of H100s over NVLink and NVSwitch, is the point-to-point test discussed on slide 21: one 128 GB/s bidirectional link against a switched fabric that gives a pair the GPU's full 450 GB/s in one direction$^\dagger$. For the all-reduce collectives that training actually uses the gap was smaller but present, and the study attributed part of it to RCCL's tuning relative to NCCL's.

The study's title, ``The CUDA moat is still alive,'' and its summary, that the MI300X was competitive hardware held back by immature software, are what the deck quotes, and they were accepted by AMD. The study also said something the chart cannot show: that AMD's development builds improved during the five months, that AMD engineers fixed several of the bugs the authors reported within days, and that the public ROCm release at the end of the study (6.2 to 6.3) was already better than the one at its start. That is the mechanism the panel's rule names.

AMD's response in 2025 and 2026 was organizational as much as technical. ROCm 6.3 (late 2024) and 7.0 (September 2025) replaced the matrix-multiply path with hipBLASLt as the default for deep-learning shapes, rewrote the attention and GEMM kernels in Composable Kernel, and added AITER, a library of hand-tuned kernels for the operations that dominate inference (attention, mixture-of-experts dispatch, quantized GEMM) with the explicit aim of matching the CUDA kernels the AI frameworks call$^\dagger$. The company moved to a six-week release cadence, so that the 7.x series ran to 7.14 by July 2026 before the renumbering to 10.0 in August (slide 24), and it began submitting to MLPerf: inference in August 2024, training in June 2025, with the MI355X results of 2025--26 landing within a few percent of the B200 on the same benchmarks (slide 30)$^\dagger$. Whether the MI355X's achieved GEMM fraction in 2026 is closer to 47\% or to 73\% is not something the deck claims to know, because no equally careful independent measurement had been published by November 2026; the MLPerf results are the nearest public evidence and they say the gap on those workloads is small.

The rule the panel ends with is the thing to teach. A peak is the product of unit count, operations per clock and clock, and it is fixed when the die is; it is also the number both vendors print largest. The achieved fraction is a property of the libraries, the compiler and the runtime, and it changes with every release; it is what the user gets. The two are related only through engineering effort, and the effort is spent over a product's life, so a GPU measured at launch is measured at its worst. This is as true of NVIDIA as of AMD (the H100's own GEMM efficiency rose with CUTLASS 3 and cuBLAS releases through 2023), and it is why the deck's comparisons of measured performance (slides 26 and 30) always carry a date. The honest comparison of two GPUs is measured throughput on the user's own workload, with each vendor's current software, on the same day. Slide 30 shows the one public venue where that happens.

A derived footnote on the arithmetic of the chart: the MI300X's measured 620 TFLOP/s on 304 compute units at 2.1 GHz corresponds to $620\times10^{12} / (304 \times 2.1\times10^{9}) = 971$ FLOP per CU per clock, against a peak of $1{,}307\times10^{12} / (304 \times 2.1\times10^{9}) = 2{,}047$, that is, 2{,}048 FLOP (1{,}024 FMA) per CU per clock for BF16, exactly the CDNA 3 whitepaper's figure; the H100's 720 TFLOP/s on 132 SMs at a tensor clock of about 1.83 GHz is $720\times10^{12} / (132 \times 1.83\times10^{9}) = 2{,}981$ FLOP per SM per clock against a peak of 4{,}096 (derived). Both machines were leaving matrix-unit cycles idle; the MI300X was leaving more, and the reasons (instruction issue, operand delivery from the LDS and registers, the GEMM kernel's tiling) are the things a BLAS library controls.

\begin{terms}
\term{achieved fraction}{Measured throughput divided by datasheet peak for a given operation; a property of the software and the date as much as of the die.}
\term{GEMM}{General matrix multiply, $C \leftarrow \alpha AB + \beta C$; the kernel that dominates transformer training and the standard test of a matrix unit and its library.}
\term{BF16 with FP32 accumulation}{The usual training arithmetic: 16-bit brain-float operands multiplied and summed in 32-bit; the mode whose dense peaks are 989 (H100) and 1,307 (MI300X) TFLOP/s.}
\term{SemiAnalysis study}{The December 2024 independent benchmark of MI300X, H100 and H200 training by Patel, Nishball and Knuhtsen; the source of the 720 and 620 TFLOP/s figures.}
\term{hipBLASLt}{AMD's lightweight BLAS library for the matrix-multiply shapes and data types of deep learning, the counterpart of cuBLASLt; made the default path in ROCm 6.3--7.0.}
\term{Composable Kernel (CK)}{AMD's C++ template library for writing tensor kernels, the counterpart of CUTLASS; rewritten for the 2025 releases.}
\term{AITER}{AMD's AI Tensor Engine for ROCm (2025), a library of hand-tuned kernels for attention, mixture-of-experts and quantized GEMM meant to match the CUDA kernels frameworks call.}
\term{RCCL and NCCL}{The AMD and NVIDIA collective-communication libraries; RCCL implements NCCL's API over Infinity Fabric and Ethernet.}
\term{release cadence}{The interval between software releases; AMD's six-week ROCm cadence from 2025 is the organizational form of the panel's rule that achieved fractions rise over time.}
\term{same-date comparison}{Comparing two GPUs' measured throughput on the same workload with each vendor's software of the same date, the only comparison the deck treats as fair.}
\end{terms}

\begin{intuition}
The datasheet tells you the size of the engine; the measurement tells you how much of it the software could use that month. In 2024 AMD's engine was larger and its software used less than half of it; the right response was neither to trust the peak nor to dismiss the part, but to measure again after the next release, which is what the MLPerf rounds on slide 30 do.
\end{intuition}
% =====================================================================
\setstage{outc}
\section{Systems and procurements (slides 27--29)}

\slide{27}{The exascale systems: Frontier, Aurora, El Capitan, JUPITER}{Four machines, three GPU vendors, three node designs; the FP64 side of the comparison, measured by HPL.}

\para{What the slide shows} A table of the four GPU-accelerated exascale systems on the June 2026 TOP500 list, one column each, and a panel describing their node designs. The rows give the site (ORNL, Argonne, LLNL, J\"ulich), the system family (three HPE Cray EX machines and one Eviden BullSequana), the node (one CPU and four MI250X; two CPUs and six Ponte Vecchio; four MI300A APUs; four GH200 superchips), the node count (9,408; 10,624; about 11,100; about 6,000), the accelerator count (37,632; 63,744; about 44,500; about 24,000), the first HPL result and its date (1.102 EFLOP/s in 2022; 0.585 in 2023; 1.742 in 2024; 0.793 in 2025), the June 2026 HPL result and rank (1.353 EFLOP/s, third; 1.012, fourth; 1.809, second; 1.000, fifth), the theoretical peak Rpeak (2.06, 1.98, 2.82, 1.23 EFLOP/s), the listed power (24.6, 38.7, 29.7, 15.8 MW), the derived HPL result per megawatt (55, 26, 61, 63 PFLOP/s per MW) and the network (Slingshot for the three Cray machines, InfiniBand for JUPITER). The panel, ``Three node designs'', describes the Frontier, Aurora, El Capitan and JUPITER nodes and closes with what HPL rewards: FP64 matrix throughput and bandwidth per watt. A footnote gives the provenance of the counts.

\para{Details} The slide is the FP64 half of the lecture's comparison, and it is deliberately placed after the fabrics and software sections, because the four machines are the only places where the three vendors' data-center GPUs have been built into systems of the same class, run the same benchmark and published the result. HPL, the High-Performance Linpack benchmark, solves a dense system of linear equations in FP64 by LU factorization with partial pivoting; its measured rate, Rmax, is the TOP500's ranking metric, and the ratio of Rmax to the theoretical peak Rpeak is the HPL efficiency. HPL rewards exactly two things: the rate at which the accelerators multiply FP64 matrices and the ability of the memory system and network to keep them fed through the panel broadcasts and trailing updates of the factorization. It is the one workload on which the GPUs of this lecture are compared at full scale on equal terms, which is why the deck gives it a slide even though no AI buyer runs it.

Frontier at Oak Ridge National Laboratory was the first machine to exceed one exaflop on HPL, with 1.102 EFLOP/s in June 2022, and it has been upgraded since: 1.194 EFLOP/s in May 2023 and 1.353 EFLOP/s from November 2024, the figure it still carries on the June 2026 list in third place$^\dagger$. It is an HPE Cray EX235a system of 9,408 compute nodes in 74 cabinets; each node holds one AMD EPYC 7A53 ``Trento'' CPU of 64 cores and four MI250X modules, joined by coherent Infinity Fabric at 36 GB/s in each direction per CPU--GPU link, with the four modules connected to one another by Infinity Fabric at 50 GB/s each way and four Slingshot-11 network interfaces at 200 Gb/s each (slide 8)$^\dagger$. Because each MI250X holds two graphics compute dies that the software sees as two GPUs, Frontier's 37,632 modules are 75,264 GPU devices to ROCm, which is why Frontier's documentation speaks of ``eight GPUs per node''. The TOP500 lists Frontier's Rpeak at 2.06 EFLOP/s and its power at 24.6 MW (24,607 kW on the June 2026 list)$^\dagger$, so its HPL efficiency is about 66\% and its HPL result per megawatt about 55 PFLOP/s (both derived).

Aurora at Argonne National Laboratory is the Intel system and the largest by GPU count: 10,624 nodes in 166 HPE Cray EX racks, each node with two Xeon Max 9470C processors (52 cores each, with 64 GB of HBM2e on the CPU package and 512 GB of DDR5) and six Data Center GPU Max 1550 (Ponte Vecchio) GPUs, for 63,744 GPUs and 21,248 CPUs, with eight Slingshot-11 interfaces per node on a dragonfly network (slide 19)$^\dagger$. Its first HPL submission, 585.34 PFLOP/s in November 2023 on 5,439 nodes, placed second while the machine was still being brought up; in June 2024 a run on 9,234 nodes reached 1.012 EFLOP/s, again second, and that figure has stood through the November 2024, 2025 and June 2026 lists while the rank slipped to third and then fourth as El Capitan and LineShine arrived$^\dagger$. The list gives Rpeak 1.980 EFLOP/s and power 38,698 kW, so Aurora runs HPL at about 51\% of peak and about 26 PFLOP/s per MW (derived), the lowest efficiency of the four. The research behind the deck notes a common error here: the 24.69 MW sometimes quoted for Aurora was the power of the partial November 2023 run, not of the full system, and the deck uses the June 2026 list's 38.7 MW. Slide 19 explains why the efficiency is low: PVC's FP64 runs at its FP32 rate of 52.4 TFLOP/s with no separate matrix path, its 600 W package was a year late against Hopper, and the 47-tile package spent power on moving data between tiles.

El Capitan at Lawrence Livermore National Laboratory is the AMD APU machine and the fastest accelerated system in the world as of June 2026: 1.742 EFLOP/s in November 2024, when it took first place, and 1.809 EFLOP/s from November 2025, which kept first place until LineShine in June 2026 pushed it to second$^\dagger$. Its node has no discrete CPU at all: four MI300A accelerated processing units, each combining 24 Zen 4 cores and 228 CDNA 3 compute units on the same 3D-stacked package with 128 GB of HBM3 shared by both (slide 9), connected to one another by Infinity Fabric and to the network by Slingshot-11$^\dagger$. The node count on the November 2024 list was 11,136 (44,544 APUs, with 43,808 used in the HPL run); LLNL's system page now lists 11,424 batch nodes plus 64 debug and 32 login nodes, 46,080 APUs in all, and a peak of 2,889 PFLOP/s, while the TOP500 entry gives Rpeak 2.821 EFLOP/s$^\dagger$. The deck therefore writes ``about 11,100 nodes'' and ``about 44,500 accelerators'' and explains the spread in its footnote. At 29.7 MW (29,685 kW listed), El Capitan's HPL efficiency is about 64\% and its result per megawatt about 61 PFLOP/s (derived), the best of the Department of Energy's three exascale machines, which is the panel's point about the APU design: with the CPU on the package, the CPU's memory traffic no longer crosses a board-level link and the memory capacity is shared rather than duplicated.

JUPITER at Forschungszentrum J\"ulich is the NVIDIA system and Europe's first exascale machine: its Booster module, built by Eviden on the BullSequana XH3000 platform, holds about 24,000 GH200 Grace Hopper superchips in nodes of four, roughly 6,000 nodes (derived from four per node; the FZJ press release gives the superchip count). Each GH200 joins a 72-core Grace CPU to an H100 GPU over NVLink-C2C at 900 GB/s, so the node has four CPUs and four GPUs with no x86 host, and the nodes are connected by NVIDIA Quantum-2 InfiniBand at 400 Gb/s$^\dagger$. JUPITER's first HPL result, 793 PFLOP/s in June 2025, placed fourth; the 1.000 EFLOP/s run announced on November 17, 2025 made it the first European system above an exaflop (fourth on the November 2025 list), and it is fifth in June 2026 at 15.8 MW (15,794 kW)$^\dagger$. With Rpeak 1.23 EFLOP/s its efficiency is about 81\% and its result per megawatt about 63 PFLOP/s (derived), the best of the four on both measures. Two things explain that: the H100's FP64 tensor cores (67 TFLOP/s per GPU) are the most efficient FP64 matrix engine of the four at the clock and power the GH200 runs at, and the system is the newest, so its HPL run benefited from the most tuning per node.

The ``Three node designs'' panel draws the architectural comparison. Frontier is the conventional design: a host CPU with four accelerators behind it, distinguished from the DGX-style design by the coherent CPU--GPU fabric, which lets the GPUs reach host memory without PCIe transfers, and by the two-die module, which gives eight devices per node. Aurora doubles the host: two CPUs with their own HBM and six GPUs, with the GPUs connected to one another by Xe Link rather than through a switch, and with eight NICs because each GPU pair gets its own path to the network. El Capitan removes the discrete CPU entirely: four APUs, each a complete node in miniature, so that CPU and GPU share one physical memory and the question ``which memory is this in?'' disappears for the programmer. JUPITER does the same thing with two chips on a board: a Grace and an H100 joined at 900 GB/s, four such superchips per node. The lesson is that in 2022--25 the three vendors converged on tightly coupled CPU--GPU designs from three directions---coherent fabric, on-package HBM for the CPU, APU, superchip---because the HPC codes these machines run are not purely GPU codes and move data between the two sides constantly.

The last row of the panel, what HPL rewards, is the caveat for the whole slide. HPL is dense FP64 matrix multiply at the highest occupancy any application ever reaches, and its efficiency is set by the FP64 matrix rate, the memory bandwidth needed to stream panels, and the network's ability to broadcast them. It does not exercise FP16 or FP8, sparse or irregular access, or the latency-bound communication of a real application, so a ranking by HPL is a ranking by one capability that the AI buyers of slides 29 and 30 do not purchase. HPCG, the conjugate-gradient benchmark that stresses memory bandwidth, and HPL-MxP, which runs the factorization in low precision and refines the result to FP64 accuracy, are the TOP500's own correctives (slide 28).

\begin{terms}
\term{HPL}{High-Performance Linpack: a dense FP64 LU factorization and solve; its measured rate, Rmax, ranks the TOP500.}
\term{Rmax and Rpeak}{The achieved HPL rate of a system and its theoretical FP64 peak; their ratio is the HPL efficiency.}
\term{exascale system}{A system exceeding $10^{18}$ FP64 FLOP/s on HPL; five were on the June 2026 list, four of them GPU-accelerated.}
\term{Frontier}{Oak Ridge's HPE Cray EX system of 9,408 nodes with one EPYC 7A53 and four MI250X each; first above an exaflop (June 2022), 1.353 EFLOP/s in 2026.}
\term{Aurora}{Argonne's HPE Cray EX system of 10,624 nodes with two Xeon Max 9470C and six Data Center GPU Max 1550 each; 1.012 EFLOP/s at 38.7 MW.}
\term{El Capitan}{Livermore's HPE Cray EX system of about 11,100 nodes of four MI300A APUs; 1.809 EFLOP/s, first in the world in 2024--25, second in June 2026.}
\term{JUPITER}{J\"ulich's Eviden BullSequana system of about 24,000 GH200 superchips; Europe's first exaflop (November 2025), 1.000 EFLOP/s at 15.8 MW.}
\term{HPE Cray EX}{HPE's liquid-cooled supercomputer platform (formerly Cray Shasta) with the Slingshot network; Frontier, Aurora and El Capitan are Cray EX machines.}
\term{Slingshot}{HPE Cray's Ethernet-based high-performance network; Slingshot-11 interfaces run at 200 Gb/s.}
\term{BullSequana}{Eviden's (formerly Atos) supercomputer platform; JUPITER's Booster is a BullSequana XH3000 system.}
\term{EPYC 7A53 Trento}{The 64-core AMD CPU made for Frontier, with coherent Infinity Fabric links to its four MI250X.}
\term{Xeon Max 9470C}{Intel's 52-core CPU with 64 GB of HBM2e on the package, two per Aurora node.}
\term{coherent CPU--GPU fabric}{A link over which CPU and GPU share one address space with hardware cache coherence, so either side can read the other's memory directly (Infinity Fabric on Frontier and El Capitan, NVLink-C2C on JUPITER).}
\term{APU}{Accelerated processing unit: AMD's name for a package combining CPU cores and GPU compute units sharing one memory; the MI300A.}
\term{dragonfly network}{A topology of all-to-all connected groups of switches, each group connected to every other; Slingshot systems use it.}
\term{HPL efficiency}{Rmax divided by Rpeak; about 66\% on Frontier, 51\% on Aurora, 64\% on El Capitan and 81\% on JUPITER (derived).}
\term{HPL per megawatt}{Rmax divided by the listed power; a crude energy-efficiency measure that the Green500 formalizes.}
\end{terms}

\begin{intuition}
Four machines bought for the same job, built three ways, measured by one benchmark. The table says that the newest design wins on efficiency (JUPITER), the APU wins on absolute result (El Capitan), and the machine with the most GPUs is the least efficient (Aurora). Before drawing a conclusion about the GPUs, ask how old each machine is, how its CPU is attached, and how much of what you care about looks like HPL.
\end{intuition}

% ---------------------------------------------------------------------
\slide{28}{TOP500 and Green500, June 2026: who runs what}{A CPU-only Chinese system leads; AMD holds four of the top eight; the four most efficient are Grace Hopper.}

\para{What the slide shows} Two tables and a panel. The upper table is the top ten of the June 2026 TOP500 list with rank, system and site, accelerator, Rmax in PFLOP/s and power in MW: LineShine (Shenzhen, CPU only, 2,198, 42.2 MW); El Capitan (LLNL, AMD MI300A, 1,809, 29.7); Frontier (ORNL, AMD MI250X, 1,353, 24.6); Aurora (Argonne, Intel Max 1550, 1,012, 38.7); JUPITER (FZJ, NVIDIA GH200, 1,000, 15.8); HPC7 (Eni, AMD MI300A, 572, 8.7); Eagle (Microsoft Azure, NVIDIA H100, 561, power not listed); HPC6 (Eni, AMD MI250X, 478, 8.5); Fugaku (RIKEN, CPU only, 442, 29.9); Alps (CSCS, NVIDIA GH200, 435, 7.1). The lower table gives five Green500 entries in GFLOP/s per watt: KAIROS (GH200, 73.3), ROMEO-2025 (GH200, 70.9), the Levante GPU extension (GH200, 69.4) and Isambard-AI phase 1 (GH200, 68.8) in the first four places, and Ouranos (MI300A, 66.5) in ninth. The panel, ``Reading the lists'', summarizes: five exascale systems and three vendors of accelerator plus one of none; AMD in four of the top eight; Intel's one system the least efficient of the exascale group; NVIDIA's entries all GH200 or H100, the Green500's top four Grace Hopper followed by three H100 systems, and no Blackwell or Rubin system in the top ten; and the HPL-MxP order.

\para{Details} The TOP500 is published twice a year, in June at the ISC conference in Germany and in November at SC in the United States, and ranks the systems whose owners choose to submit an HPL result. Its June 2026 edition is the most recent at the time of these notes (the November 2026 list is released at SC26 in the second half of the month), and it is the first list led by a Chinese system since 2017. LineShine, at the National Supercomputing Center in Shenzhen, reached 2.198 EFLOP/s on HPL with CPUs alone---the list names the processor as the LX2 with 304 cores---at 42.2 MW (42,220 kW), and it also leads the HPCG list at 22.0 PFLOP/s and places fourth on HPL-MxP at 7.92 EFLOP/s$^\dagger$. Chinese centers had largely stopped submitting new systems after 2021 while exascale-class machines were described in conference papers; LineShine's return to the list is a political event as much as a technical one, and for this lecture its significance is negative: the fastest HPL machine in the world has no GPU in it, which is a reminder that the benchmark measures FP64 matrix arithmetic, not the capability the rest of this deck is about.

The next four systems are the four of slide 27, and the panel's first observation follows from them: five systems above an exaflop, built on three vendors' accelerators and one vendor's CPUs, no two alike. AMD's position is the strongest it has held on any list: El Capitan (MI300A) second, Frontier (MI250X) third, and Eni's two machines in Italy, HPC7 on MI300A at 571.5 PFLOP/s (sixth) and HPC6 on MI250X at 477.9 PFLOP/s (eighth), give AMD four of the top eight$^\dagger$. All four are HPE Cray EX systems, and all four were bought for FP64 simulation---nuclear stockpile stewardship at Livermore, open science at Oak Ridge, seismic imaging and reservoir modeling at Eni---which is the pattern the deck's AMD section predicted: CDNA's doubled matrix FP64 (slide 12) and the coherent CPU--GPU designs (slide 9) were built to the specifications of exactly these buyers. Eni's two machines deserve a note because they are commercial, not governmental: an energy company bought the sixth- and eighth-fastest computers in the world for its own exploration workloads, both on AMD, at 8.7 and 8.5 MW$^\dagger$.

Intel's position is the opposite. Aurora is its one system in the top ten and the least efficient of the exascale group at 26 PFLOP/s per MW (derived), and no Ponte Vecchio system has been added to the list since the two of November 2023, Dawn at Cambridge (1,024 GPUs, 19.5 PFLOP/s, 41st) and SuperMUC-NG Phase 2 at LRZ (960 GPUs, 17.2 PFLOP/s, 52nd)$^\dagger$. Slide 20 explains why there will be no more: Rialto Bridge and Falcon Shores were cancelled and the GPU D-series is an inference card without FP64 ambitions. NVIDIA's position is in between. Its three top-ten systems are JUPITER and Alps on GH200 and Eagle on H100; Eagle, Microsoft Azure's 561.2 PFLOP/s cloud system, is the only hyperscaler machine on the list and reports no power figure$^\dagger$. The panel's observation that no Blackwell or Rubin system is in the top ten is the most important fact on the slide for a reader of 2026: eighteen months after the first GB200 NVL72 racks and three months after the first Vera Rubin racks, none of their large owners has put one near the top of the list, and slide 12 says why they would not want to (GB300 at 1.4 TFLOP/s FP64 per GPU gives about 100 TFLOP/s per rack).

The Green500 ranks the same submitted systems by HPL result per watt, using a power measurement made during the HPL run under the list's rules (which permit measuring a representative fraction of the system). Its June 2026 top four are all GH200 systems: KAIROS at 73.28 GFLOP/s per watt, ROMEO-2025 at 70.91, the Levante GPU extension at 69.43 and Isambard-AI phase 1 at 68.84; places five to seven are H100 SXM systems (Otus at 68.18, Capella, SSC-24) and eighth is another GH200 system$^\dagger$. The best AMD entries are two MI300A systems, Ouranos at 66.46 (ninth) and Portage at 66.28 (tenth); El Capitan itself is 28th at 60.94 and Frontier 20th$^\dagger$. The ordering says three things. First, the Grace Hopper superchip's H100, with its FP64 tensor cores at 67 TFLOP/s and a CPU that draws tens rather than hundreds of watts, is the most efficient FP64 engine on the list, by about 10\% over the MI300A (73 against 66, derived); the APU's advantage in El Capitan's absolute result does not carry over to efficiency, partly because the Green500 leaders are small systems whose power measurement excludes much of the infrastructure a 30 MW machine must count. Second, the small systems at the top of the Green500 and the large systems at the top of the TOP500 differ by about 9\% in efficiency for the same chip (Isambard-AI 68.8 against JUPITER's 63.3 PFLOP/s per MW from the TOP500 power column, derived), which is the cost of scale: more network, more cooling, more idle margin. Third, no system of 2025--26 silicon appears on either list, so the efficiency comparison is between the 2022--23 generations of NVIDIA and AMD and does not include Blackwell, CDNA 4 or Rubin.

The panel's last line concerns HPL-MxP, the mixed-precision variant of HPL in which the factorization is done in a low precision (FP16 or FP8 on tensor and matrix cores) and the solution is refined iteratively to FP64 accuracy. Its June 2026 order is El Capitan first, Frontier second, Aurora third (11.6 EFLOP/s, from its November 2024 run on 9,500 nodes) and LineShine fourth at 7.92 EFLOP/s$^\dagger$. HPL-MxP rewards the matrix units at their fast formats and the software that drives them, and the fact that Aurora places third on it with a system that is fourth on HPL and far down on efficiency shows that PVC's XMX engines, a generation behind on paper (slide 18), are usable at scale when Intel's libraries are tuned for the benchmark. The deck uses the HPL-MxP order as a bridge to slide 30: once the arithmetic is done in low precision, the ranking changes, and the question becomes which vendor's software extracts its hardware's peak.

A reader should also keep in mind what the lists leave out. Submission is voluntary and HPL is expensive to run at full scale, so the list undercounts the cloud providers and AI laboratories entirely (Eagle is the exception), undercounts China, and overcounts public laboratories. Power figures are self-reported and some systems report none. And the list is a record of machines bought two to four years before the list date, so the June 2026 list describes the 2022--24 procurement cycle; the systems of slide 29 will appear, if they appear at all, in 2027--28.

\begin{terms}
\term{TOP500}{The twice-yearly ranking of submitted systems by HPL result; June (ISC) and November (SC).}
\term{Green500}{The TOP500's companion list ranking the same systems by HPL result per watt of measured power.}
\term{HPCG}{High-Performance Conjugate Gradients: a sparse iterative benchmark that stresses memory bandwidth and network latency; LineShine led it in June 2026 at 22.0 PFLOP/s.}
\term{HPL-MxP}{Mixed-precision HPL: LU factorization in low precision on matrix units with iterative refinement to FP64 accuracy; El Capitan, Frontier and Aurora led it in June 2026.}
\term{LineShine}{The CPU-only system at the National Supercomputing Center in Shenzhen that led the June 2026 TOP500 at 2.198 EFLOP/s and 42.2 MW; the first Chinese number one since 2017.}
\term{HPC6 and HPC7}{Eni's two HPE Cray EX systems in Italy, on MI250X (478 PFLOP/s, eighth) and MI300A (572 PFLOP/s, sixth) in June 2026.}
\term{Eagle}{Microsoft Azure's H100 system, 561 PFLOP/s, the only hyperscaler machine in the June 2026 top ten.}
\term{Fugaku}{RIKEN's CPU-only system on Fujitsu A64FX processors, first in 2020--21, ninth in June 2026 at 442 PFLOP/s.}
\term{Alps}{The CSCS system in Lugano on GH200 superchips, 435 PFLOP/s, tenth in June 2026.}
\term{KAIROS, ROMEO-2025, Levante, Isambard-AI}{The four GH200 systems that led the June 2026 Green500 at 68.8--73.3 GFLOP/s per watt.}
\term{Ouranos and Portage}{The two MI300A systems that were AMD's best Green500 entries in June 2026, ninth and tenth at about 66 GFLOP/s per watt.}
\term{Dawn and SuperMUC-NG Phase 2}{The two Ponte Vecchio systems besides Aurora on the TOP500 (Cambridge, 1,024 GPUs; LRZ, 960 GPUs), both from November 2023.}
\term{GFLOP/s per watt}{The Green500's metric: HPL Rmax divided by the power measured during the run.}
\end{terms}

\begin{intuition}
The lists are a census of one kind of buyer. Read the TOP500 as ``who bought FP64 machines two to four years ago'' and the Green500 as ``whose 2022--23 chip does FP64 most efficiently,'' and both are informative; read either as a ranking of today's GPUs for AI and both mislead, because the machines that would top such a ranking were never submitted.
\end{intuition}

% ---------------------------------------------------------------------
\slide{29}{The next machines and the gigawatt deals}{DOE's 2026--28 systems split three ways; the AI commitments of 2025--26 are measured in gigawatts, not GPUs.}

\para{What the slide shows} Two tables and a panel. The upper table lists eight HPC systems announced for 2026--28 with site, accelerator and date: Doudna (NERSC, NVIDIA Vera Rubin built by Dell, late 2026); Lux (ORNL, AMD MI355X built by HPE, early 2026); Equinox and Solstice (Argonne with Oracle, NVIDIA Blackwell, 10,000 and 100,000 GPUs, 2026 onward); Discovery (ORNL, AMD MI430X with Venice CPUs, HPE, 2028); Alice Recoque (CEA/GENCI, MI430X with Venice, Eviden, 2026--27); LUMI-AI (EuroHPC, MI430X, second half of 2027); Blue Lion (LRZ, Vera Rubin, HPE, 2027); and Herder (HLRS, AMD, the successor to Hunter, 2027). The lower table lists six AI commitments with what was agreed and when: OpenAI--AMD, 6 GW of Instinct from the second half of 2026 with a warrant for up to 160 million AMD shares (October 2025); Oracle--AMD, about 50,000 MI450 (October 2025); Meta--AMD, a custom MI450 part, up to 6 GW, with a warrant for up to 160 million shares (February 2026); Anthropic--AMD, up to 2 GW of MI450-series GPUs with AMD investing up to \$5 billion in Anthropic (July 2026); Stargate Abilene, about 450,000 GB200 at 1.2 GW (2025); and NVIDIA--Intel, a \$5 billion stake and x86 CPUs with NVLink (September 2025). The panel, ``What the procurements say'', reads the two tables: DOE split its bet three ways; the AI deals are sized in power, and 6 GW is two to four million GPUs depending on the kW per GPU assumed (derived); and secondary estimates still put NVIDIA at 75--80\% of 2026 shipments.

\para{Details} The upper table is the procurement cycle that will produce the TOP500 lists of 2027--28, and the panel's first sentence is its summary: the U.S.\ Department of Energy, which bought AMD for Frontier and El Capitan and Intel for Aurora, has divided its next round. NERSC's Doudna (also called NERSC-10) is a Dell system on Vera Rubin, and Berkeley Lab took delivery of its early-access system, Cech, in January 2026: 72 Grace CPUs and 144 Blackwell GPUs, rated by the laboratory at 5.76 PFLOP/s FP64 and 1.44 EFLOP/s NVFP4, with the full system due in late 2026$^\dagger$. Argonne's Equinox and Solstice, announced by NVIDIA, Oracle and DOE on October 28, 2025, are AI systems of 10,000 and 100,000 Blackwell GPUs, described together as 2,200 exaflops of AI performance without a stated precision, with Equinox due in the first half of 2026$^\dagger$; they are the first DOE machines sized in AI exaflops rather than FP64. Oak Ridge went the other way: Lux, announced October 27, 2025, is an HPE ProLiant XD685 cluster of MI355X GPUs with EPYC CPUs and Pensando networking, built with Oracle Cloud and due in early 2026 as an AI cluster, and Discovery, the successor to Frontier, is an HPE Cray GX5000 system on AMD's next-generation Venice EPYC CPUs and MI430X GPUs for delivery in 2028$^\dagger$. The MI430X is the HPC variant of the MI400 family with ``true hardware FP64'' at 288 TFLOP/s and 432 GB of HBM4, shipping in the first half of 2027 (slide 12)$^\dagger$.

Europe's choices follow the same split. Alice Recoque, the French exascale system at CEA for GENCI, is an Eviden BullSequana XH3500 of 94 racks on Venice and MI430X with the BXI interconnect, a EUR 554 million contract targeting more than an exaflop on HPL, for 2026--27$^\dagger$; LUMI-AI, the EuroHPC successor to LUMI at CSC in Finland, is an MI430X system for the second half of 2027$^\dagger$; Blue Lion at the Leibniz Supercomputing Centre in Garching is an HPE system on Vera Rubin for 2027$^\dagger$; and Herder at HLRS in Stuttgart, the 2027 successor to the MI300A-based Hunter (752 APUs, in operation since February 2025), is an AMD system whose GPU has not been named$^\dagger$. Intel appears nowhere in either table: after Falcon Shores (slide 20) it has no announced next system, and the panel says so.

The division is not arbitrary, and the panel's reading of it is the deck's thesis about HPC procurement. The centers whose mission is FP64 simulation, and whose procurement rules weight it---Oak Ridge, CEA, CSC---chose the vendor that kept building FP64 hardware, AMD, whose MI430X promises 288 TFLOP/s per GPU against Rubin's 33 (slide 12); the centers adding AI capacity on short timelines, NERSC's user community and Argonne's AI program, chose NVIDIA, whose rack and software were available first. Both bets are visible in the same laboratory: Oak Ridge bought an MI355X AI cluster and an MI430X leadership machine, Argonne bought a Blackwell AI machine after an Intel FP64 one. The procurements also show the CPU side of the platform contest: Discovery and Alice Recoque pair MI430X with AMD's Venice, Doudna and Blue Lion pair Rubin with NVIDIA's Vera, and no one pairs a GPU with the other vendor's CPU, because the coherent CPU--GPU link of slide 27 is now proprietary on both sides.

The lower table is a different kind of procurement and the panel's second observation is about its units. The AI commitments of 2025--26 are stated in gigawatts of data-center power rather than in GPUs or FLOP/s. The OpenAI--AMD agreement of October 6, 2025 commits OpenAI to deploy 6 GW of Instinct GPUs over several years, beginning with 1 GW of MI450 in the second half of 2026, and grants OpenAI a warrant for up to 160 million AMD shares that vests as the deployments and AMD's share price reach milestones$^\dagger$; Oracle announced about 50,000 MI450 from the third quarter of 2026$^\dagger$; Meta agreed on February 24, 2026 to a custom MI450-based part for up to 6 GW, with the same warrant structure for up to 160 million AMD shares$^\dagger$; Anthropic committed on July 22, 2026 to up to 2 GW of MI450-series GPUs from the first half of 2027, with AMD investing up to \$5 billion in Anthropic$^\dagger$; and Microsoft Azure is a Helios customer$^\dagger$. On the NVIDIA side, Stargate Abilene, the OpenAI and Oracle site in Texas, is planned for about 450,000 GB200 GPUs across eight buildings at 1.2 GW, with two buildings operating in September 2025$^\dagger$. The NVIDIA--Intel agreement of September 18, 2025 is of a third kind: NVIDIA bought \$5 billion of Intel stock at \$23.28 a share, and the two will build custom x86 CPUs with NVLink for NVIDIA's racks and x86 processors with RTX graphics chiplets for PCs$^\dagger$. For this lecture its meaning is the one slide 20 gives: Intel's route back into the AI data center is as a CPU supplier to NVLink racks, not as a GPU vendor.

Converting gigawatts to GPUs is the panel's derived number and the arithmetic should be made explicit, because the answer depends on what the gigawatt counts. A 2026 accelerator draws 1.2--1.4 kW (GB200 1,200 W; GB300 and MI355X 1,400 W; Rubin and MI455X not published), and the CPUs, switches, NICs, storage, power conversion and cooling of the rack and facility add to it. At 1.5 kW per GPU, 6 GW is 4 million GPUs; at 2 kW all-in, 3 million; and if the gigawatts are Helios racks at their reported 225--245 kW for 72 GPUs (3.1--3.4 kW per GPU at the rack, before facility overhead), about 1.8 million, or 1.5 million with a power usage effectiveness of 1.2. The deck's ``two to four million GPUs over several years'' brackets that range, and the order of magnitude is the point: the OpenAI commitment alone is comparable to AMD's entire Instinct shipments to date, and Meta's is of the same size. By comparison, Stargate Abilene's 450,000 GB200 at 1.2 GW is 2.7 kW per GPU, consistent with a GB200 NVL72 rack at about 120 kW (1.67 kW per GPU) plus facility overhead of about 1.6$\times$ (derived). The commitments are also conditional in a way the HPC procurements are not: they are capacity to be delivered over years, and the equity runs from AMD to the customer, not the other way: AMD granted OpenAI and Meta warrants for up to 160 million AMD shares each that vest as the gigawatts are deployed, and it is investing up to \$5 billion in Anthropic; Oracle simply buys racks. The panel is careful to say that these are commitments, not yet deployments. The first evidence that they are real came on November 3, 2026, when AMD's 8-K reported that the MI450 series had launched in the third quarter (``shipping in volume'', on the call) with first Helios revenue, that Helios racks were running at Anthropic, Meta, Microsoft, OpenAI and Oracle, and that its data-center segment revenue had grown 97\% year on year to \$8.5 billion$^\dagger$.

The panel's third observation is about share, and it is marked as secondary because no primary source exists: market-research estimates compiled in 2026 put NVIDIA at about 87\% of data-center AI accelerator revenue in 2024, about 80\% in 2025 and 75--80\% in 2026, AMD at 5\% rising to 7--10\%, and Intel below 1\%$^\dagger$. Those estimates predate most of the deliveries in the lower table, and the deck's claim is a conditional one: these are the first commitments large enough to move the share, if they are delivered. Whether they are is the open question of slide 32.

\begin{formulas}
GPUs from power: $N \approx P_{\text{facility}} / p_{\text{all-in}}$ with $p_{\text{all-in}} = \text{TDP} \times (1 + \text{rack overhead}) \times \text{PUE}$. At TDP 1.4 kW, 25\% rack overhead and PUE 1.2, $p_{\text{all-in}} \approx 2.1$ kW, so 6 GW $\approx 2.9$ million GPUs; at 1.5 kW all-in, 4 million; from the Helios rack figure, $245\ \text{kW}/72 = 3.4$ kW per GPU at the rack, so $6\ \text{GW}/(3.4\ \text{kW}\times 1.2) \approx 1.5$ million (all derived). Check against Abilene: $1.2\ \text{GW}/450{,}000 = 2.7$ kW per GB200, against $120\ \text{kW}/72 = 1.67$ kW at the rack, implying a facility overhead factor of about 1.6 (derived).
\end{formulas}

\begin{terms}
\term{Doudna}{NERSC's next system (NERSC-10), a Dell machine on Vera Rubin due in late 2026; its early-access system Cech (72 Grace, 144 Blackwell) arrived in January 2026.}
\term{Lux}{Oak Ridge's AI cluster of MI355X GPUs in HPE ProLiant XD685 servers, built with Oracle Cloud, announced October 2025 for early 2026.}
\term{Equinox and Solstice}{Argonne's Blackwell AI systems of 10,000 and 100,000 GPUs, announced by NVIDIA, Oracle and DOE on October 28, 2025.}
\term{Discovery}{Oak Ridge's successor to Frontier: an HPE Cray GX5000 system on AMD Venice CPUs and MI430X GPUs for 2028.}
\term{Alice Recoque}{France's exascale system at CEA for GENCI: an Eviden BullSequana XH3500 of 94 racks on Venice and MI430X, targeting more than an exaflop on HPL, 2026--27.}
\term{LUMI-AI}{The EuroHPC successor to LUMI at CSC in Finland, on MI430X, due in the second half of 2027.}
\term{Blue Lion}{LRZ's next system in Garching, an HPE machine on Vera Rubin for 2027.}
\term{Hunter and Herder}{HLRS's MI300A system (752 APUs, 2025) and its 2027 AMD successor.}
\term{MI430X}{AMD's HPC variant of the MI400 family: 288 TFLOP/s FP64 and 432 GB HBM4, shipping in the first half of 2027.}
\term{Venice}{AMD's sixth-generation EPYC CPU (Zen 6), paired with MI430X in Discovery and Alice Recoque and with MI455X in Helios.}
\term{OpenAI--AMD agreement}{The October 6, 2025 commitment to deploy 6 GW of Instinct GPUs from the second half of 2026, with a warrant for up to 160 million AMD shares.}
\term{Stargate Abilene}{The OpenAI and Oracle data-center site in Texas planned for about 450,000 GB200 GPUs at 1.2 GW.}
\term{NVIDIA--Intel agreement}{NVIDIA's \$5 billion investment in Intel (September 18, 2025) and the plan for custom x86 CPUs with NVLink and x86 processors with RTX graphics.}
\term{gigawatt}{$10^9$ W; the unit in which the AI deployments of 2025--26 are stated, each gigawatt corresponding to roughly half a million GPUs all-in (derived).}
\term{warrant}{A right to buy shares at a set price; the OpenAI--AMD agreement pays OpenAI partly in such rights as it deploys AMD GPUs.}
\term{market share estimate}{A figure from market-research firms, not from the vendors; the deck marks all such figures as secondary.}
\end{terms}

\begin{intuition}
Two kinds of buyer, two kinds of contract. A laboratory buys one machine for seven years and names the GPU, the CPU and the HPL target; an AI company buys power in gigawatts over several years and may be paid in shares to do it. The first tells you who builds FP64; the second tells you who might move the market share, once the racks are running.
\end{intuition}
% =====================================================================
\setstage{vc}
\section{Benchmarks (slide 30)}

\slide{30}{MLPerf: where AMD and NVIDIA meet on the same benchmark}{AMD entered in 2024 (inference) and 2025 (training); by 2026 the MI355X trades blows with the B200.}

\para{What the slide shows} A horizontal bar chart and a panel. The chart plots the MI355X's per-GPU result relative to the B200's (B200 $= 1$, marked by a dashed line) on five MLPerf tests: Inference 6.0, Llama 2 70B in the offline scenario (1.00); the same model in the server scenario (0.97); the same in the interactive scenario (1.19); Training 5.1, Llama 2 70B LoRA fine-tuning (0.97); and Training 5.1, Llama 3.1 8B pretraining (0.94). A caption explains that the training ratios are the B200's time divided by the MI355X's (10.18 against 9.85 minutes; 99.7 against 93.7 minutes), that the inference ratios are as AMD states them, and that against the B300 the three inference ratios are 0.92, 0.93 and 1.04. The panel, ``The rounds'', lists the MLPerf rounds from Inference 4.1 (August 2024) to Training 6.0 (June 2026) with AMD's entries and the context, and closes with Intel: Gaudi 2 submitted in 2023--24, and no Gaudi 3 submission is on record.

\para{Details} MLPerf is the benchmark suite of MLCommons, an industry consortium founded in 2018 (as MLPerf) and incorporated in 2020, and it is the only venue in which AMD and NVIDIA submit results for the same machine-learning workloads under the same rules, with the results reviewed by the other submitters before publication. It has two suites relevant here. MLPerf Training measures the wall-clock time to train a reference model to a specified quality on a specified dataset; the result is minutes, lower is better, and a submission names the system (accelerator type and count, host CPUs, software). MLPerf Inference measures throughput on a trained model under scenarios: offline (all queries available at once, maximizing throughput), server (queries arrive at random with a latency bound), and, since 2025, interactive (server with tighter time-to-first-token and time-per-output-token bounds that mimic a chat application); the result is samples or tokens per second, higher is better, and every result must meet an accuracy target (99\% or 99.9\% of the reference model's). Both suites have a Closed division, in which the model, dataset and quality target are fixed and only the implementation may change, and an Open division, in which the model may be altered. All the results on the slide are Closed-division results, which is what makes the comparison meaningful.

The panel's history is AMD's entry into this arena. In Inference 4.1 (results published August 28, 2024) AMD submitted for the first time: eight MI300X in a server, running Llama 2 70B in FP8, came within 2--3\% of an eight-GPU DGX H100 on the offline and server scenarios$^\dagger$. That result mattered because it was the first time AMD's claimed parity with the H100 (slide 9) had been measured by a third party under rules NVIDIA had agreed to; it also came eight months after the MI300X launched, and the ROCm work needed to get there is the subject of slide 26. Inference 5.0 (April 2, 2025) added the first MI325X results, and Training 5.0 (June 4, 2025) brought AMD's first training submission ever, with MI300X and MI325X systems, in the same round in which the GB200 made its debut$^\dagger$. Inference 5.1 (September 9, 2025) brought the first MI355X results and the first results in MXFP4 from any vendor: on Llama 2 70B the MI355X delivered 2.7 times the MI325X's throughput, and an eight-node submission scaled to 7.9 times a single node$^\dagger$. Training 5.1 (November 2025) is the source of the two training bars: MI355X fine-tuned Llama 2 70B with LoRA in 10.18 minutes against the B200's average of 9.85 and the B300's 9.59, and pretrained Llama 3.1 8B in 99.7 minutes against the B200's 93.7 and the B300's 95.1, all at eight GPUs$^\dagger$. Inference 6.0 (April 1, 2026) is the source of the three inference bars: a multi-node MI355X submission of 87 GPUs passed one million tokens per second on Llama 2 70B offline (1,042,110), a single MI355X reached 100,282 tokens per second in the server scenario, and AMD's per-GPU comparison with the B200 gave a tie offline, 97\% in server and 119\% in interactive, with 92\%, 93\% and 104\% against the B300$^\dagger$. Training 6.0 (June 16, 2026) set the scale record: NVIDIA trained the DeepSeek-V3 671B benchmark in 2.02 minutes on 8,192 GB300 GPUs and Llama 3.1 405B in 7.07 minutes on 8,192 GB200, with GB300 1.3--1.6 times faster than GB200 at equal counts; neither Rubin nor MI455X appeared in the round$^\dagger$.

How the ratios are computed is worth stating, because MLPerf publishes absolute results and the ratios are the vendors' or the deck's. For training, the natural comparison is time, so the ratio is the B200's time divided by the MI355X's: $9.85/10.18 = 0.97$ and $93.7/99.7 = 0.94$ (derived); a value below one means the MI355X was slower. For inference the comparison is throughput per GPU, and MLPerf systems differ in accelerator count, host CPU and software, so AMD divides each system's throughput by its accelerator count before comparing; the 1.00, 0.97 and 1.19 are AMD's figures from its results blog, and the deck marks them with a dagger. The interactive scenario's 1.19 is the most interesting number on the chart: in the scenario with the tightest latency bounds, which stresses memory bandwidth and the serving software's scheduling rather than peak arithmetic, the MI355X's 8 TB/s of HBM3e and 288 GB (against the B200's 7.7 TB/s and 180 GB, slide 4) show up as a 19\% lead, while in the throughput-bound offline scenario the two are even. Against the B300, which has the same 288 GB and 8 TB/s as the MI355X but 50\% more FP4 arithmetic (15 against 10 PFLOP/s dense), the MI355X trails by 7--8\% in offline and server and leads by 4\% in interactive$^\dagger$. The numbers are consistent with the datasheets: at equal memory the GPU with more arithmetic wins the throughput scenarios, and where memory decides, the two are even.

The two training results say something different. LoRA fine-tuning of Llama 2 70B at eight GPUs is a short, communication-light task, and the MI355X's 3\% deficit is within the spread between NVIDIA's own submissions (the B200's 9.85 is an average over several systems). Pretraining Llama 3.1 8B is longer and exercises the collective communication between the eight GPUs, where the MI355X's Infinity Fabric at $7 \times 144$ GB/s against the B200's NVLink 5 at 1.8 TB/s per GPU (slide 21) is at a disadvantage, and the deficit grows to 6\%. Both are eight-GPU results. The absence of any AMD submission at hundreds or thousands of GPUs in Training 6.0, where NVIDIA's 8,192-GPU runs set the headline, is the gap the deck emphasizes: MLPerf shows the MI355X competitive with the B200 inside a node, and shows nothing about AMD at the scale at which frontier models are trained, because AMD has not submitted there. Helios and the MI455X (slide 11) are AMD's answer to that scale, and neither had a result at the time of these notes.

Intel's line in the panel closes the vendor comparison. Habana's Gaudi 2 submitted to MLPerf Training (v3.0, v3.1, v4.0) and Inference (v3.1, v4.0) in 2023--24 with results that were competitive with the A100 and, on some tests, the H100; no Gaudi 3 submission is on record, and after 2024 Intel stopped submitting accelerator results$^\dagger$. No Ponte Vecchio result was ever submitted, and no result for any Intel data-center GPU exists in the 2025--26 rounds. Slide 16 explains the first part (Gaudi's sales never justified the engineering) and slide 20 the second (there is no current Intel data-center GPU to submit).

Three cautions apply to every MLPerf comparison. The results are for the Closed division's reference models, which lag the frontier by a year or two (Llama 2 70B dates from 2023); the submissions are made by the vendors on hand-tuned software stacks, so they measure the best each vendor can do rather than what a customer's code achieves on the first day; and a per-GPU ratio compares systems that differ in CPU, memory and network, so differences of a few percent are not significant. What MLPerf does establish, and what no datasheet can, is that by 2026 the MI355X and the B200 are the same class of machine on the workloads measured, and that the gap between them is in the software and the scale, not the silicon.

\begin{formulas}
Training ratio (B200 $= 1$): $r = t_{\text{B200}}/t_{\text{MI355X}}$; LoRA $9.85/10.18 = 0.97$; Llama 3.1 8B $93.7/99.7 = 0.94$ (derived). Inference ratio: $r = (T_{\text{MI355X}}/n_{\text{MI355X}})/(T_{\text{B200}}/n_{\text{B200}})$ with $T$ the system throughput and $n$ the accelerator count; AMD's values 1.00, 0.97, 1.19. Multi-node scaling: $1{,}042{,}110$ tokens/s on 87 GPUs is about 12{,}000 tokens/s per GPU offline (derived), against 100{,}282 for one GPU in the server scenario, two different scenarios on different accuracy and latency terms.
\end{formulas}

\begin{terms}
\term{MLPerf}{The benchmark suite of MLCommons for machine-learning training and inference, with peer-reviewed submissions published in rounds twice a year for each suite.}
\term{MLCommons}{The industry consortium (founded as MLPerf in 2018, incorporated 2020) that defines and publishes the MLPerf benchmarks.}
\term{MLPerf Training}{The suite measuring wall-clock time to train a reference model to a target quality; results in minutes, lower is better.}
\term{MLPerf Inference}{The suite measuring throughput on trained models under latency scenarios, subject to accuracy targets; results in samples or tokens per second.}
\term{offline, server, interactive}{The inference scenarios: all queries at once; random arrivals under a latency bound; random arrivals under tight time-to-first-token and time-per-output-token bounds.}
\term{Closed and Open divisions}{Closed fixes the model and quality target and allows only implementation changes; Open permits model changes. The slide's results are Closed.}
\term{Llama 2 70B}{Meta's 70-billion-parameter language model (2023), the MLPerf Inference reference model on which AMD and NVIDIA have met since 2024.}
\term{LoRA}{Low-rank adaptation: fine-tuning by training small low-rank matrices added to a frozen model; an MLPerf Training benchmark on Llama 2 70B.}
\term{Llama 3.1 8B and 405B}{Meta's 8- and 405-billion-parameter models (2024); the 8B pretraining benchmark entered MLPerf Training in 2025, the 405B in 2024.}
\term{DeepSeek-V3}{The 671-billion-parameter mixture-of-experts model (December 2024) that became an MLPerf Training benchmark in 2026; 2.02 minutes on 8,192 GB300.}
\term{MXFP4}{The Open Compute Project's block-scaled 4-bit format (32 values per 8-bit power-of-two scale); first used in MLPerf by AMD in Inference 5.1.}
\term{per-GPU ratio}{A system result divided by its accelerator count, compared across systems; the deck's inference ratios are AMD's per-GPU figures.}
\term{Gaudi in MLPerf}{Habana's Gaudi 2 submitted training and inference results in 2023--24; no Gaudi 3 submission is on record, and Intel has submitted none since.}
\end{terms}

\begin{intuition}
A datasheet is a promise; a benchmark is a measurement under rules the competitor agreed to. On the measurements that exist, the MI355X and B200 are even inside a node and the B300 is ahead where arithmetic decides. On the measurement that matters most for frontier training, thousands of GPUs, only one vendor has shown up. Read the absences as carefully as the bars.
\end{intuition}
% =====================================================================
\setstage{ink}
\section{Wrap-up (slides 31--33)}

\slide{31}{The platforms in one table}{Seven accelerators from three vendors, 2023--2026, as their datasheets state them.}

\para{What the slide shows} A single table of seven accelerators in columns---H100, GB300 and Rubin for NVIDIA; MI300X, MI355X and MI455X for AMD; the Ponte Vecchio Max 1550 for Intel---and sixteen rows: vendor; year first shipped (2022, 2025, 2026; 2023, 2025, 2026; 2023); compute dies and process (1 die on 4N; 2 on 4NP; 2 on 3 nm; 8 XCDs on N5, N3P and N2; 16 tiles on N5); transistors in billions (80, 208, 336; 153, 185, 320; over 100); HBM capacity in GB (80, 288, 288; 192, 288, 432; 128); HBM bandwidth in TB/s (3.35, 8.0, 19.2; 5.3, 8.0, 23.3; 3.28); on-chip cache (50 MB L2; not published; not published; 256 MB Infinity Cache; 256 MB Infinity Cache; 192 MB L2; 408 MB L2); FP64 vector / matrix in TFLOP/s (34 / 67; 1.4; 33; 81.7 / 163; 78.6; not published; 52.4); dense FP16 (989; 2,500; 4,000; 1,307; 2,500; 5,030; 839); dense FP8 (1,979; 5,000; 17,500; 2,615; 5,000; 20,130; none); dense FP4 (none; 15,000; 35,000; none; 10,000; 40,260; none); scale-up bandwidth per GPU in GB/s (900; 1,800; 3,000; 896; 1,008; 3,600; Xe Link); scale-up domain (8; 72; 72; 8; 8; 72; 6, with 8 the Xe Link maximum); power in watts (700; 1,400; not published; 750; 1,400; not published; 600); and software (CUDA; ROCm; oneAPI). A footnote states the conventions: vendor figures, dense tensor peaks, Rubin from the September 2026 revision, MI455X compute in MXFP8 and MXFP4, PVC's FP16 as Intel's figure consistent with its confirmed TF32 rate, ``n.p.'' for not published as of November 2026, and IC for Infinity Cache.

\para{Details} The table is the lecture's scorecard, and its value is in reading it row by row with the earlier slides in mind. The columns pair up by date: H100 against MI300X and the Max 1550 (2022--23, slide 18), GB300 against MI355X (2025, slides 10 and 22), Rubin against MI455X (2026, slides 11 and 22). Each pair is a generation in which the two main vendors shipped within months of each other, and the third vendor appears once.

The dies-and-process row records three packaging strategies (slide 3). NVIDIA kept the largest possible monolithic die for as long as it could---the H100's 814 mm$^2$ GH100 on TSMC's 4N process---and went to two reticle-sized dies joined by a 10 TB/s die-to-die link on Blackwell (4NP) and Rubin (3 nm). AMD went to chiplets with CDNA 2 and to 3D-stacked chiplets with CDNA 3: eight accelerator complex dies on four I/O dies for the MI300X (N5 on N6), eight on two for the MI355X (N3P on N6) and eight on two I/O plus two cache dies for the MI455X (N2 on N3P). Intel's Ponte Vecchio took 47 active tiles from five process nodes, sixteen of them compute tiles on TSMC N5 (slide 17). The transistor counts follow: 80 billion on one die against 153 on twelve, and in 2026, 336 billion on two dies against 320 on twelve; the deck notes elsewhere that NVIDIA's two-die counts and AMD's twelve-die counts include different amounts of I/O and cache silicon, so the row compares packages, not compute.

The memory rows are the row AMD won from the MI250X on, with a tie in 2025 (slide 5). Capacity per GPU: 80 GB against 192 in 2023 (2.4$\times$), 288 against 288 in 2025, 288 against 432 in 2026 (1.5$\times$); Intel's 128 GB led both in early 2023. Bandwidth: 3.35 TB/s against 5.3 (1.6$\times$), 8.0 against 8.0, 19.2 against 23.3 (1.2$\times$), with PVC's 3.28 TB/s equal to the H100's (all ratios derived). The two 2025 parts are the same memory---eight stacks of 12-high HBM3e---and the two 2026 parts differ in stack count, eight stacks of HBM4 for Rubin and twelve for the MI455X, which is also why the MI455X package is wider.

The cache row shows a design divergence that the earlier slides explain. The H100 has 50 MB of L2; NVIDIA has not published the L2 of Blackwell or Rubin (the deck writes n.p.). AMD put a 256 MB Infinity Cache on the I/O dies of CDNA 3 and CDNA 4 as a memory-side last-level cache shared by all compute dies, then removed it on CDNA 5 in favor of 192 MB of L2 shared across the package. Intel's 408 MB on PVC, the largest of the seven, is spread over the base tiles and the RAMBO cache tiles (slide 17), and the deck notes that Intel's own documents split it inconsistently.

The FP64 row is the lecture's longest argument in one line (slide 12). NVIDIA: 34 TFLOP/s vector and 67 on tensor cores for the H100; 1.4 for the GB300, which removed most of its FP64 units; 33 for Rubin, which restored them without a separate matrix rate. AMD: 81.7 vector and 163 matrix on the MI300X, the highest of any GPU to date; 78.6 on the MI355X, where CDNA 4 halved the matrix FP64 rate per compute unit so that vector and matrix coincide; and no published figure for the MI455X, a part built for AI whose HPC sibling, the MI430X, is promised at 288. Intel: 52.4 for PVC, equal to its FP32 rate. The consequence in slide 27 is that every exascale machine runs on AMD or Intel FP64 except JUPITER, and JUPITER runs on the H100 generation that NVIDIA has since cut.

The three low-precision rows show the convergence (slide 6). Dense FP16: 989 against 1,307 against 839 in the first generation, 2,500 against 2,500 in the second, 4,000 against 5,030 in the third. Dense FP8: 1,979 against 2,615, 5,000 against 5,000, 17,500 against 20,130. Dense FP4: 15,000 against 10,000 in 2025 (the GB300's extra FP4 is Blackwell Ultra's redesigned tensor core, slide 10), 35,000 against 40,260 in 2026. The MI455X figures are in the MX block-scaled formats and Rubin's in NVFP4 and FP8, formats that differ in block size and scale encoding, so the row compares the vendors' own headline figures, not identical arithmetic. PVC has neither FP8 nor FP4, and its FP16 figure is Intel's own, consistent with its confirmed 419 TFLOP/s of TF32 and with the per-core throughput table in Intel's architecture material (4,096 FP16 operations per clock per Xe-core); Gaudi, Intel's other accelerator, is not in the table because it is not a GPU.

The scale-up rows are the row NVIDIA won at every generation (slide 21). Per GPU: 900 GB/s of NVLink 4 against 896 of Infinity Fabric in 2023, a tie in bandwidth that the switch turned into a rout, because the H100's 900 GB/s reaches every other GPU in the domain through NVSwitch while the MI300X's seven links of 128 GB/s each reach one neighbor; 1,800 against 1,008 in 2025; 3,000 against 3,600 in 2026, where AMD's UALink-over-Ethernet figure exceeds NVIDIA's for the first time on paper. Domain size: 8 GPUs for the HGX and UBB boards, 72 for the NVL72 racks since 2024 and for Helios since 2026, six for an Aurora node's Xe Link mesh (Xe Link supports up to eight GPUs per node). The deck gives Rubin's 3,000 as revised in September 2026 from the 3,600 announced in January, so the two 2026 figures were equal when announced.

The power row records the convergence of the thermal envelope: 700 and 750 W in 2023, 1,400 W for both in 2025, and no published figure for either 2026 part as of November 2026; PVC's 600 W in 2023 was the highest of its generation. The software row names the three stacks of slide 24, and the deck's footnote on the table is the lecture's method in one sentence: vendor figures, dense peaks, the revision named, the unconfirmed figure flagged, and ``n.p.'' where a number does not exist rather than an estimate in its place.

\begin{formulas}
Generation ratios (AMD over NVIDIA, derived): HBM capacity $192/80 = 2.4$, $288/288 = 1.0$, $432/288 = 1.5$; HBM bandwidth $5.3/3.35 = 1.6$, $8.0/8.0 = 1.0$, $23.3/19.2 = 1.2$; dense FP8 $2{,}615/1{,}979 = 1.3$, $5{,}000/5{,}000 = 1.0$, $20{,}130/17{,}500 = 1.15$; dense FP4 $10{,}000/15{,}000 = 0.67$, $40{,}260/35{,}000 = 1.15$; FP64 matrix $163/67 = 2.4$, $78.6/1.4 = 56$; scale-up per GPU $896/900 = 1.0$, $1{,}008/1{,}800 = 0.56$, $3{,}600/3{,}000 = 1.2$. Transistors per compute die (derived, package totals divided by compute dies, I/O silicon included): H100 80B on 1; GB300 104B on each of 2; MI455X 320B over 8 XCDs plus 4 base dies.
\end{formulas}

\begin{terms}
\term{scorecard table}{The deck's one-page comparison of seven accelerators on sixteen rows, all from vendor material, with the conventions stated in its footnote.}
\term{accelerator complex die}{AMD's name (XCD) for a CDNA compute chiplet; eight per MI300X, MI355X and MI455X.}
\term{I/O die}{A chiplet carrying memory controllers, cache and fabric links on which compute dies are stacked (AMD IOD; Intel base tile).}
\term{Infinity Cache}{AMD's memory-side last-level cache on the I/O dies of CDNA 3 and 4 (256 MB); removed on CDNA 5 in favor of a 192 MB shared L2.}
\term{RAMBO tile}{Ponte Vecchio's cache tile (Random Access Memory, Bandwidth Optimized), eight per package, part of its 408 MB of L2.}
\term{dense peak}{A tensor-core peak counted without structured sparsity; NVIDIA's headline figures are 2:4-sparse and twice the dense ones.}
\term{MX formats}{The Open Compute Project's block-scaled formats (MXFP8, MXFP6, MXFP4) in which AMD states the MI455X's compute peaks.}
\term{n.p.}{Not published: the deck's mark for a figure the vendor had not disclosed as of November 2026, used in place of an estimate.}
\end{terms}

\begin{intuition}
Read the table in pairs of columns and ask one question per row: who led, and by how much. Memory and FP64 say AMD; the fabric and the domain say NVIDIA; low-precision arithmetic says neither, by 2025. Then recall that the table is datasheets, and that slides 26 and 30 measured what fraction of each number a customer gets.
\end{intuition}

% ---------------------------------------------------------------------
\slide{32}{Summary: what to take from this lecture}{The dies converged, the platforms did not; the fabric, the software and the buying model decide.}

\para{What the slide shows} Two panels. ``Eight things to remember'' lists: (1) on paper the compute GPUs converged, with 10 PFLOP/s of dense FP4 from both AMD and NVIDIA in 2025 and 35--40 in 2026; (2) AMD led on memory per GPU from 2021, tied in 2025; (3) FP64 diverged, with AMD building it up to the MI430X's 288 TFLOP/s, NVIDIA cutting it, and Intel running it at full rate once; (4) the scale-up fabric is NVIDIA's widest lead; (5) packaging, with Intel's 47 tiles, AMD's chiplets and NVIDIA's dies; (6) HPC is portable and buys FP64 per watt, so the exascale list is AMD, Intel, one NVIDIA and one CPU-only machine; (7) AI is CUDA-first and buys by the quarter, so NVIDIA's share follows from software and fabric; (8) measure, do not read: MLPerf, HPL and microbenchmarks. ``Open questions for 2027'' asks whether Helios delivers at multi-gigawatt scale, whether Rubin's lead holds against Rubin Ultra and the MI500 when both vendors' 2026 specifications moved after announcement, whether Intel returns through the GPU D-series and the NVIDIA deal, what HPC does with 288 TFLOP/s of FP64 from AMD against 33 with emulation from NVIDIA, and whether the standards (UALink silicon, Ultra Ethernet NICs, OCP racks) bite or NVLink Fusion absorbs them.

\para{Details} The eight points are the lecture's argument compressed, and each has a slide behind it. The first is slide 6: in 2020 the MI100's 185 TFLOP/s of FP16 trailed the A100's 312 by 40\%, and the MI250X's 383 trailed the H100's 989 by 60\% a year later; by 2025 the MI355X and GB200 both delivered 10 PFLOP/s of dense FP4 and 5 of FP8, and in 2026 the MI455X's 40.3 PFLOP/s exceeds Rubin's 35 (vendor figures, different block formats). The convergence has a cause the deck names: both vendors buy the same TSMC processes and the same HBM, both have reached the reticle limit and gone to multi-die packages, and the arithmetic of a tensor core at a given format is not a secret, so the peak per package is set by die area and power, which are the same for both.

The second point is slide 5 and the third is slide 12, and together they are the deck's account of why the HPC market and the AI market bought different GPUs. AMD's memory lead---128 GB against 80 in 2021, 192 against 80 in 2023, 432 against 288 in 2026---came from designing for the FP64 buyers, whose codes are bandwidth-bound and whose procurements specify memory per node, and from the chiplet package, which could carry more stacks. AMD's FP64 came from the same buyers: the 2019 CORAL-2 contracts for Frontier and El Capitan specified FP64 performance, CDNA 2 and 3 doubled the matrix FP64 rate to meet them, and the MI430X continues the line for Discovery and Alice Recoque. NVIDIA designed Blackwell Ultra and Rubin for inference, where FP64 is dead weight: the GB300 kept 1.4 TFLOP/s and Rubin 33, and NVIDIA offers emulation on INT8 tensor cores in its place. Intel's Ponte Vecchio ran FP64 at the full FP32 rate because Aurora's contract required it, and Intel then left the market.

The fourth and fifth points are slides 21--23 and 3. The scale-up fabric is NVIDIA's widest lead because it is the one component of the platform that no standard has yet replaced: NVLink with NVSwitch since 2018 gives every GPU in the domain full bandwidth to every other, the domain grew from 8 to 72 GPUs in 2024, and AMD's answer, UALink over Ethernet in Helios, shipped in the third quarter of 2026 at the same domain size and comparable bandwidth but with no measured results yet. Packaging is the reverse case: Intel built the most complex package ever shipped and gained nothing in the market from it; AMD's chiplets gave it the memory lead and the APU; NVIDIA's two dies were enough. The lesson the deck draws is that packaging is a means and the fabric is an end, because the fabric sets what a customer can buy as one unit.

The sixth and seventh points are slide 25 and the lecture's thesis. HPC codes are portable by construction because the Department of Energy and its European counterparts funded Kokkos, RAJA, OpenMP and SYCL ports for a decade, so an HPC center can buy whichever vendor offers the most FP64 per watt and per dollar, and the June 2026 exascale list (slide 28) is the result: El Capitan and Frontier on AMD, Aurora on Intel, JUPITER on NVIDIA, LineShine on CPUs. AI code is CUDA-first because the kernels that set performance---FlashAttention, the Triton compiler, the serving engines---are written against CUDA and ported later, and an AI laboratory buying capacity by the quarter cannot wait for the port, so NVIDIA's share follows from the software and the fabric rather than from the datasheet. The eighth point is slides 26 and 30: a peak is a property of the die and the achieved fraction a property of the libraries, and in 2024 independent measurement put the MI300X at 47\% of its GEMM peak against the H100's 73\%; two years later MLPerf shows the MI355X and B200 even inside a node. The instruction ``measure, do not read'' means: compare MLPerf, HPL and microbenchmark results at the same date, not datasheets.

The open questions are the five places where the deck's account could be overturned within a year. Helios: AMD's first rack-scale product began shipping in the third quarter of 2026 to five named customers, and the gigawatt commitments of slide 29 assume that UALink over Ethernet and ROCm hold up at thousands of GPUs, where no AMD result has yet been measured. Rubin's lead: both vendors' 2026 specifications moved after announcement (Rubin's HBM4 bandwidth and NVLink down in September, the MI455X's HBM4 bandwidth up in July), and both have 2027 parts promised (Rubin Ultra with four dies and HBM4e, the MI500 on the ``Helios 500'' rack), so the 2026 comparison may not hold in 2027. Intel: the GPU D-series is an inference card on LPDDR5X sampling in late 2026, Jaguar Shores has no date, and the NVIDIA investment makes Intel a CPU supplier to NVLink racks; whether that is a return to the data-center GPU market or an exit from it is not yet decidable. HPC: with AMD offering 288 TFLOP/s of hardware FP64 per GPU and NVIDIA 33 plus emulation, the 2028 machines (Discovery, the post-exascale European systems) will show whether the FP64 market is large enough to sustain a dedicated part or whether emulation on tensor cores is good enough. Standards: UALink silicon from Astera, Marvell and others was announced for 2026, Ultra Ethernet NICs are shipping, and the OCP rack designs are in production at AMD and others; against them, NVLink Fusion admits partner chips to NVLink domains on NVIDIA's terms, and which of the two models the hyperscalers choose for their own accelerators will decide whether the scale-up fabric stays proprietary.

A student should leave with a method rather than a verdict. For any accelerator comparison: identify the level (die, package, board, fabric, rack, software; slide 3); state whether each figure is dense or sparse, vector or matrix, announced or shipped; find the measured result nearest in date (MLPerf, HPL, a microbenchmark paper); and ask which buyer the part was designed for, because the answer to ``which is better'' has been ``for whom'' at every generation in this deck.

\begin{terms}
\term{platform contest}{The deck's frame: the competition is decided separately at each level (die, package, board, fabric, rack, software), and the winners differ.}
\term{CUDA-first}{The property of the AI software ecosystem that performance-critical code is written for CUDA and ported to other stacks later.}
\term{portable by construction}{The property of HPC codes written to Kokkos, RAJA, OpenMP or SYCL, which run on all three vendors' GPUs from one source.}
\term{buying by the quarter}{The AI procurement model: capacity bought on short timelines from whichever vendor can deliver, against the HPC model of one machine for seven years.}
\term{Rubin Ultra and MI500}{The 2027 parts promised by NVIDIA (four dies, HBM4e, the Kyber rack) and AMD (CDNA next, the Helios 500 rack).}
\term{FP64 emulation}{Computing double-precision matrix products from integer tensor-core products with error correction (the Ozaki scheme); NVIDIA's substitute for FP64 hardware on Blackwell Ultra and Rubin.}
\term{UALink silicon}{Switch and controller chips implementing the UALink specification, announced by Astera Labs, Marvell and others for 2026.}
\term{NVLink Fusion}{NVIDIA's program (May 2025) licensing NVLink to partners' CPUs and accelerators so that they can join an NVLink domain.}
\end{terms}

\begin{intuition}
Convergent silicon, divergent platforms. If two vendors buy the same process and the same memory and both hit the reticle, their dies will match; what they build around the die---fabric, rack, software, and the customers they build it for---is where they differ, and that is where the decade's outcomes were decided.
\end{intuition}

% ---------------------------------------------------------------------
\slide{33}{Sources}{Whitepapers and datasheets, papers and talks, lists and reporting; $^\dagger$ marks vendor material.}

\para{What the slide shows} Two columns of sources grouped under six headings. AMD: the CDNA 1 through 4 architecture whitepapers (2020--25); the Instinct MI200, MI210, MI300A, MI300X, MI325X and MI355X datasheets and briefs; ROCm documentation and release history; Smith and Alla's Hot Chips 36 paper on the MI300X (2024); AMD's ISSCC 2024 MI300 paper; the Hot Chips 37 CDNA 4 talk (2025); the Advancing AI 2025 and 2026 announcements (MI400, MI455X, Helios, MI430X); the Supermicro Helios datasheet; AMD's third-quarter 2026 results and call; the AMD--OpenAI, AMD--Eviden and ORNL press releases; and AMD's MLPerf blogs. Intel: Junkins's Gen9 compute architecture whitepaper (2015); the Gen11 whitepaper (2019); the Xe-HPG introduction and the oneAPI GPU optimization guide; the Max 1550 specifications; Gomes et al.\ on Ponte Vecchio at ISSCC 2022; the Hot Chips 34 PVC talk; Kaplan's Gaudi 3 talk at Hot Chips 2024; Intel newsroom items; Intel's third-quarter 2026 results; the Intel--NVIDIA announcement; ALCF Aurora documentation; Goto et al.\ on HPL and HPL-MxP on Aurora (2026); the AERIS paper (2025); HPCwire's ``Aurora the survivor'' (2023); and reporting on Rialto Bridge, Falcon Shores, Crescent Island and the GPU D-series. NVIDIA and platforms: the HGX A100, HGX H100 and H200 datasheets; the GB200 NVL72, GB300 NVL72 and Vera Rubin NVL72 product pages; the NVLink Fusion release; and the previous lecture's sources. Fabrics and standards: the UALink 1.0 white paper; the Ultra Ethernet Consortium 1.0 announcement; the OCP OAM and UBB specifications; the Frontier architecture overview; LLNL's El Capitan page; the Next Platform on El Capitan; the LUMI, Adastra and Hunter pages. Software and portability: Reguly (SC-W 2023); Davis et al.\ (ICS 2025); ORNL/TM-2025/4398; the Khronos SYCL 2020 release; Codeplay's oneAPI plugins; the UXL Foundation; SYCLomatic documentation; SemiAnalysis on MI300X against H100 (2024). Lists and benchmarks: the TOP500, HPCG and Green500 lists of June 2026; HPE on the June 2026 list; MLCommons MLPerf Training 5.0 and 6.0 and Inference 6.0; TrendForce (2026); and secondary market-share estimates.

\para{Details} The slide sorts the lecture's evidence by kind, and the kinds carry different weight. Architecture whitepapers and conference papers are primary for the microarchitecture: AMD's CDNA whitepapers give compute-unit counts, cache sizes, link counts and peak rates with the clocks they assume, and the ISSCC and Hot Chips papers on MI300 and Ponte Vecchio give die areas, tile counts, transistor counts and bandwidths that no datasheet states; Junkins's Gen9 whitepaper is the canonical description of Intel's execution unit and slice and the reason the Gen section of the deck can give FLOP-per-clock figures. Datasheets and product pages are primary for what was sold---capacity, bandwidth, TDP, the headline peaks---but they are marketing documents, they change (the deck's handling of the Vera Rubin page's September 2026 revision and the Helios page's bandwidth change is the example), and they quote sparse figures where it flatters; the deck marks every number from them with a dagger. Press releases and earnings calls are primary for dates and commercial facts (shipping quarters, customer commitments, revenue) and for nothing technical. The TOP500 family of lists and MLCommons are primary for measured performance, with the caveats of slides 28 and 30. Research papers on portability (Reguly; Davis et al.; the ORNL report) and independent measurement (SemiAnalysis) are the only sources in the deck for achieved rather than peak performance, and the notes rely on them for slides 25 and 26. Press reporting is used where nothing better exists---Helios rack power, the cancellation of Falcon Shores, market share---and is marked as such.

Three sources deserve particular mention because the deck's conclusions depend on them. The MI300X Hot Chips 36 and ISSCC 2024 papers are the basis for the AMD packaging slides, and the deck's research found a discrepancy between them and the CDNA 3 whitepaper on the Infinity Cache bandwidth (17.2 against 14.7 TB/s) and HBM bandwidth (5.3 against 5.2 TB/s), resolved in favor of the whitepaper. Gomes et al.'s ISSCC 2022 paper on Ponte Vecchio is the source for the tile census (47 active tiles, 63 with thermal tiles, 11 EMIB bridges), which press accounts garble; the deck's figures follow the paper. And the Vera Rubin NVL72 product page exists in at least two versions, the January 2026 announcement figures (22 TB/s HBM4, 3.6 TB/s NVLink 6) and the September 2026 revision (19.2 and 3.0), with NVIDIA's October 29 statement explaining the difference; the deck uses the revised figures throughout and names the revision in every table that carries them.

The full reading list, with one entry per source, a line on why each matters and a link, is Appendix~D of these notes and the paper kit that accompanies them; the kit's script fetches the papers that are freely available and lists the rest with their landing pages.

\begin{terms}
\term{primary source}{A document produced by the party responsible for the fact: an architecture whitepaper or paper for a design, a datasheet for a product, a list or benchmark publication for a measurement.}
\term{architecture whitepaper}{A vendor's technical description of a GPU generation (NVIDIA's Hopper and Blackwell papers, AMD's CDNA papers, Intel's Gen and Xe documents); primary for microarchitecture, silent on prices and dates.}
\term{ISSCC and Hot Chips}{The International Solid-State Circuits Conference and the Hot Chips symposium, where vendors present die-level details (areas, tile counts, transistor counts) absent from datasheets.}
\term{SemiAnalysis}{An independent analysis firm whose December 2024 measurements of MI300X against H100 and H200 are the deck's source for achieved-versus-peak figures.}
\term{paper kit}{The companion archive to these notes: the reading list as CSV and BibTeX, a fetch script for the open-access items, and the appendix in LaTeX.}
\end{terms}

\begin{intuition}
Trust the document closest to the fact: the paper for the die, the datasheet for the box, the list for the measurement, the press release for the date. When two disagree, the deck says which it chose and why; when none exists, it says ``n.p.'' Do the same in your own work.
\end{intuition}

% =====================================================================
\clearpage
\appendix
\setstage{ink}
\section{Notation}

\begin{center}\scriptsize\setlength{\tabcolsep}{4pt}\renewcommand{\arraystretch}{1.05}
\begin{tabular}{@{}lp{4.7cm}p{7.0cm}@{}}
\toprule
symbol or term & meaning & notes\\
\midrule
SM, CU, Xe-core & the compute unit of NVIDIA, AMD CDNA and Intel Xe & 132 SMs on H100, 148 on GB200; 304 CUs on MI300X, 256 on MI355X (256 WGPs on MI455X); 128 Xe-cores on the Max 1550\\
EU, subslice, slice & Intel's Gen-era units & EU: 7 threads, two SIMD-4 FPUs, 16 FP32 FLOP per clock; subslice: 8 EUs; slice: 3 subslices (Gen9); renamed XVE, Xe-core, Xe-slice in Xe\\
XVE, SIMD, partition & the vector units inside a compute unit & Xe-HPC: 8 XVEs of 16 lanes; CDNA: 4 SIMDs of 16 lanes over 4 cycles (64-wide wavefront); Hopper: 4 partitions of 32 lanes (32-wide warp)\\
tensor core, matrix core, XMX & the matrix multiply-accumulate units & 4 per SM; 4 per CU; 8 per Xe-HPC core (16 per Xe-HPG core)\\
$N$, $L$, $f$ & number of compute units, FP32 lanes per unit, clock & vector peak $= 2\,N\,L\,f$; MI300X $2\times 304\times 64\times 2.1$ GHz $= 81.7$ TFLOP/s (FP64 vector; FP32 with packed math is twice)\\
$M$ & matrix FMAs per clock per compute unit at a format & matrix peak $= 2\,N\,M\,f$; MI300X FP16 $M = 1{,}024$: $2\times 304\times 1{,}024\times 2.1 = 1{,}307$ TFLOP/s\\
dense, sparse & tensor peaks without and with 2:4 structured sparsity & sparse $= 2\times$ dense; NVIDIA headlines sparse, AMD and Intel dense; all deck figures dense\\
vector / matrix FP64 & double precision on the vector lanes and on the matrix units & H100 34 / 67; MI300X 81.7 / 163.4; MI355X 78.6 / 78.6; PVC 52.4 (no matrix FP64); GB200 40; GB300 1.4; Rubin 33; MI430X 288 (announced)\\
$B$ & HBM bandwidth per GPU, bytes per second & 3.35 TB/s (H100), 5.3 (MI300X), 3.28 (PVC), 8.0 (GB300, MI355X), 19.2 (Rubin), 23.3 (MI455X)\\
$C$ & HBM capacity per GPU & 80 GB (H100), 192 (MI300X), 128 (PVC), 288 (GB300, MI355X, Rubin), 432 (MI455X)\\
$I^\ast$ & ridge point $= P/B$ at the fastest dense format & 591 FLOP/B (H100 FP8), 493 (MI300X FP8), 256 (PVC FP16), 1{,}250 (MI355X FP4), 1{,}823 (Rubin), 1{,}728 (MI455X)\\
XCD, IOD, FCD & AMD's accelerator complex die, I/O die, cache die & MI300X: 8 XCD + 4 IOD; MI355X: 8 + 2; MI455X: 8 XCD + 2 IOD + 2 FCD\\
GCD & graphics compute die & the two compute dies of the MI250X, seen as two GPUs by software\\
tile, base tile, RAMBO & Ponte Vecchio's chiplets & 47 active tiles: 16 compute, 8 RAMBO cache, 2 base, 2 Xe Link, 8 HBM2e, 11 EMIB bridges\\
Foveros, EMIB, CoWoS & Intel's 3D stacking and 2.5D bridge; TSMC's interposer packaging & PVC uses Foveros and EMIB; MI300 and Blackwell use CoWoS variants with hybrid bonding on MI300\\
IC, L2, LDS, SLM & Infinity Cache; level-2 cache; AMD local data share; Intel shared local memory & 256 MB IC (CDNA 3, 4); 192 MB L2 (MI455X); 408 MB L2 (PVC); LDS 64 KB (CDNA 3), 160 KB (CDNA 4); SLM 512 KB per Xe-HPC core\\
NVLink, IF, Xe Link, UALink & the scale-up fabrics & per GPU 900 GB/s (NVLink 4), 896 (MI300X IF), 1{,}800 (NVLink 5), 1{,}008 (MI355X), 3{,}000 (NVLink 6), 3{,}600 (MI455X UALoE)\\
UALoE & UALink over Ethernet & Helios's scale-up fabric: UALink protocol carried on Ethernet switches, 72 lanes per GPU\\
OAM, UBB, SXM, HGX & module and baseboard form factors & OAM/UBB are OCP standards used by AMD and Intel; SXM/HGX are NVIDIA's\\
NVL72, Helios & 72-GPU racks & NVIDIA: 72 GPUs, 36 CPUs, NVLink; AMD: 72 MI455X, 18 Venice, 18 Salina DPUs, UALoE\\
Rmax, Rpeak & HPL result and theoretical FP64 peak & efficiency $= R_{\max}/R_{\text{peak}}$: 66\% Frontier, 51\% Aurora, 64\% El Capitan, 81\% JUPITER (derived)\\
GFLOP/s per W & the Green500 metric & 73.3 (KAIROS, GH200), 66.5 (Ouranos, MI300A), June 2026\\
TDP & thermal design power per GPU & 700 W (H100), 750 (MI300X), 600 (PVC), 1{,}400 (GB300, MI355X); Rubin, MI455X n.p.\\
n.p. & not published & a figure the vendor had not disclosed as of November 2026\\
$^\dagger$ & secondary source & datasheet, product page, press release, benchmark list, press\\
\bottomrule
\end{tabular}
\end{center}

\clearpage
\section{Formulas at a glance}

\begin{description}[leftmargin=0pt,labelindent=0pt,itemsep=6pt,font=\normalfont\bfseries\sffamily]\sloppy
\item[Vector peak.] $P = 2\,N\,L\,f$ with $N$ compute units, $L$ FP32 lanes per unit and $f$ the clock; the 2 counts a fused multiply-add as two operations. MI300X FP64 vector: $2\times 304\times 64\times 2.1\,\mathrm{GHz} = 81.7$ TFLOP/s (CDNA does FP64 at the FP32 lane rate; packed FP32 doubles it to 163.4). MI355X: $2\times 256\times 64\times 2.4 = 78.6$. Max 1550: $2\times 128\times 128\times 1.6 = 52.4$ TFLOP/s for FP32 and for FP64, since Intel rates each XVE at 16 FP64 FMAs per clock, as for FP32. H100 FP32: $2\times 132\times 128\times 1.98 = 66.9 \approx 67$; FP64 at half the lanes, 34.
\item[FLOP per clock per unit.] CU: $4\ \text{SIMDs}\times 16\ \text{lanes}\times 2 = 128$ FP32 FLOP per clock (the 64-wide wavefront issues over four cycles). SM (Hopper): $4\times 32\times 2 = 256$. Xe-HPC core: $8\times 16\times 2 = 256$. Gen9 EU: $2\ \text{FPUs}\times 4\ \text{lanes}\times 2 = 16$; a 24-EU GT2 slice, 384. Whole chips: MI300X $304\times 128 = 38{,}912$; H100 $132\times 256 = 33{,}792$; PVC $128\times 256 = 32{,}768$ FLOP per clock; multiplied by 2.1, 1.98 and 1.6 GHz these give 81.7, 66.9 and 52.4 TFLOP/s (derived).
\item[Matrix peak.] $P_{\text{matrix}} = 2\,N\,M\,f$ with $M$ the dense FMAs per clock per unit at the format. MI300X FP16: $M = 1{,}024$, $2\times 304\times 1{,}024\times 2.1 = 1{,}307$ TFLOP/s; FP8 at $M = 2{,}048$, 2{,}615; FP64 matrix at $M = 128$, 163.4. MI355X FP16: $2\times 256\times M\times 2.4 = 2{,}500$ gives $M = 2{,}034 \approx 2{,}048$ (CDNA 4 doubled the FP16 rate per CU); FP4 at $M \approx 8{,}192$, 10{,}000; FP64 matrix at $M = 64$, 78.6 (halved from CDNA 3). PVC TF32: $1{,}024\ \text{XMX}\times 128\ \text{FMA}\times 2\times 1.6 = 419$ TFLOP/s; FP16 at twice the FMA rate, 839 (Intel's figure); INT8, 1{,}678. H100 FP16: $2\times 132\times 2{,}048\times 1.83 = 989$.
\item[Per-generation ratios.] AMD over NVIDIA at matched dates (derived): HBM capacity 2.4, 1.0, 1.5 (2023, 2025, 2026); HBM bandwidth 1.6, 1.0, 1.2; dense FP8 1.3, 1.0, 1.15; dense FP4 ---, 0.67, 1.15; FP64 matrix 2.4 (2023), 56 (2025: 78.6 against 1.4); scale-up per GPU 1.0, 0.56, 1.2.
\item[HBM per stack.] $B/\text{stacks}$: MI300X $5.3/8 = 0.66$ TB/s (HBM3); H100 $3.35/5 = 0.67$ (HBM3, five of six stacks active); PVC $3.28/8 = 0.41$ (HBM2e); GB300 and MI355X $8.0/8 = 1.0$ (HBM3e, 12-high); Rubin $19.2/8 = 2.4$ (HBM4); MI455X $23.3/12 = 1.94$ (HBM4). Capacity per stack: 24 GB (MI300X), 36 GB (12-high HBM3e and HBM4). All derived.
\item[Ridge point and decode time.] $I^\ast = P/B$ at the fastest dense format: H100 $1{,}979/3.35 = 591$ FLOP/B; MI300X $2{,}615/5.3 = 493$; PVC $839/3.28 = 256$; GB300 $15{,}000/8 = 1{,}875$; MI355X $10{,}000/8 = 1{,}250$; Rubin $35{,}000/19.2 = 1{,}823$; MI455X $40{,}260/23.3 = 1{,}728$. One decoding step of a model with $W$ bytes of weights takes at least $W/B$: a 70-billion-parameter model in FP8 (70 GB) on an MI300X, $70/5.3 = 13$ ms; on an H100, 21 ms; on a Rubin in NVFP4 (39 GB), 2 ms (derived).
\item[Scale-up arithmetic.] Infinity Fabric per GPU $= N_{\text{links}}\times b_{\text{link}}$: MI300X $7\times 128 = 896$ GB/s (bidirectional, each link to one neighbor); MI355X $7\times 144 = 1{,}008$. UALink per GPU $= \text{lanes}\times 200\ \mathrm{Gb/s}\times 2\ \text{directions}/8$: MI455X $72\times 200\times 2/8 = 3{,}600$ GB/s. NVLink per GPU: Blackwell $18\times 100 = 1{,}800$; Rubin 3{,}000 (revised from 3{,}600). Measured pair bandwidth against nominal: SemiAnalysis found 64 GB/s per MI300X pair (one link, one direction: $128/2$) against 450 for NVLink ($900/2$), both half the bidirectional figure, but NVLink's 450 reaches every GPU through the switch while the MI300X's 64 reaches one.
\item[Rack totals.] A rack figure is 72 times the per-GPU figure. Helios: HBM $72\times 432\,\mathrm{GB} = 31.1$ TB; bandwidth $72\times 23.3 = 1{,}678$ TB/s $\approx 1.7$ PB/s; dense FP4 $72\times 40.26 = 2{,}899$ PFLOP/s $\approx 2.9$ EF; FP8 $72\times 20.13 = 1{,}449 \approx 1.4$ EF; scale-up $72\times 3.6 = 259 \approx 260$ TB/s. Vera Rubin NVL72: $72\times 288 = 20.7$ TB; $72\times 19.2 = 1{,}382 \approx 1{,}400$ TB/s; $72\times 35 = 2{,}520$ PFLOP/s; $72\times 3.0 = 216$ TB/s; FP64 $72\times 33 = 2{,}376 \approx 2{,}400$ TFLOP/s. GB300 NVL72 FP64: $72\times 1.4 \approx 100$ TFLOP/s. Helios over Vera Rubin: HBM $31/20.7 = 1.5$; FP4 $2.9/2.52 = 1.15$ (derived).
\item[HPL efficiency and efficiency per watt.] $\eta = R_{\max}/R_{\text{peak}}$: Frontier $1{,}353/2{,}055 = 66\%$; Aurora $1{,}012/1{,}980 = 51\%$; El Capitan $1{,}809/2{,}821 = 64\%$; JUPITER $1{,}000/1{,}226 = 81\%$. HPL per megawatt $= R_{\max}/P$: 55, 26, 61, 63 PFLOP/s per MW; LineShine $2{,}198/42.2 = 52$; HPC7 $572/8.7 = 66$; Alps $435/7.1 = 61$; Fugaku $442/29.9 = 15$. The Green500 metric is the same quantity in GFLOP/s per W from a power measurement made during the run (all derived from the June 2026 list).
\item[Peak from counts.] Aurora's June 2024 run used $9{,}234\times 6 = 55{,}404$ GPUs; at the datasheet 52.4 TFLOP/s the GPU peak alone would be 2.90 EFLOP/s, but the listed Rpeak is 1.98 EFLOP/s, which implies about 35.7 TFLOP/s per GPU ignoring the CPUs, or a clock near 1.1 GHz against the 1.6 GHz behind the datasheet figure (derived): the submitted peak uses the clock the GPUs sustain under power limits. El Capitan: $44{,}544\times 61.3 = 2{,}731$ PFLOP/s from the APUs' FP64 vector rate; the listed 2{,}821 includes the Zen 4 cores (derived).
\item[MLPerf ratios.] Training: $r = t_{\text{B200}}/t_{\text{MI355X}}$; $9.85/10.18 = 0.97$, $93.7/99.7 = 0.94$. Inference: $r = (T/n)_{\text{MI355X}}/(T/n)_{\text{B200}}$ with $T$ system throughput and $n$ accelerator count; 1.00, 0.97, 1.19 (AMD's figures). Scaling: $1{,}042{,}110/87 \approx 12{,}000$ tokens per second per GPU offline.
\item[Achieved fraction.] $\alpha = P_{\text{measured}}/P_{\text{peak}}$: MI300X BF16 GEMM $620/1{,}307 = 47\%$; H100 $720/989 = 73\%$ (SemiAnalysis, December 2024); Aurora HPL 51\%; Frontier 66\%. A comparison of two GPUs at a date is $\alpha_1 P_1$ against $\alpha_2 P_2$, not $P_1$ against $P_2$: in December 2024, $0.47\times 1{,}307 = 614$ against $0.73\times 989 = 722$, so the GPU with the higher peak was the slower one (derived).
\item[Gigawatts to GPUs.] $N = P_{\text{facility}}/p_{\text{all-in}}$, $p_{\text{all-in}} = \text{TDP}\times(1 + \text{rack overhead})\times\text{PUE}$. At 1.4 kW, 25\% and 1.2: 2.1 kW, so 6 GW $\approx 2.9$ million GPUs; from the Helios rack figure ($245\ \mathrm{kW}/72 = 3.4$ kW) and PUE 1.2, about 1.5 million; at a bare 1.5 kW, 4 million. Abilene: $1.2\ \mathrm{GW}/450{,}000 = 2.7$ kW per GB200 against $120/72 = 1.67$ kW at the rack, a facility factor of 1.6 (all derived).
\item[Announced versus shipped.] $(\text{shipped} - \text{announced})/\text{announced}$: Rubin HBM4 $(19.2 - 22)/22 = -13\%$; NVLink 6 $(3.0 - 3.6)/3.6 = -17\%$; MI455X HBM4 $(23.3 - 19.6)/19.6 = +19\%$; Helios HBM bandwidth $(1.7 - 1.4)/1.4 = +21\%$. The MI455X-to-Rubin bandwidth ratio moved from $19.6/22 = 0.89$ at announcement to $23.3/19.2 = 1.21$ as shipped (derived).
\item[Gen to Xe scaling.] FLOP per clock $= \text{EUs}\times 16$ (Gen) or $\text{Xe-cores}\times 256$ (Xe-HPC): Gen9 GT2 $24\times 16 = 384$; Xe-LP $96\times 16 = 1{,}536$; PVC $128\times 256 = 32{,}768$, 21 times Xe-LP per clock and 85 times Gen9 GT2 (derived); at 1.15, 1.35 and 1.6 GHz, about 0.44, 2.07 and 52.4 TFLOP/s.
\end{description}

\clearpage
\section{Exercises}

\begin{enumerate}[leftmargin=1.6em,itemsep=6pt]
\item \textbf{Three compute units.} Using the table on slide 14, compute the FP32 FLOP per clock of one CDNA 3 compute unit, one Hopper SM and one Xe-HPC core, then the whole-chip figures for the MI300X (304 CUs), H100 (132 SMs) and Max 1550 (128 cores), and multiply by 2.1, 1.98 and 1.6 GHz. Which of the three vector peaks on slide 31 do you reproduce, and which needs a note? \emph{Answer: 128, 256 and 256 FLOP per clock; 38,912, 33,792 and 32,768 per clock; 81.7, 66.9 and 52.4 TFLOP/s. The MI300X figure is its FP64 vector peak, because CDNA runs FP64 at the lane rate and reaches 163.4 for FP32 only with packed math; the H100 reproduces its 67; PVC its 52.4 for both FP32 and FP64.}
\item \textbf{The CDNA 4 halving.} From the MI300X (304 CUs, 2.1 GHz, 163.4 TFLOP/s matrix FP64) and the MI355X (256 CUs, 2.4 GHz, 78.6), compute the matrix FP64 FMAs per clock per CU of each generation. Do the same for dense FP16 (1,307 and 2,500 TFLOP/s). What did CDNA 4 trade for what? \emph{Answer: $163.4\times 10^{12}/(2\times 304\times 2.1\times 10^9) = 128$ FMA per clock per CU on CDNA 3 and $78.6\times 10^{12}/(2\times 256\times 2.4\times 10^9) = 64$ on CDNA 4, halved; FP16 1,024 to 2,034 (about 2,048), doubled. The matrix unit's area went to low precision; FP64 fell to the vector rate.}
\item \textbf{Dense and sparse, three vendors.} An NVIDIA page states ``20 PFLOPS FP4'' for the GB200, an AMD page ``10 PFLOPS FP4'' for the MI355X and ``20 PFLOPS with sparsity'', and Intel's Gaudi 3 brief ``1,835 TFLOPS FP8''. Put the three on a dense basis and state which convention each vendor uses in its headline. \emph{Answer: GB200 10 PFLOP/s dense (NVIDIA headlines sparse); MI355X 10 dense (AMD headlines dense and gives the sparse figure separately); Gaudi 3 1.835 PFLOP/s dense (Intel gives one figure for BF16 and FP8, with no sparsity mode). The GB200 and MI355X are equal in dense FP4; the GB300's 15 is the 2025 lead.}
\item \textbf{HBM per stack.} From slide 31, compute the bandwidth per HBM stack for the MI300X (8 stacks), H100 (5 of 6 stacks active), Max 1550 (8), MI355X (8), Rubin (8) and MI455X (12). Which two 2026 parts use HBM4, and why do their per-stack figures differ? \emph{Answer: 0.66, 0.67, 0.41, 1.0, 2.4 and 1.94 TB/s per stack. Rubin and the MI455X both use HBM4; Rubin's eight stacks run at a higher per-stack rate (2.4 TB/s, a faster pin speed or wider interface) while AMD reaches a higher total with twelve stacks at 1.94. The capacities per stack are both 36 GB ($288/8$ and $432/12$).}
\item \textbf{Ridge points.} Compute $I^\ast = P/B$ at the fastest dense format for the H100, MI300X, Max 1550, MI355X, Rubin and MI455X. For a decoding workload with intensity of order 2 FLOP/B at batch 1, which GPU is fastest per token, and by what measure? \emph{Answer: 591, 493, 256, 1,250, 1,823 and 1,728 FLOP/B. All six are bandwidth-bound at that intensity, so time per token is $W/B$ and the MI455X at 23.3 TB/s is fastest by the ratio of bandwidths (1.21 over Rubin, 4.4 over the MI300X); the FLOP/s peaks do not enter.}
\item \textbf{Capacity against bandwidth.} A 405-billion-parameter model in FP8 (405 GB) is to be served with the weights split across the fewest GPUs that hold it. How many MI300X (192 GB, 5.3 TB/s) and how many H100 (80 GB, 3.35 TB/s) are needed, and what is the minimum time per token at batch 1 on each (ignore the KV cache and communication)? \emph{Answer: 3 MI300X (576 GB) against 6 H100 (480 GB); per GPU 135 GB at 5.3 TB/s is 25 ms, and 67.5 GB at 3.35 TB/s is 20 ms. The H100 set is about 25\% faster per token with twice the GPUs; per GPU the MI300X is 1.6 times as productive, which is the memory argument of slide 5.}
\item \textbf{Two fabrics, one number.} The MI300X has seven Infinity Fabric links of 128 GB/s and the H100 900 GB/s of NVLink 4; SemiAnalysis measured about 64 GB/s between a pair of MI300X and 450 between a pair of H100. Express both measurements as fractions of the nominal figures and explain why the two fractions are the same yet the systems differ. \emph{Answer: $64/128 = 0.5$ and $450/900 = 0.5$: each measurement is one direction of a bidirectional figure, so both fabrics deliver their nominal one-way rate. The difference is topology: an MI300X pair shares one 128 GB/s link, while an H100 pair communicates through NVSwitch and can use all 900 GB/s; for an all-to-all collective across eight GPUs the H100's advantage is about $7\times$, not 1.}
\item \textbf{UALink lanes.} The MI455X has 72 UALink lanes at 200 Gb/s each. Compute its scale-up bandwidth in GB/s (bidirectional) and the aggregate for a 72-GPU Helios rack, and compare with the Vera Rubin NVL72's 3.0 TB/s per GPU and 216 TB/s. \emph{Answer: $72\times 200\ \mathrm{Gb/s}\times 2/8 = 3{,}600$ GB/s per GPU; $72\times 3.6 = 259 \approx 260$ TB/s per rack; 1.2 times Rubin's per-GPU and rack figures as revised in September 2026 (equal to them as announced in January).}
\item \textbf{Rack totals.} From the per-GPU figures of the MI455X (432 GB, 23.3 TB/s, 40.26 PFLOP/s MXFP4, 20.13 MXFP8) compute the Helios rack totals and compare with AMD's stated 31 TB, 1.7 PB/s, 2.9 EF and 1.4 EF. Then compute the ratio of Helios to Vera Rubin NVL72 on HBM capacity and dense FP4. \emph{Answer: 31.1 TB, 1,678 TB/s, 2,899 PFLOP/s, 1,449 PFLOP/s; all agree to rounding. Helios over Vera Rubin: $31.1/20.7 = 1.5$ on HBM, $2.9/2.52 = 1.15$ on FP4.}
\item \textbf{HPL efficiency.} From slide 27, compute $R_{\max}/R_{\text{peak}}$ and $R_{\max}$ per megawatt for Frontier, Aurora, El Capitan and JUPITER. Then rank the four by each measure and explain why the rankings agree on the first and last places but the explanations differ. \emph{Answer: efficiency 66\%, 51\%, 64\%, 81\%; per MW 55, 26, 61, 63 PFLOP/s. JUPITER is first on both (the newest machine, the H100's FP64 tensor cores, the lowest listed power) and Aurora last on both (PVC's FP64 at the FP32 vector rate with no matrix path, a late and power-hungry package). In between, El Capitan beats Frontier on both: the APU's shared memory and CDNA 3's doubled matrix FP64.}
\item \textbf{Peak from counts.} Aurora's exaflop run used 9,234 nodes of six Max 1550 GPUs. Compute the GPU count, the GPU peak at the datasheet 52.4 TFLOP/s, and the per-GPU FP64 rate implied by the listed Rpeak of 1.98 EFLOP/s (ignore the CPUs). What clock does the implied rate correspond to, and what does the exercise say about datasheet peaks? \emph{Answer: 55,404 GPUs; 2.90 EFLOP/s; $1{,}980/55{,}404 = 35.7$ TFLOP/s per GPU, which with 16,384 lanes is $35.7\times 10^{12}/(2\times 16{,}384) = 1.09$ GHz against the 1.6 GHz behind 52.4. The submitted Rpeak uses the clock the GPUs sustain at their power limit; against the datasheet peak Aurora's HPL efficiency would be 35\%, against the listed Rpeak 51\%. (Derived; the CPUs account for a few percent of Rpeak.)}
\item \textbf{Green500 against the TOP500.} Isambard-AI phase 1 (GH200) is 68.8 GFLOP/s per W on the Green500; JUPITER (GH200) is 1,000 PFLOP/s at 15.8 MW on the TOP500. Compute JUPITER's efficiency in the same unit and the ratio of the two. Do the same for Ouranos (MI300A, 66.5) against El Capitan (1,809 PFLOP/s, 29.7 MW). What does the pair of ratios say about scale? \emph{Answer: JUPITER 63.3 GFLOP/s per W, ratio $68.8/63.3 = 1.09$; El Capitan 60.9, ratio $66.5/60.9 = 1.09$. The same chip is about 9\% less efficient at exascale than in a small system on both vendors' silicon, the cost of network, cooling and the Green500's measurement rules; the vendor gap (GH200 over MI300A, about 1.04--1.10) is of the same size as the scale effect.}
\item \textbf{MLPerf ratios.} From the Training 5.1 times (MI355X 10.18 and 99.7 minutes; B200 9.85 and 93.7; B300 9.59 and 95.1), compute the MI355X's ratio to each NVIDIA part on each test, and explain why the B300 is faster than the B200 on LoRA but slower on Llama 3.1 8B pretraining. \emph{Answer: against B200 0.97 and 0.94; against B300 0.94 and 0.95. The B300 has 50\% more FP4 arithmetic but the same HBM and NVLink as the B200 and these benchmarks run in BF16 or FP8, so its advantage is small; the 8B pretraining result ($95.1 > 93.7$) is within run-to-run spread between different submitters' systems, which is also why differences of a few percent between vendors should not be read as rankings.}
\item \textbf{Achieved against peak.} In December 2024 the MI300X reached about 620 of 1,307 TFLOP/s on large BF16 matrix multiplies and the H100 about 720 of 989. Compute the achieved fractions, the ratio of achieved rates, and the ratio of peaks. If AMD's libraries later reached the H100's fraction, what would the MI300X achieve? \emph{Answer: 47\% and 73\%; achieved ratio $620/720 = 0.86$ (MI300X slower); peak ratio $1{,}307/989 = 1.32$ (MI300X faster); at 73\% the MI300X would reach $0.73\times 1{,}307 = 954$ TFLOP/s, 1.32 times the H100. The gap between the two ratios is the software gap of slide 26, and MLPerf 6.0 (slide 30) suggests that the MI355X had closed most of it by 2026.}
\item \textbf{Gigawatts.} The OpenAI--AMD agreement is for 6 GW. Estimate the number of GPUs it corresponds to under three assumptions: 1.5 kW per GPU all-in; a TDP of 1.4 kW with 25\% rack overhead and a PUE of 1.2; and Helios racks at 245 kW for 72 GPUs with a PUE of 1.2. Then compute how many Helios racks 6 GW is under the third assumption. \emph{Answer: 4.0 million; $6\times 10^9/(1.4\times 1.25\times 1.2\times 10^3) = 2.9$ million; $6\times 10^9/(3.4\times 1.2\times 10^3) = 1.5$ million GPUs, or about 20,400 racks. The spread is a factor of 2.7, which is why the deck says ``on the order of three million'' and why such commitments are better read in gigawatts than in GPUs.}
\item \textbf{Announced against shipped, both vendors.} Vera Rubin's HBM4 bandwidth went from 22 TB/s (January 2026) to 19.2 (September 2026) and its NVLink 6 from 3.6 to 3.0 TB/s; the MI455X's HBM4 bandwidth went from 19.6 TB/s (2025) to 23.3 (July 2026) and the Helios rack's from 1.4 to 1.7 PB/s. Compute the percentage changes and the MI455X-to-Rubin bandwidth ratio before and after. Using slide 31, which other 2026 figures are still unpublished, and how should a comparison written in November 2026 treat them? \emph{Answer: $-13\%$, $-17\%$, $+19\%$, $+21\%$; the ratio moved from 0.89 to 1.21. Unpublished: Rubin's and the MI455X's power, the MI455X's FP64, Helios's FP64, and Rubin's L2. Write ``n.p.'' and compare only the rows both vendors have stated; a specification that moved twice in a year is a target until a measurement replaces it.}
\end{enumerate}
%=====
\clearpage
\setstage{ink}
\section{Reading list}
\label{app:reading}
The papers, whitepapers, datasheets and reports behind each part of the deck, in the order they are discussed, with one line on why each is worth reading. Entries with an arXiv identifier are included in the paper archive of the kit; entries without one are listed with a URL. Items marked \emph{secondary source} are press coverage, analyst notes or commentary; vendor material is named as such.

\subsection*{1. Overview: three vendors, six levels (slides 1-3)}
\begin{itemize}[leftmargin=1.2em,itemsep=3pt]\sloppy
\item \textbf{SemiAnalysis InferenceX (2026).} \textit{GPU specifications table (InferenceX)}. \url{https://inferencex.semianalysis.com/gpu-specs}. Analyst compilation, secondary source. Cross-vendor dense tensor peaks on one page; the sanity check for the three-vendor chart. Slide 2.
\item \textbf{Silicon Analysts (2026).} \textit{AMD vs NVIDIA AI GPU market share 2026}. \url{https://siliconanalysts.com/analysis/amd-vs-nvidia-ai-gpu-market-share-2026}. Analyst estimate, secondary source. The share figures the deck quotes as estimates only (NVIDIA about 80 percent, AMD under 10, Intel under 1). Slides 2 and 29.
\item \textbf{TrendForce (2026).} \textit{Rubin faces delays; Blackwell over 70 percent of 2026 shipments (April 2026)}. \url{https://www.trendforce.com/presscenter/news/20260408-13003.html}. Analyst press release, secondary source. The 2026 NVIDIA product mix behind the share discussion. Slides 2 and 33.
\item \textbf{Lockwood (2026).} \textit{Rubin R200 (processor garden page)}. \url{https://glennklockwood.com/garden/processors/R200}. Independent, open. A sourced compilation of every Rubin number with its date, including the NVL144 to NVL72 renaming; the provenance check for the Rubin column throughout. Slides 2, 4 and 31.
\item \textbf{Open Compute Project (2019).} \textit{OCP Accelerator Module (OAM) wiki and specifications}. \url{https://www.opencompute.org/wiki/Server/OAM}. Standards body, open. The OAM module and UBB baseboard specifications that AMD and Intel accelerators ship on; the board level of the six-level table. Slides 3 and 23.
\end{itemize}
\subsection*{2. The NVIDIA line as platforms (slides 4-6)}
\begin{itemize}[leftmargin=1.2em,itemsep=3pt]\sloppy
\item \textbf{NVIDIA (2020).} \textit{HGX A100 datasheet}. \url{https://www.nvidia.com/content/dam/en-zz/Solutions/Data-Center/HGX/pdf/nvidia-hgx-a100-datasheet.pdf}. Vendor datasheet. The 8-GPU board: 40/80 GB per GPU, 600 GB/s NVLink 3 per GPU, 4.8 TB/s aggregate through six NVSwitches; the first row of the platform table. Slides 4 and 21.
\item \textbf{NVIDIA (2023).} \textit{H100 Tensor Core GPU datasheet}. \url{https://openzeka.com/wp-content/uploads/2023/12/h100-datasheet-1.pdf}. Vendor datasheet (mirror). The shipping H100 peaks with the sparsity footnote: TF32 494/989, FP16 989/1,979, FP8 1,979/3,958 dense/sparse; 80 GB at 3.35 TB/s, 700 W. Slides 4, 6 and 18.
\item \textbf{CoreWeave documentation (2024).} \textit{NVIDIA HGX H100 (platform documentation)}. \url{https://docs.coreweave.com/coreweave-machine-learning-and-ai/nvidia-hgx-h100}. Cloud-provider documentation, secondary source. The 8-GPU HGX H100 totals (640 GB, 900 GB/s NVLink 4 per GPU) as deployed. Slide 4.
\item \textbf{NVIDIA (2023).} \textit{NVIDIA supercharges Hopper, the world's leading AI computing platform (press release, November 13, 2023)}. \url{https://nvidianews.nvidia.com/news/nvidia-supercharges-hopper-the-worlds-leading-ai-computing-platform}. Vendor press release. The H200: 141 GB of HBM3e at 4.8 TB/s on the unchanged GH100 die, available Q2 2024; the memory refresh in the platform table and the memory chart. Slides 4 and 5.
\item \textbf{Lenovo Press (LP1944) (2025).} \textit{NVIDIA H200 product guide}. \url{https://lenovopress.lenovo.com/lp1944.pdf}. Vendor documentation. The H200 dense and sparse peak table side by side, the clearest statement that the H200 compute equals the H100's. Slide 4.
\item \textbf{NVIDIA (2025).} \textit{GB200 NVL72 (product page)}. \url{https://www.nvidia.com/en-us/data-center/gb200-nvl72/}. Vendor page. The shipping rack with dense and sparse footnotes: 1,440/720 PFLOP/s FP4, 2,880 TFLOP/s FP64, 13.4 TB HBM3e at 576 TB/s, 130 TB/s NVLink; per GPU 186 GB at 8 TB/s. Slides 4, 6 and 22.
\item \textbf{The Register (2024).} \textit{Nvidia DGX GB200 NVL72: 120 kW and 1,200 W per GPU (March 2024)}. \url{https://www.theregister.com/2024/03/21/nvidia_dgx_gb200_nvk72}. Press, secondary source. The rack power of about 120 kW and 1,200 W per B200 in the GB200 superchip, figures NVIDIA's pages do not state. Slides 4 and 22.
\item \textbf{NVIDIA (2025).} \textit{GB300 NVL72 (product page)}. \url{https://www.nvidia.com/en-gb/data-center/gb300-nvl72}. Vendor page. 1,440/1,080 PFLOP/s FP4 sparse/dense, 100 TFLOP/s FP64 per rack (1.4 per GPU), 288 GB per GPU, 20 TB HBM3e, 1,400 W. Slides 4, 6, 12 and 22.
\item \textbf{Tom's Hardware (2025).} \textit{NVIDIA's newest top-tier AI supercomputers deployed for the first time at CoreWeave (July 2025)}. \url{https://www.tomshardware.com/tech-industry/artificial-intelligence/nvidias-newest-top-tier-ai-supercomputers-deployed-for-the-first-time-grace-blackwell-ultra-superchip-systems-deployed-at-coreweave}. Press, secondary source. The first GB300 NVL72 racks, July 3, 2025: the shipment date in the rack table. Slide 22.
\item \textbf{Tom's Hardware (2024).} \textit{Dell ships the industry's first GB200 NVL72 racks to CoreWeave (November 18, 2024)}. \url{https://www.tomshardware.com/tech-industry/artificial-intelligence/dell-reaches-milestone-with-industrys-first-enterprise-ready-nvidia-blackwell-poweredge-xe9712-server-racks}. Press, secondary source. The first complete GB200 NVL72 racks (Dell PowerEdge XE9712) delivered in November 2024; the first-shipment date in the rack table. Slide 22.
\item \textbf{Tom's Hardware (2026).} \textit{Nvidia launches Vera Rubin NVL72 at CES (January 2026)}. \url{https://www.tomshardware.com/pc-components/gpus/nvidia-launches-vera-rubin-nvl72-ai-supercomputer-at-ces-promises-up-to-5x-greater-inference-performance-and-10x-lower-cost-per-token-than-blackwell-coming-2h-2026}. Press, secondary source. The January 2026 numbers (288 GB HBM4 at 22 TB/s, 3.6 TB/s NVLink 6, 260 TB/s per rack) that the September 2026 revision later reduced. Slides 4 and 22.
\item \textbf{NVIDIA (2026).} \textit{Vera Rubin NVL72 (product page, revised September 2026)}. \url{https://www.nvidia.com/en-us/data-center/vera-rubin-nvl72}. Vendor page. The current Rubin specifications used throughout: 19.2 TB/s HBM4 per GPU, 3.0 TB/s NVLink 6, 216 TB/s and 1,400 TB/s per rack, 33 TFLOP/s FP64, 35/50 PFLOP/s NVFP4. Slides 4, 5, 6, 12, 22 and 31.
\item \textbf{Hardware Busters (2026).} \textit{Nvidia quietly lowers Vera Rubin HBM4 bandwidth and the HGX NVL8 loses a quarter of it (October 2026)}. \url{https://hwbusters.com/news/nvidia-quietly-lowers-vera-rubin-hbm4-bandwidth-and-the-hgx-nvl8-loses-a-quarter-of-it/}. Press, secondary source. The old-versus-new table of the September 2026 revision, dated with the Wayback Machine; why the deck gives the revised Rubin values. Slides 4 and 22.
\item \textbf{Wccftech (2026).} \textit{NVIDIA's official statement on the Vera Rubin spec change (October 29, 2026)}. \url{https://wccftech.com/nvidia-official-statement-on-vera-rubin-spec-change/}. Press, secondary source. NVIDIA's statement that the earlier figures were preliminary architectural targets and the new ones are what qualified production sustains. Slides 4 and 32.
\end{itemize}
\subsection*{3. AMD Instinct and CDNA (slides 7-12)}
\begin{itemize}[leftmargin=1.2em,itemsep=3pt]\sloppy
\item \textbf{AMD (2020).} \textit{Introducing AMD CDNA architecture (MI100 whitepaper)}. \url{https://www.amd.com/content/dam/amd/en/documents/instinct-business-docs/white-papers/amd-cdna-white-paper.pdf}. Vendor whitepaper. The CDNA 1 compute unit with matrix cores, 120 CUs, 11.5 TFLOP/s FP64, 32 GB HBM2 at 1.2 TB/s, three Infinity Fabric links at 92 GB/s; the GCN-to-CDNA split. Slides 7 and 12.
\item \textbf{AMD (2021).} \textit{AMD CDNA 2 architecture (MI200 whitepaper)}. \url{https://www.amd.com/content/dam/amd/en/documents/instinct-business-docs/white-papers/amd-cdna2-white-paper.pdf}. Vendor whitepaper. Two graphics compute dies per module, 400 GB/s in-package Infinity Fabric, packed FP32, FP64 at the FP32 rate; 47.9/95.7 TFLOP/s FP64 vector/matrix. Slides 7, 8 and 12.
\item \textbf{AMD (2023).} \textit{AMD CDNA 3 architecture whitepaper}. \url{https://www.amd.com/content/dam/amd/en/documents/instinct-tech-docs/white-papers/amd-cdna-3-white-paper.pdf}. Vendor whitepaper. The 3D chiplet package (8 XCDs on 4 IODs), 153 billion transistors, 256 MB Infinity Cache at 17.2 TB/s, 896 GB/s of peer links; MI300X and MI300A peaks. Slides 7, 9 and 12.
\item \textbf{AMD (2025).} \textit{AMD CDNA 4 architecture whitepaper}. \url{https://www.amd.com/content/dam/amd/en/documents/instinct-tech-docs/white-papers/amd-cdna-4-architecture-whitepaper.pdf}. Vendor whitepaper. The MI355X: N3P XCDs on N6 IODs, 185 billion transistors, MXFP4 and MXFP6 rates, matrix FP64 halved per CU, Infinity Fabric 7 x 144 GB/s. Slides 7, 10 and 12.
\item \textbf{AMD (2021).} \textit{Instinct MI200 series datasheet}. \url{https://www.amd.com/content/dam/amd/en/documents/instinct-tech-docs/instinct-mi200-datasheet.pdf}. Vendor datasheet. MI250X and MI250 peaks, 128 GB HBM2e at 3.2 TB/s, eight Infinity Fabric links totaling 800 GB/s, 500-560 W. Slides 6, 8 and 18.
\item \textbf{AMD (2022).} \textit{Instinct MI210 brochure}. \url{https://www.amd.com/content/dam/amd/en/documents/instinct-business-docs/product-briefs/instinct-mi210-brochure.pdf}. Vendor brochure. The single-die PCIe CDNA 2 part, for the CDNA 2 product range. Slide 8.
\item \textbf{AMD (2023).} \textit{Instinct MI300A datasheet}. \url{https://www.amd.com/content/dam/amd/en/documents/instinct-tech-docs/data-sheets/amd-instinct-mi300a-data-sheet.pdf}. Vendor datasheet. The APU: 228 CUs and 24 Zen 4 cores sharing 128 GB of HBM3, 61.3/122.6 TFLOP/s FP64, 550-760 W; El Capitan's processor. Slides 9 and 27.
\item \textbf{AMD (2023).} \textit{Instinct MI300X accelerator product brief}. \url{https://www.amd.com/content/dam/amd/en/documents/partner-hub/instinct/instinct-mi300x-accelerator-product-briefs.pdf}. Vendor product brief. The MI300X on the 8-GPU UBB 2.0 platform: 304 CUs, 192 GB at 5.3 TB/s, 750 W, the peak table. Slides 6 and 9.
\item \textbf{AMD (2024).} \textit{Instinct MI325X datasheet}. \url{https://www.amd.com/content/dam/amd/en/documents/instinct-tech-docs/product-briefs/instinct-mi325x-datasheet.pdf}. Vendor datasheet. 256 GB of HBM3e at 6 TB/s and 1,000 W on the MI300X silicon; the answer to the H200 in the memory chart. Slides 5 and 9.
\item \textbf{AMD (2025).} \textit{Instinct MI355X GPU brochure}. \url{https://www.amd.com/content/dam/amd/en/documents/instinct-tech-docs/product-briefs/amd-instinct-mi355x-gpu-brochure.pdf}. Vendor brochure. 78.6 TFLOP/s FP64 vector and matrix, 10 PFLOP/s dense FP4, 288 GB at 8 TB/s, 1,400 W liquid cooled. Slides 6, 10 and 12.
\item \textbf{AMD ROCm documentation (2024).} \textit{MI300 series microarchitecture}. \url{https://instinct.docs.amd.com/latest/_sources/gpu-arch/mi300.md}. Vendor documentation. The XCD has 40 CUs with 38 active; the FLOP per clock per CU table behind the CDNA 3 compute-unit diagram and the FP32 doubling. Slides 7 and 9.
\item \textbf{AMD ROCm documentation (2026).} \textit{MI250 microarchitecture}. \url{https://rocm.docs.amd.com/en/docs-7.14.1/_sources/reference/gpu-arch/mi250.md}. Vendor documentation. The two GCDs and the links between them; why software sees two GPUs per module. Slide 8.
\item \textbf{AMD ROCm documentation (2026).} \textit{MI100 microarchitecture}. \url{https://rocm.docs.amd.com/en/docs-7.14.1/reference/gpu-arch/mi100.html}. Vendor documentation. Eight shader engines of fifteen CUs: the MI100's 120 compute units. Slide 7.
\item \textbf{Smith and Alla (AMD) (2024).} \textit{MI300X generative AI accelerator and platform architecture (Hot Chips 36 slides)}. \url{https://www.hc2024.hotchips.org/assets/program/conference/day1/23_HC2024.AMD.MI300X.ASmith(MI300X).v1.Final.20240817.pdf}. Primary conference slides. The MI300X package, hybrid bonding, the Infinity Cache and the 8-GPU platform as AMD presented them. Slide 9.
\item \textbf{ServeTheHome (2025).} \textit{AMD dives deep on CDNA 4 architecture and MI350 at Hot Chips 2025}. \url{https://www.servethehome.com/amd-dives-deep-on-cdna-4-architecture-and-mi350-accelerator-at-hot-chips-2025/}. Press, secondary source. The Hot Chips 37 CDNA 4 talk: N3P XCDs on N6 IODs, the LDS growth to 160 KB, the FP64 decision. Slides 7 and 10.
\item \textbf{AMD (ISSCC 2024, reference page) (2024).} \textit{AMD Instinct MI300 series modular chiplet package: HPC and AI accelerator for exa-class systems (ISSCC 2024)}. \url{https://www.aihardware.mit.edu/amd-instincttm-mi300-series-modular-chiplet-package-hpc-and-ai-accelerator-for-exa-class-systems/}. Conference paper reference (title and abstract page). The ISSCC paper on the MI300 package; the primary source for the chiplet construction. Slide 9.
\item \textbf{Lam and Cozma (Chips and Cheese) (2023).} \textit{AMD's CDNA 3 compute architecture}. \url{https://chipsandcheese.com/p/amds-cdna-3-compute-architecture}. Independent analysis, secondary source. The CDNA 3 compute unit, cache hierarchy and the Infinity Cache measured and explained. Slides 7 and 9.
\item \textbf{Chips and Cheese (2025).} \textit{SC25: estimating MI430X FP64 and Discovery}. \url{https://old.chipsandcheese.com/?p=37639}. Independent analysis, secondary source. How the 288 TFLOP/s FP64 of the MI430X and the Discovery system size follow from AMD's statements. Slides 12 and 29.
\item \textbf{Phoronix (2025).} \textit{AMD previews Instinct MI400 series and Helios AI rack (June 2025)}. \url{https://www.phoronix.com/news/AMD-Instinct-MI400-Preview}. Press, secondary source. The June 12, 2025 numbers for MI400 and Helios (432 GB HBM4, 19.6 TB/s, 40 PFLOP/s FP4) before the 2026 launch revised them. Slide 11.
\item \textbf{Tom's Hardware (2025).} \textit{AMD debuts Helios rack-scale platform at OCP Global Summit 2025}. \url{https://www.tomshardware.com/tech-industry/amd-debuts-helios-rack-scale-ai-hardware-platform-at-ocp-global-summit-2025-promises-easier-serviceability-and-50-percent-more-memory-than-nvidias-vera-rubin}. Press, secondary source. Helios on the OCP Open Rack Wide form factor, double-wide for serviceability, and the memory comparison with Vera Rubin. Slides 11, 22 and 23.
\item \textbf{Tom's Hardware (2026).} \textit{AMD takes the wraps off Instinct MI455X (CDNA 5) and the Helios rack (July 2026)}. \url{https://www.tomshardware.com/pc-components/gpus/amd-takes-the-wraps-off-its-instinct-mi455x-ai-accelerator-cdna-5-and-helios-rack-scale-architecture-combine-to-take-the-fight-to-nvidia-in-the-data-center}. Press, secondary source. The MI455X launch of July 23, 2026: N2 XCDs, 320 billion transistors, 192 MB L2 with no Infinity Cache, 23.3 TB/s HBM4, 256 work-group processors. Slides 7, 11 and 31.
\item \textbf{StorageReview (2026).} \textit{AMD MI455X and Helios: 432 GB HBM4, 72-GPU racks (July 2026)}. \url{https://www.storagereview.com/news/amd-mi455x-and-helios-432gb-hbm4-72-gpu-racks-and-a-real-answer-to-vera-rubin}. Press, secondary source. UALink over Ethernet at 3.6 TB/s per GPU, 260 TB/s per rack, 225-245 kW; the Helios figures in the rack table. Slides 11, 21 and 22.
\item \textbf{Supermicro (2026).} \textit{AMD Helios rack datasheet}. \url{https://www.supermicro.com/datasheet/datasheet_AMD_Helios.pdf}. Vendor (integrator) datasheet. The rack bill of materials: 72 MI455X, 18 Venice CPUs, 18 Salina DPUs, Vulcano NICs, 31 TB HBM4 (listed at 1.4 PB/s). Slides 11 and 22.
\item \textbf{ServeTheHome (2026).} \textit{AMD Advancing AI 2026 keynote live coverage}. \url{https://www.servethehome.com/?p=102176}. Press, secondary source. The MI430X at 288 TFLOP/s FP64 and the MI500 (2027) and MI600 (2028) roadmap as announced. Slides 11, 12 and 29.
\item \textbf{Tech Wire Asia (2026).} \textit{AMD Advancing AI 2026: Helios, OpenAI, Meta, Anthropic}. \url{https://techwireasia.com/2026/07/amd-advancing-ai-2026-helios-openai-meta-anthropic/}. Press, secondary source. The customer commitments announced with Helios (OpenAI, Meta, Anthropic, Oracle). Slides 11 and 29.
\item \textbf{AMD (2025).} \textit{Instinct MI430X: powering the next wave of AI and scientific discovery}. \url{https://www.amd.com/en/blogs/2025/amd-instinct-mi430x-powering-the-next-wave-of-ai.html}. Vendor blog. The HPC SKU of the MI400 family, its hardware FP64 and the systems that chose it (Discovery, Alice Recoque). Slides 12 and 29.
\item \textbf{AMD (SEC Form 8-K) (2026).} \textit{AMD reports third quarter 2026 financial results}. \url{https://www.sec.gov/Archives/edgar/data/0000002488/000000248826000188/q32026991.htm}. Primary filing. Helios shipping, first Helios revenue, data-center segment of 8.5 billion dollars and 97 percent growth. Slides 11 and 29.
\item \textbf{AMD Q3 2026 earnings call (transcript) (2026).} \textit{AMD Q3 2026 earnings call transcript}. \url{https://singjupost.com/amd-advanced-micro-devices-q3-2026-earnings-call-transcript/}. Transcript, secondary source. The statement that the MI455X is shipping in volume and the customers named. Slide 11.
\item \textbf{StorageNewsletter (2026).} \textit{AMD fiscal 4Q25 and FY25 financial results}. \url{https://www.storagenewsletter.com/2026/02/18/amd-fiscal-4q25-and-fy25-financial-results/}. Press, secondary source. Data-center revenue of 16.6 billion dollars in fiscal 2025, the base for the 2026 growth figures. Slide 29.
\item \textbf{AMD investor relations (2025).} \textit{AMD and OpenAI announce strategic partnership to deploy 6 gigawatts of AMD GPUs (October 2025)}. \url{https://ir.amd.com/news-events/press-releases/detail/1260/amd-and-openai-announce-strategic-partnership-to-deploy-6-gigawatts-of-amd-gpus}. Vendor press release. The 6 GW commitment from the second half of 2026 and the warrant for up to 160 million shares. Slides 3, 11 and 29.
\end{itemize}
\subsection*{4. Intel: the Gen families, Xe and Gaudi (slides 13-16)}
\begin{itemize}[leftmargin=1.2em,itemsep=3pt]\sloppy
\item \textbf{Junkins (Intel) (2015).} \textit{The compute architecture of Intel processor graphics Gen9}. \url{https://www.intel.com/content/dam/develop/external/us/en/documents/the-compute-architecture-of-intel-processor-graphics-gen9-v1d0-843410.pdf}. Vendor whitepaper. The execution unit (7 threads, two SIMD-4 FPUs, 16 FLOP per clock), the subslice of 8 EUs with sampler and 64 KB shared local memory, the slice with 768 KB L3; the Gen9 diagram and table. Slide 13.
\item \textbf{Intel (2019).} \textit{Intel processor graphics Gen11 architecture}. \url{https://www.intel.com/content/dam/develop/external/us/en/documents/the-architecture-of-intel-processor-graphics-gen11-r1new.pdf}. Vendor whitepaper. The 64-EU Gen11 GPU of Ice Lake, 1,024 FLOP per clock, and the changes to the slice. Slide 13.
\item \textbf{Tom's Hardware (2020).} \textit{Intel drops Xe LP graphics specifications (August 2020)}. \url{https://www.tomshardware.com/news/intel-xe-lp-graphics-specs}. Press, secondary source. Xe-LP's 96 EUs and 768 ALUs in Tiger Lake and the DG1, the last Gen-style design and the first Xe. Slide 13.
\item \textbf{Intel (2022).} \textit{Introduction to the Xe-HPG architecture}. \url{https://www.intel.com/content/www/us/en/developer/articles/technical/introduction-to-the-xe-hpg-architecture.html}. Vendor documentation. The Xe-core, Xe Vector Engine and XMX naming and the Xe-HPG core (16 XVE, 16 XMX, ray-tracing unit) in the two-core diagram. Slide 14.
\item \textbf{Intel (2024).} \textit{oneAPI GPU optimization guide: Xe architecture}. \url{https://www.intel.com/content/www/us/en/docs/oneapi/optimization-guide-gpu/2024-1/xe-arch.html}. Vendor documentation. The Xe-HPC hierarchy (XVE, Xe-core, slice, stack) with lane counts, the 512 KB L1 and SLM per core, and the cross-vendor level table. Slides 14 and 18.
\item \textbf{Tom's Hardware (2022).} \textit{Intel reveals specifications for Arc Alchemist desktop GPUs (September 2022)}. \url{https://www.tomshardware.com/news/intel-reveals-specifications-for-arc-alchemist-desktop-gpus}. Press, secondary source. The ACM-G10 table: 32 Xe-cores, 21.7 billion transistors on 406 square millimeters of TSMC N6, 17.2 TFLOP/s, 16 GB at 560 GB/s, 225 W. Slide 15.
\item \textbf{Tom's Hardware (2024).} \textit{Intel Arc B580 review (December 2024)}. \url{https://www.tomshardware.com/pc-components/gpus/intel-arc-b580-review-the-new-usd249-gpu-champion-has-arrived}. Press review, secondary source. BMG-G21 on TSMC N5: 20 Xe2 cores with SIMD16 XVEs, 19.6 billion transistors on 272 square millimeters, 12 GB at 456 GB/s, 190 W. Slide 15.
\item \textbf{Tom's Hardware (2025).} \textit{Intel launches Arc Pro B50 and B60 and Project Battlematrix workstations (May 2025)}. \url{https://www.tomshardware.com/pc-components/gpus/intel-launches-usd299-arc-pro-b50-with-16gb-of-memory-project-battlematrix-workstations-with-24gb-arc-pro-b60-gpus}. Press, secondary source. The workstation Xe2 parts with 16 and 24 GB and the eight-B60 Battlematrix box for local inference. Slide 15.
\item \textbf{HotHardware (2026).} \textit{Intel CES 2026: Panther Lake (Core Ultra series 3) is a go}. \url{https://hothardware.com/news/intel-ces-2026-panther-lake-is-a-go}. Press, secondary source. The first Xe3 GPU: 12 Xe3 cores on Intel 18A, about 120 TOPS on the GPU, in a laptop processor. Slide 15.
\item \textbf{Kaplan (Intel) (2024).} \textit{Gaudi 3 (Hot Chips 2024 slides)}. \url{https://www.hc2024.hotchips.org/assets/program/conference/day1/60_HC2024.Intel.RomanKaplan.Gaudi3-0826.pdf}. Primary conference slides. Two dies, 8 matrix engines and 64 tensor processor cores, 128 GB HBM2e at 3.67 TB/s, 24 x 200 GbE on the die; the Gaudi 3 diagram. Slide 16.
\item \textbf{The Register (2022).} \textit{Intel's Habana unit reveals Gaudi 2 (May 2022)}. \url{https://www.theregister.com/2022/05/10/intels_habana_unit_reveals_new/}. Press, secondary source. Gaudi 2: 7 nm, 24 TPCs, 96 GB HBM2e at 2.45 TB/s, 24 x 100 GbE, 600 W. Slide 16.
\item \textbf{Tom's Hardware (2024).} \textit{Intel says it will miss its 500 million dollar Gaudi goal (October 2024)}. \url{https://www.tomshardware.com/tech-industry/artificial-intelligence/intel-says-it-will-miss-its-ai-goals-with-gaudi-3-unbaked-software-leaves-intels-usd500-million-ai-goal-unachievable-as-competitors-rake-in-billions}. Press, secondary source. The withdrawn 2024 revenue target and the software reasons Intel gave; the market outcome of the Gaudi line. Slides 16 and 20.
\end{itemize}
\subsection*{5. Xe-HPC: Ponte Vecchio, Aurora and after (slides 17-20)}
\begin{itemize}[leftmargin=1.2em,itemsep=3pt]\sloppy
\item \textbf{Intel (2023).} \textit{Intel Data Center GPU Max 1550 specifications (ARK)}. \url{https://www.intel.com/content/www/us/en/products/sku/232873/intel-data-center-gpu-max-1550/specifications.html}. Vendor specification page. 128 Xe-cores, 128 GB HBM2e at 3,276.8 GB/s, 600 W OAM, launch date January 2023. Slides 17, 18 and 31.
\item \textbf{Intel Labs (2022).} \textit{Intel at ISSCC 2022: Ponte Vecchio (Gomes et al.) and other papers}. \url{https://intel.com/content/www/us/en/research/blogs/isscc-2022.html}. Vendor research blog. The ISSCC 2022 paper on Ponte Vecchio as a multi-tile 3D stacked processor; the primary source for the tile census. Slide 17.
\item \textbf{TechPowerUp (2022).} \textit{Intel details Ponte Vecchio: 63 tiles, 600 W TDP and lots of bandwidth (ISSCC 2022)}. \url{https://techpowerup.com/292250/intel-details-ponte-vecchio-accelerator-63-tiles-600-watt-tdp-and-lots-of-bandwidth}. Press, secondary source. The tile census (47 active, 63 with thermal tiles), the five process nodes and the 408 MB of L2 as presented at ISSCC. Slide 17.
\item \textbf{The Register (2022).} \textit{Intel Ponte Vecchio at Hot Chips 34 (August 2022)}. \url{https://www.theregister.com/2022/08/23/intel_ponte_vecchio/}. Press, secondary source. 52 TFLOP/s FP64 equal to FP32 and 419 TFLOP/s TF32 on the XMX engines, the figures the deck uses for the Max 1550. Slides 17 and 18.
\item \textbf{Tom's Hardware (2021).} \textit{Intel Ponte Vecchio and Xe-HPC architecture: built for big data (Architecture Day 2021)}. \url{https://www.tomshardware.com/features/intel-ponte-vecchio-and-xe-hpc-architecture-built-for-big-data}. Press, secondary source. The stack composition (8 compute tiles, 4 RAMBO cache tiles, a base tile, Xe Link tiles) and the Xe-HPC core as first described. Slides 14 and 17.
\item \textbf{TechPowerUp GPU database (2023).} \textit{Intel Data Center GPU Max 1550}. \url{https://www.techpowerup.com/gpu-specs/data-center-gpu-max-1550.c4068}. Database, secondary source. The 52.43 TFLOP/s FP32 and FP64 figures, clock and lane counts behind the derived arithmetic. Slide 18.
\item \textbf{Lockwood (2025).} \textit{Ponte Vecchio (processor garden page)}. \url{https://glennklockwood.com/garden/processors/PVC}. Independent, open. The Xe-HPC hierarchy with sources and measured DGEMM on Aurora; the provenance check for the PVC column. Slides 17, 18 and 31.
\item \textbf{Phoronix (2022).} \textit{Intel introduces Xeon Max and Data Center GPU Max (November 2022)}. \url{https://www.phoronix.com/review/intel-xeon-max-dc/3}. Press, secondary source. The Max series SKUs as announced (1100, 1350, 1550) and the Xeon Max with HBM used in Aurora. Slides 17 and 19.
\item \textbf{Argonne Leadership Computing Facility (2026).} \textit{Aurora system overview (machine overview)}. \url{https://docs.alcf.anl.gov/aurora/hardware-overview/machine-overview}. Facility documentation. The node (two Xeon Max 9470C, six Max 1550, eight Slingshot NICs), 10,624 nodes in 166 racks, 63,744 GPUs. Slides 19 and 27.
\item \textbf{Goto et al. (2026).} \textit{HPL and HPL-MxP at exascale on Aurora}. arXiv:2604.09517. The run campaigns behind Aurora's TOP500 entries: 0.585 EFLOP/s on 5,439 nodes, 1.012 on 9,234, and 11.64 EFLOP/s HPL-MxP. Slide 19.
\item \textbf{HPCwire (2023).} \textit{Aurora the survivor: exascale supercomputer arrives after eight years (November 2023)}. \url{https://www.hpcwire.com/2023/11/13/aurora-the-survivor-exascale-supercomputer-arrives-after-eight-years-of-doom/}. Press, secondary source. The 2015 to 2023 history: Knights Hill, the 2017 re-plan, the 7 nm slip, installation in 2023. Slide 19.
\item \textbf{Intel (2023).} \textit{Intel systems in the November 2023 TOP500 (newsroom)}. \url{https://www.intel.com/content/www/us/en/newsroom/news/top500-supercomputers-nov-2023.html}. Vendor newsroom. Aurora's first entry and the other PVC systems (Dawn at Cambridge, SuperMUC-NG Phase 2 at LRZ). Slide 19.
\item \textbf{Argonne Leadership Computing Facility (2025).} \textit{Argonne researchers and supercomputers contribute to four Gordon Bell prize finalist projects (2025)}. \url{https://alcf.anl.gov/news/argonne-researchers-and-supercomputers-contribute-four-projects-named-gordon-bell-prize}. Facility news. The four 2025 Gordon Bell finalists run on Aurora, the application-scale evidence for the machine. Slide 19.
\item \textbf{Hatanpaa et al. (2025).} \textit{AERIS: Argonne Earth Systems Model for reliable and skillful predictions}. arXiv:2509.13523. The weather and climate model that sustained 10.21 EFLOP/s in mixed precision on 10,080 Aurora nodes. Slide 19.
\item \textbf{insideHPC (2023).} \textit{Intel alters HPC-AI roadmap: Rialto Bridge GPU discontinued (March 2023)}. \url{https://insidehpc.com/2023/03/intel-alters-hpc-ai-roadmap-rialto-bridge-gpu-discontinued/}. Press, secondary source. The cancellation of Ponte Vecchio's successor and the roadmap reset toward Falcon Shores. Slides 3 and 20.
\item \textbf{Tom's Hardware (2025).} \textit{Intel cancels Falcon Shores GPU for AI workloads; Jaguar Shores to be successor (January 2025)}. \url{https://www.tomshardware.com/tech-industry/artificial-intelligence/intel-cancels-falcon-shores-gpu-for-ai-workloads-jaguar-shores-to-be-successor}. Press, secondary source. The January 30, 2025 decision to keep Falcon Shores as an internal test chip and the naming of Jaguar Shores. Slides 3 and 20.
\item \textbf{Phoronix (2025).} \textit{Intel announces Crescent Island (October 2025)}. \url{https://www.phoronix.com/review/intel-crescent-island}. Press, secondary source. The Xe3P inference GPU with 160 GB of LPDDR5X announced at OCP 2025. Slides 15 and 20.
\item \textbf{StorageReview (2026).} \textit{Intel GPU D-series: 480 GB LPDDR5X, 1.3 PFLOP/s FP8, 350 W PCIe card (October 2026)}. \url{https://www.storagereview.com/news/intel-gpu-d-series-480gb-lpddr5x-1-3-pflops-fp8-350w-pcie-card}. Press, secondary source. The GPU D-series as named and specified in October 2026 (32 Xe3P cores, 1.5 TB/s, sampling); its current status. Slides 20 and 32.
\item \textbf{Intel (2026).} \textit{Intel reports third-quarter 2026 financial results}. \url{https://www.intc.com/news-events/press-releases/detail/1786/intel-reports-third-quarter-2026-financial-results}. Primary filing (press release). Data center and AI revenue of 7.3 billion dollars, growth on CPUs, the D-series as the only accelerator mentioned. Slides 16 and 20.
\item \textbf{Intel (2025).} \textit{NVIDIA and Intel to develop AI infrastructure and personal computing products (September 2025)}. \url{https://www.intc.com/news-events/press-releases/detail/1750/}. Vendor press release. The 5 billion dollar NVIDIA investment and the plan for x86 CPUs with NVLink. Slides 20, 23 and 29.
\end{itemize}
\subsection*{6. Fabrics, racks and standards (slides 21-23)}
\begin{itemize}[leftmargin=1.2em,itemsep=3pt]\sloppy
\item \textbf{UALink Consortium (2025).} \textit{UALink 1.0 white paper (April 2025)}. \url{https://ualinkconsortium.org/wp-content/uploads/2025/04/UALink-1.0-White_Paper_FINAL.pdf}. Standards body white paper. The 200G 1.0 specification: 200 Gb/s per lane, up to 1,024 accelerators, memory-semantic operations; the open scale-up fabric. Slides 21 and 23.
\item \textbf{Ultra Ethernet Consortium (Linux Foundation) (2025).} \textit{UEC launches specification 1.0 (June 2025)}. \url{https://www.linuxfoundation.org/press/uec-launches-spec-1.0}. Standards body press release. Ultra Ethernet 1.0, the scale-out counterpart of UALink and the alternative to InfiniBand. Slide 23.
\item \textbf{NVIDIA (2025).} \textit{NVIDIA unveils NVLink Fusion (press release, May 2025)}. \url{https://investor.nvidia.com/news/press-release-details/2025/NVIDIA-Unveils-NVLink-Fusion-for-Industry-to-Build-Semi-Custom-AI-Infrastructure-With-NVIDIA-Partner-Ecosystem/default.aspx}. Vendor press release. NVLink opened to partners' CPUs and accelerators on NVIDIA's terms; the answer to UALink in the standards timeline. Slide 23.
\item \textbf{Abraham (OLCF) (2023).} \textit{Frontier architecture overview}. \url{https://www.olcf.ornl.gov/wp-content/uploads/Frontier-Architecture-Overview_Abraham.pdf}. Facility presentation. The Frontier node: one Trento CPU, four MI250X, coherent Infinity Fabric CPU-GPU links, Slingshot-11 NICs; the node diagram. Slides 8, 21 and 27.
\item \textbf{Lawrence Livermore National Laboratory (2026).} \textit{El Capitan system page}. \url{https://hpc.llnl.gov/node/1269}. Facility documentation. Node counts, APU counts and the 2,889 PFLOP/s peak of the MI300A system. Slides 9 and 27.
\item \textbf{Morgan (The Next Platform) (2024).} \textit{El Capitan blazes the trail for converged CPU-GPU compute (November 2024)}. \url{https://www.nextplatform.com/2024/11/18/el-capitan-supercomputer-blazes-the-trail-for-converged-cpu-gpu-compute/}. Press, secondary source. 11,136 nodes and 44,544 APUs at the first TOP500 entry, and the case for the APU node. Slides 9 and 27.
\end{itemize}
\subsection*{7. Software stacks and portability (slides 24-26)}
\begin{itemize}[leftmargin=1.2em,itemsep=3pt]\sloppy
\item \textbf{AMD (2026).} \textit{ROCm release history (versions page)}. \url{https://rocm.docs.amd.com/en/latest/release/versions.html}. Vendor documentation. Dated ROCm releases from 5.0 to 10.x, the source of the dates in the three-stacks panel. Slide 24.
\item \textbf{AMD (2026).} \textit{HIPIFY documentation}. \url{https://rocm.docs.amd.com/projects/HIPIFY/en/latest}. Vendor documentation. hipify-clang and hipify-perl, the CUDA-to-HIP source translators behind the porting-path discussion. Slide 24.
\item \textbf{Khronos Group (2021).} \textit{Khronos releases SYCL 2020 final specification (February 2021)}. \url{https://www.khronos.org/news/press/khronos-releases-sycl-2020-final-specification}. Standards body press release. SYCL 2020 with unified shared memory and generic back ends; the specification Intel's DPC++ implements. Slide 24.
\item \textbf{Phoronix (2020).} \textit{Intel oneAPI 1.0 officially released (September 2020)}. \url{https://www.phoronix.com/review/intel-oneapi-10}. Press, secondary source. The oneAPI 1.0 specification release and the first toolkits that followed in December 2020. Slides 13 and 24.
\item \textbf{TechCrunch (2023).} \textit{The Unified Acceleration Foundation wants an open standard for accelerator programming (September 2023)}. \url{https://techcrunch.com/2023/09/19/the-unified-acceleration-foundation-wants-to-create-an-open-standard-for-accelerator-programming/}. Press, secondary source. oneAPI's move into the Linux Foundation's UXL, the governance point in the three-stacks panel. Slide 24.
\item \textbf{Codeplay (2022).} \textit{Codeplay announces oneAPI for NVIDIA and AMD GPU hardware (December 2022)}. \url{https://codeplay.com/portal/press-releases/2022/12/16/codeplay-announces-oneapi-for-nvidia-and-amd-gpu-hardware}. Vendor press release. The SYCL plugins that run DPC++ code on CUDA and HIP back ends; the portability-layer claim. Slides 24 and 25.
\item \textbf{Intel (2026).} \textit{Migrate from CUDA to C++ with SYCL (SYCLomatic and the DPC++ compatibility tool)}. \url{https://www.intel.com/content/www/us/en/developer/tools/oneapi/training/migrate-from-cuda-to-cpp-with-sycl.html}. Vendor documentation. The CUDA-to-SYCL migration tool and Intel's claim that 80 to 90 percent of code migrates automatically. Slide 24.
\item \textbf{Reguly (2023).} \textit{Evaluating the performance portability of SYCL across CPUs and GPUs from different vendors}. arXiv:2309.10075. SC-W 2023 (P3HPC). SYCL against native on Max 1100, A100 and MI250X: DPC++ within 1.2 percent of CUDA on the A100, 15.9 percent behind HIP on the MI250X, weaker on CPUs. Slide 25.
\item \textbf{Davis et al. (2025).} \textit{Taking GPU programming models to task for performance portability}. arXiv:2402.08950. ICS 2025. Seven programming models and five proxy applications across five GPUs: portability layers reach most of native speed and the spread sits in compilers and runtimes. Slide 25.
\item \textbf{ORNL and SNL (ORNL/TM-2025/4398) (2025).} \textit{Experiences with SYCL on AMD GPUs with Kokkos}. \url{https://www.osti.gov/biblio/3016989}. Laboratory technical report. Kokkos with DPC++ and AdaptiveCpp back ends on AMD GPUs; DPC++ less mature on AMD than on Intel, HIP the reference. Slide 25.
\item \textbf{Patel, Nishball and Knuhtsen (SemiAnalysis) (2024).} \textit{MI300X vs H100 vs H200 benchmark part 1: training (December 2024)}. \url{https://newsletter.semianalysis.com/p/mi300x-vs-h100-vs-h200-benchmark-part-1-training}. Analyst report, secondary source. Achieved BF16 GEMM against peak (about 720 of 989 TFLOP/s on the H100, about 620 of 1,307 on the MI300X), 64 against 450 GB/s between GPUs; the software gap in numbers. Slides 21 and 26.
\end{itemize}
\subsection*{8. Systems and procurements (slides 27-29)}
\begin{itemize}[leftmargin=1.2em,itemsep=3pt]\sloppy
\item \textbf{TOP500 (2026).} \textit{TOP500 list, June 2026}. \url{https://top500.org/lists/2026/06}. Benchmark list, open. The June 2026 top ten with accelerators, Rmax, Rpeak and power; the exascale and TOP500 tables. Slides 2, 12, 27 and 28.
\item \textbf{TOP500 (2026).} \textit{HPCG results, June 2026}. \url{https://top500.org/lists/hpcg/2026/06/}. Benchmark list, open. The HPCG ranks that complement HPL for the same machines. Slide 28.
\item \textbf{TOP500 (2026).} \textit{Green500 list, June 2026}. \url{https://top500.org/lists/green500/green500-june-2026/}. Benchmark list, open. GFLOP/s per watt: the Grace Hopper systems at the top and the MI300A systems near 66. Slide 28.
\item \textbf{HPE (2026).} \textit{HPE delivers six of the ten most powerful supercomputers (June 2026)}. \url{https://www.hpe.com/us/en/newsroom/blog-post/2026/06/hpe-delivers-six-out-of-ten-of-worlds-most-powerful-supercomputers.html}. Vendor blog. The HPE Cray EX share of the top ten and the HPL-MxP order (El Capitan, Frontier, Aurora). Slides 27 and 28.
\item \textbf{CSC (LUMI) (2021).} \textit{LUMI's full system architecture revealed}. \url{https://www.lumi-supercomputer.eu/lumis-full-system-architecture-revealed/}. Facility news. 2,978 GPU nodes of four MI250X each; the European Frontier sibling. Slides 8 and 28.
\item \textbf{CINES (2024).} \textit{Adastra's architecture}. \url{https://dci.dci-gitlab.cines.fr/webextranet/architecture/index.html}. Facility documentation. MI250X and MI300A partitions in one system; another AMD HPC deployment. Slide 8.
\item \textbf{HLRS (2025).} \textit{Hunter}. \url{https://www.hlrs.de/solutions/systems/hunter}. Facility page. 752 MI300A, the transition system before Herder. Slides 8 and 29.
\item \textbf{Oak Ridge National Laboratory (2025).} \textit{ORNL, AMD and HPE to deliver DOE's newest AI supercomputers, Discovery and Lux}. \url{https://www.ornl.gov/news/ornl-amd-and-hpe-deliver-does-newest-ai-supercomputers-discovery-and-lux}. Laboratory press release. Lux (MI355X, early 2026) and Discovery (MI430X with Venice on HPE Cray GX5000, 2028). Slides 10, 12 and 29.
\item \textbf{AMD and Eviden (2025).} \textit{AMD and Eviden to power Europe's new exascale supercomputer, the first in France}. \url{https://newsroom.amd.com/news/amd-and-eviden-to-power-europe-s-new-exascale-supe/}. Vendor press release. Alice Recoque: MI430X with Venice CPUs in 94 racks on the BXI network. Slides 12 and 29.
\end{itemize}
\subsection*{9. Benchmarks (slide 30)}
\begin{itemize}[leftmargin=1.2em,itemsep=3pt]\sloppy
\item \textbf{AMD (2024).} \textit{Engineering insights: unveiling MLPerf results on AMD Instinct MI300X}. \url{https://www.amd.com/en/blogs/2024/engineering-insights-unveiling-mlperf-results-on.html}. Vendor blog. AMD's first MLPerf submission (Inference 4.1, August 2024): eight MI300X within 2 to 3 percent of a DGX H100 on Llama 2 70B. Slide 30.
\item \textbf{AMD ROCm blog (2025).} \textit{MI325X accelerates MLPerf Inference v5.0}. \url{https://rocm.blogs.amd.com/artificial-intelligence/mi325x-accelerates-mlperf-inference/README.html}. Vendor blog. The first MI325X results, April 2025. Slide 30.
\item \textbf{AMD ROCm blog (2025).} \textit{Technical dive into AMD's MLPerf Inference v5.1 submission}. \url{https://rocm.blogs.amd.com/artificial-intelligence/mlperf-inference-v5.1/README.html}. Vendor blog. The first MI355X and first MXFP4 results (September 2025), Llama 2 70B at 2.7 times the MI325X. Slides 10 and 30.
\item \textbf{AMD (2026).} \textit{AMD delivers breakthrough MLPerf Inference 6.0 results: MI355X past one million tokens per second}. \url{https://amd.com/en/blogs/2026/amd-delivers-breakthrough-mlperf-inference-6-0-results.html}. Vendor blog. The April 2026 ratios against the B200 (1.00, 0.97, 1.19) and B300 (0.92, 0.93, 1.04) on Llama 2 70B, and the 87-GPU million-token run. Slides 10 and 30.
\item \textbf{MLCommons (2025).} \textit{MLPerf Training v5.0 results (June 2025)}. \url{https://mlcommons.org/2025/06/mlperf-training-v5-0-results}. Benchmark consortium, open. AMD's first training submission and the GB200's debut in the same round. Slide 30.
\item \textbf{AMD (2025).} \textit{Accelerating AI training: MI350 series in MLPerf Training v5.1}. \url{https://www.amd.com/en/blogs/2025/accelerating-ai-training.html}. Vendor blog. Llama 2 70B LoRA in 10.18 minutes against 9.85 (B200) and 9.59 (B300), and Llama 3.1 8B at 99.7 against 93.7 minutes. Slide 30.
\item \textbf{NVIDIA (2026).} \textit{NVIDIA Blackwell tops MLPerf Training 6.0 (developer blog, June 2026)}. \url{https://developer.nvidia.com/blog/nvidia-blackwell-tops-mlperf-training-6-0-with-industry-leading-scale-and-performance/}. Vendor blog. The 8,192-GPU GB300 results (DeepSeek-V3 671B in 2.02 minutes) that set the scale of Training 6.0, with no AMD result at that size. Slide 30.
\end{itemize}
\subsection*{10. Wrap-up and the record (slides 31-33)}
\begin{itemize}[leftmargin=1.2em,itemsep=3pt]\sloppy
\item \textbf{Lockwood (2026).} \textit{GTC 2026 recap (blog, March 2026)}. \url{https://blog.glennklockwood.com/2026/03/gtc-2026-recap.html}. Independent analyst notes, open. The NVIDIA roadmap changes of 2026 (Rubin CPX dropped, Kyber shown less dense, Feynman) against which the AMD and Intel roadmaps are read. Slide 32.
\item \textbf{ServeTheHome (2026).} \textit{NVIDIA Vera Rubin NVL72 rack at Hot Chips 2026}. \url{https://www.servethehome.com/nvidia-vera-rubin-nvl72-rack-at-hot-chips-2026/}. Press, secondary source. The rack as presented in the same Hot Chips session as AMD's MI400 talks: NVLink 6 switch with in-network compute, 800 V DC, cable-free trays. Slides 22 and 31.
\item \textbf{Wikipedia (2026).} \textit{AMD Instinct}. \url{https://en.wikipedia.org/wiki/AMD_Instinct}. Tertiary source. A dated pointer page for the Instinct generations; useful for finding primary sources, not for quoting. Slide 31.
\end{itemize}
%=====
\clearpage
\printindex
\end{document}
